% =====================================================================
%  《打通工科数学任督二脉》/ Math Rendu
%  Bridging the Threads of Engineering Mathematics
%  A Fun Lecture Note from High School to Graduate Level
%  Li Zhou — 2027
%  LaTeX source — MIT License
%  仿《打通物理任督二脉》V5 风格
% =====================================================================

\documentclass[11pt,a4paper,openany]{book}

% ---------- 中文 ----------
\usepackage{xeCJK}
\setCJKmainfont[ItalicFont=Droid Sans Fallback,BoldFont=Droid Sans Fallback]{Droid Sans Fallback}
\setCJKsansfont{Droid Sans Fallback}
\setCJKmonofont{Droid Sans Fallback}
% 西文字体——DejaVu Serif 有 em-dash 字形
\setmainfont{DejaVu Serif}
\setsansfont{DejaVu Sans}
\setmonofont{DejaVu Sans Mono}
% 把 U+2014（em-dash ——）路由到 Latin 字体——破折号显示为"——"
\xeCJKDeclareCharClass{HalfLeft}{"2014}
% 行距、段距——适中
\linespread{1.32}
\setlength{\parskip}{0.4em}
\setlength{\parindent}{2em}
% 中文章节命名——\thechapter 只放数字，"章"字由 titlesec/书签格式控制
\renewcommand{\chaptername}{第}
\renewcommand{\contentsname}{目录}
\renewcommand{\figurename}{图}
\renewcommand{\tablename}{表}
\renewcommand{\appendixname}{附录}
\renewcommand{\partname}{第}
% 让 thepart 在引用时输出"X 部分"——但 part 标题的"第 X 部分"由 titlesec 处理
\renewcommand{\thepart}{\Roman{part}}

% ---------- 页面 ----------
\usepackage[a4paper,margin=2.5cm,headheight=14pt]{geometry}

% ---------- 数学 ----------
\usepackage{amsmath,amssymb,amsthm,mathtools,bm,esint}

% ---------- 图表 ----------
\usepackage{graphicx,booktabs,array,xcolor,multirow}

% ---------- 色板 ----------
\definecolor{sectionblue}{RGB}{32,52,90}
\definecolor{sectiongold}{RGB}{180,140,40}

% ---------- 盒子 ----------
\definecolor{storyboxbg}{RGB}{255,247,224}
\definecolor{storyboxframe}{RGB}{180,140,40}
\definecolor{essbg}{RGB}{230,243,255}
\definecolor{essframe}{RGB}{40,90,180}
\definecolor{ahabg}{RGB}{235,255,235}
\definecolor{ahaframe}{RGB}{40,140,60}
\usepackage[most]{tcolorbox}
\tcbset{
  storybox/.style={colback=storyboxbg, colframe=storyboxframe, boxrule=0.6pt, arc=2pt, left=8pt, right=8pt, top=6pt, bottom=6pt, fonttitle=\bfseries, title={#1}},
  essbox/.style={colback=essbg, colframe=essframe, boxrule=0.6pt, arc=2pt, left=8pt, right=8pt, top=6pt, bottom=6pt, fonttitle=\bfseries, title={#1}},
  ahabox/.style={colback=ahabg, colframe=ahaframe, boxrule=0.6pt, arc=2pt, left=8pt, right=8pt, top=6pt, bottom=6pt, fonttitle=\bfseries, title={#1}},
}

% ---------- 列表 ----------
\usepackage{enumitem}
\setlist[itemize]{itemsep=0.1em, topsep=0.2em, parsep=0.05em, partopsep=0pt}
\setlist[enumerate]{itemsep=0.1em, topsep=0.2em, parsep=0.05em, partopsep=0pt}

% ---------- 代码 ----------
\usepackage{listings}
\lstset{
  basicstyle=\ttfamily\small,
  breaklines=true,
  numbers=left,
  numberstyle=\tiny\color{gray},
  keywordstyle=\color{blue}\bfseries,
  commentstyle=\color{olive}\itshape,
  stringstyle=\color{purple},
  showstringspaces=false,
  frame=single,
  framerule=0.3pt,
  rulecolor=\color{gray},
}

% ---------- 链接 ----------
\usepackage[bookmarks=true,bookmarksopen=false,unicode=true]{hyperref}
\hypersetup{colorlinks=true,linkcolor=blue,urlcolor=blue,citecolor=blue,
  pdftitle={打通工科数学任督二脉},
  pdfauthor={Li Zhou}}

% ---------- V5 风格章节标题 ----------
\usepackage{titlesec}
\titleformat{\chapter}[display]
  {\normalfont\huge\bfseries\color{sectionblue}}
  {\color{sectiongold}\Large 第 \thechapter\ 章}{14pt}{\Huge}
  [\vspace{6pt}{\color{sectionblue}\hrule height 1pt}]
\titleformat{\section}
  {\Large\bfseries\color{sectionblue}}
  {\raisebox{0.05em}{\textcolor{sectionblue}{\rule{0.75em}{0.75em}}}\;\thesection}{0.5em}{}
  [\vspace{2pt}{\color{sectionblue}\hrule height 0.5pt}]
\titleformat{\subsection}
  {\large\bfseries}
  {\color{sectionblue}\thesubsection}{0.5em}{}
\titlespacing*{\chapter}{0pt}{20pt}{20pt}
\titlespacing*{\section}{0pt}{1.5em}{0.6em}
\titlespacing*{\subsection}{0pt}{1em}{0.3em}
\titleformat{\part}[display]
  {\centering\normalfont\Huge\bfseries\color{sectionblue}}
  {\Large\color{sectiongold}\partname\ \thepart}{20pt}{\Huge}
  [\vspace{8pt}{\color{sectiongold}\hrule height 1pt}]

% ---------- 定理 ----------
\newtheorem{theorem}{定理}[chapter]
\newtheorem{lemma}[theorem]{引理}
\newtheorem{proposition}[theorem]{命题}
\newtheorem{corollary}[theorem]{推论}
\theoremstyle{definition}
\newtheorem{definition}[theorem]{定义}
\newtheorem{example}[theorem]{例}
\theoremstyle{remark}
\newtheorem{remark}[theorem]{注}

% =====================================================================
\begin{document}

% ---------- 封面（V5 风格） ----------
\definecolor{titleline}{RGB}{32,52,90}
\definecolor{titlegold}{RGB}{180,140,40}

\begin{titlepage}
\centering
\vspace*{3.2cm}

{\color{titleline}\hrule height 0.9pt}
\vspace{1.0cm}

{\Huge\bfseries\color{titleline} 打通工科数学任督二脉 \par}
\vspace{0.6cm}

{\color{titlegold}\hrule height 1.2pt width 6cm}
\vspace{0.8cm}

{\Large 高中到研究生的数学贯通课 \par}
\vspace{0.5cm}

{\Large\itshape Bridging the Threads of Engineering Mathematics \par}
\vspace{0.15cm}
{\large A Fun Lecture Note from High School to Graduate Level \par}

\vspace{2.5cm}

{\Large\color{titlegold}\bfseries Li Zhou \par}

\vspace{1.2cm}

{\normalsize 第一版 \quad 2027 \par}

\vfill

{\color{titleline}\hrule height 0.9pt}
\vspace{0.8cm}

{\normalsize\itshape 数学是活的——但它需要有人来贯通 \par}

\vspace{1cm}
\end{titlepage}

% ---------- 版权页 ----------
\clearpage
\thispagestyle{empty}
\noindent\textbf{《打通工科数学任督二脉——高中到研究生的数学贯通课》}\\[4pt]
Author: Li Zhou\\
Email: \texttt{lizhou\_alfred2011@hotmail.com}\\
First edition: 2027\\[8pt]
本书采用 MIT 许可证开源发布。配套 Python 代码模块亦遵循 MIT 许可证。\\
您可以自由阅读、复制、分发、修改本书内容，但请保留原作者署名。\\[16pt]

文责自负。

% ---------- 引言句 ----------
\clearpage
\thispagestyle{empty}
\vspace*{4cm}
\begin{center}
\itshape
\itshape
"数学家最大的敌人不是公式——是孤立。\\[6pt]
教材把它们分开教——研究把它们拼起来用。\\[6pt]
这本书想做的是后者。"\\[20pt]
\normalfont 写在前面
\end{center}

\clearpage
\tableofcontents

\mainmatter

% =====================================================================
%  第 0 章 · 序言 / Preface
%  Why Bridge Engineering Mathematics
% =====================================================================

\chapter*{序言：为什么要打通工科数学}
\addcontentsline{toc}{chapter}{序言：为什么要打通工科数学}
\markboth{序言}{为什么要打通工科数学}

\section*{一段个人前情}

我读大一的时候——第一学期的高等数学，期中考完，全班 36 个人，挂了 11 个。

老师讲极限的 $\varepsilon\text{-}\delta$，讲了三周。讲完之后我问同寝室的舍友：你听懂了吗？他说没有。我说我也没有。我们互相安慰：先背下来再说，反正期末考的不是这个，考的是题目。

后来我才明白——那个学期我们浪费的不是三周，是\textbf{整本高等数学的入场券}。我们没听懂极限——所以导数只是"那个 $\frac{f(x+h)-f(x)}{h}$ 的极限"——一个我们也不真懂的东西，再用一个不真懂的东西取极限——能信吗？信不了——但考试还是过了。

\medskip

到了大三做毕业设计，我要解一个二阶 ODE。我打开《高等数学》——翻到 ODE 那一章——里面写了"分离变量法、待定系数法、特征方程法"——三种方法都给了。我试着套——套不上。

为什么？因为教材的例题是 $y'' + 4y = 0$，参数是漂亮的整数。我的方程是 $y'' + 0.83 y' + 12.7 y = \sin(3.1 t)$。教材说"令 $y = e^{rt}$，代入求 $r$"——可我代入完之后，$r$ 是个复数——教材没继续讲了。

我那天晚上盯着这个方程看了三个小时——最后用 Python 的 \texttt{scipy.integrate.odeint} 一行代码把它数值解了。

\textbf{大学四年的数学——在毕业设计前那个晚上——被一行 \texttt{scipy} 取代。}

\medskip

这是个隐喻——但它不只是隐喻。\textbf{中国工科生的数学教育——和真实工程世界——存在三道断层}。这本书就是为这三道断层而写。

\section*{工科数学的三道断点}

\subsection*{断点 1：从高中到大学}

高中数学像"算术"——\textbf{操作明确、答案唯一、训练题海}。

大学数学突然变成"\textbf{$\varepsilon\text{-}\delta$ + 链式法则 + 多重积分 + 矩阵}"——\textbf{概念抽象、计算冗长、考题逼人}。

大一下学期——\textbf{$30\%\sim 40\%$ 的学生数学就掉队了}——剩下的 60\% 也只是在"硬刷题"。等到大二学《线性代数》——他们已经不知道矩阵是什么——只知道"上三角，下三角，行列式按某一行展开"。

\textbf{断点 1 的本质}——不是难——是"\textbf{从直觉学习到符号操作的切换}"——而国内教材几乎\textbf{完全跳过了直觉这一步}。

\subsection*{断点 2：从大学到研究生 / 工程岗}

大学数学讲的是"\textbf{解出闭合解}"——研究生 / 工程岗用的是"\textbf{数值算}"。

\textbf{大学教 $y'' + 4y = \cos(2t)$}——这种方程有共振——很美——但工程里 \textbf{99\% 的方程没有共振}——\textbf{99\% 的方程没有闭合解}——必须数值。

\textbf{大学教矩阵的本征值 $A\bm{v}=\lambda\bm{v}$}——给你一个 $3\times 3$ 矩阵让你手算 $\det(A-\lambda I)=0$。

\textbf{研究生 / 工业用的是 $1000\times 1000$ 矩阵}——\textbf{你必须用 \texttt{numpy.linalg.eig}}——\textbf{手算的方法你一辈子都不会再用}。

\textbf{断点 2 的本质}——大学数学和工业数学——\textbf{是两套不同的数学}——但教材里只有第一套——\textbf{第二套你必须自学}。

\subsection*{断点 3：经典数学 vs 现代工程}

\textbf{现代工程的数学发生了一次革命}——大约从 \textbf{2010 年}（深度学习兴起）开始——\textbf{加速到 2020 年（ChatGPT 训练）}。

具体表现：

\begin{itemize}
\item \textbf{线性代数从"辅助课"变成"主菜"}——21 世纪所有 AI / ML / 数据科学都跑在矩阵上
\item \textbf{优化从"运筹学小话题"变成"核心引擎"}——SGD / Adam / 损失函数曲面成为日常
\item \textbf{概率从"实验数据分析"变成"模型本身"}——贝叶斯、生成模型、扩散模型
\item \textbf{数值方法从"工程师小工具"变成"科研主力"}——\texttt{numpy} + \texttt{scipy} + \texttt{PyTorch}
\end{itemize}

但\textbf{中国大学的工科数学教材}——\textbf{基本停留在 1990 年代}——同济高数、线代、概率论——\textbf{三本书加起来——讲 SVD 的不超过一段话——讲梯度下降的没有一章——讲 PageRank 的根本没有}。

\textbf{断点 3 的本质}——\textbf{教材跟不上工业}——\textbf{老师跟不上时代}——学生学完之后\textbf{对现代工程世界几乎一无所知}——必须补课。

\medskip

\textbf{这本书的目的——就是打通这三个断点。}

\section*{为什么书名叫"打通任督二脉"}

\textbf{中国读者都懂这个梗}——\textbf{武侠小说里}——\textbf{一旦"打通任督二脉"——内力大增、武功跨越式提升}。\textbf{数学也是这样}——\textbf{一旦你"打通"了几条核心脉络——所有应用数学都豁然开朗}。

\textbf{但"任督二脉"四个字}——\textbf{字面上就告诉你}：\textbf{是\textbf{两条}脉}——\textbf{不是六条}——\textbf{不是十条}——\textbf{就两条}。

\textbf{数学的这两条脉是什么}？——\textbf{大学数学教学早就给了答案}：

\begin{center}
\fbox{\parbox{0.85\textwidth}{\centering
\textbf{\large 任脉 = 分析 Analysis} \quad / \quad \textbf{\large 督脉 = 代数 Algebra}\\[6pt]
\textit{连续 vs 离散 \quad / \quad 极限 vs 结构 \quad / \quad 变化率 vs 变换}
}}
\end{center}

\textbf{任脉}贯穿\textbf{Ch 1-5, 9-11}——\textbf{覆盖所有"连续"的数学}：微积分、矢量分析、ODE/PDE、复变、概率统计。

\textbf{督脉}贯穿\textbf{Ch 6-8}——\textbf{覆盖所有"线性结构"的数学}：矩阵入门、本征值核心、SVD/PCA/PageRank 应用。

\textbf{Ch 12 数值方法}与\textbf{Ch 13 优化}——\textbf{是任督相通的两章}——\textbf{ML 训练公式 $\theta_{t+1}=\theta_t-\eta\nabla L(\theta_t)$ 就在这里}——\textbf{用分析的梯度求解代数的最小化}。

\textbf{为什么是分析 / 代数的二分}——\textbf{有四个理由}：

\begin{enumerate}
\item \textbf{数学史的天然分界}——\textbf{17 世纪 Newton/Leibniz 开启分析}；\textbf{19 世纪中叶 Cayley 开启线代}——\textbf{各自独立发展 200 年}
\item \textbf{教学路径的天然分界}——\textbf{大一上 = 高数（分析）}；\textbf{大一下 / 大二上 = 线代（代数）}——\textbf{中国所有工科客观就是两段}
\item \textbf{研究语言的天然分界}——\textbf{所有应用数学问题}——\textbf{要么用分析（极限 / 积分 / 微分方程）}——\textbf{要么用代数（矩阵 / 本征值）}——\textbf{要么两者混合（数值线代）}
\item \textbf{"二脉"字面}——\textbf{书名说二就是二}——\textbf{列六条线从根上和书名打架}
\end{enumerate}

\textbf{打通任脉}——\textbf{把任何工程问题翻译成微分方程或积分}。\textbf{打通督脉}——\textbf{把任何高维问题翻译成矩阵和本征值}。\textbf{打通"任督二脉"}——\textbf{知道它们在 Ch 12/13 怎么相遇}——\textbf{这就是 ML、AI、数据科学的全部数学骨架}。

\section*{这本书的方法论：7 件套 + 贯通}

每一章——我用同一个\textbf{7 件套结构}：

\begin{enumerate}
\item \textbf{写在前面}——为什么要学这一章 / 我个人的体会
\item \textbf{一句话本质}——这一章如果只能记一句话，记哪句
\item \textbf{教材里你看到的}——同济 / Stewart / Zill 怎么讲——它们的优缺点
\item \textbf{直觉}——这个数学到底在做什么——几何 / 物理 / 工程意义
\item \textbf{数学}——必要的公式、定理、推导——\textbf{节制}
\item \textbf{Aha 例子}——一个工业级真实场景——让你说"原来如此"
\item \textbf{代码与思考题}——Python 模块的 demo + 6 道思考题
\end{enumerate}

\textbf{加一个"小故事盒"}——每章半页——讲这个数学是怎么诞生的——\textbf{Newton 在瘟疫之年}、\textbf{Fourier 被法国学院拒稿三次}、\textbf{Larry Page 的博士论文}——数学不是凭空冒出来的——\textbf{是某些人在某个时代，解决某个问题}。

\section*{风格定调：参考吴军《数学之美》}

\textbf{我不想写}：

\begin{itemize}
\item 同济风格——题海 + 应试 + 缺乏直觉
\item Zill 风格——英文 + 偏理论 + 翻译腔
\item 网红科普——标题党 + 浅薄 + 段子
\item 哲学随笔——务虚 + 不接地气
\end{itemize}

\textbf{我想写的是}（参考吴军《数学之美》）：

\begin{itemize}
\item 从工程问题切入——不是"概念→应用"，而是"问题→数学"
\item 公式量节制——必要才出现，不堆砌
\item 故事 + 人物 + 历史——数学不是凭空发明
\item 接地气——中国工程师的语言习惯
\item 克制的优雅——不浮夸、不嬉笑
\item 应用导向——Google / 微软 / AI 公司的真实场景
\end{itemize}

气质参照——\textbf{像一位老工程师在咖啡馆给后辈讲数学}——严肃但不严厉——\textbf{偶尔幽默}——但每章只一两处——\textbf{多了就过度}。

\section*{这本书的目标读者}

\textbf{核心读者}：

\begin{itemize}
\item 工科本科生（机械、电子、土木、化工、计算机、材料、能源……）
\item 跨学科研究生（物理转 AI、化学转计算、生物转生信……）
\item 在职工程师（基础不牢 / 想转 ML / 想做科研）
\item 数学焦虑者（高考数学就没学好——但工作里又躲不掉）
\end{itemize}

\textbf{不适合的读者}：

\begin{itemize}
\item 想看严格证明的人——请去读 Apostol / Rudin / Folland
\item 想刷考研真题的人——请去买李永乐 / 张宇
\item 想纯学应用的人——请去看 \texttt{numpy} 文档
\end{itemize}

\textbf{本书的定位}——\textbf{介于"严格"与"应用"之间——介于"理论"与"代码"之间}——一本\textbf{贯通式}的工科数学讲义。

\section*{与《打通物理任督二脉》的关系}

这是"\textbf{打通系列}"的第二弹。

\begin{itemize}
\item \textbf{品牌呼应}——同样的"7 件套" + 数值实验 + 双语 + GitHub 开源
\item \textbf{风格继承}——同一位作者——同样的"老工程师在咖啡馆"气质
\item \textbf{内容独立}——本书数学为主体——\textbf{物理只作为应用之一}——\textbf{你不需要先读物理版}
\item \textbf{交叉引用}——本书 Ch 8 量子比特 $\leftrightarrow$ 物理版 Ch 17；本书 Ch 13 优化 $\leftrightarrow$ 物理版 Ch 21
\end{itemize}

\textbf{建议阅读顺序}：

\begin{itemize}
\item 工科生 / 工程师——\textbf{直接读本书}
\item 物理学生——先读物理版 Ch 0（数学预备）——再来读本书 Ch 6-8 + Ch 13
\item AI / ML 从业者——重点读本书 Ch 6-8 + Ch 10-11 + Ch 13
\end{itemize}

\section*{怎么用这本书}

\textbf{完整阅读（4-6 周）}——Ch 0 $\to$ Ch 1-13——按顺序——每章 1-2 个晚上。

\textbf{速成路径}：

\begin{tabular}{ll}
\toprule
\textbf{目标} & \textbf{重点章节} \\
\midrule
补线代 & Ch 6-8（3 章） \\
补 ML 数学 & Ch 6-8 + Ch 10-11 + Ch 13 \\
补现代数学 & Ch 8 + Ch 10 + Ch 12-13 \\
补 PDE & Ch 4-5 + Ch 12 \\
补优化 & Ch 13 \\
\bottomrule
\end{tabular}

\textbf{代码并行}——每章 §7 给出 Python 模块的 5-6 个 demo——\textbf{边读边跑}——\textbf{不跑就白读}。

\textbf{思考题}——每章 6 道——\textbf{留给自己想}——书末\textbf{不给答案}——这是工程师的习惯：\textbf{答案在代码里、在论文里、在和同事的讨论里}——\textbf{不在书后的"答案部分"}。

\section*{诚实声明：这本书不是什么}

\begin{itemize}
\item \textbf{不是}严格的数学教材——想看证明，请读 Apostol、Rudin、Folland
\item \textbf{不是}考研辅导书——不会教你怎么对付李永乐
\item \textbf{不是}百科全书——\textbf{有意省略很多内容}：测度论、张量分析、群论、变分法、随机过程、信息论——这些留给\textbf{后续专题书}
\item \textbf{不是}"零基础学数学"——你至少需要\textbf{高中数学的基础}
\item \textbf{不是}"一周速成数学"——\textbf{不存在这种东西}——\textbf{6 周是合理的}
\end{itemize}

\section*{最后一句}

这本书在 2027 年上半年完成。

\medskip

\textbf{数学是活的}——\textbf{但它需要有人来贯通}。

\medskip

\hfill Li Zhou
\hfill 2027 春

\clearpage


% 全书脉络图——分析/代数 二脉框架（紧接序言之后）
% =====================================================================
%  数学版全书脉络图 / Math Threads Map
%  二脉 = 分析（任脉）+ 代数（督脉）
% =====================================================================

\cleardoublepage
\thispagestyle{empty}

\vspace*{1cm}

\begin{center}
{\huge\bfseries 全书脉络图 \par}
\vspace{0.2cm}
{\large\itshape The Map of Two Meridians of Engineering Math \par}
\end{center}

\vspace{0.8cm}

\section*{两条脉——分析与代数}

\textbf{这本书叫}"\textbf{打通工科数学任督二脉}"——\textbf{"二脉"不是修辞}——\textbf{工科数学的两条根本脉络}：

\begin{center}
\fbox{\parbox{0.85\textwidth}{\centering
\textbf{\Large 任脉 = 分析 Analysis} \quad \quad \textbf{\Large 督脉 = 代数 Algebra} \\[8pt]
\textit{连续 vs 离散 \quad / \quad 极限 vs 结构 \quad / \quad 变化率 vs 变换}
}}
\end{center}

\vspace{0.5cm}

\textbf{为什么是这两条}——\textbf{回到数学史}：

\textbf{17 世纪}——\textbf{Newton + Leibniz 发明微积分}——\textbf{开启分析}（任脉的起点）。

\textbf{19 世纪}——\textbf{Cayley + Hamilton + Sylvester 发明矩阵代数}——\textbf{开启线性代数}（督脉的起点）。

\textbf{20 世纪}——\textbf{Hilbert 把两条脉拼起来}——\textbf{泛函分析 = 分析 × 代数}——\textbf{量子力学的数学骨架}。

\textbf{21 世纪}——\textbf{机器学习 = 优化（分析）× 矩阵（代数）}——\textbf{两条脉在同一个公式 $\theta_{t+1} = \theta_t - \eta \nabla L(\theta_t)$ 里相遇}。

\section*{任脉 / 分析 — 五条支脉}

\textbf{任脉}贯穿\textbf{Ch 1-5, 9-11}——\textbf{覆盖工科数学里所有"连续"的部分}：

\begin{center}
\begin{tabular}{|c|p{4.5cm}|p{5.5cm}|}
\hline
\textbf{支脉} & \textbf{主题} & \textbf{覆盖章节} \\
\hline
任 · 一 & 极限与微积分（一元） & Ch 1 微积分入门 \\
\hline
任 · 二 & 多元微积分 & Ch 2 多元微积分 / Ch 3 矢量分析 \\
\hline
任 · 三 & 微分方程 & Ch 4 ODE / Ch 5 PDE \\
\hline
任 · 四 & 复变 + Fourier & Ch 9 复变函数 \\
\hline
任 · 五 & 概率与不确定性 & Ch 10 概率 / Ch 11 统计 \\
\hline
\end{tabular}
\end{center}

\section*{督脉 / 代数 — 三条支脉}

\textbf{督脉}贯穿\textbf{Ch 6-8}——\textbf{覆盖所有"线性结构"的部分}：

\begin{center}
\begin{tabular}{|c|p{4.5cm}|p{5.5cm}|}
\hline
\textbf{支脉} & \textbf{主题} & \textbf{覆盖章节} \\
\hline
督 · 一 & 线代入门 & Ch 6（矩阵 = 线性变换） \\
\hline
督 · 二 & 线代核心 & Ch 7（本征值 + 谱定理） \\
\hline
督 · 三 & 线代应用 & Ch 8（SVD + PCA + PageRank） \\
\hline
\end{tabular}
\end{center}

\section*{交汇——任督相通的两章}

\textbf{Ch 12 数值方法}：\textbf{把分析问题（积分、ODE、根）转化为代数运算（矩阵求逆、本征值）}——\textbf{这是分析与代数相遇的工程层}。

\textbf{Ch 13 优化}：\textbf{用分析的工具（梯度、Hessian）求解代数的问题（二次型、最小化）}——\textbf{这是 21 世纪工程数学的核心}——\textbf{ML 训练的全部数学都在这里}。

\section*{二脉 × 章节对照表}

\textbf{下表标出每章属于哪条脉}——\textbf{$\bullet$ 是主脉}——\textbf{$\circ$ 是支援}——\textbf{$\star$ 是任督相通}：

\vspace{0.3cm}

\begin{center}
\small
\begin{tabular}{|c|p{2.6cm}|c|c|c|c|c|c|c|c|c|c|c|c|c|c|}
\hline
\textbf{脉} & \textbf{支脉} & 1 & 2 & 3 & 4 & 5 & 6 & 7 & 8 & 9 & 10 & 11 & 12 & 13 & 14 \\
\hline
\multirow{5}{*}{\parbox{0.9cm}{\centering\textbf{任·分析}}}
& 微积分一元   & $\bullet$ &   &   &   &   &   &   &   &   &   &   & $\circ$ & $\circ$ & \\
\cline{2-16}
& 多元/矢量    & & $\bullet$ & $\bullet$ &   &   &   &   &   &   &   &   & $\circ$ & $\circ$ & \\
\cline{2-16}
& ODE/PDE      & &   &   & $\bullet$ & $\bullet$ &   &   &   &   &   &   & $\circ$ &   & \\
\cline{2-16}
& 复变/Fourier & &   &   &   & $\circ$ &   &   &   & $\bullet$ &   &   & $\circ$ &   & \\
\cline{2-16}
& 概率/统计    & &   &   &   &   &   &   & $\circ$ &   & $\bullet$ & $\bullet$ &   & $\circ$ & $\circ$ \\
\hline
\multirow{3}{*}{\parbox{0.9cm}{\centering\textbf{督·代数}}}
& 线代入门     & &   &   &   &   & $\bullet$ &   &   &   &   &   & $\circ$ &   & \\
\cline{2-16}
& 线代核心     & &   &   &   &   &   & $\bullet$ & $\circ$ &   & $\circ$ &   & $\circ$ & $\circ$ & \\
\cline{2-16}
& 线代应用     & &   &   &   &   &   &   & $\bullet$ &   & $\circ$ & $\circ$ & $\circ$ & $\circ$ & \\
\hline
\multirow{2}{*}{\parbox{0.9cm}{\centering\textbf{相通}}}
& 数值方法     & &   &   &   &   &   &   &   &   &   &   & $\star$ & $\circ$ & \\
\cline{2-16}
& 优化         & & $\circ$ &   &   &   &   & $\circ$ & $\circ$ &   & $\circ$ &   & $\circ$ & $\star$ & \\
\hline
\end{tabular}
\end{center}

\vspace{0.3cm}

\begin{center}
\footnotesize
\textbf{$\bullet$} = 这章是该支脉的核心打通点 \qquad
$\circ$ = 这章涉及该支脉 \qquad
$\star$ = 任督相通点
\end{center}

\section*{阅读建议}

\textbf{第一遍阅读}（按章号顺序）：\textbf{Ch 1 $\to$ Ch 13}——\textbf{遇到 $\bullet$ 就重点学}——\textbf{遇到 $\circ$ 注意它和前面 $\bullet$ 的呼应}。

\textbf{第二遍阅读}（追一条脉）：

\begin{itemize}
\item \textbf{追任脉分析}：Ch 1 $\to$ Ch 2 $\to$ Ch 4 $\to$ Ch 5 $\to$ Ch 9 $\to$ Ch 10——\textbf{从一元导数到 PDE 到 Fourier 到概率}
\item \textbf{追督脉代数}：Ch 6 $\to$ Ch 7 $\to$ Ch 8——\textbf{从矩阵到本征值到 SVD/PCA/PageRank}——\textbf{是个完整三部曲}
\item \textbf{追任督相通}：Ch 2 + Ch 6 $\to$ Ch 7 + Ch 9 $\to$ Ch 12 $\to$ Ch 13——\textbf{Jacobian + 厄米矩阵 + 数值线代 + 优化 = ML 的全部数学}
\end{itemize}

\section*{读法速成路线}

\begin{itemize}
\item \textbf{补线代（最常见需求）}：Ch 6-8 三章 + Ch 13 优化
\item \textbf{补 ML 数学}：Ch 2-3（梯度）+ Ch 6-8（矩阵）+ Ch 10-11（概率）+ Ch 13（SGD/Adam）
\item \textbf{补 PDE（科学计算）}：Ch 1-5 + Ch 12 数值
\item \textbf{补现代数学（数据科学方向）}：Ch 8 SVD + Ch 10 概率 + Ch 13 优化
\end{itemize}

\vspace{0.5cm}

\textbf{打通任脉}——\textbf{你看懂了所有"连续 + 变化率"的工程数学}。

\textbf{打通督脉}——\textbf{你看懂了所有"高维 + 结构"的工程数学}。

\textbf{打通任督二脉}——\textbf{两条脉相通的地方都熟}——\textbf{整本工科数学的"任督二脉"由此通畅}——\textbf{再看 PyTorch / numpy / scipy 的源码}——\textbf{你能识别每一行用的是哪条脉的工具}。

\clearpage


\part{微积分与场 / Calculus and Fields}
% =====================================================================
%  第 1 章 · 微积分入门 / Calculus I
%  极限、导数、积分、应用
%  小故事：Newton 与瘟疫之年
%  Aha：短信流量预测
% =====================================================================

\chapter{微积分入门 / Calculus I}
\label{ch:calc1}

\begin{quote}\itshape
1665 年，伦敦闹瘟疫。剑桥大学停课。年轻的牛顿回到林肯郡的老家躲避。\\
这个 22 岁的青年——闲在家里——在那一年发明了三件大事：\\
\textbf{微积分、万有引力、光学}。\\
我们这一章——就讲他在那个瘟疫之年想清楚的第一件事：\textbf{导数}。
\end{quote}

\section{写在前面 / Preface}

\textbf{我刚学微积分时}——和大多数人一样——\textbf{我背的是公式}。

$\sin'(x) = \cos(x)$。$\int x^n \, dx = \frac{x^{n+1}}{n+1}$。链式法则、分部积分、换元。

\textbf{大一上半学期}——\textbf{我整整考了 91 分}——但你问我导数是什么——我说不上来。我会算——但我\textbf{不懂}。

到了\textbf{大三}——我做毕业设计——\textbf{需要预测一个温度信号的"上升速度"}——我打开 Python：

\begin{lstlisting}[language=Python]
import numpy as np
dy_dx = np.gradient(y, x)
\end{lstlisting}

\textbf{一行代码——做完了导数}。

那一刻我才反应过来——\textbf{我大一背的所有求导公式}——在工业里——\textbf{我一个都没用过}。

\textbf{真正用的是}：

\begin{itemize}
\item 导数是\textbf{变化率}——\textbf{$\Delta y / \Delta x$ 的极限}
\item 导数是\textbf{斜率}——\textbf{切线的斜率}
\item 导数是\textbf{速度}——位置 $\to$ 速度 $\to$ 加速度
\item 导数告诉我们\textbf{何时增、何时减、何时到顶}——这就是\textbf{极值}
\end{itemize}

\textbf{这一章就把这四件事说清楚}——\textbf{再多的求导公式}——\textbf{都是这四件事的变体}。

\section{一句话本质 / The Essence in One Sentence}

\begin{tcolorbox}[essbox={一句话本质}]
\textbf{中文}：\textbf{导数是变化率——积分是累积量——它们互为逆运算（微积分基本定理）}。\\[6pt]
\textbf{English}: The derivative is a rate of change; the integral is an accumulated quantity; they are inverse operations (the Fundamental Theorem of Calculus).
\end{tcolorbox}

记住这一句话——\textbf{这一章已经讲完了一半}。

\section{教材里你看到的 / What the Textbook Tells You}

\subsection{同济《高等数学》}

同济大学应用数学系主编的\textbf{《高等数学》（第七版）}——是中国 90\% 工科生的\textbf{第一本大学数学教材}。

它的章节是这样的：

\begin{itemize}
\item 第一章：函数与极限（$\varepsilon\text{-}\delta$ 定义）
\item 第二章：导数与微分
\item 第三章：微分中值定理与导数的应用
\item 第四章：不定积分
\item 第五章：定积分
\item 第六章：定积分的应用
\end{itemize}

\textbf{优点}：完整、严格、覆盖国内考研所有考点。

\textbf{缺点}：

\begin{enumerate}
\item \textbf{$\varepsilon\text{-}\delta$ 放在第一章}——把 90\% 的学生劝退在起跑线
\item \textbf{题量巨大 + 偏怪}——大量"凑微分"、"三角换元"——\textbf{工业里这辈子都用不到}
\item \textbf{应用部分少}——只有"求面积、求体积、求功"——\textbf{现代工程一个都不沾}
\end{enumerate}

\subsection{Stewart《Calculus》}

\textbf{Stewart} 的《Calculus》是\textbf{美国大学的国民微积分}——\textbf{1200 页}——\textbf{彩色印刷}——\textbf{每章都有 100 道题}。

它比同济好的地方：\textbf{图多、应用多、写得活}。

它的缺点：\textbf{1200 页太长了}——\textbf{你读不完}——\textbf{而且它讲的是美国人的应用（高尔夫球弹道、糖果销量）}——\textbf{中国工程师看着隔膜}。

\subsection{Zill《Advanced Engineering Mathematics》}

\textbf{Zill} 是工科研究生用的——\textbf{偏理论 + 偏 PDE}——\textbf{不适合入门}。

\subsection{我的策略}

\textbf{本章不会重复同济和 Stewart 的内容}——\textbf{它们都已经够全了}。

\textbf{本章只做三件事}：

\begin{enumerate}
\item \textbf{把直觉讲清楚}——为什么导数 = 切线斜率——为什么积分 = 面积——\textbf{这是同济不讲的}
\item \textbf{把基本定理讲透}——\textbf{微分和积分是逆运算}——\textbf{这是整本微积分的中心}
\item \textbf{把 Aha 例子讲到能跑代码}——\textbf{短信流量预测}——\textbf{用 Python 的 \texttt{numpy.gradient}}
\end{enumerate}

\begin{tcolorbox}[storybox={\Large 小故事：Newton 与瘟疫之年}]
\textbf{1665 年}——\textbf{伦敦闹瘟疫}——\textbf{剑桥大学停课}。
22 岁的 \textbf{Isaac Newton}（刚拿到学士学位）——\textbf{回到林肯郡 Woolsthorpe 的老家}——\textbf{在那里待了 18 个月}。

\medskip

那 18 个月——他\textbf{独自一人}——\textbf{想清楚了三件事}：

\begin{enumerate}
\item \textbf{微积分}（他叫"流数法"，fluxions）——\textbf{为了算行星运动}
\item \textbf{万有引力}——\textbf{为了解释 Kepler 三大定律}
\item \textbf{光的分解}（用棱镜实验）——\textbf{为了解释颜色}
\end{enumerate}

\medskip

\textbf{在德国}——\textbf{Gottfried Leibniz}——\textbf{独立发明了同样的东西}——\textbf{用了不同的符号}（$dx, dy, \int$——这是我们今天用的版本）。

\textbf{Newton 和 Leibniz 后来为"谁先发明的"吵了 30 年}——\textbf{学术界的第一场国际口水仗}。

\medskip

\textbf{微积分不是凭空发明的}——\textbf{它是为了解决"行星怎么绕太阳运动"这个 2000 年没人能算清的问题}。

\textbf{17 世纪没有 Python}——\textbf{没有计算器}——\textbf{只有羽毛笔}——他们\textbf{用纯思想}搞出了让\textbf{人类登月}的数学工具。

\medskip

你今天学的每一招——\textbf{背后都有这样一群人}。
\end{tcolorbox}

\section{直觉 / Intuition}

\subsection{极限：让"无穷小"变成"任意小"}

\textbf{极限的核心思想——一句话}：

\begin{quote}
\textbf{"我想要的精度——我能做到。"}\\[4pt]
\itshape "Whatever precision you want, I can achieve."
\end{quote}

举例：$\lim_{x \to 2} x^2 = 4$。这句话的意思——\textbf{不是} "$x^2$ 在 $x=2$ 处等于 4"（虽然确实等于）——\textbf{而是} "无论你想让 $x^2$ 离 4 多近——我都能让 $x$ 离 2 足够近，就达到了"。

\textbf{这就是 $\varepsilon\text{-}\delta$ 的本质}——\textbf{你给我 $\varepsilon$（你想要的精度）——我给你 $\delta$（我能达到的范围）}——\textbf{这是一种"挑战与回应"的游戏}。

\textbf{同济教材的错误}——把 $\varepsilon\text{-}\delta$ \textbf{当成一个"严格定义"塞给学生}——\textbf{没说它是干什么的}——\textbf{学生当然懵}。

\textbf{真正的故事是}：

\begin{itemize}
\item Newton / Leibniz 时代——大家用"无穷小量 $dx$"算微积分——但\textbf{无穷小到底是什么——没人说得清}
\item Bishop Berkeley 1734 年写文章批评——\textbf{"你们的 dx 一会儿是 0 一会儿不是 0——这是数学还是巫术？"}
\item 整整 100 年——数学家无法回应这个质疑——他们\textbf{知道公式对、不知道为什么对}
\item 直到 1820 年代——\textbf{Cauchy} 引入 $\varepsilon\text{-}\delta$——\textbf{Bolzano 和 Weierstrass 完善}——\textbf{终于把"无穷小"这个概念用"任意小"代替}
\end{itemize}

\textbf{$\varepsilon\text{-}\delta$ 不是为了折磨大一新生发明的}——\textbf{是为了回应 Berkeley 的质疑、为了让 200 年来"半懂半不懂"的微积分获得严格基础}。

\subsection{导数 = 切线斜率 / 变化率}

\textbf{导数的定义}：

\begin{equation}
f'(x_0) = \lim_{\Delta x \to 0} \frac{f(x_0 + \Delta x) - f(x_0)}{\Delta x}
\end{equation}

\textbf{这个公式的几何意义}（图意，请脑补）：

\begin{itemize}
\item 取 $x_0$ 附近的两点 $(x_0, f(x_0))$ 和 $(x_0 + \Delta x, f(x_0 + \Delta x))$
\item 它们的连线——\textbf{割线}——斜率为 $\frac{f(x_0+\Delta x) - f(x_0)}{\Delta x}$
\item 让 $\Delta x \to 0$——\textbf{两点合并成一点}——\textbf{割线变成切线}
\item \textbf{切线的斜率——就是导数 $f'(x_0)$}
\end{itemize}

\textbf{物理意义}：\textbf{位置 $\to$ 速度 $\to$ 加速度}——\textbf{这是微积分发明的第一个动机}。

\begin{example}[导数的多重身份]
\textbf{同一个公式 $\frac{dy}{dx}$——在不同领域有不同名字}：
\begin{itemize}
\item 数学：\textbf{斜率 / 变化率}
\item 物理：\textbf{速度 / 加速度 / 力}（如果 $y$ 是位置）
\item 化学：\textbf{反应速率 / 浓度变化率}
\item 经济：\textbf{边际成本 / 边际收益}
\item 信号处理：\textbf{瞬时频率 / 上升沿斜率}
\item ML：\textbf{梯度 / 损失下降方向}
\end{itemize}
\textbf{它们是同一个东西}——\textbf{不同领域起了不同名字}——\textbf{这就是数学的力量}。
\end{example}

\subsection{积分 = 累积量 / 面积}

\textbf{积分的两个面孔}：

\textbf{不定积分}：$\int f(x) \, dx = F(x) + C$——\textbf{"反求导"}——找一个 $F$ 使得 $F'=f$。

\textbf{定积分}：$\int_a^b f(x) \, dx = \lim_{n\to\infty} \sum_{i=1}^n f(x_i^*) \Delta x$——\textbf{"加起来取极限"}——直观上是\textbf{$f(x)$ 在 $[a,b]$ 上与 $x$ 轴围的有向面积}。

\textbf{为什么这两个东西用同一个符号 $\int$ ？}——\textbf{因为它们等价}——\textbf{这就是微积分基本定理}。

\subsection{微积分基本定理 / Fundamental Theorem of Calculus}

\begin{theorem}[微积分基本定理]
若 $f$ 在 $[a,b]$ 上连续，$F$ 是 $f$ 的原函数（即 $F'=f$），则
\begin{equation}
\int_a^b f(x) \, dx = F(b) - F(a)
\end{equation}
\end{theorem}

\textbf{这是整本微积分的"任督二脉"}——\textbf{它把"求面积（积分）"和"求斜率（导数）"连成了一件事}。

\textbf{为什么这件事这么重要}？

\textbf{在 Newton 之前}——\textbf{阿基米德}（公元前 250 年）就会算抛物线下面的面积——\textbf{用穷竭法}——\textbf{每算一个图形要花几页纸的工作}。

\textbf{Newton 和 Leibniz 之后}——\textbf{只需要找到原函数 $F$——做一个减法 $F(b)-F(a)$——就完了}。

\textbf{2000 年的难题——被一个定理简化成一个减法}——\textbf{这就是它的伟大}。

\section{数学 / The Math}

\subsection{极限的 $\varepsilon\text{-}\delta$ 定义}

\begin{definition}[函数极限]
设 $f(x)$ 在 $x_0$ 的某个去心邻域内有定义。称 $\lim_{x\to x_0} f(x) = L$，若
\[
\forall \varepsilon > 0,\ \exists \delta > 0,\ \text{使得 } 0<|x-x_0|<\delta \Rightarrow |f(x)-L|<\varepsilon.
\]
\end{definition}

\textbf{我不展开这个定义的应用}——\textbf{在工程里你几乎用不到它}——\textbf{你需要知道的只是：极限是"逼近"——但不一定相等}。

\textbf{重要极限}（背熟这几个就够了）：
\begin{align}
\lim_{x \to 0} \frac{\sin x}{x} &= 1 \\
\lim_{x \to 0} \frac{1 - \cos x}{x^2} &= \frac{1}{2} \\
\lim_{x \to \infty} \left(1 + \frac{1}{x}\right)^x &= e
\end{align}

\subsection{导数的运算法则}

\textbf{记住四件事——你能求 95\% 工程函数的导数}：

\begin{enumerate}
\item \textbf{加法}：$(f+g)' = f' + g'$
\item \textbf{乘法}：$(fg)' = f'g + fg'$
\item \textbf{链式法则}：$\frac{d}{dx}f(g(x)) = f'(g(x)) \cdot g'(x)$ —— \textbf{ML 反向传播的核心}
\item \textbf{除法}：$\left(\frac{f}{g}\right)' = \frac{f'g - fg'}{g^2}$
\end{enumerate}

\textbf{基本函数的导数}（这些是肌肉记忆——必须熟）：
\begin{align}
\frac{d}{dx}x^n &= n x^{n-1} \\
\frac{d}{dx}e^x &= e^x \\
\frac{d}{dx}\ln x &= \frac{1}{x} \\
\frac{d}{dx}\sin x &= \cos x \\
\frac{d}{dx}\cos x &= -\sin x
\end{align}

\textbf{隐函数求导}：方程 $F(x,y)=0$——直接对 $x$ 两边求导——把 $y$ 当成 $x$ 的函数——\textbf{$y' = -F_x/F_y$}。

\subsection{积分技巧（节制版）}

\textbf{同济花了 50 页讲积分技巧}——\textbf{我只讲三个}——\textbf{其余都让 \texttt{sympy.integrate} 去算}。

\textbf{换元法}：$\int f(g(x)) g'(x) \, dx = \int f(u) \, du$（其中 $u = g(x)$）

\textbf{分部积分}：$\int u \, dv = uv - \int v \, du$ —— \textbf{记口诀"反对幂指三"}（反三角、对数、幂函数、指数、三角函数——按这个顺序选 $u$）。

\textbf{有理函数}：用部分分式分解 —— \textbf{现代版}：直接 \texttt{sympy.apart()}。

\begin{remark}[工业现实]
\textbf{99\% 的工程积分都是数值积分}——\textbf{\texttt{scipy.integrate.quad}——一行代码——比手算靠谱十倍}。\\
\textbf{手算积分的能力——只在两个地方有用}：
\begin{itemize}
\item \textbf{考试}（包括考研、博士入学）
\item \textbf{你要发现一个新公式}——这时数值算不出"封闭形式"——必须手算
\end{itemize}
\textbf{大部分工程师——一辈子都不会再用第二种}。
\end{remark}

\subsection{微积分基本定理：完整版}

\begin{theorem}[微积分基本定理（两部分）]
设 $f$ 在 $[a,b]$ 上连续。
\begin{enumerate}
\item（第一部分）令 $F(x) = \int_a^x f(t) \, dt$，则 $F'(x) = f(x)$。
\item（第二部分）若 $G$ 是 $f$ 的任一原函数，则 $\int_a^b f(x) \, dx = G(b) - G(a)$。
\end{enumerate}
\end{theorem}

\textbf{证明的核心思想（不严格）}：

考虑 $F(x+h) - F(x) = \int_x^{x+h} f(t) \, dt$。当 $h$ 小时，$f$ 在 $[x, x+h]$ 上几乎是常数 $f(x)$——所以这个积分 $\approx f(x) \cdot h$。两边除以 $h$ 并取极限——得到 $F'(x) = f(x)$。

\textbf{这个证明的画面——记下来——比公式重要}。

\section{Aha 例子：短信流量预测 / Aha: SMS Traffic Forecasting}

\textbf{场景}：你在中国移动总部做数据工程师。每天早上 9 点——\textbf{运营部门给你一组数据}：\textbf{昨天每分钟的短信发送量}。

数据是这样的（部分）：

\begin{center}
\begin{tabular}{cc}
\toprule
时间（分钟） & 短信量 \\
\midrule
08:59 & 12,034 \\
09:00 & 12,847 \\
09:01 & 14,201 \\
09:02 & 16,932 \\
09:03 & 19,500 \\
$\vdots$ & $\vdots$ \\
\bottomrule
\end{tabular}
\end{center}

\textbf{问题}：

\begin{enumerate}
\item \textbf{今天的高峰会出现在什么时候？}（导数的应用——找极值）
\item \textbf{今天一共会发多少条短信？}（积分的应用——累积量）
\item \textbf{流量上升最快的时刻是什么时候？}（二阶导数——拐点）
\end{enumerate}

\textbf{为什么这三个问题用微积分？}——\textbf{因为它们都是"变化率"或"累积量"的问题}——\textbf{这就是微积分的母语}。

\begin{tcolorbox}[ahabox={Aha: 微积分就是数据工程师的日常}]
\textbf{你以为微积分是"算 $\int \sin^3 x \cos^2 x \, dx$"}——\textbf{那是同济在折磨你}。

\medskip

\textbf{真实的工业微积分是这样的}：

\begin{enumerate}
\item 给你一个\textbf{离散的时间序列}（每分钟的数据）
\item 你需要找\textbf{极值}（高峰期 / 低谷期）——\textbf{用一阶导数}
\item 你需要找\textbf{拐点}（增速最快的时刻）——\textbf{用二阶导数}
\item 你需要算\textbf{累积量}（今日总流量）——\textbf{用数值积分}
\end{enumerate}

\medskip

\textbf{核心 API 就两个}——\textbf{\texttt{numpy.gradient}} 和 \textbf{\texttt{numpy.trapz}}——\textbf{你大学学的 300 个求导公式}——\textbf{在这两个 API 面前——一个都用不上}。

\medskip

但这不是说大学学的没用——\textbf{你必须知道"导数 = 变化率"——你才能反应过来"找高峰用一阶导数"}——\textbf{API 不会自己告诉你这件事}。

\medskip

\textbf{这就是"打通"的意义——把课本里的概念，和工业的 API，连成同一件事}。
\end{tcolorbox}

\textbf{解决问题的 Python 代码}：

\begin{lstlisting}[language=Python]
import numpy as np

# 1. 读数据（假设已加载为 t, sms_count）
t = np.arange(0, 1440)              # 一天的每分钟 (0-1439)
sms = load_sms_data()               # 长度 1440 的数组

# 2. 一阶导数：变化率（每分钟流量的变化速度）
rate = np.gradient(sms, t)

# 3. 找高峰：rate == 0 且 rate 由正变负 → 极大值
peaks = np.where(np.diff(np.sign(rate)) < 0)[0]
print("高峰时刻 (分钟):", peaks)
print("高峰流量:", sms[peaks])

# 4. 二阶导数：增速变化（找上升最快的时刻）
acceleration = np.gradient(rate, t)
fastest_rise = np.argmax(acceleration)
print("上升最快的时刻:", fastest_rise, "分钟")

# 5. 数值积分：今日总流量
total = np.trapz(sms, t)
print("今日总短信量:", int(total))
\end{lstlisting}

\textbf{这段代码——\textbf{20 行}——\textbf{解决了一个真实的运营问题}}。

它背后用到的微积分概念：

\begin{itemize}
\item \texttt{np.gradient}——\textbf{中心差分逼近导数}——\textbf{原理是导数定义的离散化}
\item \texttt{np.trapz}——\textbf{梯形法逼近积分}——\textbf{原理是 Riemann 和 + 取极限}
\item \textbf{找极值}——\textbf{用一阶导数变号}——\textbf{这是大一就学的内容}
\end{itemize}

\textbf{你看}——\textbf{大学学的没浪费}——\textbf{它在 \texttt{numpy} 的源码里}——\textbf{在你做决策的脑子里}——\textbf{只是不在你算 $\int \sin^3 x \, dx$ 的草稿纸上了}。

\section{现代视角：自动微分与 ML / Modern: Autodiff and ML}

\textbf{微积分的现代形态——叫"自动微分"（automatic differentiation, autodiff）}。

\textbf{这是 21 世纪的微积分革命}——\textbf{比 Newton 那次更"工业"}。

\subsection{什么是自动微分}

考虑函数
\[
f(x) = \sin(x^2) \cdot e^x
\]

用手算它的导数——\textbf{要用乘法法则 + 链式法则}——3 行公式：
\[
f'(x) = 2x \cos(x^2) e^x + \sin(x^2) e^x
\]

\textbf{自动微分的思路}——\textbf{不用人推公式}——\textbf{让计算机用链式法则一步步往下传}：

\begin{lstlisting}[language=Python]
import torch
x = torch.tensor(1.5, requires_grad=True)
f = torch.sin(x**2) * torch.exp(x)
f.backward()
print(x.grad)  # 直接给出 f'(1.5)
\end{lstlisting}

\textbf{两行——完事}。

\subsection{为什么这件事改变了一切}

\textbf{深度学习里的损失函数}——是\textbf{几百万个参数的函数}——\textbf{手算导数完全不可能}。

\textbf{有了 autodiff}——\textbf{反向传播}——\textbf{随机梯度下降}——\textbf{才有了 ChatGPT 这样的模型}。

\textbf{这就是为什么我说——21 世纪的微积分发生了第二次革命}：

\begin{itemize}
\item \textbf{Newton 1665}——\textbf{发明了微积分}
\item \textbf{Cauchy 1820}——\textbf{严格化了微积分}
\item \textbf{2010 年后}——\textbf{自动化了微积分}——\textbf{让 1 万亿参数的模型可以训练}
\end{itemize}

\textbf{你今天打开 PyTorch 调一个 \texttt{loss.backward()}}——\textbf{背后是 Newton + Leibniz + Cauchy + 现代 PL 设计的合谋}。

\section{代码与思考 / Code and Exercises}

\subsection{配套模块 \texttt{calculus\_intro.py}}

GitHub 上的 \texttt{modules/calculus\_intro.py} 包含\textbf{6 个 demo}：

\begin{enumerate}
\item \texttt{demo\_limit()}——可视化 $\lim_{x\to 0} \sin(x)/x = 1$ 的"逼近"过程
\item \texttt{demo\_derivative()}——比较解析导数 vs 数值导数（\texttt{np.gradient}）
\item \texttt{demo\_integral()}——比较解析积分 vs 数值积分（trapz vs Simpson vs quad）
\item \texttt{demo\_FTC()}——验证微积分基本定理：$\frac{d}{dx}\int_a^x f(t)\,dt = f(x)$
\item \texttt{demo\_sms\_traffic()}——本章 Aha 例子：完整的短信流量分析
\item \texttt{demo\_autodiff()}——用 PyTorch 做自动微分——和 \texttt{sympy} 的解析结果对比
\end{enumerate}

\textbf{使用}：

\begin{lstlisting}[language=bash]
git clone https://github.com/inaturelz-coder/math-rendu
cd math-rendu
pip install -r requirements.txt
python modules/calculus_intro.py
\end{lstlisting}

\subsection{思考题 / Exercises}

\begin{enumerate}
\item \textbf{[直觉]} 用一句话向你高中的同桌解释"导数"——不许用公式。
\item \textbf{[计算]} 求 $f(x) = x^x$ 的导数。提示：先两边取对数。
\item \textbf{[应用]} 一辆车在 $[0, 10]$ 秒内位置为 $s(t) = 3t^2 + 5t$（米）。它在第 5 秒的速度和加速度各是多少？
\item \textbf{[代码]} 写一个 Python 函数 \texttt{numerical\_derivative(f, x, h=1e-5)}——用中心差分实现。和 \texttt{numpy.gradient} 比较精度。
\item \textbf{[历史]} 查一查：为什么 Leibniz 的 $dx, dy, \int$ 符号最终胜过了 Newton 的"流数"记号？
\item \textbf{[现代]} 用 PyTorch 计算 $f(x) = x^4 - 3x^3 + 2x$ 在 $x=1.5$ 的一阶和二阶导数。和手算结果对比。
\end{enumerate}

\section{本章小结 / Chapter Summary}

\begin{tcolorbox}[essbox={本章核心}]
\begin{itemize}
\item \textbf{极限}——"我想要的精度——我能做到"——$\varepsilon\text{-}\delta$ 是"挑战与回应"的游戏
\item \textbf{导数}——变化率——切线斜率——速度——梯度——\textbf{它们是同一个东西}
\item \textbf{积分}——累积量——面积——\textbf{它和导数是逆运算}（微积分基本定理）
\item \textbf{现代}——\textbf{99\% 的工程微积分是数值微积分}——\textbf{\texttt{np.gradient + np.trapz}}
\item \textbf{未来}——自动微分（autodiff）+ 反向传播——\textbf{ML 的核心数学}
\end{itemize}
\end{tcolorbox}

\textbf{下一章——多元微积分}——\textbf{我们从"一维"走到"多维"}——\textbf{并第一次遇到"梯度"——这个 ML 训练的核心概念}。

\clearpage

% =====================================================================
%  第 2 章 · 多元微积分 / Calculus II
%  偏导、全微分、重积分、Lagrange 乘数法
%  小故事：Lagrange 与拿破仑时代的优化问题
%  Aha：外卖配送优化（多变量梯度下降）
% =====================================================================

\chapter{多元微积分 / Calculus II}
\label{ch:calc2}

\begin{quote}\itshape
\textbf{Joseph-Louis Lagrange} 是\textbf{法国大革命之后、拿破仑时代}的数学家。
他活到 1813 年——\textbf{见证了大革命、督政府、拿破仑称帝、滑铁卢的前夜}。

在那个动荡的年代——他写了一本《\textit{解析力学}》——\textbf{把整个牛顿力学，用变分原理 + 多元微积分重新写了一遍}。

我们这一章——讲的就是他用过的工具：\textbf{偏导、梯度、Lagrange 乘数法}。
\end{quote}

\section{写在前面 / Preface}

\textbf{多元微积分是工程师的"日常数学"}——\textbf{比一元微积分日常 100 倍}。

为什么？因为\textbf{真实世界没有"一元"问题}：

\begin{itemize}
\item \textbf{温度} 是 $(x, y, z, t)$ 的函数——\textbf{四元}
\item \textbf{应力} 是 $(x, y, z)$ 的张量——\textbf{每个分量都是三元}
\item \textbf{神经网络} 的损失 $L(\theta_1, \theta_2, \ldots, \theta_N)$——\textbf{$N=10^9$ 元}
\end{itemize}

\textbf{一元微积分是给数学家的}——\textbf{多元微积分才是给工程师的}。

\textbf{这一章的核心}——\textbf{把"导数"推广到"梯度"——把"求面积"推广到"求体积 / 求质量 / 求电荷"}。

\textbf{并且——我们第一次遇到\textbf{梯度} $\nabla f$}——\textbf{这是接下来所有章节的"母语"}——\textbf{Ch 3 矢量分析、Ch 13 优化、整个 ML——全都建立在这上面}。

\section{一句话本质 / The Essence in One Sentence}

\begin{tcolorbox}[essbox={一句话本质}]
\textbf{中文}：\textbf{梯度 $\nabla f$——是函数上升最快的方向}；\textbf{多重积分——是把"累积"从线段推广到区域、体积、高维空间}。\\[6pt]
\textbf{English}: The gradient $\nabla f$ points in the direction of steepest ascent; multiple integrals generalize accumulation from intervals to regions, volumes, and higher-dimensional domains.
\end{tcolorbox}

\section{教材里你看到的 / What the Textbook Tells You}

\subsection{同济《高等数学（下册）》}

第八章到第十一章——共\textbf{120 页}——讲多元微积分。

\textbf{优点}：覆盖全——\textbf{偏导、全微分、重积分、曲线积分、曲面积分}全都有。

\textbf{缺点}：

\begin{enumerate}
\item \textbf{符号繁琐}——$\frac{\partial}{\partial x}, \frac{\partial}{\partial y}, \frac{\partial^2}{\partial x \partial y}$——\textbf{学生看着头大}
\item \textbf{几乎不提\textbf{梯度}}——只在习题里偶尔出现"方向导数"——\textbf{这是最严重的遗漏}
\item \textbf{Lagrange 乘数法\textbf{用 1 页讲完}}——\textbf{然后给 5 道考研真题}——\textbf{几何意义完全没讲}
\end{enumerate}

\subsection{Stewart 的版本}

Stewart 把这部分写得\textbf{活灵活现}——\textbf{有大量"温度场""压力场"的可视化}——\textbf{这是同济不及的}。

\subsection{我的策略}

\begin{enumerate}
\item \textbf{把梯度 $\nabla f$ 提到第一位}——\textbf{它是这一章的灵魂}
\item \textbf{Lagrange 乘数法用几何直觉讲}——\textbf{不只是公式}
\item \textbf{Aha 例子用"外卖配送优化"}——\textbf{多变量梯度下降的真实场景}
\end{enumerate}

\begin{tcolorbox}[storybox={\Large 小故事：Lagrange 与拿破仑时代}]
\textbf{Joseph-Louis Lagrange}（1736-1813）——\textbf{意大利人——成名于柏林——晚年定居巴黎}。

\medskip

他在 1788 年出版《\textbf{解析力学}》——\textbf{这是力学史上的一座丰碑}。\textbf{Hamilton} 后来说："这本书是数学的\textbf{史诗}。"

\medskip

\textbf{Lagrange 解决了什么问题}？

\medskip

\textbf{Newton 力学}的写法是——\textbf{对每一个粒子写 $F=ma$}——\textbf{对 $N$ 个粒子要写 $3N$ 个方程}——\textbf{遇到约束（绳子、铰链）要单独处理}。

\medskip

\textbf{Lagrange 的洞察}——\textbf{用"广义坐标" $q_1, q_2, \ldots$ 代替笛卡尔坐标}——\textbf{用一个标量函数 $L = T - V$（拉格朗日量）代替力的矢量分析}——\textbf{从此约束变成"消元"——力学变成"求极值"}。

\medskip

\textbf{这背后用到的核心工具——就是"多元函数求极值——可能带约束"}——\textbf{也就是 Lagrange 乘数法}。

\medskip

\textbf{Lagrange 一生没结婚}（晚年才在拿破仑撮合下结了第二次）——\textbf{他在大革命中差点被杀}（数学家拉瓦锡就被送上断头台）——\textbf{他靠拿破仑的赏识活到 77 岁}——\textbf{死后葬入先贤祠}。

\medskip

\textbf{拿破仑说过}：\textit{"Lagrange 是数学界的金字塔。"}

\medskip

\textbf{你今天用的 Lagrangian、Lagrange 乘数法、Euler-Lagrange 方程}——\textbf{都源于这位 18 世纪末的意大利人——在巴黎的小公寓里——一个人想了几十年}。
\end{tcolorbox}

\section{直觉 / Intuition}

\subsection{多元函数：等高线与曲面}

考虑二元函数 $z = f(x, y)$。\textbf{它的"图"是一个曲面}——而不是一条曲线。

\textbf{怎么画一个三维曲面在二维纸上}？——\textbf{用等高线}：

\begin{itemize}
\item 把曲面切成"水平片" $z = c$（$c$ 是常数）
\item 每一片在 $xy$ 平面上的投影——\textbf{就是一条等高线}
\item 等高线\textbf{越密}——\textbf{曲面越陡}
\end{itemize}

\textbf{地形图、温度图、电势图}——\textbf{都是等高线}——\textbf{这是工程师最常见的"多元函数可视化"}。

\subsection{偏导数：把其他变量当常数}

\begin{definition}[偏导数]
\[
\frac{\partial f}{\partial x}(x_0, y_0) = \lim_{h \to 0} \frac{f(x_0+h, y_0) - f(x_0, y_0)}{h}
\]
\end{definition}

\textbf{直觉}：\textbf{把 $y$ 固定}——\textbf{$f$ 就变成 $x$ 的一元函数}——\textbf{对它求导}。

\textbf{几何意义}：\textbf{$\partial f / \partial x$ 是"沿 $x$ 方向切曲面"得到的曲线在那一点的斜率}。

\subsection{梯度 \texorpdfstring{$\nabla f$}{∇f}：本章的灵魂}

\begin{definition}[梯度]
$n$ 元函数 $f(x_1, \ldots, x_n)$ 的梯度是矢量：
\[
\nabla f = \left( \frac{\partial f}{\partial x_1}, \frac{\partial f}{\partial x_2}, \ldots, \frac{\partial f}{\partial x_n} \right)
\]
\end{definition}

\textbf{三大性质}（每一条都要记住）：

\begin{enumerate}
\item \textbf{方向}：$\nabla f$ 指向 $f$ \textbf{上升最快}的方向
\item \textbf{大小}：$|\nabla f|$ 是 $f$ 在那个方向的\textbf{变化率}
\item \textbf{正交}：$\nabla f$ 在每一点上\textbf{垂直于过那点的等高线}
\end{enumerate}

\textbf{第三条特别重要}——\textbf{这是 Lagrange 乘数法的几何基础}。

\subsection{方向导数}

沿单位矢量 $\bm{u}$ 方向的导数：
\[
D_{\bm{u}} f = \nabla f \cdot \bm{u}
\]

\textbf{这告诉我们一件事}：\textbf{所有方向的变化率——都是梯度在那个方向的投影}。\textbf{梯度是"方向变化率的母体"}。

\subsection{全微分}

\begin{equation}
df = \frac{\partial f}{\partial x} dx + \frac{\partial f}{\partial y} dy + \cdots = \nabla f \cdot d\bm{r}
\end{equation}

\textbf{物理意义}（热力学里你天天见）：

\[
dU = T \, dS - P \, dV \quad \text{（热力学第一定律）}
\]

\textbf{这是 $U(S, V)$ 的全微分}——\textbf{$T = \partial U / \partial S$（温度）}——\textbf{$P = -\partial U / \partial V$（压强）}——\textbf{热力学的整个机器就在这里}。

\subsection{Lagrange 乘数法：约束下的极值}

\textbf{问题}：求 $f(x,y)$ 的极值——\textbf{在约束 $g(x,y) = 0$ 下}。

\textbf{几何直觉}：

\begin{itemize}
\item $f$ 的等高线——\textbf{一族曲线}
\item 约束曲线 $g=0$——\textbf{一条曲线}
\item \textbf{在 $g=0$ 上的极值点}——\textbf{必然是 $f$ 的等高线与 $g=0$ 相切的点}
\item \textbf{相切——意味着两条曲线的法向量平行}——\textbf{$\nabla f = \lambda \nabla g$}
\end{itemize}

\textbf{这就是 Lagrange 乘数法的本质}——\textbf{把"约束极值"变成"梯度平行"}。

\textbf{解法}：解方程组
\[
\nabla f = \lambda \nabla g, \quad g = 0
\]

\section{数学 / The Math}

\subsection{链式法则（多元版）}

设 $z = f(x, y)$，$x = x(t), y = y(t)$。那么
\[
\frac{dz}{dt} = \frac{\partial f}{\partial x} \frac{dx}{dt} + \frac{\partial f}{\partial y} \frac{dy}{dt}
\]

\textbf{矩阵版}：若 $\bm{z} = \bm{f}(\bm{x})$ 是 $\mathbb{R}^n \to \mathbb{R}^m$ 的映射，则\textbf{Jacobian 矩阵}
\[
J = \frac{\partial \bm{z}}{\partial \bm{x}} = \begin{pmatrix}
\frac{\partial z_1}{\partial x_1} & \cdots & \frac{\partial z_1}{\partial x_n} \\
\vdots & & \vdots \\
\frac{\partial z_m}{\partial x_1} & \cdots & \frac{\partial z_m}{\partial x_n}
\end{pmatrix}
\]

\textbf{ML 的反向传播——本质上就是 Jacobian 矩阵的乘法链}。

\subsection{重积分：从面积到体积}

\textbf{二重积分}：
\[
\iint_D f(x,y) \, dA = \lim_{\|P\|\to 0} \sum f(x_i, y_i) \Delta A_i
\]

\textbf{物理意义}：\textbf{$f$ 是面密度——$\iint f \, dA$ 是总质量}。

\textbf{计算手段}——\textbf{化为累次积分}（Fubini 定理）：
\[
\iint_D f \, dA = \int_a^b \left( \int_{y_1(x)}^{y_2(x)} f(x, y) \, dy \right) dx
\]

\textbf{坐标变换}（极坐标、球坐标）——\textbf{必须乘上 Jacobian}：
\[
dx \, dy = r \, dr \, d\theta \quad \text{(极坐标)}
\]

\subsection{Hessian 矩阵：二阶导数的推广}

对 $f(x_1, \ldots, x_n)$，定义\textbf{Hessian 矩阵}
\[
H_{ij} = \frac{\partial^2 f}{\partial x_i \partial x_j}
\]

\textbf{它告诉我们函数曲面在某点的"曲率"}：

\begin{itemize}
\item $H$ 正定 $\Rightarrow$ \textbf{局部极小值}（碗形）
\item $H$ 负定 $\Rightarrow$ \textbf{局部极大值}（馒头形）
\item $H$ 不定 $\Rightarrow$ \textbf{鞍点}（马鞍形）
\end{itemize}

\textbf{Newton 优化法的核心}：用 $H^{-1} \nabla f$ 作为"二阶下降方向"——\textbf{见第 13 章}。

\section{Aha 例子：外卖配送优化 / Aha: Food Delivery Optimization}

\textbf{场景}：你在美团做算法工程师。系统有 $N$ 个骑手——每个骑手位置 $(x_i, y_i)$——\textbf{你要决定每个骑手该往哪个新订单冲}。

\textbf{目标}：\textbf{最小化总配送时间}。

\textbf{建模}：

\begin{itemize}
\item 订单位置 $(a, b)$
\item 第 $i$ 个骑手到订单的时间 $t_i = d_i / v_i$（距离/速度）
\item 总成本（简化）：$C = \sum_i t_i \cdot p_i$，其中 $p_i$ 是\textbf{你分配给骑手 $i$ 的概率}
\item 约束：$\sum p_i = 1, \quad p_i \geq 0$
\end{itemize}

\textbf{这是一个带约束的多元优化问题}——\textbf{Lagrange 乘数法上场}。

\begin{tcolorbox}[ahabox={Aha: 多元微积分就是 ML 的全部}]
\textbf{ML 的训练过程——一句话总结}：

\medskip

\textbf{找一个 $\theta$，使损失函数 $L(\theta_1, \theta_2, \ldots, \theta_N)$ 最小。}

\medskip

\textbf{怎么找}？——\textbf{梯度下降}：
\[
\theta_{t+1} = \theta_t - \eta \nabla L(\theta_t)
\]

\medskip

\textbf{这里的 $\nabla L$——就是本章讲的"梯度"}——\textbf{$\eta$ 是学习率}——\textbf{你大学学的偏导数、链式法则、Jacobian——全在这里}。

\medskip

\textbf{所以你大学的多元微积分——不是"为了考试发明的"}——\textbf{是 ML 的全部数学根基}。

\medskip

\textbf{你不学好这一章——你就看不懂 PyTorch 的源码}——\textbf{你不学好这一章——你就理解不了为什么 Adam 优化器要"自适应学习率"}——\textbf{你不学好这一章——21 世纪的工程对你就是个黑盒}。

\medskip

\textbf{这就是为什么我说——多元微积分比一元微积分日常 100 倍}。
\end{tcolorbox}

\textbf{解决方法（简化版）}：

设我们想最小化 $C(p_1, \ldots, p_N) = \sum_i t_i p_i$，约束 $\sum p_i = 1$。

\textbf{Lagrangian}：
\[
\mathcal{L} = \sum_i t_i p_i + \lambda \left( \sum_i p_i - 1 \right)
\]

\textbf{对每个 $p_i$ 求偏导}：$\partial \mathcal{L} / \partial p_i = t_i + \lambda = 0$——\textbf{所有 $t_i$ 相等}。

\textbf{现实里}——$t_i$ 通常不等——\textbf{所以最优解是 $p_i = \delta(i = i^*)$}——\textbf{把订单全分给时间最短的骑手}。

\textbf{这听起来废话}——但\textbf{真实工业场景里}——\textbf{$t_i$ 是不确定的（堵车、骑手繁忙度）}——\textbf{所以要加入概率/期望}——\textbf{这就是为什么美团的算法不是"最短距离原则"——而是一个复杂的多元优化}。

\textbf{Python 实现}（用 \texttt{scipy.optimize.minimize}）：

\begin{lstlisting}[language=Python]
from scipy.optimize import minimize
import numpy as np

N = 10
t = np.random.uniform(5, 30, N)  # 10 个骑手的预估时间

# 目标函数
def cost(p):
    return np.sum(t * p)

# 约束：和为 1
cons = {'type': 'eq', 'fun': lambda p: np.sum(p) - 1}
# 非负约束
bnds = [(0, 1) for _ in range(N)]

p0 = np.ones(N) / N  # 初始猜测
res = minimize(cost, p0, constraints=cons, bounds=bnds)
print("最优分配:", res.x)
print("最优成本:", res.fun)
\end{lstlisting}

\textbf{20 行——解决了一个真实的运筹问题}。

\section{现代视角：神经网络与梯度 / Modern: NN and Gradients}

\subsection{神经网络的损失函数：百万维的山地}

一个中等规模的神经网络——\textbf{参数数 $N \approx 10^7 \sim 10^9$}。

它的损失函数 $L(\theta_1, \ldots, \theta_N)$——\textbf{是一个 $10^7$ 维空间中的函数}。

\textbf{我们无法可视化它}——\textbf{但我们知道它的梯度 $\nabla L$ 仍然存在}——\textbf{它仍然指向上升最快的方向}。

\textbf{梯度下降}——\textbf{每次沿 $-\nabla L$ 走一小步}——\textbf{走到一个"局部低点"}——\textbf{这就是训练}。

\subsection{Adam 优化器：自适应学习率}

\textbf{Adam}（2014, Kingma \& Ba）——\textbf{现代 ML 训练最常用的优化器}——\textbf{它在普通梯度下降上加了两个改进}：

\begin{itemize}
\item \textbf{Momentum}——用过去几步梯度的\textbf{加权平均}——避免"震荡"
\item \textbf{自适应学习率}——根据每个参数的\textbf{梯度方差}调整步长
\end{itemize}

\textbf{Adam 的"自适应"——就是多元微积分的"分量"概念}——\textbf{每个 $\theta_i$ 都有自己的学习率 $\eta_i$}。

\textbf{Adam 训练了 ChatGPT}——\textbf{训练了 AlphaGo}——\textbf{训练了所有你听说过的 AI 系统}——\textbf{它本质上——是多元微积分的现代化}。

\section{代码与思考 / Code and Exercises}

\subsection{配套模块 \texttt{multi\_calc.py}}

\begin{enumerate}
\item \texttt{demo\_partial\_derivatives()}——可视化 $\partial f / \partial x$ 和 $\partial f / \partial y$
\item \texttt{demo\_gradient\_field()}——画 $f(x,y) = x^2 + y^2$ 的等高线与梯度场
\item \texttt{demo\_total\_differential()}——验证热力学全微分 $dU = TdS - PdV$
\item \texttt{demo\_lagrange()}——用 Lagrange 乘数法求约束极值——和 scipy 对比
\item \texttt{demo\_delivery\_optimization()}——本章 Aha 例子
\item \texttt{demo\_double\_integral()}——数值二重积分（\texttt{scipy.integrate.dblquad}）
\end{enumerate}

\subsection{思考题 / Exercises}

\begin{enumerate}
\item \textbf{[直觉]} 用一句话解释：\textbf{为什么 $\nabla f$ 垂直于等高线？}（不用公式）
\item \textbf{[计算]} 求 $f(x,y,z) = x^2 y - yz + e^{xz}$ 在点 $(1, 2, 0)$ 的梯度。
\item \textbf{[应用]} 用 Lagrange 乘数法求：在椭圆 $x^2/4 + y^2/9 = 1$ 上离原点最远的点。
\item \textbf{[代码]} 写一个 \texttt{gradient\_descent} 函数——对 $f(x,y) = (x-3)^2 + (y+1)^2$ 找极小值。
\item \textbf{[历史]} 查一查：为什么 Lagrange 不用"乘数法"这个名字——这个名字是后人给的？
\item \textbf{[ML]} 用 PyTorch 计算 $f(x, y) = \sin(xy) + x^2$ 在 $(1, \pi/2)$ 的梯度。和手算对比。
\end{enumerate}

\section{本章小结 / Chapter Summary}

\begin{tcolorbox}[essbox={本章核心}]
\begin{itemize}
\item \textbf{偏导} = 固定其他变量、对单变量求导
\item \textbf{梯度 $\nabla f$} = 上升最快方向 = 等高线的法向 = ML 训练的方向
\item \textbf{全微分} = $df = \nabla f \cdot d\bm{r}$ = 热力学第一定律的写法
\item \textbf{重积分} = 把累积从一维推广到多维 = 总质量 / 总电荷 / 总能量
\item \textbf{Lagrange 乘数法} = 约束极值 = $\nabla f = \lambda \nabla g$ = 等高线与约束曲线相切
\item \textbf{Hessian} = 多元二阶导数 = 曲面的"曲率" = 决定极小 / 极大 / 鞍点
\end{itemize}
\end{tcolorbox}

\textbf{下一章——矢量分析}——\textbf{我们让梯度遇到"散度"和"旋度"——这是电磁学、流体、弹性的共用语言}。

\clearpage

% =====================================================================
%  第 3 章 · 矢量分析 / Vector Analysis
%  梯度、散度、旋度、Gauss、Stokes
%  小故事：Maxwell 与 Gibbs 的矢量发明
%  Aha：手机信号场强
% =====================================================================

\chapter{矢量分析 / Vector Analysis}
\label{ch:vec}

\begin{quote}\itshape
\textbf{Maxwell 1865 年}发表《电磁场的动力学理论》——里面是\textbf{20 个偏微分方程}——\textbf{用四元数（quaternions）写的}——\textbf{没人看得懂}。

\medskip

\textbf{Heaviside 和 Gibbs}\textbf{在 1880 年代}独立改写了它——\textbf{用我们今天熟悉的"矢量"语言}——\textbf{20 个方程缩成 4 个}。

\medskip

我们这一章——讲的就是 Heaviside 和 Gibbs 发明的"矢量分析"——它是电磁、流体、弹性、传热的\textbf{共用语言}。
\end{quote}

\section{写在前面 / Preface}

\textbf{我读大二上《电动力学》}——\textbf{第一节课老师写了一个公式}：
\[
\nabla \cdot \bm{E} = \rho / \varepsilon_0
\]
\textbf{我懵了 10 分钟}。

那个倒三角是什么？点乘是什么？$\rho / \varepsilon_0$ 我懂——电荷密度除以真空介电常数——但等号左边是个什么"操作"——我说不上来。

\textbf{我翻同济《高等数学下册》}——\textbf{第十一章}——\textbf{讲过"散度"——总共一页半}——\textbf{讲完就给习题}。

\textbf{我翻 Griffiths《Introduction to Electrodynamics》}——\textbf{它专门花了一章讲矢量分析}——\textbf{我跟读了三天}——\textbf{才反应过来}：

\textbf{$\nabla$ 不是"数学家发明的折磨工具"}——\textbf{它是"工程师发明的简写"}。

它简化了什么？——\textbf{它把"对每个分量分别写偏导"的繁琐版}——\textbf{压缩成了一个符号}。

\textbf{这一章的目的}——\textbf{把 $\nabla \cdot, \nabla \times, \nabla^2$ 这三个符号——讲到你看 Maxwell 方程组不再懵}。

\section{一句话本质 / The Essence in One Sentence}

\begin{tcolorbox}[essbox={一句话本质}]
\textbf{中文}：\textbf{$\nabla$ 是一个"矢量算符"——它的点乘是散度（衡量源）、叉乘是旋度（衡量旋转）、自乘是 Laplacian（衡量曲率）}。\\[6pt]
\textbf{English}: $\nabla$ is a vector operator——its dot product gives divergence (measuring sources), its cross product gives curl (measuring rotation), and its square gives the Laplacian (measuring curvature).
\end{tcolorbox}

\section{教材里你看到的 / What the Textbook Tells You}

同济下册有一章\textbf{"曲线积分与曲面积分"}——\textbf{讲了 Gauss 定理、Stokes 定理}——\textbf{但只当成"计算工具"}——\textbf{没讲它们的物理意义}。

\textbf{物理系的教材}（比如梁昆淼《数学物理方法》）——\textbf{讲得更深}——\textbf{但工科生很少看}。

\textbf{我的策略}——\textbf{把"梯度、散度、旋度"当成三个\textbf{独立的几何对象}}——\textbf{每一个都配一个物理画面}：

\begin{itemize}
\item \textbf{梯度} = 山的最陡上升方向
\item \textbf{散度} = "流出量"（水从源头喷出）
\item \textbf{旋度} = "转动量"（水的小漩涡）
\end{itemize}

\begin{tcolorbox}[storybox={\Large 小故事：Maxwell, Gibbs, Heaviside 的矢量发明}]
\textbf{1865 年}——\textbf{James Clerk Maxwell}（苏格兰人）——\textbf{发表《电磁场的动力学理论》}——\textbf{20 个偏微分方程}——\textbf{用 Hamilton 的"四元数"写的}。

\medskip

\textbf{四元数}是什么？——\textbf{它是 1843 年 Hamilton 发明的"3D 复数"}——\textbf{$i^2 = j^2 = k^2 = ijk = -1$}——\textbf{乘法不交换}——\textbf{学起来比矢量分析难十倍}。

\medskip

\textbf{Maxwell 的方程组用四元数写出来——非常对称、非常美}——\textbf{但工程师看不懂}。

\medskip

\textbf{1880 年代}——\textbf{两个"实用派"独立行动}：

\begin{itemize}
\item \textbf{Josiah Willard Gibbs}（耶鲁大学的化学家——后来成为物理学家）
\item \textbf{Oliver Heaviside}（英国电信工程师——自学成才——脾气怪——晚年贫困）
\end{itemize}

\medskip

\textbf{他们独立做了同一件事}——\textbf{把 Maxwell 的四元数版}——\textbf{改写成"矢量分析"版}——\textbf{20 个方程 → 4 个方程}。

\medskip

\textbf{Hamilton 的拥护者大怒}——\textbf{发起"四元数 vs 矢量分析"的口水战}——\textbf{持续了 20 年}。

\medskip

\textbf{最后矢量分析赢了}——\textbf{因为它更直观、更易教、更适合工程}。

\medskip

\textbf{我们今天看到的 $\nabla \cdot E = \rho / \varepsilon_0$、$\nabla \times B = \mu_0 J + \mu_0 \varepsilon_0 \partial E / \partial t$——都是 Heaviside 和 Gibbs 的简化版}。

\medskip

\textbf{Maxwell 本人没活到看见这场胜利}（他 1879 年去世——48 岁）。
\end{tcolorbox}

\section{直觉 / Intuition}

\subsection{矢量场 vs 标量场}

\textbf{标量场}——\textbf{每点一个数}——温度 $T(\bm{r})$、电势 $\phi(\bm{r})$、海拔 $h(x,y)$。

\textbf{矢量场}——\textbf{每点一个矢量}——速度场 $\bm{v}(\bm{r})$、电场 $\bm{E}(\bm{r})$、磁场 $\bm{B}(\bm{r})$。

\textbf{我们要做的事}——\textbf{对"场"求导}——\textbf{得到三种不同的"变化率"}：梯度、散度、旋度。

\subsection{$\nabla$ 算符：矢量微分}

\textbf{$\nabla$（"nabla"）}定义为：
\[
\nabla = \left( \frac{\partial}{\partial x}, \frac{\partial}{\partial y}, \frac{\partial}{\partial z} \right)
\]

\textbf{它是个"矢量"}——\textbf{但分量是"偏导操作"而不是数}——所以叫\textbf{矢量算符}。

\textbf{它能做四件事}（前三件本章讲，最后一件是 Ch 5 PDE 的核心）：

\begin{enumerate}
\item \textbf{作用于标量场}：$\nabla f$ = \textbf{梯度}（结果是矢量）
\item \textbf{点乘矢量场}：$\nabla \cdot \bm{F}$ = \textbf{散度}（结果是标量）
\item \textbf{叉乘矢量场}：$\nabla \times \bm{F}$ = \textbf{旋度}（结果是矢量）
\item \textbf{自身点乘}：$\nabla^2 f = \nabla \cdot \nabla f$ = \textbf{Laplacian}（结果是标量）
\end{enumerate}

\subsection{散度：源还是汇}

\textbf{物理画面}——\textbf{想象一个水流场 $\bm{v}(\bm{r})$}——\textbf{每点的水流速度}。

\textbf{散度 $\nabla \cdot \bm{v}$ 告诉你}：\textbf{这一点周围——是水的源头（喷出来）还是汇（吸进去）}？

\begin{itemize}
\item $\nabla \cdot \bm{v} > 0$：\textbf{源}（这一点向外喷水）
\item $\nabla \cdot \bm{v} < 0$：\textbf{汇}（这一点吸水）
\item $\nabla \cdot \bm{v} = 0$：\textbf{无源无汇}（水进多少出多少）
\end{itemize}

\textbf{Maxwell 第一方程 $\nabla \cdot \bm{E} = \rho / \varepsilon_0$ 的物理意义}：\textbf{电场的源是电荷}——\textbf{电荷在哪里——电场就从哪里"喷出来"}。

\textbf{Maxwell 第二方程 $\nabla \cdot \bm{B} = 0$ 的物理意义}：\textbf{磁场无源}——\textbf{没有"单极磁荷"}——\textbf{磁感线必须闭合}。

\subsection{旋度：漩涡}

\textbf{物理画面}——\textbf{把一个小风车放进水流场}——\textbf{看它转不转}。

\textbf{旋度 $\nabla \times \bm{v}$ 告诉你}：\textbf{这一点的水流——有没有"旋转"}？

\begin{itemize}
\item $\nabla \times \bm{v} \neq \bm{0}$：\textbf{有旋}（小风车会转）——\textbf{方向沿转轴}
\item $\nabla \times \bm{v} = \bm{0}$：\textbf{无旋}（小风车不转）
\end{itemize}

\textbf{Maxwell 第三方程 $\nabla \times \bm{E} = -\partial \bm{B}/\partial t$ 的物理意义}（Faraday 电磁感应）：\textbf{磁场的变化——会产生有旋的电场}。

\textbf{Maxwell 第四方程 $\nabla \times \bm{B} = \mu_0 \bm{J} + \mu_0 \varepsilon_0 \partial \bm{E}/\partial t$}（Ampère-Maxwell）：\textbf{电流 + 电场变化——会产生有旋的磁场}。

\subsection{Laplacian：曲率}

\[
\nabla^2 f = \frac{\partial^2 f}{\partial x^2} + \frac{\partial^2 f}{\partial y^2} + \frac{\partial^2 f}{\partial z^2}
\]

\textbf{物理画面}——\textbf{这一点的 $f$，比周围邻居们的平均值多了多少}？

\begin{itemize}
\item $\nabla^2 f > 0$：\textbf{这点比邻居低}（凹陷处）
\item $\nabla^2 f < 0$：\textbf{这点比邻居高}（凸起处）
\item $\nabla^2 f = 0$：\textbf{这点等于邻居的平均}——\textbf{Laplace 方程}——\textbf{第 5 章主角}
\end{itemize}

\section{数学 / The Math}

\subsection{矢量恒等式：必背的五个}

\begin{align}
\nabla \times (\nabla f) &= \bm{0} \quad \text{(梯度场无旋)} \\
\nabla \cdot (\nabla \times \bm{F}) &= 0 \quad \text{(旋度场无源)} \\
\nabla \times (\nabla \times \bm{F}) &= \nabla(\nabla \cdot \bm{F}) - \nabla^2 \bm{F} \\
\nabla \cdot (f \bm{F}) &= f \nabla \cdot \bm{F} + \nabla f \cdot \bm{F} \\
\nabla \times (f \bm{F}) &= f \nabla \times \bm{F} + \nabla f \times \bm{F}
\end{align}

\textbf{前两个特别重要}——\textbf{它们告诉我们"梯度场无旋"、"旋度场无源"}——\textbf{这是 Helmholtz 分解定理的基础}。

\subsection{Gauss 定理（散度定理）}

\begin{theorem}[Gauss 散度定理]
设 $V$ 是 $\mathbb{R}^3$ 中体积——$\partial V$ 是其边界曲面——$\bm{F}$ 是 $V$ 上的可微矢量场。则
\[
\iiint_V \nabla \cdot \bm{F} \, dV = \oiint_{\partial V} \bm{F} \cdot d\bm{S}
\]
\end{theorem}

\textbf{物理意义}：\textbf{体积内的"总源"——等于边界的"总流量"}——\textbf{水进多少——必须流出多少}（无积累）。

\textbf{Maxwell 第一方程的积分版}：
\[
\oiint \bm{E} \cdot d\bm{S} = Q_{\text{内}} / \varepsilon_0
\]
\textbf{封闭曲面的电通量 = 内部总电荷 / $\varepsilon_0$}——\textbf{Gauss 定律}。

\subsection{Stokes 定理}

\begin{theorem}[Stokes 定理]
设 $S$ 是有向曲面——$\partial S$ 是其边界曲线（按右手定则）——$\bm{F}$ 是 $S$ 邻域上的可微矢量场。则
\[
\iint_S (\nabla \times \bm{F}) \cdot d\bm{S} = \oint_{\partial S} \bm{F} \cdot d\bm{r}
\]
\end{theorem}

\textbf{物理意义}：\textbf{曲面内的"总旋转"——等于边界的"环流量"}——\textbf{所有小漩涡的总和——等于宏观的循环}。

\textbf{Maxwell 第三方程的积分版}：
\[
\oint \bm{E} \cdot d\bm{r} = -\frac{d\Phi_B}{dt}
\]
\textbf{电场沿闭合回路的积分 = 磁通量变化的负值}——\textbf{Faraday 电磁感应定律}。

\subsection{Helmholtz 分解定理}

\textbf{任何（性质足够好的）矢量场}——\textbf{可以唯一分解为}：
\[
\bm{F} = -\nabla \phi + \nabla \times \bm{A}
\]
即\textbf{一个"无旋"分量 + 一个"无源"分量}。

\textbf{在电磁学里}——\textbf{$\phi$ 是电势——$\bm{A}$ 是磁矢势}——\textbf{它们是 Maxwell 方程组的"原函数"}。

\section{Aha 例子：手机信号场强 / Aha: Mobile Signal Strength}

\textbf{场景}：你在华为做天线工程师。要预测一个新基站的\textbf{信号覆盖图}——\textbf{每点的信号强度}。

\textbf{物理}：基站发射电磁波——电磁波的能量密度 $u = \frac{1}{2}(\varepsilon_0 E^2 + B^2/\mu_0)$。

\textbf{建模}（远场近似）：基站位置 $(0,0,h)$，频率 $f$，功率 $P$。则在地面 $(x, y, 0)$ 处的信号场强（近似）：
\[
E(x, y) = \frac{C \sqrt{P}}{r}, \quad r = \sqrt{x^2 + y^2 + h^2}
\]
$C$ 是常数。

\textbf{问题}：

\begin{enumerate}
\item \textbf{信号最弱的方向在哪}（梯度方向）
\item \textbf{封闭区域里的总发射能量是多少}（Gauss 定理）
\item \textbf{多个基站的信号叠加}（矢量场的线性叠加）
\end{enumerate}

\begin{tcolorbox}[ahabox={Aha: Maxwell 方程组的几何意义}]
\textbf{你以为 Maxwell 方程组是物理系学的"高深公式"}——\textbf{它其实是工程师每天都在用的工具}。

\medskip

\textbf{4G/5G 基站设计}——\textbf{$\nabla \cdot E = \rho/\varepsilon_0$}——\textbf{算电荷布局}\\
\textbf{Wi-Fi 路由器天线}——\textbf{$\nabla \times B = \mu_0 J + \mu_0 \varepsilon_0 \partial E/\partial t$}——\textbf{算电磁波辐射}\\
\textbf{电力变压器}——\textbf{$\nabla \times E = -\partial B/\partial t$}——\textbf{Faraday 电磁感应}\\
\textbf{MRI 核磁共振}——\textbf{整个 Maxwell 方程组的工程实现}

\medskip

\textbf{这些方程不是物理学家的私产}——\textbf{是 21 世纪工程的骨架}。

\medskip

\textbf{矢量分析——是它们的母语}。
\end{tcolorbox}

\textbf{Python 实现}：

\begin{lstlisting}[language=Python]
import numpy as np
import matplotlib.pyplot as plt

# 基站参数
h = 30        # 高度 30m
P = 100       # 功率 100W
C = 1.0       # 比例常数

# 场强计算
X, Y = np.meshgrid(np.linspace(-200, 200, 100),
                    np.linspace(-200, 200, 100))
r = np.sqrt(X**2 + Y**2 + h**2)
E = C * np.sqrt(P) / r

# 等高线 + 梯度
plt.contourf(X, Y, E, 20, cmap='viridis')
plt.colorbar(label='|E| field strength')
gx, gy = np.gradient(E)
plt.quiver(X[::5, ::5], Y[::5, ::5], gx[::5, ::5], gy[::5, ::5])
plt.title('Mobile Signal Coverage Map')
plt.savefig('signal_map.png')
\end{lstlisting}

\section{现代视角：Maxwell, 流体, ML / Modern}

\textbf{矢量分析的现代应用}：

\begin{itemize}
\item \textbf{Maxwell 方程组}——所有电磁工程的根
\item \textbf{Navier-Stokes 方程}——流体的"Maxwell"——$\partial \bm{v} / \partial t + (\bm{v} \cdot \nabla) \bm{v} = \cdots$
\item \textbf{弹性力学}——应力张量的散度 = 体力 + 惯性
\item \textbf{ML 里的"Hamiltonian flow"}——\textbf{Hamilton Monte Carlo}——\textbf{把矢量分析用到 Bayesian inference}
\end{itemize}

\textbf{你掌握了 $\nabla$——你掌握了 21 世纪工程数学的"通用语"}。

\section{代码与思考 / Code and Exercises}

\subsection{配套模块 \texttt{vector\_analysis.py}}

\begin{enumerate}
\item \texttt{demo\_gradient\_field()}——画 $\nabla(x^2+y^2)$ 的矢量场
\item \texttt{demo\_divergence()}——比较点源 vs 旋转场的散度
\item \texttt{demo\_curl()}——可视化旋转场的旋度
\item \texttt{demo\_gauss\_theorem()}——数值验证 Gauss 定理
\item \texttt{demo\_stokes\_theorem()}——数值验证 Stokes 定理
\item \texttt{demo\_signal\_coverage()}——本章 Aha 例子
\end{enumerate}

\subsection{思考题 / Exercises}

\begin{enumerate}
\item \textbf{[直觉]} 用一句话区分"散度"和"旋度"——不许用公式。
\item \textbf{[计算]} 求 $\bm{F} = (x^2 y, y^2 z, z^2 x)$ 的散度和旋度。
\item \textbf{[应用]} 一个 $\bm{B}$ 场 $\bm{B} = B_0 \hat{\bm{z}}$（均匀磁场）——它的散度和旋度各是什么？
\item \textbf{[代码]} 用 Python 数值验证 $\nabla \times \nabla f = \bm{0}$（任取一个 $f$）。
\item \textbf{[历史]} 查一查：Heaviside 为什么晚年贫困？他的工作为什么被低估？
\item \textbf{[ML]} 在 PyTorch 里实现一个简单的"矢量场"——计算它的 Jacobian——和"散度 + 旋度"对比。
\end{enumerate}

\section{本章小结 / Chapter Summary}

\begin{tcolorbox}[essbox={本章核心}]
\begin{itemize}
\item \textbf{$\nabla$} = $(\partial_x, \partial_y, \partial_z)$ = 矢量微分算符
\item \textbf{梯度 $\nabla f$} = 最陡上升方向（结果：矢量）
\item \textbf{散度 $\nabla \cdot \bm{F}$} = 源/汇的强度（结果：标量）
\item \textbf{旋度 $\nabla \times \bm{F}$} = 漩涡的强度（结果：矢量）
\item \textbf{Laplacian $\nabla^2 f$} = 偏离邻居平均的量（结果：标量）
\item \textbf{Gauss 定理} = 体积内的源 = 边界的流量
\item \textbf{Stokes 定理} = 曲面内的旋 = 边界的环流
\end{itemize}
\end{tcolorbox}

\textbf{下一章——常微分方程 ODE}——\textbf{我们让"变化率"出现在方程里}——\textbf{进入"动力系统"的世界}。

\clearpage

% =====================================================================
%  第 4 章 · 常微分方程 / Ordinary Differential Equations
%  一阶 ODE、二阶线性、系统与相图、数值解 (Euler/RK4)、应用
%  小故事：Bernoulli 一家与悬链线之争
%  Aha：阻尼振子 与 RLC 电路的同构 + SIR 传播模型
% =====================================================================

\chapter{常微分方程 / Ordinary Differential Equations}
\label{ch:ode}

\begin{quote}\itshape
\textbf{1696 年}——\textbf{Johann Bernoulli} 在《教师学报》上向全欧洲数学家挑战：
"一根均匀的链子两端固定挂着——它形成的曲线是什么？"

\medskip

他自己以为答案是抛物线——\textbf{猜错了}。他哥哥 \textbf{Jakob} 也猜错了。\\
真正解出来的是\textbf{年轻的 Leibniz、Huygens 和 Johann 自己}——答案是\textbf{双曲余弦} $y = a \cosh(x/a)$。

\medskip

这场\textbf{悬链线之争}（\textit{catenary controversy}），\textbf{是 ODE 的第一次"工业测试"}——\textbf{它告诉人类：很多自然形状不是代数曲线、而是微分方程的解}。

\medskip

我们这一章——讲的就是 Bernoulli 兄弟那一代人发明的工具：\textbf{常微分方程}。
\end{quote}

\section{写在前面 / Preface}

\textbf{我大二下学《工程数学》}——\textbf{ODE 那一章是最让我崩溃的}。

老师在黑板上一口气写了\textbf{六种"形式"}——可分离变量、齐次、一阶线性、Bernoulli、Riccati、全微分——\textbf{每种形式配一个"对应解法"}——\textbf{再每种配一道难题}。

我抄了\textbf{18 页笔记}——\textbf{期末考了 87 分}——\textbf{然后转身就忘了}。

\textbf{真正让我反应过来 ODE 是什么的}——是大三做\textbf{电力电子的课程设计}——\textbf{RLC 串联电路}：

\[
L \frac{d^2 i}{dt^2} + R \frac{di}{dt} + \frac{1}{C} i = \frac{dV}{dt}
\]

\textbf{我盯着这个方程看了 5 秒}——\textbf{突然想起来}——\textbf{这不就是高中物理里的\textbf{弹簧——阻尼——质量}系统？}

\[
m \ddot{x} + c \dot{x} + k x = F(t)
\]

\textbf{两个方程长得一模一样}——\textbf{只是物理量换了名字}。

那一刻我才意识到——\textbf{ODE 不是"六种形式 + 六种解法"}——\textbf{ODE 是"自然界写给工程师的母语"}——\textbf{电路、振动、化学反应、人口、传染病、金融——它们说的都是同一种话}。

\textbf{这一章的目的}——\textbf{把那 18 页笔记}——\textbf{压缩成 3 张"通用图"}：

\begin{itemize}
\item \textbf{一阶 ODE}——"\textbf{当下变化率 = 当下状态}"——\textbf{指数增长 / 衰减}
\item \textbf{二阶线性 ODE}——"\textbf{加速度 = 弹力 + 阻力 + 外力}"——\textbf{振动 / 电路}
\item \textbf{ODE 系统}——"\textbf{多个状态互相影响}"——\textbf{捕食者——猎物 / SIR 传染病}
\end{itemize}

\textbf{再加一招——\texttt{scipy.integrate.solve\_ivp}}——\textbf{你几乎所有的工程 ODE 都能算出来}。

\section{一句话本质 / The Essence in One Sentence}

\begin{tcolorbox}[essbox={一句话本质}]
\textbf{中文}：\textbf{ODE 是"用变化率描述未来"——给定当下状态——它告诉你下一刻去哪里}。\\[6pt]
\textbf{English}: An ODE describes the future through rates of change——given the present state, it tells you where the next instant goes.
\end{tcolorbox}

\textbf{再说一遍}——\textbf{ODE 不是"解方程"}——\textbf{ODE 是"写一条规则——让时间帮你演化"}。

\section{教材里你看到的 / What the Textbook Tells You}

\subsection{同济《高等数学下册》第十二章}

同济把 ODE 放在\textbf{下册的倒数第二章}——\textbf{30 多页}——\textbf{讲了 8 节}：

\begin{itemize}
\item 微分方程的基本概念
\item 可分离变量的微分方程
\item 齐次方程
\item 一阶线性方程
\item 可降阶的高阶方程
\item 高阶线性方程
\item 常系数齐次线性方程
\item 常系数非齐次线性方程
\end{itemize}

\textbf{优点}：\textbf{完整}——所有"分类与解法"都覆盖了——\textbf{考研够用}。

\textbf{缺点}：

\begin{enumerate}
\item \textbf{只讲"解法"——不讲"为什么这么解"}——\textbf{学生背了 6 套口诀——不知道哪一套对哪一种物理}
\item \textbf{没有相图}——没有"\textbf{用几何看 ODE}"——这是\textbf{现代动力系统的入门视角}
\item \textbf{没有数值方法}——\textbf{Euler 法、RK4——这是工业里 99\% 的 ODE 解法}——同济一字不提
\item \textbf{应用部分孤立}——讲了"\textbf{LR 电路}"——但\textbf{没讲它和"弹簧振子"是同一件事}
\end{enumerate}

\subsection{Boyce \& DiPrima《Elementary Differential Equations》}

\textbf{美国的国民 ODE 教材}——\textbf{900 页}——\textbf{讲得活、图多、相图齐全}。

它比同济好的地方：\textbf{有相图、有数值方法（一整章 Euler + RK4）、有 Laplace 变换的完整故事}。

它的缺点：\textbf{900 页太长}——\textbf{中国工程师没时间通读}。

\subsection{Strogatz《Nonlinear Dynamics and Chaos》}

\textbf{Strogatz 这本是非线性 ODE 的"必读"}——\textbf{用"几何 + 故事"讲分岔、混沌、Lorenz 吸引子}——\textbf{读起来像小说}。

\textbf{它不教解法}——\textbf{它教你"看相图"的本能}。\textbf{推荐每个工程师都翻一遍}。

\subsection{我的策略}

本章只做三件事：

\begin{enumerate}
\item \textbf{把"六种解法"压缩成两条主线}——\textbf{一阶 = 指数家族}——\textbf{二阶线性 = 振动家族}
\item \textbf{第一次引入相图}——\textbf{让你"看见" ODE 的解——而不只是"算出"它}
\item \textbf{把数值方法说透到一行代码}——\textbf{\texttt{solve\_ivp(...)} 背后是 Euler + RK4 的现代化}
\end{enumerate}

\begin{tcolorbox}[storybox={\Large 小故事：Bernoulli 一家与悬链线之争}]
\textbf{17 世纪末的瑞士巴塞尔}——\textbf{有一个特殊的家族}——\textbf{Bernoulli 家族}——\textbf{三代人里出了八位数学家}——\textbf{史无前例}。

\medskip

\textbf{Jakob Bernoulli}（1654--1705）——\textbf{大哥}——\textbf{巴塞尔大学数学教授}——\textbf{发明了"Bernoulli 数"和"大数定律"}。

\textbf{Johann Bernoulli}（1667--1748）——\textbf{弟弟}——\textbf{比 Jakob 小 13 岁}——\textbf{脾气暴躁、好胜心强}——\textbf{后来教出了 Euler}。

\textbf{Daniel Bernoulli}（1700--1782）——\textbf{Johann 的儿子}——\textbf{流体力学里"Bernoulli 方程"就是他的工作}。

\medskip

\textbf{1690 年}——\textbf{Jakob 在《教师学报》（Acta Eruditorum）上提了一个问题}：

\medskip

\centerline{\itshape "一根均匀柔软的链子，两端固定，自然下垂——它的形状是什么曲线？"}

\medskip

\textbf{Galileo 在 1638 年猜过}——\textbf{他说是\textbf{抛物线}}——\textbf{听起来很合理}（抛物线也长那样）。

\medskip

\textbf{1691 年}——\textbf{Leibniz、Huygens、Johann Bernoulli 三个人独立解出来了}——\textbf{答案不是抛物线}——\textbf{是}：

\[
y = a \cosh\!\left(\frac{x}{a}\right) = \frac{a}{2}\left(e^{x/a} + e^{-x/a}\right)
\]

\medskip

这个曲线叫\textbf{悬链线（catenary）}——\textbf{源自拉丁语 catena（链子）}。

\medskip

\textbf{怎么解出来的？}——\textbf{他们写下了一个二阶 ODE}：

\[
y'' = \frac{1}{a}\sqrt{1 + (y')^2}
\]

\textbf{背后的物理}——\textbf{链子上每一小段——受重力 + 两端张力}——\textbf{力平衡——化成微分方程}——\textbf{再积分}。

\medskip

\textbf{这是 ODE 的"第一次公开胜利"}——\textbf{它解出了一个\textbf{2000 年没人能用代数搞定的几何问题}}。

\medskip

\textbf{Johann 后来在书信里偷偷得意}——\textbf{"我用了\textbf{一晚上}就解出来了——我哥哥 Jakob 想了\textbf{一整年}才追上"}——\textbf{兄弟俩从此结下梁子——一辈子打数学口水仗}。

\medskip

\textbf{今天}——\textbf{St. Louis Gateway Arch}（圣路易斯大拱门）——\textbf{Antoni Gaud{\'i} 的圣家堂}——\textbf{都是按悬链线（倒过来）设计的}——\textbf{因为这是\textbf{纯压缩——没有弯矩}的天然结构}。

\medskip

\textbf{Bernoulli 兄弟 330 年前为之吵架的曲线}——\textbf{今天还在撑着 192 米高的拱门}。
\end{tcolorbox}

\section{直觉 / Intuition}

\subsection{ODE 的本质——"局部规则"决定"全局轨迹"}

\textbf{一句话——ODE 是这样工作的}：

\begin{quote}
\textbf{"你告诉我现在在哪——我告诉你下一刻往哪走。"}\\[4pt]
\itshape "Tell me where you are now——I'll tell you where to go next."
\end{quote}

\textbf{每个 ODE 都长这样}：
\[
\frac{dy}{dt} = f(t, y)
\]

\textbf{这个 $f$——就是"游戏规则"}——\textbf{它告诉时间——在每个状态 $y$ 下——下一刻该怎么变}。

\textbf{ODE 不像代数方程"求未知数"}——\textbf{ODE 是"求未知函数"}——\textbf{它的答案是\textbf{一条曲线 $y(t)$}}——\textbf{不是\textbf{一个数}}。

\subsection{斜率场 / Slope Field——把 ODE 画出来}

\textbf{斜率场是 ODE 最直观的可视化}：

\begin{itemize}
\item 在平面 $(t, y)$ 上的\textbf{每一点}——\textbf{画一根短斜线}——\textbf{斜率 $= f(t, y)$}
\item 一个\textbf{初值 $(t_0, y_0)$}——\textbf{决定了一条曲线}——\textbf{沿着斜率场"流"下去}
\end{itemize}

\textbf{初值不同——曲线不同}——\textbf{这就是"初值问题"（IVP）的画面}。

\subsection{相图 / Phase Portrait——把"轨迹"画出来}

对\textbf{二阶 ODE 或 ODE 系统}——\textbf{斜率场不够用}——\textbf{我们画相图}：

\begin{itemize}
\item 横轴 $x$、纵轴 $\dot x$——\textbf{每一个状态是平面上的一个点}
\item \textbf{解 $\to$ 平面上的一条轨迹}——\textbf{时间是"运动方向"}
\end{itemize}

\textbf{相图能一眼看出}：

\begin{itemize}
\item 系统是\textbf{稳定的}吗？（所有轨迹收敛到某点 $\to$ \textbf{吸引子}）
\item 系统会\textbf{振荡}吗？（轨迹绕圈 $\to$ \textbf{极限环}）
\item 系统会\textbf{发散}吗？（轨迹飞出去 $\to$ \textbf{不稳定}）
\end{itemize}

\textbf{这是同济不讲的"几何视角"}——\textbf{却是动力系统的现代主流语言}。

\section{数学 / The Math}

\subsection{一阶 ODE / First-Order ODEs}

\textbf{一般形式}：
\[
\frac{dy}{dx} = f(x, y)
\]

\textbf{三种最常见的可解形式}——\textbf{背一遍——其余交给数值}：

\medskip

\textbf{(1) 可分离变量 / Separable}：

\[
\frac{dy}{dx} = g(x) h(y) \;\Longrightarrow\; \int \frac{dy}{h(y)} = \int g(x) \, dx + C
\]

\textbf{例}：$\dfrac{dy}{dx} = -k y$（指数衰减）$\;\Rightarrow\;y = y_0 e^{-kx}$。

\textbf{应用}：\textbf{放射性衰变、电容放电、Newton 冷却定律}。

\medskip

\textbf{(2) 一阶线性 / First-Order Linear}：

\[
\frac{dy}{dx} + P(x) y = Q(x)
\]

\textbf{解法——乘"积分因子" $\mu(x) = e^{\int P(x) dx}$}：

\[
y = \frac{1}{\mu(x)}\left[\int \mu(x) Q(x) \, dx + C\right]
\]

\textbf{应用}：\textbf{RC 电路的充电过程、混合溶液的浓度变化}。

\medskip

\textbf{(3) Bernoulli 方程 / Bernoulli Equation}（\textbf{1695 年由 Jakob Bernoulli 提出}）：

\[
\frac{dy}{dx} + P(x) y = Q(x) y^n \quad (n \neq 0, 1)
\]

\textbf{解法——令 $z = y^{1-n}$——化成一阶线性}。

\textbf{应用}：\textbf{Logistic 增长模型}（$n=2$）——\textbf{种群生态、SIR 传染病}。

\begin{example}[Newton 冷却定律 / Newton's Law of Cooling]
一杯 90°C 的咖啡放在 20°C 的房间——\textbf{温度变化率与温差成正比}：
\[
\frac{dT}{dt} = -k(T - T_{\text{环境}}), \quad T(0) = 90
\]
\textbf{这是可分离 + 线性双重身份}。解：
\[
T(t) = T_{\text{环境}} + (T_0 - T_{\text{环境}}) e^{-kt} = 20 + 70 e^{-kt}
\]
\textbf{这就是为什么咖啡总是先快、后慢地冷下来}——\textbf{指数衰减是物理界的"母模板"}。
\end{example}

\subsection{二阶线性 ODE / Second-Order Linear ODEs}

\textbf{一般形式}：
\[
a y'' + b y' + c y = f(x)
\]

\textbf{解的结构}（\textbf{这是同济讲对了的地方}）：

\[
y(x) = \underbrace{y_h(x)}_{\text{齐次解}} + \underbrace{y_p(x)}_{\text{特解}}
\]

\subsubsection{齐次方程 / Homogeneous}

\[
a y'' + b y' + c y = 0
\]

\textbf{试探解 $y = e^{\lambda x}$}——代入——得\textbf{特征方程}：
\[
a \lambda^2 + b \lambda + c = 0
\]

\textbf{三种情况}：

\begin{enumerate}
\item \textbf{两个不同实根} $\lambda_1, \lambda_2$：$y_h = C_1 e^{\lambda_1 x} + C_2 e^{\lambda_2 x}$ —— \textbf{过阻尼振动}
\item \textbf{重根} $\lambda$：$y_h = (C_1 + C_2 x) e^{\lambda x}$ —— \textbf{临界阻尼}
\item \textbf{共轭复根} $\alpha \pm i\beta$：$y_h = e^{\alpha x}(C_1 \cos\beta x + C_2 \sin\beta x)$ —— \textbf{欠阻尼振动}
\end{enumerate}

\textbf{这三种情况——对应物理上的三种"振动行为"}——\textbf{这是 ODE 和物理的第一次"语义对齐"}。

\subsubsection{非齐次：待定系数法 / Method of Undetermined Coefficients}

对 $a y'' + b y' + c y = f(x)$ —— \textbf{$f(x)$ 形式决定 $y_p$ 形式}：

\begin{center}
\begin{tabular}{ll}
\toprule
\textbf{$f(x)$ 的形式} & \textbf{试探 $y_p$} \\
\midrule
$P_n(x)$（多项式） & $Q_n(x)$（同次多项式） \\
$e^{\alpha x}$ & $A e^{\alpha x}$ \\
$\sin\beta x, \cos\beta x$ & $A\cos\beta x + B\sin\beta x$ \\
$e^{\alpha x}\sin\beta x$ 等 & $e^{\alpha x}(A\cos\beta x + B\sin\beta x)$ \\
\bottomrule
\end{tabular}
\end{center}

\textbf{口诀}：\textbf{"什么形式来——什么形式去"}——\textbf{除非和齐次解共振——共振要乘 $x$}。

\subsubsection{常数变易法 / Variation of Parameters}

\textbf{若 $f(x)$ 形式怪}——\textbf{用常数变易}：
\[
y_p(x) = -y_1(x) \int \frac{y_2(x) f(x)}{a W(x)} dx + y_2(x) \int \frac{y_1(x) f(x)}{a W(x)} dx
\]
其中 $W = y_1 y_2' - y_1' y_2$ 是\textbf{Wronskian 行列式}——\textbf{它衡量两个解"线性无关到什么程度"}。

\subsection{ODE 系统 / Systems of ODEs}

\textbf{一般形式}（用向量）：
\[
\frac{d\bm{y}}{dt} = \bm{f}(t, \bm{y}), \quad \bm{y} \in \mathbb{R}^n
\]

\textbf{重要事实}——\textbf{任何 $n$ 阶 ODE 都能化成 $n$ 维的一阶 ODE 系统}：

二阶 ODE $y'' + p y' + q y = 0$ —— 令 $y_1 = y, y_2 = y'$ —— 化成：
\[
\begin{pmatrix} \dot y_1 \\ \dot y_2 \end{pmatrix} = \begin{pmatrix} 0 & 1 \\ -q & -p \end{pmatrix} \begin{pmatrix} y_1 \\ y_2 \end{pmatrix}
\]

\textbf{这就是数值求解器只支持一阶系统的原因}——\textbf{所有高阶 ODE 先化成一阶}。

\subsubsection{线性系统的稳定性 / Stability of Linear Systems}

\textbf{$\dot{\bm{y}} = A \bm{y}$——稳定性由 $A$ 的特征值决定}：

\begin{itemize}
\item \textbf{所有特征值实部 $< 0$}：\textbf{稳定}（轨迹收敛到原点）
\item \textbf{有特征值实部 $> 0$}：\textbf{不稳定}（轨迹发散）
\item \textbf{纯虚根}：\textbf{中心}（轨迹绕圆/椭圆）
\end{itemize}

\textbf{这是\textbf{线性代数和 ODE 的"接口"}}——\textbf{为什么我们要学特征值？因为它告诉我们一个系统会不会"爆"}。

\subsection{数值解法 / Numerical Methods}

\subsubsection{Euler 法 / Euler's Method}

\textbf{最简单的数值解法}——\textbf{Newton 1687 年就用过}：

\[
y_{n+1} = y_n + h \cdot f(t_n, y_n)
\]

\textbf{几何意义}：\textbf{在当前点画切线——走一步}。

\textbf{误差}：\textbf{每步 $O(h^2)$——总误差 $O(h)$}——\textbf{$h$ 减半——精度只翻一倍}——\textbf{慢}。

\subsubsection{Runge--Kutta 4 阶法 / RK4}

\textbf{工业界默认的数值求解器}——\textbf{由 Carl Runge 和 Wilhelm Kutta 在 1900 年前后发明}：

\begin{align}
k_1 &= f(t_n, y_n) \\
k_2 &= f(t_n + h/2, y_n + h k_1 / 2) \\
k_3 &= f(t_n + h/2, y_n + h k_2 / 2) \\
k_4 &= f(t_n + h, y_n + h k_3) \\
y_{n+1} &= y_n + \frac{h}{6}(k_1 + 2k_2 + 2k_3 + k_4)
\end{align}

\textbf{误差}：\textbf{每步 $O(h^5)$——总误差 $O(h^4)$}——\textbf{$h$ 减半——精度翻 16 倍}——\textbf{这就是为什么 RK4 是默认选择}。

\textbf{现代版}：\textbf{\texttt{scipy.integrate.solve\_ivp(method='RK45')}}——\textbf{它是带自适应步长的 5 阶 Runge--Kutta}——\textbf{你只需要写微分方程的右端、给初值、范围——剩下全交给它}。

\begin{remark}[工业现实]
\textbf{99\% 的工程 ODE 都是数值解}——\textbf{解析解只在两个地方有用}：
\begin{itemize}
\item \textbf{教材考试}（包括考研）
\item \textbf{你要做\textbf{参数研究}}——比如"阻尼比 $\zeta$ 从 0.1 到 1.0 变化时，振荡频率怎么变"——\textbf{这时解析解更直观}
\end{itemize}
\textbf{其余}——\textbf{打开 SciPy——10 行搞定}。
\end{remark}

\section{Aha 例子：阻尼振子 与 RLC 电路的同构 / Aha: Damped Oscillator ≡ RLC Circuit}

\textbf{场景}——\textbf{你是机械工程师}——\textbf{设计一辆车的悬架}：

\[
m \ddot{x} + c \dot{x} + k x = F(t)
\]

\textbf{$m$ 是车身质量、$c$ 是减震器阻尼、$k$ 是弹簧刚度、$F(t)$ 是路面冲击}。

\medskip

\textbf{同时}——\textbf{你的同事是电气工程师}——\textbf{设计一个 RLC 滤波器}：

\[
L \ddot{q} + R \dot{q} + \frac{1}{C} q = V(t)
\]

\textbf{$L$ 是电感、$R$ 是电阻、$C$ 是电容、$V(t)$ 是输入电压、$q$ 是电荷}。

\medskip

\textbf{两个方程长得一模一样}——\textbf{这不是巧合}——\textbf{这是 ODE 在告诉你}：

\medskip

\begin{center}
\begin{tabular}{ccc}
\toprule
\textbf{机械} & \textbf{电气} & \textbf{角色} \\
\midrule
质量 $m$ & 电感 $L$ & 储存\textbf{动能 / 磁能}（"惯性"） \\
阻尼 $c$ & 电阻 $R$ & 耗散能量 \\
弹簧 $k$ & 电容倒数 $1/C$ & 储存\textbf{势能 / 电场能}（"恢复力"） \\
位移 $x$ & 电荷 $q$ & 状态变量 \\
速度 $\dot x$ & 电流 $\dot q = i$ & 状态变量的变化率 \\
外力 $F$ & 电压 $V$ & 外部驱动 \\
\bottomrule
\end{tabular}
\end{center}

\begin{tcolorbox}[ahabox={Aha: 同一个 ODE——同时是机械振动 和 电路滤波}]
\textbf{你以为电路和机械是两门课}——\textbf{老师把它们分开教}——\textbf{教材封面不一样、单位不一样、符号不一样}。

\medskip

\textbf{但它们的"数学骨架"是同一个二阶线性 ODE}：

\[
\ddot{y} + 2\zeta \omega_n \dot{y} + \omega_n^2 y = u(t)
\]

\textbf{其中}：
\begin{itemize}
\item $\omega_n = \sqrt{k/m} = 1/\sqrt{LC}$ —— \textbf{自然频率}
\item $\zeta = \dfrac{c}{2\sqrt{km}} = \dfrac{R}{2}\sqrt{\dfrac{C}{L}}$ —— \textbf{阻尼比}
\end{itemize}

\medskip

\textbf{$\zeta$ 这一个数——决定了系统所有"性格"}：

\begin{itemize}
\item $\zeta = 0$：\textbf{无阻尼——永远振荡}（理想 LC 振荡器、单摆）
\item $0 < \zeta < 1$：\textbf{欠阻尼——振荡衰减}（汽车悬架、铃铛）
\item $\zeta = 1$：\textbf{临界阻尼——最快回归}（关门器、相机镜头电机）
\item $\zeta > 1$：\textbf{过阻尼——缓慢回归}（液压门、太软的弹簧）
\end{itemize}

\medskip

\textbf{这就是 ODE 的"母语威力"}——\textbf{一个方程——同时描述\textbf{机械、电气、声学、生物}的振动}——\textbf{你学会了一个——就学会了一切}。

\medskip

\textbf{大学课表把它分成 6 门课}——\textbf{每门课用不同的符号}——\textbf{但骨子里是同一个 ODE}——\textbf{这就是"打通"要做的事}。
\end{tcolorbox}

\textbf{解决问题的 Python 代码}：

\begin{lstlisting}[language=Python]
import numpy as np
from scipy.integrate import solve_ivp
import matplotlib.pyplot as plt

# 物理参数 (车辆悬架)
m, c, k = 300.0, 1200.0, 30000.0   # kg, N·s/m, N/m
omega_n = np.sqrt(k/m)             # 自然频率
zeta    = c / (2*np.sqrt(k*m))     # 阻尼比
print(f"omega_n = {omega_n:.2f} rad/s,  zeta = {zeta:.3f}")

# 写成一阶系统：y = [x, x_dot]
def rhs(t, y, F):
    x, xdot = y
    xddot = (F(t) - c*xdot - k*x) / m
    return [xdot, xddot]

# 外部冲击：t=0 时 1000 N 阶跃
F = lambda t: 1000.0 if t > 0 else 0.0

# 数值求解：RK45 + 自适应步长
sol = solve_ivp(rhs, t_span=(0, 2.0), y0=[0, 0],
                args=(F,), max_step=0.001, dense_output=True)

t = np.linspace(0, 2, 2000)
x = sol.sol(t)[0]
plt.plot(t, x*1000); plt.xlabel('t [s]'); plt.ylabel('x [mm]')
plt.title(f'Damped oscillator: zeta = {zeta:.3f}'); plt.grid(True)
\end{lstlisting}

\textbf{把这段代码改成"RLC 电路"}——\textbf{只换一行}：

\begin{lstlisting}[language=Python]
# 同样的代码——换成 RLC：只改物理参数
L, R, C = 0.1, 50.0, 10e-6        # H, Ohm, F
omega_n = 1/np.sqrt(L*C)
zeta    = (R/2) * np.sqrt(C/L)
# rhs 也是 [q, q_dot]：L*q_ddot + R*q_dot + q/C = V(t)
\end{lstlisting}

\textbf{这就是"同构"的力量}——\textbf{一个求解器——通吃两个学科}。

\subsection{Aha 续：SIR 模型与 COVID-19}

\textbf{2020 年初}——\textbf{所有人突然变成了"流行病学家"}——\textbf{每天讨论 R0、群体免疫}。

背后的数学——\textbf{是 1927 年 Kermack \& McKendrick 提出的 SIR 模型}——\textbf{一个 3 维 ODE 系统}：

\begin{align}
\dot{S} &= -\beta \, S I / N \\
\dot{I} &= \beta \, S I / N - \gamma I \\
\dot{R} &= \gamma I
\end{align}

\textbf{$S$ = 易感者、$I$ = 感染者、$R$ = 康复者、$N = S + I + R$ 总人口}。

\textbf{$\beta$ = 接触传播率、$\gamma$ = 康复率}——\textbf{基本再生数 $R_0 = \beta/\gamma$}。

\begin{lstlisting}[language=Python]
def sir(t, y, beta, gamma, N):
    S, I, R = y
    dS = -beta * S * I / N
    dI =  beta * S * I / N - gamma * I
    dR =  gamma * I
    return [dS, dI, dR]

# COVID-19 早期估计 (示意值)
N, beta, gamma = 1e7, 0.35, 0.10
R0 = beta / gamma           # 基本再生数 ≈ 3.5

sol = solve_ivp(sir, (0, 200), [N-1, 1, 0],
                args=(beta, gamma, N), dense_output=True)
\end{lstlisting}

\textbf{这就是\textbf{2020 年 3 月各国政府决策依据的"母方程"}}——\textbf{它是 ODE 系统——不是单个 ODE}——\textbf{它有相图}——\textbf{它告诉你"封城（降低 $\beta$）会让峰值推后并压低"}。

\textbf{一个 1927 年的 ODE}——\textbf{在 2020 年决定了全球 70 亿人的生活}。

\section{现代视角：神经 ODE 与可微仿真 / Modern: Neural ODEs and Differentiable Simulation}

\textbf{ODE 在 21 世纪有了第二次青春}——\textbf{它和深度学习"联姻"了}。

\subsection{Neural ODE（2018 年 NeurIPS 最佳论文）}

\textbf{ResNet} 里残差块的形式 $y_{n+1} = y_n + f(y_n, \theta_n)$ —— 像不像 Euler 法 $y_{n+1} = y_n + h \cdot f(y_n)$？

\textbf{Chen 等人 2018 年的洞察}——\textbf{把网络层数 $\to \infty$、步长 $h \to 0$}——\textbf{ResNet 就变成 ODE}：

\[
\frac{d\bm{h}(t)}{dt} = f(\bm{h}(t), t, \theta)
\]

\textbf{学习这个"右端 $f$"——就是 Neural ODE}。\textbf{优点}：连续深度、内存恒定（用伴随方法反传梯度）。

\subsection{可微物理仿真}

\textbf{\texttt{diffrax}（JAX）、\texttt{torchdiffeq}（PyTorch）}——\textbf{把 ODE 求解器变成"可求导的"}——\textbf{这样你可以\textbf{用梯度下降优化物理系统的参数}}：

\begin{itemize}
\item 用\textbf{真实测量数据}——\textbf{反演弹簧刚度 $k$}
\item 用\textbf{目标轨迹}——\textbf{优化无人机控制律}
\item 用\textbf{化学反应曲线}——\textbf{反演反应速率常数}
\end{itemize}

\textbf{这是\textbf{科学计算 + 机器学习}的结合点}——\textbf{现在叫"科学机器学习（SciML）"}——\textbf{是 2030 年代最热的方向之一}。

\subsection{时间线}

\begin{itemize}
\item \textbf{1690}：\textbf{Bernoulli 兄弟}——\textbf{ODE 的几何起源}（悬链线）
\item \textbf{1736}：\textbf{Euler}——\textbf{发明 Euler 法}
\item \textbf{1900}：\textbf{Runge \& Kutta}——\textbf{RK4 现代数值方法}
\item \textbf{1927}：\textbf{Kermack \& McKendrick}——\textbf{SIR 模型}
\item \textbf{1963}：\textbf{Lorenz}——\textbf{偶然发现混沌}（蝴蝶效应）
\item \textbf{2018}：\textbf{Neural ODE}——\textbf{深度学习与 ODE 合体}
\end{itemize}

\textbf{ODE 不是死的}——\textbf{它每隔一代人就被"重新发明"一次}。

\section{代码与思考 / Code and Exercises}

\subsection{配套模块 \texttt{ode.py}}

GitHub 上的 \texttt{modules/ode.py} 包含\textbf{6 个 demo}：

\begin{enumerate}
\item \texttt{demo\_first\_order()}——可视化 Newton 冷却定律——\textbf{解析解 vs 数值解}
\item \texttt{demo\_damped\_oscillator()}——\textbf{扫描 $\zeta$ 从 0.1 到 2.0}——画出\textbf{欠/临界/过阻尼}三种行为
\item \texttt{demo\_rlc\_vs\_mass\_spring()}——\textbf{本章 Aha 例子}：同一个求解器跑两个物理系统——证明它们同构
\item \texttt{demo\_phase\_portrait()}——\textbf{画相图}：用 \texttt{matplotlib.quiver} 画矢量场 + 几条典型轨迹
\item \texttt{demo\_sir\_covid()}——\textbf{SIR 模型}：扫描 $R_0$ 从 1.5 到 4.0——\textbf{演示"封城（降低 $\beta$）"的效果}
\item \texttt{demo\_euler\_vs\_rk4()}——\textbf{比较 Euler 法和 RK4}的精度——\textbf{画 log--log 误差——步长曲线}
\end{enumerate}

\textbf{使用}：

\begin{lstlisting}[language=bash]
git clone https://github.com/inaturelz-coder/math-rendu
cd math-rendu
pip install -r requirements.txt
python modules/ode.py
\end{lstlisting}

\subsection{思考题 / Exercises}

\begin{enumerate}
\item \textbf{[直觉]} 用一句话解释为什么\textbf{Newton 冷却定律}产生的是\textbf{指数衰减}——不许用积分。
\item \textbf{[计算]} 解一阶线性方程 $y' + 2y = e^{-x}$，初值 $y(0)=1$。再用 \texttt{solve\_ivp} 数值验证。
\item \textbf{[计算]} 求 $y'' - 3y' + 2y = e^{3x}$ 的通解（待定系数法）。
\item \textbf{[物理]} 单摆方程 $\ddot{\theta} + (g/L)\sin\theta = 0$ \textbf{不是线性 ODE}——小角度近似下退化为 SHM。用 \texttt{solve\_ivp} 比较\textbf{非线性精确解}和\textbf{线性近似解}在初角 $\theta_0 = 60°$ 时的差异。
\item \textbf{[代码]} 实现你自己的 RK4 求解器（不用 SciPy），求解 $\dot y = -y + \sin t$，$y(0)=0$，与 SciPy 的 \texttt{RK45} 对比误差。
\item \textbf{[应用]} 用 SIR 模型——给定 $R_0 = 2.5, \gamma = 1/14$——\textbf{画出疫情曲线}——并\textbf{找到峰值时刻和峰值感染人数}。\textbf{若 30 天后突然封城（让 $\beta$ 减半）}——\textbf{画对比曲线}。
\end{enumerate}

\section{本章小结 / Chapter Summary}

\begin{tcolorbox}[essbox={本章核心}]
\begin{itemize}
\item \textbf{ODE 的本质}——"\textbf{用变化率描述未来}"——给当下、定演化
\item \textbf{一阶}——\textbf{指数家族}——衰变、冷却、Logistic 增长
\item \textbf{二阶线性}——\textbf{振动家族}——\textbf{特征方程 $a\lambda^2+b\lambda+c=0$ 的三种根 = 三种振动行为}
\item \textbf{系统}——任何高阶 ODE 都化成一阶系统——\textbf{稳定性看 $A$ 的特征值}
\item \textbf{相图}——\textbf{ODE 的"几何视角"}——一眼看出稳定性、吸引子、极限环
\item \textbf{数值}——\textbf{Euler 易懂 / RK4 实用}——现代版 \texttt{solve\_ivp(method='RK45')}
\item \textbf{同构}——\textbf{机械振动 ≡ RLC 电路}——\textbf{ODE 是工程师的母语}
\item \textbf{现代}——\textbf{Neural ODE + 可微仿真——ODE 与深度学习联姻}
\end{itemize}
\end{tcolorbox}

\textbf{下一章——偏微分方程}——\textbf{我们从"时间一维"走到"时间+空间多维"}——\textbf{见到\textbf{热方程、波方程、Laplace 方程}}——\textbf{并理解为什么 Fourier 变换是 PDE 的"瑞士军刀"}。

\clearpage

% =====================================================================
%  第 5 章 · 偏微分方程 / Partial Differential Equations
%  椭圆、抛物、双曲——分离变量、Fourier、有限差分、有限元
%  小故事：Fourier 被法兰西学院拒稿三次
%  Aha：手机芯片散热仿真
% =====================================================================

\chapter{偏微分方程 / Partial Differential Equations}
\label{ch:pde}

\begin{quote}\itshape
\textbf{1807 年}——\textbf{Joseph Fourier} 把一篇关于\textbf{热传导}的论文交给法兰西学院。\\
评审是 \textbf{Lagrange} 和 \textbf{Laplace}——\textbf{当时数学界的两座大山}。\\
他们说：\textbf{"你这个'三角级数'——一会儿连续一会儿不连续——不严格——拒。"}\\
Fourier 改了再交——\textbf{再拒}。再改再交——\textbf{再拒}。\\
直到 1822 年——他干脆\textbf{自己出书}——《\textbf{热的解析理论}》——\textbf{现代 PDE 和信号处理的奠基之作}。\\[4pt]
我们这一章——讲的就是 Fourier 用一辈子才让世人接受的那套方法——\textbf{偏微分方程 + 分离变量 + 三角级数}。
\end{quote}

\section{写在前面 / Preface}

\textbf{我刚接触 PDE 是大三《传热学》}——\textbf{老师在黑板上写了热传导方程}：
\[
\frac{\partial u}{\partial t} = \alpha \frac{\partial^2 u}{\partial x^2}
\]
\textbf{然后说}：\textbf{"这就是描述温度随时间和空间变化的方程——大家回去自己看。"}

\textbf{我懵了整整一周}。

那个左边的 $\partial u / \partial t$——我懂——温度对时间的变化率。
那个右边的 $\partial^2 u / \partial x^2$——我也勉强懂——温度对空间的二阶导。
但\textbf{"温度的变化率 等于 温度的二阶导"}——\textbf{这是什么物理}？

\textbf{这一周我跑去蹭了《数学物理方法》的课}——\textbf{讲到分离变量法}——\textbf{$u(x,t) = X(x) T(t)$}——\textbf{老师写满了三块黑板的求解步骤}——\textbf{我抄了一份回去}——\textbf{依然不懂}。

\textbf{直到我用 Python 写了 10 行有限差分代码}——\textbf{把一根铁棒切成 100 段}——\textbf{每个时间步让每一段的温度 = 自己 + $\alpha\Delta t / \Delta x^2 \times$（左右邻居 - 2 倍自己）}——\textbf{我看到一个初始的"尖峰"温度——慢慢摊平、扩散、变成平台}。

\textbf{那一刻我才反应过来}：

\begin{quote}
\textbf{"二阶导 $\partial^2 u / \partial x^2$——就是\textbf{邻居平均 - 自己}"}——\textbf{所以热传导方程在说一件最朴素的事}——\textbf{每个点的温度——在向它的邻居平均靠拢}——\textbf{这就是"扩散"}。
\end{quote}

\textbf{PDE 不是用来折磨工科生的}——\textbf{它是用来描述"场"的语言}。温度场、电场、磁场、压力场、速度场——\textbf{凡是"在空间里铺开 + 随时间演化"的东西}——\textbf{都得用 PDE}。

\textbf{这一章只做四件事}：

\begin{enumerate}
\item \textbf{讲清 PDE 三大类型}——\textbf{椭圆、抛物、双曲}——\textbf{它们对应三种完全不同的物理}
\item \textbf{讲清分离变量法}——\textbf{它为什么会引出 Fourier 级数}
\item \textbf{讲清有限差分}——\textbf{它就是把 PDE "搬"到 numpy 数组上的方法}
\item \textbf{Aha 例子}——\textbf{手机芯片散热——20 行 Python——你能看见温度在硅片上扩散}
\end{enumerate}

\section{一句话本质 / The Essence in One Sentence}

\begin{tcolorbox}[essbox={一句话本质}]
\textbf{中文}：\textbf{偏微分方程描述"场"如何随空间和时间演化——三大类型分别对应\textbf{扩散}（抛物）、\textbf{振动}（双曲）、\textbf{平衡}（椭圆）}。\\[6pt]
\textbf{English}: A PDE describes how a field evolves in space and time. The three canonical types correspond respectively to \emph{diffusion} (parabolic), \emph{wave propagation} (hyperbolic), and \emph{equilibrium} (elliptic).
\end{tcolorbox}

\textbf{记住这一句话——这一章已经讲完了三分之一}。

\section{教材里你看到的 / What the Textbook Tells You}

\subsection{梁昆淼《数学物理方法》}

中国物理系学生人手一本的\textbf{梁昆淼《数学物理方法》}——\textbf{下册整本讲 PDE}——\textbf{讲得非常深}——\textbf{从分离变量到 Sturm-Liouville 到 Green 函数到积分变换}。

\textbf{优点}：\textbf{把"为什么"讲到了}——\textbf{比工科教材深一层}。

\textbf{缺点}：\textbf{物理系思维}——\textbf{偏理论 + 偏特殊函数}（Bessel、Legendre、Hermite 各占一章）——\textbf{工科生看着觉得"和我的电机/芯片/管道有什么关系"}。

\subsection{Zill《Advanced Engineering Mathematics》}

\textbf{Zill 第 12--15 章}——\textbf{专门讲 PDE}——\textbf{从分离变量讲到 Fourier 变换到 Laplace 变换}——\textbf{有不少应用例子}。

\textbf{优点}：\textbf{工科口味}——\textbf{每个方法配一个例子}。

\textbf{缺点}：\textbf{没有数值方法}——\textbf{停在解析解}——\textbf{现实问题 90\% 没有解析解}。

\subsection{Strauss《Partial Differential Equations: An Introduction》}

\textbf{Strauss} 是\textbf{美国数学系学生的标准 PDE 入门}——\textbf{写得清晰、例子丰富}——\textbf{但偏纯数学}。

\subsection{LeVeque《Finite Difference Methods for ODEs and PDEs》}

\textbf{真正的 PDE 工程师工作手册}——\textbf{告诉你怎么把 PDE 写成代码、怎么分析稳定性、怎么估误差}——\textbf{是 ANSYS / COMSOL / 任何 CAE 软件背后的数学}。

\subsection{我的策略}

\textbf{本章不会重复梁昆淼的"特殊函数大全"}——\textbf{也不会重复 Strauss 的"严格证明"}——\textbf{本章把"三大类型 + 分离变量 + 有限差分"讲到}：

\begin{itemize}
\item 你看到一个 PDE——\textbf{能立刻判断它是哪一类}（看二阶导的"判别式"）
\item 你拿到一个最简单的 PDE（热、波、Laplace）——\textbf{能徒手用分离变量法求解析解}
\item 你拿到一个有点复杂的 PDE——\textbf{能写 20 行 Python 用有限差分跑出数值解}
\end{itemize}

\begin{tcolorbox}[storybox={\Large 小故事：Fourier 与三次被拒的论文}]
\textbf{1807 年}——\textbf{Joseph Fourier}（41 岁——\textbf{拿破仑远征埃及时的科学顾问}——刚回到法国）——\textbf{把一篇 234 页的论文}投给\textbf{法兰西学院}——题目：\textbf{《热在固体中的传播》}。

\medskip

论文里他做了三件\textbf{当时的人觉得不可思议}的事：

\begin{enumerate}
\item \textbf{推导出热传导方程}：$\partial u / \partial t = \alpha \partial^2 u / \partial x^2$
\item \textbf{用分离变量法求解}——\textbf{把解写成 $\sum_n B_n \sin(n\pi x/L) e^{-\alpha (n\pi/L)^2 t}$}
\item \textbf{断言}：\textbf{\textbf{任意函数都可以写成三角级数 $\sum a_n \cos(nx) + b_n \sin(nx)$}}
\end{enumerate}

\medskip

\textbf{评审委员会是当时数学界的"三巨头"}：\textbf{Lagrange、Laplace、Legendre}——\textbf{再加 Monge 和 Lacroix}。

\medskip

\textbf{Lagrange 反对最激烈}——\textbf{他在 50 年前研究过弦振动}——\textbf{当时和 Euler 已经吵过"任意函数能不能写成三角级数"}——\textbf{他坚信不能}——\textbf{尤其是有跳跃的函数（比如方波）}。

\medskip

\textbf{Laplace 比较温和}——\textbf{但他也觉得 Fourier 的"收敛性"没说清}。

\medskip

\textbf{结果}：\textbf{拒稿}。\textbf{Fourier 修改后再投——再拒——再投——又被搁置}。

\medskip

\textbf{1812 年}——\textbf{法兰西学院终于把数学大奖给了 Fourier}——\textbf{但评语里写}：\textbf{"虽然推导和应用都很有价值——但严格性不够——\textbf{不予正式出版}"}。

\medskip

\textbf{Fourier 气坏了}——\textbf{1822 年他干脆自己掏钱出书}——\textbf{《Théorie analytique de la chaleur》}——\textbf{《热的解析理论》}。

\medskip

\textbf{这本书后来成了现代物理和工程的奠基书之一}——\textbf{爱因斯坦读过、Maxwell 读过、Kelvin 读过}——\textbf{Kelvin 说}：\textbf{"Fourier 的定理——不仅是这个时代最美的数学结果——还是解决几乎所有自然哲学难题的不可缺少的工具。"}

\medskip

\textbf{今天}——\textbf{Fourier 级数 / Fourier 变换 是}：\textbf{JPEG 图像压缩、MP3 音频压缩、5G 通信、医学 CT、地震信号、量子力学、信号处理}——\textbf{无处不在}。

\medskip

\textbf{Lagrange 和 Laplace 当年的担心}——\textbf{严格性问题}——\textbf{直到 1966 年 Carleson 才彻底解决}（\textbf{L\textsuperscript{2} 函数的 Fourier 级数几乎处处收敛}）——\textbf{Carleson 拿了 2006 年的 Abel 奖}。

\medskip

\textbf{从 1807 年的"被拒稿"到 2006 年的"Abel 奖"——这件事整整跨了 200 年}。\textbf{科学不怕一时不被理解——只怕你自己中途放弃}。
\end{tcolorbox}

\section{直觉：三大类型 / Three Types}

\subsection{PDE 的"判别式"——和二次曲线一样}

\textbf{二阶线性 PDE 的一般形式}：
\begin{equation}
A u_{xx} + 2B u_{xy} + C u_{yy} + (\text{低阶项}) = 0
\end{equation}

\textbf{看判别式} $\Delta = B^2 - AC$：

\begin{itemize}
\item $\Delta < 0$——\textbf{椭圆型}（elliptic）——\textbf{对应"$x^2 + y^2 = 1$ 椭圆"}——例：\textbf{Laplace 方程} $u_{xx} + u_{yy} = 0$
\item $\Delta = 0$——\textbf{抛物型}（parabolic）——\textbf{对应"$y = x^2$ 抛物线"}——例：\textbf{热传导方程} $u_t = u_{xx}$（这里"$t$"扮演 $y$ 的角色，注意 $A=1, B=0, C=0$）
\item $\Delta > 0$——\textbf{双曲型}（hyperbolic）——\textbf{对应"$x^2 - y^2 = 1$ 双曲线"}——例：\textbf{波动方程} $u_{tt} = c^2 u_{xx}$
\end{itemize}

\textbf{这个分类不是数学家凭空造的}——\textbf{它和"二次曲线分类"用的是同一套代数}——\textbf{是因为 PDE 在变换下的"主部"恰好是一个二次型}。

\subsection{三种物理画面}

\textbf{记住这三句话——比记任何公式都管用}：

\begin{tcolorbox}[colback=blue!4,colframe=blue!50,boxrule=0.6pt,arc=2pt]
\textbf{抛物型（热）}：\textbf{"今天的温度——是昨天温度被邻居拉平之后的样子"}——\textbf{扩散、磨平、走向平衡——\textbf{时间不可逆}}。

\medskip

\textbf{双曲型（波）}：\textbf{"扰动以有限速度传播——左右都跑"}——\textbf{振动、回声、能量守恒——\textbf{时间可逆}}。

\medskip

\textbf{椭圆型（Laplace）}：\textbf{"每个点的值 = 邻居的平均值"}——\textbf{已达稳态——\textbf{没有时间}}。
\end{tcolorbox}

\begin{example}[一句话记住三大类型]
\textbf{把同一片湖比作三种物理}：
\begin{itemize}
\item \textbf{热}：你往湖里倒一杯热水——温度会慢慢"摊开、降低"——\textbf{抛物}
\item \textbf{波}：你往湖里扔一个石子——涟漪向外传播——\textbf{双曲}
\item \textbf{Laplace}：湖面静止——每个点都等于周围点的平均——\textbf{椭圆}
\end{itemize}
\end{example}

\section{§1 PDE 三大类型 / Three Types: Math}

\subsection{Laplace 方程（椭圆）}
\begin{equation}
\nabla^2 u = u_{xx} + u_{yy} + u_{zz} = 0
\end{equation}
\textbf{物理}：稳态温度场、静电势、稳态流场、引力势。

\textbf{关键性质}：\textbf{均值定理}——\textbf{$u$ 在球心的值 = $u$ 在球面上的平均}——\textbf{所以椭圆方程的解\textbf{极其光滑}}（你不能在 Laplace 解里出现尖角，否则均值定理不成立）。

\subsection{热传导方程（抛物）}
\begin{equation}
\frac{\partial u}{\partial t} = \alpha \nabla^2 u
\end{equation}
\textbf{物理}：温度随时间扩散；浓度扩散；金融期权定价（Black-Scholes 的核心是抛物型）。

\textbf{关键性质}：\textbf{最大值原理}——\textbf{$u$ 在内部任意点 $\le$ 它在初始时刻 + 边界上的最大值}——\textbf{所以"温度只会被磨平、不会自己冒尖"}。\textbf{扩散是不可逆的}——\textbf{把方程反向 $t \to -t$——它变成"反扩散"——数值上完全不稳定}（这就是为什么我们不能"从今天反推昨天的温度分布"）。

\subsection{波动方程（双曲）}
\begin{equation}
\frac{\partial^2 u}{\partial t^2} = c^2 \nabla^2 u
\end{equation}
\textbf{物理}：弦振动、声波、电磁波、地震波。

\textbf{关键性质}：\textbf{有限传播速度 $c$}——\textbf{$t$ 时刻能影响到 $u(x_0, t)$ 的"初始数据"只在以 $x_0$ 为中心、半径 $ct$ 的范围内}（"光锥"）——\textbf{这和热传导的"瞬时无穷远"完全相反}。\textbf{能量守恒}——\textbf{$E = \int (\tfrac{1}{2} u_t^2 + \tfrac{1}{2} c^2 u_x^2) dx$ 不变}。

\section{§2 分离变量法 / Separation of Variables}

\textbf{这是 Fourier 1807 年发明的核心技术}——\textbf{把一个 PDE 拆成两个 ODE}。

\subsection{基本想法}

\textbf{猜}：解能写成\textbf{一个只依赖 $x$ 的函数 $\times$ 一个只依赖 $t$ 的函数}：
\[
u(x, t) = X(x) \, T(t)
\]

\textbf{代入热传导方程 $u_t = \alpha u_{xx}$}：
\[
X(x) T'(t) = \alpha X''(x) T(t)
\]
两边除以 $\alpha X T$：
\[
\frac{T'(t)}{\alpha T(t)} = \frac{X''(x)}{X(x)} = -\lambda
\]
\textbf{左边只依赖 $t$，右边只依赖 $x$——\textbf{它们相等的唯一可能是}：\textbf{两边都等于同一个常数 $-\lambda$}}（习惯写负号——因为我们期待振荡解）。

\subsection{两个 ODE}

\begin{align}
X''(x) + \lambda X(x) &= 0 \quad \text{（空间方程）}\\
T'(t) + \alpha \lambda T(t) &= 0 \quad \text{（时间方程）}
\end{align}

\textbf{加边界条件}（例如 $X(0) = X(L) = 0$——两端温度恒为 0）——\textbf{空间方程变成一个本征值问题}：
\[
X_n(x) = \sin\!\left(\frac{n\pi x}{L}\right), \quad \lambda_n = \left(\frac{n\pi}{L}\right)^2, \quad n = 1, 2, 3, \dots
\]

\textbf{时间方程的解}：
\[
T_n(t) = e^{-\alpha \lambda_n t}
\]

\subsection{叠加 = Fourier 级数}

\textbf{每个 $n$ 都给出一个"基本模式"} $u_n(x,t) = \sin(n\pi x/L) \, e^{-\alpha(n\pi/L)^2 t}$。

\textbf{线性 PDE 满足叠加原理}——\textbf{所以一般解是它们的和}：
\begin{equation}
u(x, t) = \sum_{n=1}^{\infty} B_n \sin\!\left(\frac{n\pi x}{L}\right) e^{-\alpha (n\pi/L)^2 t}
\end{equation}

\textbf{系数 $B_n$ 怎么定}？——\textbf{从初始条件 $u(x, 0) = f(x)$}——\textbf{这就是\textbf{把 $f(x)$ 展开成 Fourier 正弦级数}}：
\begin{equation}
B_n = \frac{2}{L} \int_0^L f(x) \sin\!\left(\frac{n\pi x}{L}\right) dx
\end{equation}

\textbf{这就是 Fourier 一辈子被人骂"不严格"的那个公式}——\textbf{今天它是所有信号处理和数据压缩的根}。

\begin{example}[为什么"高频模式衰减得最快"]
\textbf{看那个指数因子}：\textbf{$e^{-\alpha (n\pi/L)^2 t}$}——\textbf{$n$ 越大——指数衰减越快}。

\textbf{物理意义}：\textbf{初始温度分布里的"尖锐细节"（高频成分）}——\textbf{很快被磨掉}；\textbf{"大块的趋势"（低频成分）}——\textbf{保留得久}。

\textbf{这就是为什么"热扩散 = 低通滤波"——和 Gaussian blur 在图像里做的是同一件事}！
\end{example}

\section{§3 热传导方程 / The Heat Equation}

\subsection{推导：从能量守恒开始}

\textbf{取一小段 $[x, x+\Delta x]$}——\textbf{它的内能 $\propto u(x, t) \cdot \Delta x$}。

\textbf{Fourier 定律}：\textbf{热流密度 $q = -k \, \partial u / \partial x$}——\textbf{热从高温流向低温——比例常数 $k$ 是热导率}。

\textbf{能量守恒}：\textbf{这一小段内能的变化率 = 左边流入 - 右边流出}：
\[
\rho c \, \Delta x \, \frac{\partial u}{\partial t} = q(x) - q(x+\Delta x) \approx -\frac{\partial q}{\partial x} \Delta x
\]
代入 Fourier 定律：
\[
\rho c \, \frac{\partial u}{\partial t} = k \, \frac{\partial^2 u}{\partial x^2}
\]
\textbf{除以 $\rho c$——记 $\alpha = k/(\rho c)$ 为\textbf{热扩散率}}：
\begin{equation}
\boxed{\, \frac{\partial u}{\partial t} = \alpha \frac{\partial^2 u}{\partial x^2} \,}
\end{equation}

\subsection{为什么"二阶导 = 邻居平均 - 自己"}

\textbf{把 $u(x \pm \Delta x)$ 做 Taylor 展开}：
\[
u(x+\Delta x) + u(x-\Delta x) = 2 u(x) + u_{xx}(x) \, \Delta x^2 + O(\Delta x^4)
\]
\textbf{整理}：
\[
u_{xx}(x) \approx \frac{1}{\Delta x^2} \left[ \underbrace{\frac{u(x+\Delta x) + u(x-\Delta x)}{2}}_{\text{邻居的平均}} - u(x) \right] \cdot 2
\]

\textbf{所以热传导方程的本质是}：\textbf{"温度变化率 $\propto$ 邻居平均 - 自己"}——\textbf{比邻居冷——升温；比邻居热——降温}。\textbf{这就是扩散}。

\section{§4 波动方程 / The Wave Equation}

\subsection{推导：从牛顿第二定律开始}

\textbf{弦上的一小段 $[x, x+\Delta x]$}——\textbf{质量 $\rho \Delta x$}——\textbf{受张力 $T$ 作用}。

\textbf{弦稍微倾斜时}——\textbf{左端的张力沿弦切向、右端也沿弦切向}——\textbf{它们的合力的"竖直分量" $\approx T(\partial u/\partial x)|_{x+\Delta x} - T(\partial u/\partial x)|_x = T u_{xx} \Delta x$}。

\textbf{牛顿第二定律}：
\[
\rho \Delta x \cdot \frac{\partial^2 u}{\partial t^2} = T \, u_{xx} \Delta x
\]
\textbf{记 $c^2 = T/\rho$}：
\begin{equation}
\boxed{\, \frac{\partial^2 u}{\partial t^2} = c^2 \, \frac{\partial^2 u}{\partial x^2} \,}
\end{equation}
\textbf{$c$ 是弦上波速}——\textbf{张力越大、密度越小——波越快}。

\subsection{d'Alembert 公式（无穷长弦）}

\textbf{对于无边界问题}——\textbf{解可以写成"两个反向行波"的叠加}：
\begin{equation}
u(x, t) = F(x - ct) + G(x + ct)
\end{equation}
\textbf{$F$ 是一个不变形地向右跑的波}——\textbf{$G$ 是向左跑的波}——\textbf{两者完全独立}。\textbf{这就是双曲方程"有限传播速度"的直接体现}。

\subsection{有界弦的振动模式（与吉他弦相同）}

\textbf{对于 $[0, L]$ 上的弦，两端固定}——\textbf{分离变量给出}：
\begin{equation}
u(x, t) = \sum_{n=1}^{\infty} \sin\!\left(\frac{n\pi x}{L}\right) \left[ A_n \cos\!\left(\frac{n\pi c t}{L}\right) + B_n \sin\!\left(\frac{n\pi c t}{L}\right) \right]
\end{equation}

\textbf{每一个 $n$ 都是一个"模式"}——\textbf{第 $n$ 模式的频率是 $f_n = nc/(2L)$}——\textbf{$n=1$ 是基频}——\textbf{$n=2, 3, \dots$ 是\textbf{泛音}}——\textbf{这就是吉他弦发声的物理}。\textbf{乐器的音色 = 各次泛音的相对强度（即 $B_n$ 的分布）}。

\section{§5 Laplace / Poisson 方程 / Elliptic}

\subsection{Laplace 方程}
\[
\nabla^2 u = 0
\]
\textbf{含义}：\textbf{稳态——没有时间——没有源}。

\subsection{Poisson 方程}
\[
\nabla^2 u = -f
\]
\textbf{含义}：\textbf{稳态——有源 $f$}。\textbf{$f$ 可以是电荷密度（电场）、热源密度（稳态传热）、流体源（不可压流）}。

\subsection{物理画面：肥皂膜}

\textbf{Laplace 方程的最直观画面}：\textbf{把一个金属丝弯成 3D 曲线——蘸肥皂水——形成的肥皂膜——\textbf{它在每一点的高度 $u(x, y)$ 满足 $\nabla^2 u = 0$}}。

\textbf{为什么}？——\textbf{肥皂膜的表面张力让它"局部尽可能平"——也就是}\textbf{"每一点都是邻居的平均"——这正是 Laplace 方程的均值性质}。

\subsection{边界条件的两种类型}

\begin{itemize}
\item \textbf{Dirichlet 条件}：\textbf{在边界上指定 $u$ 的值}（例：边界温度恒定）
\item \textbf{Neumann 条件}：\textbf{在边界上指定 $\partial u / \partial n$}（例：边界绝热——法向导热 = 0）
\end{itemize}

\textbf{椭圆方程的强大之处}：\textbf{边界条件一旦给定——内部解就唯一确定}——\textbf{这就是"稳态"的含义}。

\section{§6 数值方法 / Numerical Methods}

\textbf{现实工程问题——99\% 没有解析解}——\textbf{所以你必须会数值方法}。\textbf{两大主流}：\textbf{有限差分（FDM）}和\textbf{有限元（FEM）}。

\subsection{有限差分法 / Finite Difference Method}

\textbf{核心思想}：\textbf{把所有偏导用差分代替}。

\textbf{常用差分公式}：
\begin{align}
u_t &\approx \frac{u^{n+1}_i - u^n_i}{\Delta t} \quad \text{（前向）}\\
u_{xx} &\approx \frac{u^n_{i+1} - 2 u^n_i + u^n_{i-1}}{\Delta x^2} \quad \text{（中心二阶）}
\end{align}

\textbf{把它们代进热传导方程}——\textbf{得显式更新公式}：
\begin{equation}
u^{n+1}_i = u^n_i + r (u^n_{i+1} - 2 u^n_i + u^n_{i-1}), \quad r = \frac{\alpha \Delta t}{\Delta x^2}
\end{equation}
\textbf{这就是}\textbf{显式格式}——\textbf{用今天的所有点——一步算出明天的每一点}。

\textbf{稳定性条件}（\textbf{CFL 条件}）：\textbf{$r \le 1/2$}——\textbf{否则数值解会爆炸}——\textbf{这告诉你 $\Delta t$ 不能太大——必须配合 $\Delta x$}。

\subsection{隐式格式（Crank-Nicolson）}

\textbf{把右边的二阶导取"今天和明天的平均"}：
\[
\frac{u^{n+1}_i - u^n_i}{\Delta t} = \frac{\alpha}{2}\left[ (u_{xx})^n_i + (u_{xx})^{n+1}_i \right]
\]
\textbf{结果}：\textbf{每一步要解一个三对角线性方程组}——\textbf{慢一点——但\textbf{无条件稳定}}——\textbf{$\Delta t$ 可以取任意大}。\textbf{这是工业 CAE 软件的默认时间积分器}。

\subsection{有限元法 / Finite Element Method（一段话概述）}

\textbf{核心思想}：\textbf{把求解区域划成小三角形（"单元"）}——\textbf{在每个单元里把 $u$ 假设成简单的多项式（线性或二次）}——\textbf{用\textbf{弱形式 / 变分原理}}把 PDE 转化成一个线性方程组 $K \bm{u} = \bm{f}$ ——\textbf{然后求解大型稀疏线性系统}。

\textbf{为什么 FEM 比 FDM 强}：

\begin{itemize}
\item \textbf{复杂几何}（任意形状的硅片、机翼、人体）——\textbf{FDM 只在规则网格上方便——FEM 在任意三角剖分上都行}
\item \textbf{自然处理边界条件}——\textbf{弱形式自动把 Neumann 条件吸收进来}
\item \textbf{自适应加密}——\textbf{应力集中处自动用更密的单元}
\end{itemize}

\textbf{ANSYS、COMSOL、Abaqus、FEniCS——它们底层都是 FEM}。

\section{Aha 例子：手机芯片散热 / Smartphone Chip Heat Dissipation}

\begin{tcolorbox}[ahabox={\Large Aha 例子：手机芯片散热}]
\textbf{一颗 5G 手机的 SoC（系统级芯片）}——\textbf{尺寸约 10 mm $\times$ 10 mm $\times$ 1 mm}——\textbf{满负荷时功耗 5 W}——\textbf{热流密度 $\approx 50$ W/cm\textsuperscript{2}}（比一块家用电熨斗的发热面还高 5 倍）。

\medskip

\textbf{芯片不能超过 100°C}——\textbf{否则会降频甚至烧毁}——\textbf{所以散热设计是手机性能的瓶颈之一}。

\medskip

\textbf{工程师怎么"算热"}？——\textbf{建模成一个 2D 热传导问题}：
\[
\rho c \frac{\partial T}{\partial t} = k \nabla^2 T + Q(x, y)
\]
\begin{itemize}
\item \textbf{$Q(x, y)$}：芯片上 CPU 核 / GPU 核 / 基带 等\textbf{发热源的位置和功率}
\item \textbf{边界条件}：\textbf{芯片表面通过石墨片散热 → Neumann 条件 $-k \partial T / \partial n = h(T - T_\infty)$（对流冷却）}
\item \textbf{初始条件}：\textbf{$T(x,y,0) = T_\infty$（室温 25°C）}
\end{itemize}

\medskip

\textbf{我们用最简单的 1D 模型}——\textbf{把芯片简化成一根长度 $L = 10$ mm 的硅棒——一端贴着 CPU 核（发热源 $Q_0$）、另一端贴着散热片（恒温 25°C）}——\textbf{用显式有限差分跑一下}：

\begin{itemize}
\item \textbf{看到温度分布如何\textbf{从均匀的 25°C 逐渐\textbf{升起}}}
\item \textbf{看到经过几十毫秒——\textbf{稳态形成}}
\item \textbf{看到稳态下\textbf{热源端最高温度}——这是芯片设计师最关心的数字}
\end{itemize}
\end{tcolorbox}

\subsection{Python 实现（仅 20 行核心代码）}

\begin{lstlisting}[language=Python]
import numpy as np
import matplotlib.pyplot as plt

# ---- 物理参数 ----
L = 1e-2            # 芯片长度 10 mm
alpha = 8.8e-5      # 硅的热扩散率 m^2/s
T0 = 25.0           # 室温
Q = 5e7             # 体热源密度 W/m^3 (CPU 核局部)

# ---- 离散参数 ----
N = 100             # 空间网格点数
dx = L / (N - 1)
dt = 0.4 * dx**2 / alpha   # 满足 CFL: r = alpha*dt/dx^2 = 0.4
n_steps = 4000

# ---- 初始 / 边界 ----
T = np.ones(N) * T0
T[-1] = T0          # 散热片端恒温 (Dirichlet)

# ---- 时间推进 (显式有限差分) ----
r = alpha * dt / dx**2
for n in range(n_steps):
    T_new = T.copy()
    T_new[1:-1] = T[1:-1] + r * (T[2:] - 2*T[1:-1] + T[:-2])
    # 内部 30% 区域是 CPU 核 - 持续加热
    src = (Q / 2330 / 700) * dt   # Q / (rho*cp) * dt
    T_new[:int(0.3*N)] += src
    T_new[-1] = T0    # 维持边界
    T = T_new

print(f"稳态最高温: {T.max():.1f} °C")
plt.plot(np.linspace(0, L*1000, N), T)
plt.xlabel('Position (mm)'); plt.ylabel('Temperature (°C)')
plt.title('Smartphone Chip 1D Heat Profile')
plt.savefig('chip_heat.png')
\end{lstlisting}

\textbf{运行结果}（典型输出）：\textbf{稳态最高温 $\approx 95$°C}——\textbf{紧贴 100°C 的红线}——\textbf{这就是为什么你打游戏久了手机会"降频"}。

\textbf{工程师的工作}——\textbf{改 $Q$ 的分布}（让发热源更分散）、\textbf{改 $\alpha$}（换更高导热的封装材料，如石墨烯）、\textbf{改边界}（加均热板 VC，增大有效散热面积）——\textbf{每改一次——重跑仿真——找到最优方案}。

\textbf{你刚才写的 20 行代码——和 ANSYS Icepak / Flotherm 干的是同一件事——只是它们做 3D + 复杂几何 + 耦合多物理场}。

\section{现代视角 / Modern Perspective}

\subsection{PDE 在 21 世纪的位置}

\begin{itemize}
\item \textbf{芯片设计}：\textbf{热传导 + 电场 Laplace + 应力分析}——\textbf{每代工艺验证都靠 PDE 仿真}
\item \textbf{电磁仿真}：\textbf{天线设计、5G/6G 信道——Maxwell 方程组 = 一组双曲 PDE}
\item \textbf{流体力学}：\textbf{飞机、汽车、火箭——Navier-Stokes（混合型 PDE）}——\textbf{是\textbf{千禧年七大难题之一}}
\item \textbf{金融工程}：\textbf{Black-Scholes 期权定价方程——是一个抛物型 PDE}
\item \textbf{医学影像}：\textbf{CT 重建、MRI——背后是 Radon 变换 + PDE}
\item \textbf{气候模拟}：\textbf{全球气候模型——耦合大气 + 海洋 PDE，运行在超级计算机上}
\end{itemize}

\subsection{PDE 与机器学习}

\textbf{2018 年以来一个迅速崛起的领域}：\textbf{\textbf{Physics-Informed Neural Networks（PINN）}}——\textbf{用神经网络逼近 PDE 的解}——\textbf{把"残差 $\| u_t - \alpha u_{xx} \|^2$"作为损失函数训练}——\textbf{$\Longrightarrow$ 可以处理\textbf{逆问题}（已知部分观测——反推参数 $\alpha$）和\textbf{高维 PDE}（$n>10$ 维度下传统 FEM 完全无能为力）}。

\textbf{\textbf{Neural Operators}（Fourier Neural Operator, DeepONet）}——\textbf{学的不是某一个 PDE 的解——而是"从初始条件到解"的算子}——\textbf{训练好之后用 $10^{-4}$ 秒得到解——比传统 FEM 快 $10^4$--$10^6$ 倍}。\textbf{这是\textbf{科学计算+ML} 最有前途的方向之一}。

\subsection{扩散模型——为什么名字里有"扩散"}

\textbf{2020 年起 AI 绘画领域爆火的\textbf{扩散模型}}（Stable Diffusion, DALL-E, Midjourney 的核心）——\textbf{它的数学基础\textbf{完全是 PDE}}——\textbf{具体是 Fokker-Planck 方程（一个广义的热传导方程）}。

\textbf{想法}：\textbf{前向过程把图像"加噪声"——本质是让像素概率分布按热传导方程演化——最终变成纯 Gaussian}。\textbf{反向过程用神经网络学习"反扩散"——从噪声里恢复图像}。

\textbf{你今天学的热传导方程——就是 OpenAI / Stability AI 的看家本领之一}。

\section{代码与思考 / Code and Exercises}

\subsection{配套模块 \texttt{pde.py}}

\begin{lstlisting}[language=Python]
"""
pde.py - 偏微分方程数值实验
"""
import numpy as np
import matplotlib.pyplot as plt

def demo_classify():
    """打印 PDE 三大类型 + 判别式"""
    cases = [
        ("Laplace  u_xx + u_yy = 0", 1, 0, 1),
        ("Heat     u_t = u_xx     ", 1, 0, 0),
        ("Wave     u_tt = c^2 u_xx", 1, 0, -1),
    ]
    for name, A, B, C in cases:
        disc = B**2 - A*C
        kind = "elliptic" if disc<0 else ("parabolic" if disc==0 else "hyperbolic")
        print(f"{name}   disc={disc:+d}  {kind}")

def demo_heat_1d(N=100, n_steps=2000):
    """1D 热传导 - 显式有限差分"""
    L, alpha = 1.0, 0.01
    dx = L/(N-1); dt = 0.4*dx**2/alpha
    x = np.linspace(0, L, N)
    u = np.exp(-200*(x-0.5)**2)   # 初始尖峰
    u[0] = u[-1] = 0
    r = alpha*dt/dx**2
    for _ in range(n_steps):
        u[1:-1] += r*(u[2:] - 2*u[1:-1] + u[:-2])
    plt.plot(x, u); plt.title('Heat after diffusion')
    plt.savefig('heat_1d.png')

def demo_wave_1d(N=200, n_steps=600):
    """1D 弦振动 - 显式有限差分 (双时间层)"""
    L, c = 1.0, 1.0
    dx = L/(N-1); dt = 0.9*dx/c   # CFL: c*dt/dx <= 1
    x = np.linspace(0, L, N)
    u_prev = np.sin(np.pi*x)      # 初始: 基模
    u_curr = u_prev.copy()
    R = (c*dt/dx)**2
    snapshots = []
    for n in range(n_steps):
        u_next = np.zeros_like(u_curr)
        u_next[1:-1] = (2*u_curr[1:-1] - u_prev[1:-1]
                        + R*(u_curr[2:] - 2*u_curr[1:-1] + u_curr[:-2]))
        u_prev, u_curr = u_curr, u_next
        if n % 100 == 0:
            snapshots.append(u_curr.copy())
    for s in snapshots:
        plt.plot(x, s)
    plt.title('Standing wave - fundamental mode')
    plt.savefig('wave_1d.png')

def demo_laplace_2d(N=50, n_iter=5000):
    """2D Laplace 方程 - Jacobi 迭代"""
    u = np.zeros((N, N))
    u[0, :] = 100   # 顶部边界 100 度
    for _ in range(n_iter):
        u[1:-1,1:-1] = 0.25*(u[2:,1:-1] + u[:-2,1:-1]
                              + u[1:-1,2:] + u[1:-1,:-2])
    plt.contourf(u, 20, cmap='hot'); plt.colorbar()
    plt.title('Laplace solution: top=100, others=0')
    plt.savefig('laplace_2d.png')

def demo_chip_heat():
    """本章 Aha: 手机芯片 1D 散热"""
    L, alpha, T0, Q = 1e-2, 8.8e-5, 25.0, 5e7
    N = 100; dx = L/(N-1); dt = 0.4*dx**2/alpha
    T = np.ones(N)*T0; T[-1] = T0
    r = alpha*dt/dx**2
    for _ in range(4000):
        T_new = T.copy()
        T_new[1:-1] = T[1:-1] + r*(T[2:] - 2*T[1:-1] + T[:-2])
        T_new[:int(0.3*N)] += (Q/2330/700)*dt
        T_new[-1] = T0
        T = T_new
    print(f"Max chip temperature: {T.max():.1f} C")

def demo_fourier_series(N_terms=20):
    """方波的 Fourier 级数 - 看 Gibbs 现象"""
    x = np.linspace(-np.pi, np.pi, 1000)
    s = np.zeros_like(x)
    for k in range(1, N_terms+1, 2):
        s += (4/np.pi) * np.sin(k*x)/k
    plt.plot(x, s); plt.title(f'Square wave: {N_terms} terms')
    plt.savefig('fourier_square.png')
\end{lstlisting}

\subsection{Aha 重点：1D 弦振动模式}

\textbf{把 \texttt{demo\_wave\_1d()} 跑起来——你会亲眼看到\textbf{驻波模式}}：

\begin{itemize}
\item \textbf{初始时}：\textbf{弦呈现 $\sin(\pi x / L)$ 形状（基模）}
\item \textbf{时间推进}：\textbf{弦"上下振动"——振幅按 $\cos(\pi c t / L)$ 周期变化}
\item \textbf{两端始终为 0}——\textbf{中点振幅最大}
\end{itemize}

\textbf{改初始条件为 $\sin(2\pi x/L)$ 或 $\sin(3\pi x/L)$}——\textbf{你看到\textbf{二次、三次模式}}——\textbf{这就是吉他、小提琴、钢琴的"泛音"由来}。

\textbf{改初始条件为\textbf{两种模式的叠加}}——\textbf{你会听到（如果转成声音）\textbf{和弦}}——\textbf{这就是音乐的物理}。

\textbf{Fourier 1822 年用一辈子争来的"任意函数 = 三角级数"——你今天 30 行 Python 就能实验它}——\textbf{改 N\_terms——看方波被逼近的过程——\textbf{你会看到不肯消失的"Gibbs 跳动"}——这正是 Lagrange 当年最反对的"不收敛"现象}。

\subsection{思考题 / Exercises}

\begin{enumerate}
\item \textbf{[直觉]} 用一句话区分"抛物"和"双曲"——不许用"判别式"三个字。
\item \textbf{[判别]} 判断 $u_{tt} - 2 u_{xt} + u_{xx} + 3 u_t = 0$ 的类型。
\item \textbf{[计算]} 用分离变量法求 $u_t = u_{xx}$，$u(0,t)=u(\pi,t)=0$，$u(x,0)=\sin(3x)$ 的解。
\item \textbf{[代码]} 在 \texttt{demo\_heat\_1d()} 里把 $r = \alpha \Delta t / \Delta x^2$ 改成 $0.6$——\textbf{看会发生什么}（CFL 违反——数值爆炸）。
\item \textbf{[代码]} 写一个 \texttt{demo\_drum\_2d()}——\textbf{解 2D 波动方程}——\textbf{看圆形鼓面的振动模式}（结果会涉及 Bessel 函数）。
\item \textbf{[物理]} 为什么"反扩散"在数学上不稳定？——\textbf{用最大值原理给个直观解释}。
\item \textbf{[ML]} 查一下\textbf{扩散模型}（Diffusion Model）和热传导方程的联系——\textbf{为什么 Stable Diffusion 的"噪声调度"本质是一个 PDE？}
\item \textbf{[工程]} 用 \texttt{demo\_chip\_heat()}——\textbf{加大 $\alpha$（模拟换石墨烯封装）}——\textbf{看最高温度降了多少}——\textbf{这是工业散热设计的真实问题}。
\item \textbf{[历史]} 查一下 Carleson 1966 年的工作——\textbf{它解决了 Fourier 200 年前留下的什么严格性问题？}
\end{enumerate}

\section{本章小结 / Chapter Summary}

\begin{tcolorbox}[essbox={本章核心}]
\begin{itemize}
\item \textbf{PDE} 描述\textbf{场}（温度、电、压力、波）\textbf{在空间和时间上的演化}
\item \textbf{三大类型}（看判别式 $\Delta = B^2 - AC$）：
   \begin{itemize}
   \item \textbf{椭圆}（$\Delta < 0$）—— Laplace / Poisson —— \textbf{稳态/平衡}
   \item \textbf{抛物}（$\Delta = 0$）—— 热传导 —— \textbf{扩散/不可逆}
   \item \textbf{双曲}（$\Delta > 0$）—— 波动 —— \textbf{振动/有限传播}
   \end{itemize}
\item \textbf{分离变量法}：\textbf{$u = X(x)T(t)$} → \textbf{两个 ODE} → \textbf{叠加成 Fourier 级数}
\item \textbf{二阶导 $\partial^2 u / \partial x^2 \approx$ 邻居平均 - 自己}——\textbf{解释了"扩散"的本质}
\item \textbf{有限差分}：\textbf{把偏导换成差分——20 行 Python 就能跑}——\textbf{注意 CFL 稳定性 $r \le 1/2$（热）或 $c\Delta t/\Delta x \le 1$（波）}
\item \textbf{有限元}：\textbf{任意几何 + 弱形式}——\textbf{ANSYS / COMSOL 的内核}
\item \textbf{现代应用}：\textbf{芯片散热、电磁仿真、流体、金融、医学影像、PINN、扩散模型}
\end{itemize}
\end{tcolorbox}

\textbf{下一章——复变函数与积分变换}——\textbf{我们让 Fourier 级数推广到 Fourier \textbf{变换}}——\textbf{进入信号处理的世界}。

\clearpage


\part{线性代数 / Linear Algebra}
% =====================================================================
%  第 6 章 · 线代入门 / Linear Algebra I
%  矢量空间、矩阵 = 线性变换、行列式、矩阵的逆、Gauss 消元
%  小故事：Cayley 与"矩阵的抽象"（1858）
%  Aha：3D 游戏旋转
% =====================================================================

\chapter{线代入门 / Linear Algebra I}
\label{ch:linalg1}

\begin{quote}\itshape
1858 年，英国剑桥。一位 37 岁的律师——白天打官司——晚上写数学。\\
他叫 \textbf{Arthur Cayley}。\\
那一年，他写下一篇 17 页的论文：\\
\textbf{《A Memoir on the Theory of Matrices》}。\\
人类历史上第一次——把"矩阵"当成一个\textbf{独立的对象}——\textbf{而不是方程组的附属品}。\\[6pt]
今天你打开 PyTorch——敲下 \texttt{torch.matmul(A, B)}——背后是这位律师 168 年前的一念之差。
\end{quote}

\section{写在前面 / Preface}

\textbf{大一下学期我学线性代数}——\textbf{同济版的浅蓝色课本}——\textbf{第一章就是"行列式"}——\textbf{第二章是"矩阵"}——\textbf{第三章是"向量"}。

\textbf{这个顺序——是我学过最反直觉的顺序之一}。

为什么？因为它\textbf{把答案放在了问题前面}——你还没明白\textbf{"矩阵是干什么的"}——就被要求\textbf{算 4 阶行列式}——\textbf{展开 + 化三角 + 按行按列}——\textbf{三套办法换着练 20 道题}。

\textbf{结果是}：我考了 88 分——\textbf{但我大三做信号处理时}——\textbf{看到 MATLAB 里的 \texttt{A \textbackslash b}——还是发愣}。

到了\textbf{研究生}——我做\textbf{机器人运动学}——\textbf{第一次需要把一个 3D 点从机器人坐标系转到世界坐标系}——\textbf{老板甩给我一个 4×4 矩阵}——说"乘上去就行"——\textbf{我懂了"矩阵是变换"的那一刻}——\textbf{比我所有大一题目加起来都管用}。

\textbf{从那以后我学线代——再也没有从行列式开始过}。

我的顺序是这样的：

\begin{itemize}
\item \textbf{第一步}：\textbf{什么是向量}——\textbf{从 2D 箭头开始}
\item \textbf{第二步}：\textbf{什么是矩阵}——\textbf{它是一个变换}——\textbf{不是一堆数}
\item \textbf{第三步}：\textbf{矩阵怎么乘}——\textbf{为什么这么乘}
\item \textbf{第四步}：\textbf{行列式是什么}——\textbf{它告诉我们"变换把面积缩放了多少倍"}
\item \textbf{第五步}：\textbf{解方程组}——\textbf{Gauss 消元 + 矩阵的逆}
\end{itemize}

\textbf{这一章——就是这个顺序}。

\textbf{这是本书第六章——也是"线性代数三部曲"的第一章}。后面两章是\textbf{第 7 章（线代核心）}讲\textbf{特征值、SVD、对角化}——\textbf{第 8 章（线代应用）}讲\textbf{PCA、PageRank、神经网络}。\textbf{三章合起来——就是整个工科+ML 用得到的线代}。

\textbf{如果整本书有一个"任督二脉"的核心}——\textbf{就是这三章}。

\section{一句话本质 / The Essence in One Sentence}

\begin{tcolorbox}[essbox={一句话本质}]
\textbf{中文}：\textbf{矩阵不是一堆数——矩阵是一个线性变换——它把向量从一个空间送到另一个空间}。\\[6pt]
\textbf{English}: A matrix is not a grid of numbers; it is a linear transformation that maps vectors from one space to another.
\end{tcolorbox}

记住这一句——\textbf{这一章已经讲完了 70\%}。剩下的 30\% 是\textbf{怎么算、怎么用}。

\section{教材里你看到的 / What the Textbook Tells You}

\subsection{同济《线性代数》}

\textbf{同济大学应用数学系编的《线性代数》（第六版）}——\textbf{大一下学期的国民教材}——\textbf{薄薄的一本——200 页左右}。

它的章节顺序：

\begin{enumerate}
\item 行列式
\item 矩阵及其运算
\item 矩阵的初等变换与线性方程组
\item 向量组的线性相关性
\item 相似矩阵及二次型
\end{enumerate}

\textbf{优点}：\textbf{薄、便宜、考研覆盖}——\textbf{中国考研 80\% 的线代题它都罩得住}。

\textbf{缺点}（致命）：

\begin{enumerate}
\item \textbf{从行列式开始}——\textbf{用一个高度抽象的"求和符号定义"砸学生的脸}——\textbf{学生根本不知道行列式有什么用}
\item \textbf{"矩阵"和"线性变换"完全脱节}——\textbf{矩阵讲了 50 页——"变换"两个字出现不到 5 次}
\item \textbf{应用部分几乎没有}——\textbf{所有题目都是"求行列式、求秩、解方程"}——\textbf{没有一张图、没有一行代码}
\end{enumerate}

\subsection{Gilbert Strang《Introduction to Linear Algebra》}

\textbf{MIT 的 Gilbert Strang}——\textbf{线性代数教科书届的"教皇"}。

他的《Introduction to Linear Algebra》——\textbf{开篇就是"列空间"和"向量的线性组合"}——\textbf{把矩阵当成"列向量的容器"}——\textbf{把矩阵乘法解释成"列的线性组合"}。

\textbf{他在 MIT 18.06 的课程视频}——\textbf{B 站、YouTube 都有}——\textbf{是全世界线代学生的圣经}。

\textbf{Strang 的优点}：\textbf{从几何讲起、图多、贯穿线性变换思想}。

\textbf{Strang 的缺点}：\textbf{英语原版稍贵}——\textbf{排版偏美式繁复}——\textbf{部分章节顺序对中国学生有点跳跃}。

\subsection{3Blue1Brown《Essence of Linear Algebra》}

\textbf{这不是一本书}——\textbf{是 YouTube 上的一个视频系列}——\textbf{Grant Sanderson 制作}——\textbf{共 15 集——每集 10 分钟左右}。

\textbf{这是 21 世纪线性代数教学最伟大的成果}——\textbf{没有之一}。

\textbf{它的核心}：\textbf{每一个矩阵的概念——都用动画表示成"空间的变形"}。\textbf{你看一遍——就懂了矩阵 = 线性变换这件事——比读 300 页课本还有用}。

\textbf{中文字幕版在 B 站上有几百万播放量}——\textbf{凡是没看过的——立刻去看}。

\subsection{我的策略}

\textbf{本章不和同济硬碰硬}——\textbf{它的计算技巧——你刷它的题就行}。

\textbf{本章只做三件事}：

\begin{enumerate}
\item \textbf{把"矩阵 = 线性变换"讲到刻进骨子里}——\textbf{这是 3Blue1Brown 和 Strang 共同的核心}
\item \textbf{把行列式的几何意义讲清楚}——\textbf{它是"变换的面积/体积缩放因子"}——\textbf{不是"按某行展开的求和"}
\item \textbf{把 Gauss 消元和矩阵的逆讲到能解 3D 游戏的旋转问题}
\end{enumerate}

\begin{tcolorbox}[storybox={\Large 小故事：Cayley 与"矩阵的抽象"（1858）}]
\textbf{Arthur Cayley}（1821--1895）——\textbf{英国数学家}——\textbf{也是律师}。

\medskip

\textbf{他的本职工作是律师}——\textbf{在伦敦执业 14 年}——\textbf{白天打官司——晚上写数学论文}——\textbf{这 14 年里他写了 250 多篇数学论文}（你没看错——250 篇）。

\medskip

\textbf{1858 年——他 37 岁那年——发表了《A Memoir on the Theory of Matrices》}——\textbf{这是人类历史上第一篇把"矩阵"作为独立研究对象的论文}。

\medskip

在他之前——人们写出一堆数排成方阵——只是为了\textbf{方便写线性方程组的系数}——\textbf{矩阵本身没有名字}——\textbf{更没有运算}。

\medskip

\textbf{Cayley 做了三件革命性的事}：

\begin{enumerate}
\item \textbf{给"矩阵"起了名字}（matrix——来自拉丁文"子宫、母体"——\textbf{矩阵孕育了向量}）
\item \textbf{定义了矩阵的加法、乘法、逆}——\textbf{特别是乘法——他第一次写下 $AB$ 不等于 $BA$}
\item \textbf{发现了 Cayley-Hamilton 定理}——\textbf{每个矩阵都满足它自己的特征方程}（这是第 7 章的内容）
\end{enumerate}

\medskip

\textbf{Cayley 那篇 17 页论文的核心}——\textbf{"我们应该把矩阵当成一种新的数——而不是数的集合"}。

\medskip

\textbf{这是数学史上著名的"抽象一跃"}——\textbf{把一个具体对象抽象成一类东西}——\textbf{从此整个代数学进入了"抽象代数"的时代}。

\medskip

\textbf{50 年后}——\textbf{Heisenberg} 在量子力学里用矩阵描述粒子状态——\textbf{他承认自己"重新发明了 Cayley 的矩阵"}——\textbf{因为他不知道 Cayley}。

\medskip

\textbf{100 年后}——\textbf{John von Neumann} 用矩阵设计第一台电子计算机——\textbf{现代 BLAS 库的核心是矩阵乘法}。

\medskip

\textbf{160 年后}——\textbf{你打开 PyTorch}——\textbf{torch.matmul(A, B)}——\textbf{每一次 GPU 上 10\^{}{12} 次矩阵乘法}——\textbf{都在向 Cayley 致敬}。
\end{tcolorbox}

\section{直觉 / Intuition}

\subsection{向量：从 2D 箭头开始}

\textbf{向量是什么}？——\textbf{最简单的回答}：\textbf{一个箭头}。

\textbf{它有两个属性}：\textbf{方向} + \textbf{长度}。

在二维平面上——\textbf{我们用一对数 $(x, y)$ 表示一个向量}：

\[
\vec{v} = \begin{pmatrix} 3 \\ 2 \end{pmatrix}
\]

\textbf{这个 $(3,2)$ 的意思}——\textbf{从原点出发——向右走 3 步——向上走 2 步——画一个箭头到终点}。

\textbf{向量的三个基本身份}（必须同时记住）：

\begin{itemize}
\item \textbf{几何上}——\textbf{它是一个箭头}
\item \textbf{代数上}——\textbf{它是一个有序的数对（或数组）}
\item \textbf{物理上}——\textbf{它是位移、速度、力——任何"有方向 + 有大小"的量}
\end{itemize}

\textbf{为什么我们要"既几何又代数"？}——\textbf{因为几何让你"看懂"——代数让你"算"}。

\textbf{这是 Descartes 1637 年发明解析几何时的核心思想}——\textbf{把几何问题转成代数问题}——\textbf{从此数学进入了新时代}。

\subsection{向量空间与基}

\textbf{把所有的 2D 向量放在一起}——\textbf{这个集合}——\textbf{就是 2D 向量空间 $\mathbb{R}^2$}。

\textbf{它有两个核心运算}：

\begin{itemize}
\item \textbf{加法}：$\vec{u} + \vec{v}$——\textbf{对应位置相加}——\textbf{几何上是"首尾相接"}
\item \textbf{数乘}：$c \vec{v}$——\textbf{每个分量都乘 $c$}——\textbf{几何上是"拉长 / 缩短 / 反向"}
\end{itemize}

\textbf{基（basis）是什么}？——\textbf{基是"标尺"}——\textbf{一组特殊的向量——你能用它们的线性组合表示空间里的任何向量}。

\textbf{标准基}（最常用）：

\[
\vec{e}_1 = \begin{pmatrix} 1 \\ 0 \end{pmatrix}, \quad
\vec{e}_2 = \begin{pmatrix} 0 \\ 1 \end{pmatrix}
\]

\textbf{任何 2D 向量 $\vec{v} = (3, 2)$ 都能写成}：

\[
\vec{v} = 3 \vec{e}_1 + 2 \vec{e}_2
\]

\textbf{这就是"线性组合"}——\textbf{每个分量乘上对应的基向量——再加起来}。

\textbf{推广到 3D}——\textbf{$\mathbb{R}^3$ 的标准基}：

\[
\vec{e}_1 = \begin{pmatrix} 1 \\ 0 \\ 0 \end{pmatrix}, \quad
\vec{e}_2 = \begin{pmatrix} 0 \\ 1 \\ 0 \end{pmatrix}, \quad
\vec{e}_3 = \begin{pmatrix} 0 \\ 0 \\ 1 \end{pmatrix}
\]

\textbf{"维度"是什么}？——\textbf{基里向量的个数}——\textbf{$\mathbb{R}^2$ 是 2 维——$\mathbb{R}^3$ 是 3 维——$\mathbb{R}^n$ 是 $n$ 维}。

\textbf{基不是唯一的}——\textbf{你可以选别的}。比如 2D 里——\textbf{$\{(1,1), (1,-1)\}$ 也是一组基}——\textbf{它们也能"撑起"整个 $\mathbb{R}^2$}。

\textbf{选不同的基——同一个向量的"坐标"会变}——\textbf{但向量本身没变}——\textbf{这是"换基"的核心}（第 7 章详细讲）。

\subsection{矩阵 = 线性变换（核心中的核心）}

\textbf{现在到这一章最重要的一句话}——\textbf{画重点——背下来}：

\begin{quote}
\textbf{矩阵不是一堆数——矩阵是一个变换——它把一个向量变成另一个向量}。
\end{quote}

\textbf{什么叫"变换"}？——\textbf{你给它一个输入向量——它吐给你一个输出向量}——\textbf{就像函数 $f(x) = 2x + 1$ 一样——只不过现在的输入和输出都是向量}。

\textbf{什么叫"线性"}？——\textbf{它满足两条性质}：

\begin{enumerate}
\item \textbf{$T(\vec{u} + \vec{v}) = T(\vec{u}) + T(\vec{v})$}——\textbf{加法保持}
\item \textbf{$T(c \vec{v}) = c T(\vec{v})$}——\textbf{数乘保持}
\end{enumerate}

\textbf{用大白话说}——\textbf{"先加再变" = "先变再加"}——\textbf{"先乘再变" = "先变再乘"}——\textbf{这就是线性}。

\textbf{核心定理}（这是整章的"任督二脉"）：

\begin{tcolorbox}[essbox={矩阵的真正定义}]
\textbf{每一个线性变换 $T: \mathbb{R}^n \to \mathbb{R}^m$——都对应唯一的一个 $m \times n$ 矩阵 $A$——使得 $T(\vec{v}) = A \vec{v}$}。\\[4pt]
\textbf{反过来——每一个 $m \times n$ 矩阵——都是从 $\mathbb{R}^n$ 到 $\mathbb{R}^m$ 的一个线性变换}。
\end{tcolorbox}

\textbf{矩阵 = 线性变换的"坐标表示"}。

\subsection{具体例子：看见变换}

\textbf{看一个矩阵——立刻想象它对 2D 空间做了什么}：

\textbf{例 1：恒等矩阵（什么也不做）}：
\[
I = \begin{pmatrix} 1 & 0 \\ 0 & 1 \end{pmatrix}
\]
\textbf{效果}：\textbf{每个向量原地不动}——\textbf{这是"恒等变换"}。

\textbf{例 2：放大 2 倍}：
\[
A = \begin{pmatrix} 2 & 0 \\ 0 & 2 \end{pmatrix}
\]
\textbf{效果}：\textbf{每个向量长度变成 2 倍}——\textbf{整个平面"撑大了"}。

\textbf{例 3：水平翻转}：
\[
A = \begin{pmatrix} -1 & 0 \\ 0 & 1 \end{pmatrix}
\]
\textbf{效果}：\textbf{x 坐标取反}——\textbf{整个平面"镜像翻转"}。

\textbf{例 4：旋转 90 度}：
\[
A = \begin{pmatrix} 0 & -1 \\ 1 & 0 \end{pmatrix}
\]
\textbf{效果}：\textbf{每个向量逆时针转 90 度}——\textbf{这是旋转矩阵}。

\textbf{例 5：剪切（shear）}：
\[
A = \begin{pmatrix} 1 & 1 \\ 0 & 1 \end{pmatrix}
\]
\textbf{效果}：\textbf{矩形被"推歪"成平行四边形}——\textbf{这就是"剪切变换"}。

\textbf{关键观察}：\textbf{矩阵的列——就是"基向量被变换后的样子"}。

\textbf{看矩阵 $A = \begin{pmatrix} a & b \\ c & d \end{pmatrix}$ 时}——\textbf{要这样读它}：

\begin{itemize}
\item \textbf{第一列 $\begin{pmatrix}a\\c\end{pmatrix}$ = $\vec{e}_1$ 被变换后的样子}
\item \textbf{第二列 $\begin{pmatrix}b\\d\end{pmatrix}$ = $\vec{e}_2$ 被变换后的样子}
\end{itemize}

\textbf{这是"看懂矩阵"的最重要技巧}——\textbf{它是 Strang 和 3Blue1Brown 共同强调的}。

\section{数学 / The Math}

\subsection{矩阵 × 向量}

矩阵乘向量的公式：

\[
\begin{pmatrix} a & b \\ c & d \end{pmatrix}
\begin{pmatrix} x \\ y \end{pmatrix}
=
\begin{pmatrix} ax + by \\ cx + dy \end{pmatrix}
\]

\textbf{这个公式不是凭空捏造的}——\textbf{它来自"线性组合"}：

\[
\begin{pmatrix} a & b \\ c & d \end{pmatrix}
\begin{pmatrix} x \\ y \end{pmatrix}
= x \begin{pmatrix} a \\ c \end{pmatrix} + y \begin{pmatrix} b \\ d \end{pmatrix}
\]

\textbf{看见了吗}？——\textbf{输出 = "新坐标 $x$" 乘"第一列"——加"新坐标 $y$" 乘"第二列"}。

\textbf{这就是变换的几何意义}——\textbf{$x$ 倍的"新 $\vec{e}_1$" + $y$ 倍的"新 $\vec{e}_2$"}。

\subsection{矩阵 × 矩阵（最容易记错的运算）}

矩阵乘矩阵——\textbf{为什么是"行乘列"}？

\textbf{答案}：\textbf{矩阵乘法 = 复合变换}。

如果 $A$ 表示一个变换——$B$ 表示另一个变换——\textbf{那么 $AB$ 表示"先做 $B$，再做 $A$"的复合变换}。

\textbf{注意顺序}——\textbf{右边的先做}——\textbf{这是函数复合 $f \circ g$ 的同一个习惯}。

\textbf{乘法公式}：若 $A$ 是 $m \times n$，$B$ 是 $n \times p$，则 $AB$ 是 $m \times p$，且

\[
(AB)_{ij} = \sum_{k=1}^{n} A_{ik} B_{kj}
\]

\textbf{记忆口诀}：\textbf{"左行右列"}——\textbf{结果第 $i$ 行第 $j$ 列 = $A$ 的第 $i$ 行 · $B$ 的第 $j$ 列}（点积）。

\textbf{关键性质}：

\begin{itemize}
\item \textbf{$AB \ne BA$}——\textbf{矩阵乘法不交换}——\textbf{这是 Cayley 1858 年的大发现}
\item \textbf{$(AB)C = A(BC)$}——\textbf{结合律成立}
\item \textbf{$A(B+C) = AB + AC$}——\textbf{分配律成立}
\item \textbf{$AI = IA = A$}——\textbf{单位矩阵 $I$ 是"乘法的 1"}
\end{itemize}

\begin{example}[$AB \ne BA$ 的反例]
取 $A = \begin{pmatrix}0&1\\0&0\end{pmatrix}$，$B = \begin{pmatrix}0&0\\1&0\end{pmatrix}$。\\
\textbf{计算}：$AB = \begin{pmatrix}1&0\\0&0\end{pmatrix}$，$BA = \begin{pmatrix}0&0\\0&1\end{pmatrix}$。\\
\textbf{完全不相等}——\textbf{这告诉你——做事的顺序——很重要}。
\end{example}

\subsection{矩阵的转置}

\textbf{转置}：\textbf{把矩阵的行变列、列变行}——\textbf{记作 $A^T$}。

\[
A = \begin{pmatrix} 1 & 2 & 3 \\ 4 & 5 & 6 \end{pmatrix}, \quad
A^T = \begin{pmatrix} 1 & 4 \\ 2 & 5 \\ 3 & 6 \end{pmatrix}
\]

\textbf{转置的性质}：

\begin{align}
(A^T)^T &= A \\
(A+B)^T &= A^T + B^T \\
(AB)^T &= B^T A^T \quad \text{(注意顺序反了)}\\
(cA)^T &= c A^T
\end{align}

\textbf{转置在工程里的两个核心用途}：

\begin{itemize}
\item \textbf{内积}：$\vec{u} \cdot \vec{v} = \vec{u}^T \vec{v}$——\textbf{把"两个向量"变成"一个数"}
\item \textbf{对称矩阵}：$A = A^T$——\textbf{协方差矩阵、Hessian 矩阵——都是对称的}
\end{itemize}

\subsection{行列式：变换的"面积缩放因子"}

\textbf{行列式（determinant）}——\textbf{记作 $\det(A)$ 或 $|A|$}——\textbf{它是一个数}（标量）。

\textbf{它的几何意义——一句话}：

\begin{quote}
\textbf{行列式告诉我们——这个矩阵把"单位正方形的面积"放大了多少倍}。
\end{quote}

\textbf{对 2 阶矩阵}：

\[
\det\begin{pmatrix} a & b \\ c & d \end{pmatrix} = ad - bc
\]

\textbf{解释}：\textbf{矩阵的两个列向量}（$\vec{e}_1$ 被变换后的样子 和 $\vec{e}_2$ 被变换后的样子）\textbf{张成的平行四边形——它的面积——就是 $|ad - bc|$}。

\textbf{绝对值是面积——符号告诉你"方向"}：

\begin{itemize}
\item \textbf{$\det > 0$}——\textbf{变换"保持方向"}（不翻面）
\item \textbf{$\det < 0$}——\textbf{变换"翻面了"}（镜像翻转）
\item \textbf{$\det = 0$}——\textbf{变换把空间"压扁"了}——\textbf{平行四边形面积为 0——意味着两个列向量在同一条线上}——\textbf{矩阵"退化"了}（不可逆）
\end{itemize}

\textbf{对 3 阶矩阵}——\textbf{它是"体积缩放因子"}——\textbf{三个列向量张成的平行六面体的体积}。

\textbf{对 $n$ 阶}——\textbf{它是"$n$ 维体积缩放因子"}。

\textbf{这才是行列式的真正定义}——\textbf{不是同济书里那个吓人的求和公式}。

\subsection{行列式的计算}

\textbf{2 阶}：\textbf{交叉相乘相减}——\textbf{$ad - bc$}——\textbf{秒算}。

\textbf{3 阶}：\textbf{"对角线法则"（Sarrus 规则）}——\textbf{只对 3 阶有效}：

\[
\det\begin{pmatrix} a & b & c \\ d & e & f \\ g & h & i \end{pmatrix}
= aei + bfg + cdh - ceg - bdi - afh
\]

\textbf{4 阶及以上}——\textbf{用"按行/列展开"或"化为三角形"}。

\textbf{现代做法——用 Gauss 消元化为上三角}——\textbf{然后对角线相乘}——\textbf{效率 $O(n^3)$}——\textbf{电脑用的就是这种方法}。

\textbf{行列式的关键性质}：

\begin{itemize}
\item \textbf{$\det(AB) = \det(A) \det(B)$}——\textbf{乘法保持}——\textbf{这告诉我们"复合变换的缩放 = 两个变换缩放的乘积"}
\item \textbf{$\det(A^T) = \det(A)$}——\textbf{转置不变}
\item \textbf{$\det(A^{-1}) = 1/\det(A)$}——\textbf{逆的行列式}
\item \textbf{$\det(I) = 1$}——\textbf{恒等变换不缩放}
\end{itemize}

\subsection{矩阵的逆}

\textbf{矩阵的逆}——\textbf{记作 $A^{-1}$}——\textbf{它"撤销" $A$ 的变换}。

\textbf{定义}：\textbf{若存在 $B$ 使得 $AB = BA = I$——则 $B = A^{-1}$}。

\textbf{逆存在的条件}：\textbf{$\det(A) \ne 0$}——\textbf{方阵 + 非退化}。

\textbf{$\det = 0$ 时为什么没有逆}？——\textbf{因为 $A$ 把空间"压扁"了}——\textbf{从压扁后的空间——无法"反推"回原来的空间}——\textbf{信息丢了——逆变换不存在}。

\textbf{2 阶逆的公式}（背熟）：

\[
\begin{pmatrix} a & b \\ c & d \end{pmatrix}^{-1}
= \frac{1}{ad - bc} \begin{pmatrix} d & -b \\ -c & a \end{pmatrix}
\]

\textbf{记忆}——\textbf{"主对角交换、副对角变号、再除以行列式"}。

\textbf{逆的性质}：

\begin{align}
(A^{-1})^{-1} &= A \\
(AB)^{-1} &= B^{-1} A^{-1} \quad \text{(顺序又反了)}\\
(A^T)^{-1} &= (A^{-1})^T
\end{align}

\subsection{Gauss 消元：解 $Ax = b$ 的"瑞士军刀"}

\textbf{线性方程组的核心问题}：

\[
\begin{cases}
2x + y + z = 5 \\
x - y + 2z = 3 \\
3x + 2y - z = 4
\end{cases}
\]

\textbf{用矩阵写出}：

\[
\underbrace{\begin{pmatrix} 2 & 1 & 1 \\ 1 & -1 & 2 \\ 3 & 2 & -1 \end{pmatrix}}_{A}
\underbrace{\begin{pmatrix} x \\ y \\ z \end{pmatrix}}_{\vec{x}}
=
\underbrace{\begin{pmatrix} 5 \\ 3 \\ 4 \end{pmatrix}}_{\vec{b}}
\]

\textbf{Gauss 消元法的三个允许操作}（行变换）：

\begin{enumerate}
\item \textbf{两行互换}
\item \textbf{某一行乘以非零常数}
\item \textbf{某一行加上另一行的常数倍}
\end{enumerate}

\textbf{目标}：\textbf{把增广矩阵 $[A | \vec{b}]$ 化为上三角形}——\textbf{然后从最后一行回代}。

\textbf{这个算法叫"Gauss-Jordan 消元"}——\textbf{Gauss 1810 年代发明}（其实中国人 2000 年前在《九章算术》"方程"章里就用过类似方法——\textbf{叫"消首法"}）。

\textbf{现代实现}——\textbf{LU 分解}：\textbf{把 $A$ 拆成 $L \cdot U$}（下三角 × 上三角）——\textbf{然后两步回代}——\textbf{这是 LAPACK / NumPy 里 \texttt{linalg.solve} 的底层算法}。

\begin{theorem}[$Ax = b$ 的解结构]
设 $A$ 是 $n \times n$ 方阵。则：
\begin{enumerate}
\item \textbf{$\det(A) \ne 0$}——\textbf{方程有唯一解}——\textbf{$\vec{x} = A^{-1} \vec{b}$}
\item \textbf{$\det(A) = 0$}——\textbf{要么无解——要么有无穷多解}（取决于 $\vec{b}$）
\end{enumerate}
\end{theorem}

\textbf{工程实践}——\textbf{永远不要计算 $A^{-1}$ 然后乘 $\vec{b}$}——\textbf{要用 \texttt{np.linalg.solve(A, b)}}——\textbf{它做 LU 分解 + 回代}——\textbf{更快、更稳定、误差更小}。

\section{Aha 例子：3D 游戏旋转 / Aha: 3D Game Rotation}

\textbf{场景}：你在做一个 3D 游戏——\textbf{《我的世界》风格的方块世界}。

玩家按下 \textbf{A 键}——\textbf{要把摄像机水平转 5 度}——\textbf{然后按 \textbf{W 键}——再前进 1 米}。

\textbf{问题}：\textbf{游戏里的每一个方块、每一个角色——它们的位置都要重新算}。

\textbf{用什么算}？——\textbf{矩阵乘法}。

\subsection{3D 旋转矩阵}

\textbf{绕 z 轴（垂直轴）旋转角度 $\theta$ 的矩阵}：

\[
R_z(\theta) = \begin{pmatrix}
\cos\theta & -\sin\theta & 0 \\
\sin\theta & \cos\theta & 0 \\
0 & 0 & 1
\end{pmatrix}
\]

\textbf{绕 x 轴}：

\[
R_x(\theta) = \begin{pmatrix}
1 & 0 & 0 \\
0 & \cos\theta & -\sin\theta \\
0 & \sin\theta & \cos\theta
\end{pmatrix}
\]

\textbf{绕 y 轴}：

\[
R_y(\theta) = \begin{pmatrix}
\cos\theta & 0 & \sin\theta \\
0 & 1 & 0 \\
-\sin\theta & 0 & \cos\theta
\end{pmatrix}
\]

\textbf{任意 3D 旋转}——\textbf{都是这三个的组合}——\textbf{乘起来就行}。

\subsection{为什么这件事是 Aha 时刻}

\begin{tcolorbox}[ahabox={Aha: 矩阵乘法不交换——是游戏的物理}]
\textbf{你在游戏里——先水平转 90 度——再向前走 1 米——你到的地方}——

\medskip

\textbf{和——先向前走 1 米——再水平转 90 度——你到的地方}——

\medskip

\textbf{完！全！不！一！样！}

\medskip

\textbf{这就是 $AB \ne BA$ 的物理意义}——\textbf{矩阵乘法的顺序——就是动作的顺序}。

\medskip

Cayley 1858 年发现的\textbf{"$AB$ 不等于 $BA$"}——\textbf{在 2026 年的每一款 3D 游戏里——每秒被 GPU 验证几亿次}——\textbf{它不是抽象的怪事——它是空间的本性}。

\medskip

\textbf{再补一刀}——\textbf{在量子力学里——位置算符 $\hat{x}$ 和动量算符 $\hat{p}$——也满足 $\hat{x}\hat{p} \ne \hat{p}\hat{x}$}——\textbf{这就是 Heisenberg 不确定性原理的数学根源}。

\medskip

\textbf{从游戏到量子力学——背后是同一个数学事实}——\textbf{矩阵乘法不交换}。
\end{tcolorbox}

\textbf{完整的旋转 + 平移 = 齐次坐标}：

在游戏里——\textbf{我们要同时做"旋转"和"平移"}——但平移不是线性变换（它不满足 $T(\vec{0}) = \vec{0}$）。

\textbf{解决方法}：\textbf{把 3D 点扩成 4D 齐次坐标 $(x, y, z, 1)$}——\textbf{然后用 4×4 矩阵处理}：

\[
\begin{pmatrix} R_{3\times3} & \vec{t} \\ \vec{0}^T & 1 \end{pmatrix}
\begin{pmatrix} \vec{p} \\ 1 \end{pmatrix}
= \begin{pmatrix} R\vec{p} + \vec{t} \\ 1 \end{pmatrix}
\]

\textbf{这就是 OpenGL / Unity / Unreal 引擎里 \texttt{ModelViewMatrix} 的本质}——\textbf{每一个像素的位置——都是一次 4×4 矩阵 × 4×1 向量}。

\subsection{Python 代码}

\begin{lstlisting}[language=Python]
import numpy as np

def rot_z(theta):
    """绕 z 轴旋转 theta 弧度"""
    c, s = np.cos(theta), np.sin(theta)
    return np.array([
        [c, -s, 0],
        [s,  c, 0],
        [0,  0, 1]
    ])

def rot_x(theta):
    c, s = np.cos(theta), np.sin(theta)
    return np.array([
        [1, 0,  0],
        [0, c, -s],
        [0, s,  c]
    ])

# 一个 3D 点
p = np.array([1, 0, 0])    # 沿 x 轴方向 1 米的点

# 绕 z 轴转 90 度
p_rotated = rot_z(np.pi / 2) @ p
print("旋转后:", p_rotated)   # [0, 1, 0]——确实转到了 y 轴

# 复合：先绕 z 转 90 度——再绕 x 转 90 度
R_combined = rot_x(np.pi / 2) @ rot_z(np.pi / 2)
p_combined = R_combined @ p
print("复合后:", p_combined)

# 验证：AB != BA
R1 = rot_x(np.pi / 2) @ rot_z(np.pi / 2)
R2 = rot_z(np.pi / 2) @ rot_x(np.pi / 2)
print("R1 == R2 吗?", np.allclose(R1, R2))   # False!

# 行列式：旋转矩阵的 det 应该 = 1（保持体积、保持方向）
print("det(R_combined) =", np.linalg.det(R_combined))   # 1.0

# 解一个线性方程组
A = np.array([[2, 1, 1], [1, -1, 2], [3, 2, -1]])
b = np.array([5, 3, 4])
x = np.linalg.solve(A, b)
print("方程的解:", x)
\end{lstlisting}

\textbf{这段代码——\textbf{30 行}——\textbf{涵盖了本章 80\% 的内容}}：

\begin{itemize}
\item \textbf{矩阵 = 旋转变换}
\item \textbf{矩阵乘法 = 复合}
\item \textbf{$AB \ne BA$}
\item \textbf{行列式 = 1（旋转不缩放）}
\item \textbf{\texttt{np.linalg.solve} = Gauss 消元的现代化}
\end{itemize}

\section{现代视角：矩阵在 ML 与图形学 / Modern: Matrices in ML and Graphics}

\subsection{深度学习 = 大型矩阵运算}

\textbf{一个最简单的神经网络层}：

\[
\vec{y} = \sigma(W \vec{x} + \vec{b})
\]

\textbf{这里}：
\begin{itemize}
\item \textbf{$\vec{x}$}——输入向量（比如 784 维——MNIST 图片展平）
\item \textbf{$W$}——权重矩阵（比如 $128 \times 784$——把 784 维映射到 128 维）
\item \textbf{$\vec{b}$}——偏置向量
\item \textbf{$\sigma$}——非线性激活函数（ReLU、sigmoid 等）
\end{itemize}

\textbf{$W \vec{x}$——就是矩阵乘向量——就是线性变换}。

\textbf{多层堆叠}——\textbf{就是多次线性变换的复合}——\textbf{再加非线性}。

\textbf{ChatGPT 的核心 Transformer——每一层都是几十个大矩阵乘法}——\textbf{一次前向传播——几百亿次浮点运算}——\textbf{全部用 GPU 上的 BLAS 库实现}。

\textbf{你今天和 ChatGPT 说"你好"——背后是几十亿次 Cayley 1858 年定义的矩阵乘法}。

\subsection{图形学 = 4×4 矩阵的世界}

\textbf{所有 3D 图形}——\textbf{从《英雄联盟》到《赛博朋克 2077》}——\textbf{都建立在 4×4 矩阵之上}：

\begin{itemize}
\item \textbf{Model Matrix}——\textbf{把模型从局部坐标系转到世界坐标系}
\item \textbf{View Matrix}——\textbf{把世界坐标系转到摄像机坐标系}
\item \textbf{Projection Matrix}——\textbf{把 3D 投影到 2D 屏幕}
\end{itemize}

\textbf{每个像素的位置 = Projection × View × Model × 顶点坐标}——\textbf{4×4 矩阵乘 4×4 矩阵乘 4×4 矩阵乘 4×1 向量}——\textbf{你每帧看的画面——是 GPU 跑了几百万次这个公式}。

\subsection{数值线性代数：工程世界的基石}

\textbf{LAPACK + BLAS}——\textbf{Linear Algebra PACKage + Basic Linear Algebra Subprograms}——\textbf{是 1970 年代以来一直在演化的"数学软件"}——\textbf{用 Fortran / C / 汇编写成}。

\textbf{你今天用的}：

\begin{itemize}
\item \textbf{NumPy} 的 \texttt{linalg.solve}——\textbf{底层是 LAPACK}
\item \textbf{PyTorch} 的 \texttt{torch.matmul}——\textbf{底层是 cuBLAS（GPU 上的 BLAS）}
\item \textbf{MATLAB} 的 \texttt{A\textbackslash b}——\textbf{底层是 LAPACK}
\end{itemize}

\textbf{所有的"高级 AI 软件"——本质上都是在调用 50 年前写的数值线性代数库}。

\textbf{这就是为什么——线代是现代工程的"任督二脉"中的"督脉"}——\textbf{它在每一处——你看不见的地方——撑着整个数字世界}。

\section{代码与思考 / Code and Exercises}

\subsection{配套模块 \texttt{linalg\_intro.py}}

GitHub 上的 \texttt{modules/linalg\_intro.py} 包含\textbf{7 个 demo}：

\begin{enumerate}
\item \texttt{demo\_vector\_basics()}——可视化 2D 向量的加法、数乘
\item \texttt{demo\_matrix\_as\_transform()}——\textbf{核心 demo}——画"单位正方形被矩阵变换后的样子"——\textbf{你会"看见"线性变换}
\item \texttt{demo\_matmul\_noncommutative()}——验证 $AB \ne BA$——用旋转 + 剪切的复合
\item \texttt{demo\_determinant()}——计算并可视化"面积缩放"
\item \texttt{demo\_gauss\_elimination()}——手工实现 Gauss 消元——和 \texttt{np.linalg.solve} 对比
\item \texttt{demo\_3d\_rotation()}——本章 Aha 例子：3D 旋转 + 平移 + 齐次坐标
\item \texttt{demo\_image\_rotation()}——把一张真实图片旋转任意角度——用矩阵实现
\end{enumerate}

\textbf{使用}：

\begin{lstlisting}[language=bash]
git clone https://github.com/inaturelz-coder/math-rendu
cd math-rendu
pip install -r requirements.txt
python modules/linalg_intro.py
\end{lstlisting}

\subsection{思考题 / Exercises}

\begin{enumerate}
\item \textbf{[直觉]} 用一句话向你的高中同桌解释"矩阵是什么"——不许说"一堆数排成方阵"。
\item \textbf{[几何]} 矩阵 $\begin{pmatrix}2&0\\0&3\end{pmatrix}$ 对单位圆做了什么变换？画出来。
\item \textbf{[计算]} 求 $A = \begin{pmatrix}1&2\\3&4\end{pmatrix}$ 的行列式、逆矩阵——并验证 $A A^{-1} = I$。
\item \textbf{[复合]} 取两个 2×2 旋转矩阵 $R(30°)$ 和 $R(60°)$——验证 $R(30°) R(60°) = R(90°)$。\textbf{为什么旋转矩阵的乘法是交换的}（提示：2D 里）？
\item \textbf{[反例]} 找出两个 2×2 矩阵 $A$、$B$——使得 $AB \ne BA$——并算出两者之差。
\item \textbf{[应用]} 一个机器人手臂的末端在它自己的坐标系里位置为 $(0.5, 0, 0)$。手臂坐标系相对世界坐标系——\textbf{原点在 $(1, 2, 0)$}——\textbf{绕 z 轴转了 30 度}。求末端在世界坐标系里的位置。
\item \textbf{[Gauss]} 手算 Gauss 消元——解
\[
\begin{cases} x+2y+z=4 \\ 2x+y-z=1 \\ x-y+3z=2 \end{cases}
\]
\textbf{然后用 \texttt{np.linalg.solve} 验证}。
\item \textbf{[历史]} 查一查：为什么 Cayley 1858 年的那篇论文起初\textbf{没什么人引用}？它在数学界"火起来"用了多少年？
\item \textbf{[现代]} 用 PyTorch 写一个最简单的"线性层"——\textbf{torch.nn.Linear(10, 5)}——\textbf{看看它内部存了什么矩阵}——\textbf{这个矩阵的 shape 是什么}？
\end{enumerate}

\section{本章小结 / Chapter Summary}

\begin{tcolorbox}[essbox={本章核心}]
\begin{itemize}
\item \textbf{向量}——\textbf{箭头 + 数组 + 物理量}——\textbf{三位一体}
\item \textbf{基}——空间的"标尺"——\textbf{基不唯一——但维度唯一}
\item \textbf{矩阵 = 线性变换}——\textbf{不是一堆数}——\textbf{是一个"映射"}
\item \textbf{矩阵的列 = 基向量被变换后的样子}——\textbf{看矩阵的标准姿势}
\item \textbf{矩阵乘法 = 变换复合}——\textbf{$AB \ne BA$ 是空间的本性}
\item \textbf{行列式 = 体积缩放因子}——\textbf{$\det = 0$ 意味着空间被压扁}
\item \textbf{逆矩阵 = 撤销变换}——\textbf{$\det \ne 0$ 时才存在}
\item \textbf{Gauss 消元 = 解 $Ax = b$ 的瑞士军刀}——\textbf{现代用 \texttt{np.linalg.solve}}
\item \textbf{现代应用}——\textbf{3D 游戏（4×4 矩阵）}——\textbf{深度学习（大矩阵乘法）}——\textbf{都是矩阵的世界}
\end{itemize}
\end{tcolorbox}

\textbf{下一章——线代核心}——\textbf{我们要进入"特征值、特征向量、SVD、对角化"}——\textbf{这是线性代数"内功"的部分}——\textbf{也是 PCA、推荐系统、量子力学共同的数学基础}。

\textbf{今天讲的"矩阵 = 变换"——下章我们要问}：\textbf{这个变换里——有没有"不变"的方向}？——\textbf{答案就是"特征向量"}。

\clearpage

% =====================================================================
%  第 7 章 · 线代核心 / Linear Algebra II
%  特征值、特征向量、对角化、内积、正交基、Gram-Schmidt、
%  对称矩阵、谱定理、厄米矩阵、二次型、正定矩阵
%  小故事：Hilbert 与"Hilbert 空间"的诞生（20 世纪初 · Göttingen）
%  Aha：人脸识别的特征向量（Eigenfaces）
% =====================================================================

\chapter{线代核心 / Linear Algebra II}
\label{ch:linalg2}

\begin{quote}\itshape
1904 年，德国 Göttingen。\\
一位 42 岁的数学家——在他的"积分方程"讲义里——\textbf{第一次写下了一个无限维空间里的"内积"}。\\
他叫 \textbf{David Hilbert}。\\
20 年后——\textbf{Heisenberg 和 Schrödinger} 用这个空间——描述了电子和原子——\textbf{量子力学诞生}。\\[6pt]
100 年后——\textbf{你打开手机的相册——它自动把你和家人分组}——\textbf{背后是同一个空间的特征向量}——\textbf{从 Hilbert 到 Eigenface——一脉相承}。
\end{quote}

\section{写在前面 / Preface}

\textbf{第 6 章我们学了"矩阵是线性变换"}——\textbf{这一章我们要回答一个更深的问题}：

\textbf{一个矩阵——它最"本质"的样子——到底是什么}？

\textbf{这个问题——是 20 世纪整个数学的中心问题之一}——\textbf{答案叫"特征值分解"}——\textbf{它是 21 世纪机器学习的根}。

\textbf{我研究生第一次真正"用"线代}——\textbf{不是解方程}——\textbf{而是做人脸识别}。

\textbf{老板给我一堆人脸照片}——\textbf{每张 92×112 像素}——\textbf{展平成一个 10304 维的向量}——\textbf{他说"算它们的协方差矩阵——然后求前 50 个特征向量"}。

\textbf{我把那 50 个特征向量画出来}——\textbf{每一张——都是一张"幽灵般的人脸"}——\textbf{他们叫 Eigenfaces}。

\textbf{我那一刻才明白}——\textbf{特征向量不是"满足 $A\bm{v} = \lambda\bm{v}$ 的奇怪向量"}——\textbf{而是"数据自己最想沿着延伸的方向"}——\textbf{是矩阵"骨架"里的"主轴"}。

\textbf{从那以后——我看任何矩阵——都先问"它的特征向量长什么样"}。

\textbf{这一章——就是讲这件事}。我们按下面这个顺序走：

\begin{itemize}
\item \textbf{第一步}：\textbf{特征值 + 特征向量}——\textbf{矩阵不会旋转的那些方向}
\item \textbf{第二步}：\textbf{对角化 + 矩阵幂}——\textbf{找最好的基——让矩阵变成对角阵}
\item \textbf{第三步}：\textbf{内积 + 正交基 + Gram-Schmidt}——\textbf{把"角度"和"长度"装进空间里}
\item \textbf{第四步}：\textbf{对称矩阵 + 谱定理}——\textbf{所有数据科学背后的那个大定理}
\item \textbf{第五步}：\textbf{厄米矩阵}——\textbf{对称矩阵的复数版——通向量子力学}
\item \textbf{第六步}：\textbf{二次型 + 正定矩阵}——\textbf{凸优化、机器学习损失函数的几何}
\end{itemize}

\textbf{这一章——可以说是整本书 + 整个 ML 时代——最重要的一章}。\textbf{读懂它——你会拿到一把万能钥匙}。

\section{一句话本质 / The Essence in One Sentence}

\begin{tcolorbox}[essbox={一句话本质}]
\textbf{中文}：\textbf{特征向量是矩阵"不旋转、只缩放"的方向——对角化就是换一组基——让矩阵在新基下变成最简单的对角形式}。\\[6pt]
\textbf{English}: Eigenvectors are the directions a matrix only stretches but never rotates; diagonalization is choosing a new basis so the matrix becomes a diagonal one.
\end{tcolorbox}

\textbf{记住这一句}——\textbf{这一章已经讲完了 60\%}。剩下的 40\% 是\textbf{对称矩阵的特殊待遇 + 内积带来的"几何"}。

\section{教材里你看到的 / What the Textbook Tells You}

\subsection{同济《线性代数》第五章}

\textbf{同济版把这一章压缩成 30 页}——\textbf{叫"相似矩阵与二次型"}——\textbf{包括特征值、对角化、对称矩阵、二次型}。

\textbf{优点}：\textbf{紧凑、考研覆盖、公式齐全}。

\textbf{缺点}（致命）：

\begin{enumerate}
\item \textbf{"特征向量"的几何意义只字未提}——\textbf{学生算了一辈子 $\det(A - \lambda I) = 0$——不知道这个 $\lambda$ 是什么}
\item \textbf{"对角化"被讲成纯粹的计算技巧}——\textbf{从来没说"对角化 = 换一组好基"}
\item \textbf{"谱定理"几乎没讲}——\textbf{只说"实对称矩阵可对角化"}——\textbf{没说为什么这件事是 20 世纪数学的奇迹}
\item \textbf{二次型完全脱离几何}——\textbf{讲到"正惯性指数"——学生彻底投降}
\end{enumerate}

\subsection{Strang《Linear Algebra and Its Applications》第 5--7 章}

\textbf{Strang 是这一章的最佳老师}。

\textbf{他的特征值章节}——\textbf{第一页就画了一个"矩阵把空间拉伸"的图}——\textbf{然后告诉你 $\bm{v}$ 是那些"只被拉、不被转"的方向}。

\textbf{他的谱定理}——\textbf{他说"这是线性代数的核心定理"}——\textbf{不止讲完——还讲了为什么}。

\subsection{Axler《Linear Algebra Done Right》}

\textbf{Sheldon Axler 的这本书}——\textbf{是数学系本科生的"硬核"教材}。

\textbf{Axler 的主张}：\textbf{把"行列式"推迟到最后}——\textbf{用"特征值"作为整本书的主线}。

\textbf{他证明谱定理的方式——比任何教科书都优雅}——\textbf{推荐数学品味强的同学读}。

\subsection{我的策略}

\textbf{这一章我们要做三件事}：

\begin{enumerate}
\item \textbf{把特征值 + 特征向量的几何意义讲到刻进骨子里}——\textbf{它们是矩阵"骨架"的主轴}
\item \textbf{把谱定理 + 对角化讲清楚}——\textbf{它是 PCA / Eigenface / 物理振动模态背后的同一个定理}
\item \textbf{把厄米矩阵和量子力学的桥搭起来}——\textbf{让你知道线代不只是工程工具——还是物理的语言}
\end{enumerate}

\begin{tcolorbox}[storybox={\Large 小故事：Hilbert 与"Hilbert 空间"的诞生（1904 · Göttingen）}]
\textbf{David Hilbert}（1862--1943）——\textbf{德国数学家}——\textbf{20 世纪上半叶最伟大的数学家——没有之一}。

\medskip

\textbf{1900 年——巴黎国际数学家大会}——\textbf{他登台演讲——提出了"23 个数学问题"}——\textbf{为整个 20 世纪数学定下了路线图}。

\medskip

\textbf{但他真正改变物理学的——是 1904 年开始研究的"积分方程"}。

\medskip

那时候——\textbf{物理学家在研究"弦的振动"、"热的传导"、"光的衍射"}——\textbf{这些问题都归结为一类奇怪的方程}——\textbf{未知数不是一个数 $x$——而是一个函数 $f(x)$}。

\medskip

\textbf{Hilbert 灵机一动}——\textbf{他说}：

\medskip

\textit{"如果我把每一个函数 $f(x)$ 当成一个 \textbf{向量}——把所有这样的函数放在一起当成一个 \textbf{无限维空间}——那么积分方程就变成了无限维空间里的线性方程组——所有线代的工具都能搬过来用！"}

\medskip

\textbf{这个无限维空间——后来被命名为 Hilbert 空间}（Hilbert Space）。

\medskip

\textbf{它的核心结构是"内积"} $\langle f, g \rangle = \int f(x)\overline{g(x)}\,dx$——\textbf{有了内积——就有了"角度"和"长度"和"正交"}——\textbf{Hilbert 把整个线代搬进了无限维}。

\medskip

\textbf{1925 年}——\textbf{Heisenberg} 在哥廷根读博士时——\textbf{用矩阵描述电子的状态}——\textbf{Schrödinger} 用波函数描述同一件事——\textbf{两个人吵了整整一年——谁对谁错}？

\medskip

\textbf{1926 年——Hilbert 的学生 John von Neumann} 站出来——\textbf{他说"你们俩都对——你们的对象都活在同一个 Hilbert 空间里——只是用了不同的基！"}

\medskip

\textbf{这个证明——是 20 世纪物理学最优雅的统一}——\textbf{量子力学从此被装进了"Hilbert 空间 + 厄米算子"的框架}——\textbf{一直用到今天}。

\medskip

\textbf{Hilbert 一生最喜欢说的一句话}：

\medskip

\textit{"Wir m\"ussen wissen —— wir werden wissen."}（\textbf{我们必须知道——我们终将知道}）

\medskip

\textbf{这句话刻在他 Göttingen 的墓碑上}。

\medskip

\textbf{今天——你打开 PyTorch——\texttt{torch.linalg.eigh(A)}——它给你的特征向量在 Hilbert 空间里"正交"}——\textbf{每一次调用——都是对这位 19 世纪老者的致敬}。
\end{tcolorbox}

\section{直觉 / Intuition}

\subsection{特征向量：矩阵不旋转的方向}

\textbf{先讲一个最直观的画面}。

\textbf{想象一个 2D 平面}——\textbf{你拿一个矩阵 $A$ 作用在每一个向量上}——\textbf{绝大多数向量}——\textbf{都会被 $A$"扭一下"+"拉一下"}——\textbf{方向变了——长度也变了}。

\textbf{但是——总有那么几个特殊的方向}——\textbf{$A$ 作用在它身上时}——\textbf{只把它"拉长或缩短"}——\textbf{方向纹丝不动}。

\textbf{这些"$A$ 不会旋转的方向"——就叫特征向量}（eigenvector）。

\textbf{被拉长的倍数——就叫特征值}（eigenvalue）。

\textbf{用一个方程写出来}：

\[
A \bm{v} = \lambda \bm{v}
\]

\textbf{读法}：\textbf{矩阵 $A$ 作用在 $\bm{v}$ 上——结果等于 $\bm{v}$ 自己乘以一个数 $\lambda$}。

\textbf{"乘以一个数" = "只缩放、不旋转"}——\textbf{这就是特征向量的全部秘密}。

\subsection{物理意义：振动模态}

\textbf{把这个画面换成物理}——\textbf{特征向量的意义就更清楚}。

\textbf{想象一根吉他弦}——\textbf{你拨它——它会振动}——\textbf{这根弦能怎么振}？

\begin{itemize}
\item \textbf{最简单的一种}——\textbf{整根弦上下抖}——\textbf{这是第一个"模态"}——\textbf{频率最低}（叫"基频"）
\item \textbf{第二种}——\textbf{弦的中间不动——左右两半反向抖}——\textbf{频率是基频的 2 倍}
\item \textbf{第三种}——\textbf{弦上有两个不动点——分成三段}——\textbf{频率是基频的 3 倍}
\end{itemize}

\textbf{每一种"振动模态"——就是这根弦的一个特征向量}——\textbf{对应的频率——就是特征值}。

\textbf{你拨弦时——不管你怎么拨}——\textbf{弦的实际运动——总是这些模态的叠加}。

\textbf{这就是 Fourier 级数的物理本源}——\textbf{把任何运动分解成"基本模态" = 把任何向量分解成"特征向量的线性组合"}。

\textbf{第 5 章我们讲 Fourier}——\textbf{它讲的就是"$d^2/dx^2$"这个微分算子的特征向量}——\textbf{答案是 $\sin(nx)$ 和 $\cos(nx)$}——\textbf{特征值是 $-n^2$}。

\textbf{所以——Fourier 级数 = 特征分解}——\textbf{两件事是一回事}。

\subsection{几何意义：椭圆的主轴}

\textbf{第三个画面}——\textbf{我最喜欢的画面}。

\textbf{拿一个 $2\times 2$ 的对称矩阵 $A$}——\textbf{比如 $A = \begin{pmatrix} 3 & 1 \\ 1 & 2 \end{pmatrix}$}——\textbf{让它作用在单位圆上的每个点}。

\textbf{单位圆会被压扁/拉长}——\textbf{变成一个椭圆}。

\textbf{这个椭圆有两条主轴}——\textbf{一长一短}——\textbf{它们就是 $A$ 的两个特征向量}——\textbf{主轴的长度——就是特征值}。

\textbf{这就是"特征向量是矩阵的主轴"这个说法的来源}——\textbf{它把"代数"和"几何"统一在了一起}。

\section{数学 / Mathematics}

\subsection{特征值的定义和求法}

\textbf{定义}：\textbf{对于一个 $n\times n$ 的矩阵 $A$——如果存在一个非零向量 $\bm{v}$ 和一个数 $\lambda$——使得}

\[
A\bm{v} = \lambda \bm{v}
\]

\textbf{那么 $\bm{v}$ 叫做 $A$ 的特征向量}——\textbf{$\lambda$ 叫做对应的特征值}。

\textbf{怎么求}？——\textbf{把方程改写一下}：

\[
(A - \lambda I)\bm{v} = \bm{0}
\]

\textbf{$\bm{v}$ 非零}——\textbf{意味着 $A - \lambda I$ 不可逆}——\textbf{意味着}：

\[
\boxed{\det(A - \lambda I) = 0}
\]

\textbf{这个方程叫"特征方程"}（characteristic equation）——\textbf{它是关于 $\lambda$ 的 $n$ 次多项式}——\textbf{解出来就是 $n$ 个特征值}。

\textbf{对每一个特征值 $\lambda_i$——回去解 $(A - \lambda_i I)\bm{v} = \bm{0}$——就得到对应的特征向量}。

\subsection{对角化}

\textbf{假设 $A$ 有 $n$ 个线性无关的特征向量 $\bm{v}_1, \dots, \bm{v}_n$——对应特征值 $\lambda_1, \dots, \lambda_n$}。

\textbf{把这些特征向量按列排起来——组成矩阵 $P$}：

\[
P = [\bm{v}_1 \;|\; \bm{v}_2 \;|\; \cdots \;|\; \bm{v}_n]
\]

\textbf{把特征值放在对角线上——组成对角矩阵 $\Lambda$}：

\[
\Lambda = \mathrm{diag}(\lambda_1, \lambda_2, \dots, \lambda_n)
\]

\textbf{那么}：

\[
\boxed{A = P \Lambda P^{-1}}
\]

\textbf{这就是对角化}（diagonalization）。

\textbf{它的真正含义}：

\textbf{$P^{-1}$——把任意向量换到"特征向量基"下表示}——\textbf{$\Lambda$——在新基下做缩放}——\textbf{$P$——再换回原来的基}。

\textbf{在新基下——矩阵就是一个简单的对角阵——只做缩放——什么都不做}。

\textbf{这就是"找最好的基"的含义}——\textbf{对角化 = 换一组让矩阵看起来最简单的基}。

\subsection{矩阵幂的爆炸式简化}

\textbf{为什么对角化这么有用}？——\textbf{看一个例子}。

\textbf{你要算 $A^{100}$}——\textbf{硬乘的话——99 次矩阵乘法——每次 $O(n^3)$ 操作}——\textbf{累死}。

\textbf{但如果 $A = P\Lambda P^{-1}$}——\textbf{那么}：

\[
A^2 = P\Lambda P^{-1} \cdot P\Lambda P^{-1} = P \Lambda^2 P^{-1}
\]

\[
A^{100} = P \Lambda^{100} P^{-1}
\]

\textbf{$\Lambda^{100}$ 是对角阵——直接把每个对角元素的 100 次方算出来——$O(n)$ 操作}。

\textbf{然后两边乘 $P$ 和 $P^{-1}$——总共 $O(n^3)$}——\textbf{原来 $99 \cdot n^3$ 的工作——变成了 $n^3$}。

\textbf{这就是为什么 Fibonacci 数列、Markov 链、PageRank 的稳态——都用对角化来算}。

\subsection{内积——把"角度"装进空间}

\textbf{光有向量空间——只能加和缩放——没有"长度"也没有"角度"}。

\textbf{要装"几何"——必须引入"内积"}（inner product）。

\textbf{对实数向量 $\bm{u}, \bm{v} \in \mathbb{R}^n$——标准内积}：

\[
\langle \bm{u}, \bm{v} \rangle = \bm{u}^\top \bm{v} = \sum_{i=1}^n u_i v_i
\]

\textbf{有了内积——立刻得到三样东西}：

\begin{itemize}
\item \textbf{长度}：\;$\|\bm{v}\| = \sqrt{\langle \bm{v}, \bm{v} \rangle}$
\item \textbf{角度}：\;$\cos\theta = \dfrac{\langle \bm{u}, \bm{v} \rangle}{\|\bm{u}\|\|\bm{v}\|}$
\item \textbf{正交}：\;$\bm{u} \perp \bm{v} \iff \langle \bm{u}, \bm{v} \rangle = 0$
\end{itemize}

\textbf{这三样——是整个几何学的语言}——\textbf{没有内积——你就活在一个"扁平"的代数世界里}。

\subsection{正交基与 Gram-Schmidt}

\textbf{一组基}\;$\bm{q}_1, \dots, \bm{q}_n$\;\textbf{叫正交基}——\textbf{如果它们两两正交}（$\langle \bm{q}_i, \bm{q}_j \rangle = 0$ 当 $i\neq j$）。

\textbf{如果还满足 $\|\bm{q}_i\| = 1$}——\textbf{就叫"标准正交基"}（orthonormal basis）。

\textbf{为什么正交基好}？——\textbf{把一个向量 $\bm{v}$ 在正交基下分解——系数直接是内积}：

\[
\bm{v} = \sum_{i=1}^n \langle \bm{v}, \bm{q}_i \rangle \bm{q}_i
\]

\textbf{不用解方程——直接算内积——这就是 Fourier 系数公式的来源}。

\textbf{Gram-Schmidt 过程}——\textbf{把任意一组线性无关的向量 $\bm{a}_1, \dots, \bm{a}_n$——加工成一组正交基 $\bm{q}_1, \dots, \bm{q}_n$}：

\begin{align*}
\bm{q}_1 &= \frac{\bm{a}_1}{\|\bm{a}_1\|} \\
\bm{q}_2 &= \frac{\bm{a}_2 - \langle \bm{a}_2, \bm{q}_1 \rangle \bm{q}_1}{\|\cdot\|} \\
\bm{q}_3 &= \frac{\bm{a}_3 - \langle \bm{a}_3, \bm{q}_1 \rangle \bm{q}_1 - \langle \bm{a}_3, \bm{q}_2 \rangle \bm{q}_2}{\|\cdot\|}
\end{align*}

\textbf{每一步——把新向量在已有基方向上的分量"减掉"——剩下的就和前面所有人都正交}。

\textbf{这就是 QR 分解}（$A = QR$）的几何来源。

\subsection{对称矩阵与谱定理}

\textbf{对称矩阵}——\textbf{$A^\top = A$ 的实矩阵}——\textbf{是线代里最特殊的一类}。

\textbf{谱定理}（Spectral Theorem）：

\begin{tcolorbox}[essbox={谱定理 / Spectral Theorem}]
\textbf{任意实对称矩阵 $A$——必然可以正交对角化}：
\[
A = Q \Lambda Q^\top
\]
\textbf{其中 $Q$ 是正交矩阵}（$Q^\top Q = I$）——\textbf{$\Lambda$ 是实对角阵}。\\[4pt]
\textbf{中文意思}：\textbf{对称矩阵的特征值全是实数——特征向量可以选成两两正交}。
\end{tcolorbox}

\textbf{这个定理有多重要}？——\textbf{看下面这一长串它的应用就知道}：

\begin{itemize}
\item \textbf{PCA 主成分分析}——\textbf{对协方差矩阵做谱分解——前 $k$ 个特征向量就是主成分}
\item \textbf{Eigenfaces 人脸识别}——\textbf{同上}
\item \textbf{物理振动模态}——\textbf{刚度矩阵是对称的——谱分解给出振动频率}
\item \textbf{Google PageRank}——\textbf{用谱方法找网页重要性}
\item \textbf{谱聚类}——\textbf{社区发现、图像分割}
\item \textbf{量子力学}——\textbf{所有观测量都是厄米算子——谱定理保证测量值是实数}
\end{itemize}

\textbf{谱定理 = 一个公式 + 一片江山}。

\subsection{厄米矩阵——通向量子力学}

\textbf{把对称矩阵扩展到复数}——\textbf{就是厄米矩阵}（Hermitian matrix）：

\[
A^\dagger = A \qquad \text{其中 } A^\dagger = \overline{A^\top}
\]

\textbf{$A^\dagger$ 叫共轭转置——也叫 Hermitian 共轭}。

\textbf{厄米矩阵的两条黄金性质}：

\begin{enumerate}
\item \textbf{特征值全是实数}
\item \textbf{特征向量两两正交}（在复内积下）
\end{enumerate}

\textbf{这两条——就是量子力学能存在的数学基础}。

\textbf{为什么}？——\textbf{在量子力学里}：

\begin{itemize}
\item \textbf{一个粒子的状态——是 Hilbert 空间里的一个向量 $\bm{\psi}$}
\item \textbf{一个可观测量}（能量、动量、位置）——\textbf{是一个厄米算子 $\hat{H}$}
\item \textbf{测量这个量得到的值——必然是 $\hat{H}$ 的某个特征值 $\lambda_i$}
\item \textbf{测量后——状态会"塌缩"到对应的特征向量上}
\end{itemize}

\textbf{我们能测出来的物理量——必须是实数}（你不会量到"3+4i 米"）——\textbf{这要求算子的特征值是实数——所以算子必须是厄米的}。

\textbf{这是 20 世纪物理最深刻的发现之一}——\textbf{"测量 = 特征值"}——\textbf{由 von Neumann 在 Hilbert 空间里写得清清楚楚}。

\subsection{二次型与正定矩阵}

\textbf{二次型}（quadratic form）——\textbf{就是用一个对称矩阵把向量"夹起来"}：

\[
Q(\bm{x}) = \bm{x}^\top A \bm{x} = \sum_{i,j} a_{ij} x_i x_j
\]

\textbf{这是关于 $\bm{x}$ 的二次多项式}——\textbf{对应的几何图形是椭圆/双曲线/抛物线（在 2D）——椭球/双曲面（在 3D）}。

\textbf{根据 $Q(\bm{x})$ 的取值——对称矩阵 $A$ 被分成四类}：

\begin{itemize}
\item \textbf{正定}（positive definite）——\textbf{对所有 $\bm{x}\neq \bm{0}$——$Q(\bm{x}) > 0$}——\textbf{$\iff$ 所有特征值 $> 0$}
\item \textbf{半正定}——\textbf{$Q(\bm{x}) \geq 0$}——\textbf{$\iff$ 所有特征值 $\geq 0$}
\item \textbf{负定}——\textbf{$Q(\bm{x}) < 0$}——\textbf{$\iff$ 所有特征值 $< 0$}
\item \textbf{不定}——\textbf{有正有负——既上凸又下凹}
\end{itemize}

\textbf{为什么"正定"这么重要}？

\textbf{因为机器学习里的"损失函数 $L(\bm{\theta})$"——在最优点附近——它的 Hessian 矩阵就是一个对称矩阵}。

\textbf{Hessian 正定 $\iff$ 当前点是局部最小 $\iff$ 梯度下降会收敛}。

\textbf{Hessian 不定 $\iff$ 鞍点——梯度下降会卡住}。

\textbf{所以——整个凸优化理论 + 二阶优化方法（牛顿法、L-BFGS）——都建立在"正定矩阵"的几何上}。

\section{Aha 例子 / The Aha Moment}

\subsection{对角化 = 找最好的基}

\textbf{我学这一章——最大的 aha 时刻——是有一天晚上明白了下面这句话}：

\begin{tcolorbox}[ahabox={Aha 时刻}]
\textbf{矩阵 $A$ 本身不是"差的"——只是\textbf{你用了"差的基"来描述它}}。\\[4pt]
\textbf{对角化的真正含义}——\textbf{是换一组"好的基"（特征向量基）}——\textbf{在新基下——$A$ 看起来就是一个简单的对角阵}——\textbf{每一个特征向量方向独立缩放——互不干扰}。\\[4pt]
\textbf{In English}: The matrix isn't ugly —— your basis is. Change to the eigenvector basis, and the matrix becomes diagonal —— each axis acting independently.
\end{tcolorbox}

\textbf{打个比方}——\textbf{你用中文写一首英文诗——别扭}——\textbf{换成英文写——流畅}。\textbf{矩阵也一样——换到它"母语"（特征向量基）下——它就变成对角阵——最纯粹的样子}。

\subsection{Eigenfaces——人脸识别的特征向量}

\textbf{现在我们来到这一章最让人震撼的应用}——\textbf{人脸识别}。

\textbf{故事是这样的}——\textbf{1991 年——MIT 的 Turk 和 Pentland 发表了一篇论文}——\textbf{《Eigenfaces for Recognition》}——\textbf{第一次把"特征值分解"用在了人脸识别上}。

\textbf{他们的算法极简}：

\begin{enumerate}
\item \textbf{收集 $N$ 张人脸照片}——\textbf{每张 92×112 像素 = 10304 个像素}——\textbf{展平成 10304 维的向量}
\item \textbf{把 $N$ 个向量减去平均脸——得到"偏离平均"的 $N$ 个向量}
\item \textbf{算它们的协方差矩阵 $C$}（对称 + 半正定——\textbf{谱定理直接生效}）
\item \textbf{对 $C$ 做特征分解}——\textbf{取前 $k$ 个最大特征值对应的特征向量——叫 Eigenfaces}
\item \textbf{每张新人脸——都可以用这 $k$ 个 Eigenfaces 的线性组合来近似表示}
\end{enumerate}

\textbf{效果}：\textbf{10304 维的人脸数据——压缩到 50 维——识别率达到 95\%}——\textbf{1991 年的硬件上跑得动}。

\textbf{Eigenfaces 长什么样}？——\textbf{每一张都像一张"幽灵脸"}——\textbf{第一张大致是"平均脸的轮廓"}——\textbf{后面几张分别捕捉"鼻子的位置"、"眉毛的角度"、"光照的方向"、"嘴角的弧度"……}

\textbf{这些"特征"——不是人类设计的}——\textbf{是数据自己"长出来"的}——\textbf{是协方差矩阵的特征向量自动告诉我们的}：\textbf{"这就是人脸数据里最重要的几个方向！"}

\textbf{2012 年深度学习出来之前}——\textbf{Eigenfaces 是人脸识别的工业标准}——\textbf{Microsoft、Sony、Canon 的相机里"自动对焦人脸"——背后都是它}。

\textbf{今天——你的手机相册"按人物分组"——本质还是同一个数学}——\textbf{虽然外面包了一层 CNN}。

\textbf{Eigenfaces——是"特征向量"在 21 世纪最伟大的应用之一}。\textbf{它也是 PCA（主成分分析）的直接前身}——\textbf{第 8 章我们会专门讲 PCA}。

\subsection{Python 实战——把 Eigenfaces 跑出来看看}

\textbf{下面用 scikit-learn 的内置人脸数据集——亲手算一遍 Eigenfaces}：

\begin{lstlisting}[language=Python, caption={Eigenfaces 最小实现}]
import numpy as np
import matplotlib.pyplot as plt
from sklearn.datasets import fetch_olivetti_faces

# 1. 加载 400 张人脸（40 人 x 10 张），每张 64x64
faces = fetch_olivetti_faces()
X = faces.data           # shape: (400, 4096)
n_samples, n_features = X.shape
print(f"数据形状: {X.shape}")

# 2. 中心化：减去"平均脸"
mean_face = X.mean(axis=0)
X_centered = X - mean_face

# 3. 算协方差矩阵 (4096 x 4096) —— 太大，用 SVD 技巧
#    更高效：直接对 X_centered 做 SVD，奇异向量就是特征向量
U, S, Vt = np.linalg.svd(X_centered, full_matrices=False)

# 4. 前 16 个特征脸
eigenfaces = Vt[:16].reshape(16, 64, 64)

# 5. 画出来
fig, axes = plt.subplots(4, 4, figsize=(8, 8))
for i, ax in enumerate(axes.flat):
    ax.imshow(eigenfaces[i], cmap='gray')
    ax.set_title(f'Eigenface #{i+1}')
    ax.axis('off')
plt.suptitle('前 16 个 Eigenfaces —— 数据自己长出来的"主方向"')
plt.tight_layout()
plt.show()

# 6. 用前 k 个 Eigenfaces 重构一张人脸
def reconstruct(face_idx, k):
    coeffs = Vt[:k] @ X_centered[face_idx]
    return mean_face + coeffs @ Vt[:k]

# 比较 k=5, 20, 50, 200 的重构效果
fig, axes = plt.subplots(1, 5, figsize=(15, 3))
axes[0].imshow(X[0].reshape(64, 64), cmap='gray'); axes[0].set_title('原图')
for i, k in enumerate([5, 20, 50, 200]):
    axes[i+1].imshow(reconstruct(0, k).reshape(64, 64), cmap='gray')
    axes[i+1].set_title(f'k={k}')
[ax.axis('off') for ax in axes]
plt.show()
\end{lstlisting}

\textbf{你跑完会看到}——\textbf{$k=5$ 时人脸还很模糊}——\textbf{$k=50$ 时已经认得出是谁了}——\textbf{$k=200$ 时几乎和原图无异}。

\textbf{4096 维的人脸——50 维就够认人了}——\textbf{压缩了 80 倍——这就是特征值分解的"魔法"}。

\section{现代视角 / Modern Perspective}

\subsection{机器学习里的"特征值"无处不在}

\textbf{我列一张表——你看一下"特征值/特征向量"在现代 ML 里的覆盖面}：

\begin{center}
\begin{tabular}{|l|l|}
\hline
\textbf{应用} & \textbf{用到的特征分解} \\
\hline
PCA 主成分分析 & 协方差矩阵的特征分解 \\
PageRank & 转移矩阵的最大特征向量 \\
谱聚类 & 图 Laplacian 的最小特征向量 \\
Spectral Embedding & 图 Laplacian 的特征向量 \\
图神经网络 GCN & 图 Laplacian 的谱分解 \\
推荐系统 SVD++ & 评分矩阵的奇异值分解 \\
LSI 潜在语义索引 & 词-文档矩阵 SVD \\
Word2Vec & 共现矩阵的低秩分解 \\
Transformer 注意力 & QK 矩阵的低秩近似 \\
量子计算 & Hamiltonian 的特征分解 \\
\hline
\end{tabular}
\end{center}

\textbf{你看}——\textbf{几乎所有"提取重要方向"的算法——都是某种意义上的特征分解}。

\textbf{这就是为什么我说——\textbf{特征值分解是 21 世纪 ML 之根}}。

\subsection{SVD：特征分解的"普适版"}

\textbf{特征分解 $A = P\Lambda P^{-1}$ 有一个限制}——\textbf{$A$ 必须是方阵}。

\textbf{对于一般的 $m\times n$ 矩阵（比如 Netflix 的用户-电影评分矩阵）——我们用 SVD}：

\[
A = U \Sigma V^\top
\]

\textbf{其中 $U \in \mathbb{R}^{m\times m}$、$V \in \mathbb{R}^{n\times n}$ 都是正交矩阵}——\textbf{$\Sigma$ 是对角阵——对角元素叫"奇异值"}。

\textbf{SVD 是任意矩阵都能做的——它是特征分解的"普适版"}。

\textbf{第 8 章我们会专门讲 SVD 在推荐系统和 PCA 里的用法}——\textbf{这里只埋个种}。

\subsection{从 Hilbert 空间到 Transformer}

\textbf{有意思的是}——\textbf{今天最火的 Transformer 模型}——\textbf{它的"注意力机制"$\mathrm{softmax}(QK^\top/\sqrt{d})V$}——\textbf{本质上是在一个 $d$ 维向量空间里算"内积"——找"最相关"的 token}。

\textbf{这个"$d$ 维向量空间 + 内积"——就是有限维的 Hilbert 空间}。

\textbf{从 1904 年 Hilbert 写下第一个无穷维内积——到 2017 年 Vaswani 写下 Attention 公式——人类用了 113 年——但是同一个数学结构}。

\textbf{这就是数学的力量}——\textbf{一个好的抽象——可以跨越一个世纪——还活着}。

\section{代码与思考 / Code \& Reflection}

\textbf{把这一章的核心操作整理成一个 Python 模块}——\textbf{文件名 \texttt{linalg\_core.py}}：

\begin{lstlisting}[language=Python, caption={\texttt{linalg\_core.py} —— 第 7 章工具箱}]
"""
linalg_core.py
第 7 章工具：特征值、对角化、Gram-Schmidt、谱分解、Eigenfaces
作者：打通工科数学任督二脉 Ch.7
"""
import numpy as np
import matplotlib.pyplot as plt


# --- 1. 特征值 + 特征向量 ---------------------------------------------
def eig_demo(A):
    """求 A 的特征值和特征向量，并验证 Av = λv"""
    vals, vecs = np.linalg.eig(A)
    for i, lam in enumerate(vals):
        v = vecs[:, i]
        print(f"λ_{i} = {lam:.3f},  ||Av - λv|| = {np.linalg.norm(A@v - lam*v):.2e}")
    return vals, vecs


# --- 2. 对角化 + 矩阵幂 -----------------------------------------------
def matrix_power_via_diag(A, n):
    """用对角化算 A^n —— O(d^3) 而不是 O(n d^3)"""
    vals, P = np.linalg.eig(A)
    Lambda_n = np.diag(vals ** n)
    return P @ Lambda_n @ np.linalg.inv(P)


def fibonacci_via_eig(n):
    """用对角化算 Fibonacci 第 n 项"""
    A = np.array([[1, 1], [1, 0]], dtype=float)
    An = matrix_power_via_diag(A, n)
    return int(round(An[0, 1]))


# --- 3. Gram-Schmidt 正交化 -------------------------------------------
def gram_schmidt(A):
    """对 A 的列做 Gram-Schmidt，返回正交基 Q"""
    A = np.array(A, dtype=float)
    m, n = A.shape
    Q = np.zeros_like(A)
    for j in range(n):
        v = A[:, j].copy()
        for i in range(j):
            v -= (Q[:, i] @ A[:, j]) * Q[:, i]
        Q[:, j] = v / np.linalg.norm(v)
    return Q


# --- 4. 对称矩阵的谱分解 ----------------------------------------------
def spectral_decomposition(A):
    """A = Q Λ Q^T  (A 必须对称)"""
    assert np.allclose(A, A.T), "矩阵必须对称"
    vals, Q = np.linalg.eigh(A)   # eigh 专门给对称/厄米用
    return Q, np.diag(vals)


# --- 5. 厄米矩阵示例（Pauli 矩阵）-------------------------------------
def pauli_matrices():
    """量子力学里最有名的三个厄米矩阵"""
    sigma_x = np.array([[0, 1], [1, 0]], dtype=complex)
    sigma_y = np.array([[0, -1j], [1j, 0]], dtype=complex)
    sigma_z = np.array([[1, 0], [0, -1]], dtype=complex)
    return sigma_x, sigma_y, sigma_z


# --- 6. 二次型 + 正定性检测 -------------------------------------------
def is_positive_definite(A):
    """判断对称矩阵 A 是否正定"""
    if not np.allclose(A, A.T):
        return False
    vals = np.linalg.eigvalsh(A)
    return np.all(vals > 0)


# --- 7. Eigenfaces (PCA 雏形) -----------------------------------------
def eigenfaces(X, k=50):
    """
    X: (n_samples, n_features) 的人脸矩阵
    返回前 k 个 Eigenfaces 和平均脸
    """
    mean_face = X.mean(axis=0)
    Xc = X - mean_face
    U, S, Vt = np.linalg.svd(Xc, full_matrices=False)
    return Vt[:k], mean_face


if __name__ == "__main__":
    # 小测试
    A = np.array([[3., 1.], [1., 2.]])
    print("--- 特征分解 ---");  eig_demo(A)
    print("\n--- Fibonacci(30) =", fibonacci_via_eig(30))
    print("\n--- 正定性 ---", is_positive_definite(A))
\end{lstlisting}

\textbf{思考题}（自己动手）：

\begin{enumerate}
\item \textbf{Markov 链稳态}——\textbf{随便写一个 $3\times 3$ 的随机转移矩阵}（每列和为 1）——\textbf{用 \texttt{matrix\_power\_via\_diag}\;算它的 100 次幂}——\textbf{观察每一列是不是趋于同一个向量}（这就是稳态分布——也是它的最大特征向量）
\item \textbf{Gram-Schmidt 实验}——\textbf{对 $[(1,1,1)^\top, (1,2,3)^\top, (1,4,9)^\top]$ 做正交化}——\textbf{验证 $Q^\top Q = I$}
\item \textbf{Pauli 矩阵的特征值}——\textbf{用 \texttt{np.linalg.eigh} 算 $\sigma_x, \sigma_y, \sigma_z$ 的特征值——你会发现都是 $\pm 1$}——\textbf{这就是"自旋 1/2"的来源}
\item \textbf{Hessian 实验}——\textbf{写一个 2D 函数 $f(x,y) = 3x^2 + 2xy + 2y^2$——手算 Hessian 矩阵——用 \texttt{is\_positive\_definite} 验证它是正定的}——\textbf{画出 $f$ 的等高线——你会看到一个椭圆——主轴方向就是 Hessian 的特征向量}
\end{enumerate}

\section{本章小结 / Summary}

\begin{tcolorbox}[essbox={本章小结}]
\textbf{我们用六步——从"特征向量"走到"量子力学 + 机器学习"}：

\medskip

\textbf{一、特征值 + 特征向量}：\textbf{$A\bm{v}=\lambda\bm{v}$}——\textbf{矩阵不旋转、只缩放的方向}。

\medskip

\textbf{二、对角化}：\textbf{$A=P\Lambda P^{-1}$}——\textbf{换到特征向量基下——矩阵变成对角阵}——\textbf{$A^n = P\Lambda^n P^{-1}$——指数级加速}。

\medskip

\textbf{三、内积 + 正交基 + Gram-Schmidt}：\textbf{内积给空间装上"长度 + 角度"}——\textbf{Gram-Schmidt 把任意基加工成正交基}——\textbf{这是 QR 分解的根}。

\medskip

\textbf{四、对称矩阵 + 谱定理}：\textbf{$A^\top = A \Rightarrow A=Q\Lambda Q^\top$}——\textbf{特征值全实——特征向量两两正交}——\textbf{这是 PCA、Eigenfaces、振动模态、PageRank 的共同根}。

\medskip

\textbf{五、厄米矩阵}：\textbf{$A^\dagger = A$}——\textbf{对称矩阵的复数版}——\textbf{量子力学的所有"可观测量"都是它}——\textbf{特征值 = 测量值}。

\medskip

\textbf{六、二次型 + 正定矩阵}：\textbf{$\bm{x}^\top A \bm{x}$}——\textbf{正定 $\iff$ 所有特征值 $>0$}——\textbf{凸优化、Hessian 检验、损失函数局部最小的判据}。

\medskip

\textbf{一句话总结}：\textbf{特征值分解 = 找最好的基 = 21 世纪 ML 之根}。

\medskip

\textbf{In one line}: Eigen-decomposition is the art of choosing the right basis —— the root of modern machine learning.
\end{tcolorbox}

\textbf{下一章}（第 8 章）——\textbf{我们把这一章的特征分解 + SVD——直接拿到现实数据里跑}——\textbf{PCA 降维、PageRank 排序、神经网络的反向传播}——\textbf{把"线代三部曲"画上句号}。

\textbf{From Cayley's 17-page memoir in 1858 —— to Hilbert's infinite-dimensional space in 1904 —— to Heisenberg's matrix mechanics in 1925 —— to Turk's Eigenfaces in 1991 —— to today's Transformers}——\textbf{the same idea has been quietly running underneath, for 168 years: \emph{find the eigenvectors, and the rest is easy}}.

\clearpage

% =====================================================================
%  第 8 章 · 线代应用 / Linear Algebra III
%  SVD、PCA、最小二乘、Markov 链 + PageRank、矩阵分解、神经网络、量子比特
%  小故事：Larry Page 的博士论文与 1.5 万亿美元的矩阵（1998）
%  Aha：PageRank——一个特征向量是如何统治互联网的
% =====================================================================

\chapter{线代应用 / Linear Algebra III}
\label{ch:linalg3}

\begin{quote}\itshape
1998 年——斯坦福大学计算机系——一个 25 岁的博士生坐在 Gates Hall 三楼的小办公室里。\\
他叫 \textbf{Larry Page}。\\
他的博士题目本来是\textbf{"超文本链接的逆向引用结构"}——\textbf{听起来很无聊}。\\
但他写下了一个方程——一个看起来人畜无害的特征向量方程——\\
\[
\bm{r} = M \bm{r}
\]
\\[6pt]
\textbf{20 年后}——这个特征向量值\textbf{1.5 万亿美元}——\textbf{它的名字叫 Google}。
\end{quote}

\section{写在前面 / Preface}

\textbf{这是"线性代数三部曲"的第三章——也是整本书的高潮章节之一}。

第 6 章我们学了\textbf{矩阵 = 线性变换}——\textbf{是骨架}。\\
第 7 章我们学了\textbf{特征值、特征向量、SVD、对角化}——\textbf{是内功}。\\
\textbf{这一章——我们终于要让线代"出招"}——\textbf{看它如何变成 PageRank、PCA、神经网络、量子计算}。

\textbf{为什么把"应用"单独写一章}？因为\textbf{在我的工科朋友里}——\textbf{90\% 的人会算特征值——但 50\% 的人不知道特征值能干什么}——\textbf{这是中国线代教育最大的断层}。

\textbf{后来我做机器人感知}——\textbf{每一篇论文里出现的"奇异值分解"——其实是大三那道"求 SVD"题目的真实模样}。\textbf{在那一刻——题目变成了刀——线代变成了武功}。

\textbf{这一章的七个应用——是工科 + ML + 物理学共用的"七把武器"}：

\begin{enumerate}
\item \textbf{SVD}——\textbf{最有用的矩阵分解——没有之一}
\item \textbf{PCA}——\textbf{数据科学的"根"}
\item \textbf{最小二乘 + 投影}——\textbf{回归与机器学习的最古老形式}
\item \textbf{Markov 链 + PageRank}——\textbf{1.5 万亿美元的特征向量}
\item \textbf{矩阵分解（LU、QR、Cholesky）}——\textbf{数值计算的瑞士军刀}
\item \textbf{神经网络与线代}——\textbf{前馈 = 矩阵乘法的链}
\item \textbf{量子比特 + 量子门}——\textbf{物理学和线代的握手}
\end{enumerate}

\textbf{读完这一章——你会明白}：\textbf{线代不是给考研用的——它是给"造世界"用的}。

\section{一句话本质 / The Essence in One Sentence}

\begin{tcolorbox}[essbox={一句话本质}]
\textbf{中文}：\textbf{现代世界的"大数据 + AI + 物理引擎"——本质上都是\textbf{在巨大的矩阵上反复做"分解 + 投影 + 迭代"}}——\textbf{而线代——就是这一切的语言}。\\[6pt]
\textbf{English}: The modern world of big data, AI and physics engines is essentially \emph{decomposing, projecting and iterating on enormous matrices} —— and linear algebra is the language in which it speaks.
\end{tcolorbox}

\textbf{记住一个公式}——\textbf{记住整个 21 世纪的计算}：

\[
A = U \Sigma V^{\!\top}
\]

\textbf{这是 SVD}——\textbf{第 7 章我们已经见过它了}——\textbf{这一章我们要看它如何"统治世界"}。

\section{教材里你看到的 / What the Textbook Tells You}

\subsection{中国教材的"应用空白"}

\textbf{同济《线性代数》第 6 版——一共 200 页}。\textbf{我把它从头翻到尾}：\textbf{"应用"两个字出现 0 次}——\textbf{"PCA"出现 0 次}——\textbf{"SVD"出现 0 次}——\textbf{"PageRank"出现 0 次}。

\textbf{它的最后一章}——\textbf{"二次型"}——\textbf{讲了 30 页矩阵的合同变换}——\textbf{但是没有一句话告诉你}：\textbf{二次型 = 一个椭球的方程——也是 PCA 的几何核心}。

\textbf{这就是中国线代教育的"两张皮"}——\textbf{学的全是技巧——用的全在外面}。

\subsection{Strang《Linear Algebra and Learning from Data》}

\textbf{2019 年——Gilbert Strang 84 岁那年}——\textbf{写了一本新书}——\textbf{《Linear Algebra and Learning from Data》}。

\textbf{他在前言里写}：\textbf{"我教了 60 年线性代数——直到机器学习出现——我才看到线代的全部威力。"}

\textbf{这本书的结构}：\textbf{第一部分讲 SVD}——\textbf{第二部分讲优化与神经网络}——\textbf{第三部分讲概率与统计}。\textbf{300 页——全部都是"线代的应用"}。

\textbf{这是世界顶级教授对中国线代教材的隐性回答}——\textbf{线代的灵魂不在行列式——在 SVD + 应用}。

\subsection{Trefethen \& Bau《Numerical Linear Algebra》}

\textbf{Cornell + Oxford 两位教授合写}——\textbf{这是数值线性代数的圣经}。

\textbf{开篇第 1 节}——\textbf{讲的就是 SVD}——\textbf{第 1 句话}：\textbf{"Singular Value Decomposition is the most important matrix decomposition."}

\textbf{这本书里}——\textbf{LU、QR、Cholesky}——\textbf{都不是\textbf{算法的细节}}——\textbf{是\textbf{"如何让浮点数不爆掉"的工程艺术}}。

\subsection{我的策略}

\textbf{本章不讲行列式——也不讲考研技巧}。\textbf{本章只做一件事}：\textbf{把第 7 章的特征值/SVD 全部"激活"——变成 7 个真实工程系统}。

\textbf{每一节——都是\textbf{"理论 1 页 + 公式 1 页 + 代码 1 页"}的三件套}。

\begin{tcolorbox}[storybox={\Large 小故事：Larry Page 与 1.5 万亿美元的矩阵（1998）}]
\textbf{1995 年秋天——斯坦福大学计算机系新生 orientation}——\textbf{一个 22 岁的密歇根男孩 Larry Page} 见到了带他参观校园的研究生 Sergey Brin。\textbf{两个人吵了一路}——\textbf{谁都觉得对方讨厌}。

\medskip

\textbf{两年后——他们成了搭档}——\textbf{合写一篇论文}——\textbf{题目叫《The Anatomy of a Large-Scale Hypertextual Web Search Engine》}——\textbf{这就是 Google 的诞生论文}。

\medskip

\textbf{核心问题}：\textbf{1998 年的搜索引擎——AltaVista、Yahoo、Lycos}——\textbf{排序的方法是"关键词匹配次数"}——\textbf{结果}：\textbf{你搜"汽车"——第一页全是堆满"汽车汽车汽车"的垃圾网页}。

\medskip

\textbf{Larry Page 的想法}：\textbf{用网页之间的"链接结构"——给每个网页打一个分}——\textbf{这个分——他叫 PageRank}（注：\textbf{"Page" 不是网页的意思——是他自己的姓}）。

\medskip

\textbf{核心思想——一句话}：\textbf{"一个网页的重要性——取决于多少重要的网页指向它"}。

\medskip

\textbf{这是一个递归定义}——\textbf{看起来无解}——\textbf{但 Larry Page 写下了它的数学形式}：

\[
\bm{r} = M \bm{r}
\]

\medskip

\textbf{这是一个特征向量方程}——\textbf{特征值 = 1}——\textbf{$M$ 是\textbf{网页之间链接关系构成的转移矩阵}}——\textbf{$\bm{r}$ 是\textbf{每个网页的"重要性分数"}}。

\medskip

\textbf{Larry Page 的这一念之差}——\textbf{把一个工程问题——变成了一个特征向量问题}——\textbf{把"我有 1 亿个网页该怎么排序"——变成了"我求这个 1 亿 × 1 亿矩阵最大特征值对应的特征向量"}。

\medskip

\textbf{1998 年 9 月 4 日——Google 公司在朋友家的车库里成立}——\textbf{创业资本 10 万美元——一台二手服务器}。

\medskip

\textbf{2004 年 IPO}——\textbf{市值 230 亿美元}。\\
\textbf{2010 年}——\textbf{1900 亿美元}。\\
\textbf{2020 年}——\textbf{1 万亿美元}。\\
\textbf{2024 年}——\textbf{2 万亿美元}（峰值）。

\medskip

\textbf{Larry Page 的博士论文最终没有完成}——\textbf{他在 1998 年休学创业}——\textbf{斯坦福至今欠他一个 PhD}——\textbf{但他用一个特征向量——造了一家\textbf{1.5 万亿美元的公司}}。

\medskip

\textbf{这是数学史上回报率最高的特征向量}——\textbf{也是线性代数最美的"出招"}。

\medskip

\textbf{这一章——我们要复现它}——\textbf{用 30 行 Python 代码——给一个小型互联网算 PageRank}。
\end{tcolorbox}

\section{直觉 / Intuition}

\subsection{SVD：把任何矩阵"拆成三件"}

\textbf{SVD 的本质——一句话}：\textbf{任何矩阵——都可以拆成"旋转 + 拉伸 + 旋转"}。

\[
A = U \Sigma V^{\!\top}
\]

\textbf{其中}：
\begin{itemize}
\item \textbf{$V^{\!\top}$}——\textbf{第一次旋转}（把输入空间转到一个"好"的方向）
\item \textbf{$\Sigma$}——\textbf{沿每个坐标轴拉伸}（拉伸的倍数 = 奇异值）
\item \textbf{$U$}——\textbf{第二次旋转}（把结果转到输出空间）
\end{itemize}

\textbf{为什么这件事重要}？因为\textbf{奇异值 $\sigma_1 \ge \sigma_2 \ge \cdots \ge \sigma_r$ 告诉你"矩阵的能量分布"}——\textbf{前几个大——后面的小或者 0}——\textbf{这就是"压缩"的物理来源}。

\textbf{记住一句话}：\textbf{SVD 是矩阵的"DNA 检测"——它告诉你矩阵的"主成分"是什么}。

\subsection{PCA：找到数据"最胖的方向"}

\textbf{PCA 的本质——一句话}：\textbf{在数据云里——找一条"穿过去最胖"的方向——把数据投到那条方向上}。

\textbf{为什么这件事重要}？因为\textbf{人脸有 10000 个像素——但"人脸的本质"可能只有 50 个方向}（眼睛大小、鼻子高度、肤色\dots）——\textbf{PCA 帮你把 10000 维压成 50 维——精度损失 5\% 不到}。

\textbf{PCA 的数学一句话}：\textbf{它就是数据协方差矩阵的 SVD}。\textbf{所以——SVD 是 PCA 的"母亲"}。

\subsection{最小二乘：找一条"离所有点都最近"的线}

\textbf{你大一物理实验}——\textbf{画 $y = ax + b$ 拟合曲线}——\textbf{你按的是 Excel 的 LINEST 函数}——\textbf{它背后是最小二乘法}：

\[
\min_{\bm{x}} \|A \bm{x} - \bm{b}\|^2
\]

\textbf{几何含义}——\textbf{把 $\bm{b}$ 投影到 $A$ 的列空间上}——\textbf{投影点 = 最佳拟合}——\textbf{投影误差 = 最小残差}。

\textbf{结果一行公式}：

\[
\bm{x}^* = (A^{\!\top} A)^{-1} A^{\!\top} \bm{b}
\]

\textbf{这就是回归分析的"祖宗公式"}——\textbf{从 Gauss 1809 年到今天——200 多年没变过}。

\subsection{Markov 链：状态在矩阵里"散步"}

\textbf{Markov 链的本质}——\textbf{当前状态只取决于上一步——和更早的状态无关}（无记忆性）。\textbf{每一步——状态向量乘上一个"转移矩阵"}：

\[
\bm{p}_{t+1} = M \bm{p}_t
\]

\textbf{反复迭代}——\textbf{状态向量收敛到一个"平稳分布"}——\textbf{这个平稳分布 = $M$ 的特征值 1 对应的特征向量}。

\textbf{PageRank——就是把"互联网"当成一条 Markov 链}——\textbf{用户在网页之间随机游走}——\textbf{平稳分布 = 每个网页的访问概率 = 网页的重要性}。

\subsection{神经网络：矩阵乘法的链}

\textbf{ChatGPT 的核心 Transformer——本质上是}：

\[
\bm{y} = f_n(W_n \, f_{n-1}(W_{n-1} \, \cdots f_1(W_1 \bm{x} + \bm{b}_1) \cdots + \bm{b}_{n-1}) + \bm{b}_n)
\]

\textbf{把激活函数 $f$ 暂时遮起来}——\textbf{它就是 $n$ 个矩阵的连乘}——\textbf{$W_n W_{n-1} \cdots W_1 \bm{x}$}。

\textbf{所有的"深度学习"——本质上都是"高维线性变换 + 一点点非线性"——堆叠几十层}。\textbf{99\% 的计算量——是矩阵乘法}。

\subsection{量子比特：单位球面上的复向量}

\textbf{一个量子比特 = $\mathbb{C}^2$ 里一个长度为 1 的复向量}：

\[
|\psi\rangle = \alpha |0\rangle + \beta |1\rangle, \quad |\alpha|^2 + |\beta|^2 = 1
\]

\textbf{量子门 = $\mathbb{C}^2$ 上的酉矩阵}（$U^\dagger U = I$）——\textbf{它把一个量子比特变成另一个量子比特}——\textbf{不增加信息——也不丢失信息}。

\textbf{量子计算的本质}——\textbf{一长串酉矩阵乘上初始状态 $|0\rangle$}——\textbf{最后做一次"测量"得到 0 或 1}。\textbf{这就是 IBM Q、Google Sycamore 的全部数学}。

\section{数学 / The Math}

\subsection{SVD 的精确陈述}

\begin{theorem}[SVD 定理]
设 $A \in \mathbb{R}^{m \times n}$——\textbf{rank} $= r$。\textbf{存在}：
\begin{itemize}
\item \textbf{$U \in \mathbb{R}^{m \times m}$}——\textbf{正交矩阵}（$U^{\!\top} U = I_m$）
\item \textbf{$V \in \mathbb{R}^{n \times n}$}——\textbf{正交矩阵}（$V^{\!\top} V = I_n$）
\item \textbf{$\Sigma \in \mathbb{R}^{m \times n}$}——\textbf{对角矩阵——对角元 $\sigma_1 \ge \sigma_2 \ge \cdots \ge \sigma_r > 0$——其余为 0}
\end{itemize}
使得
\[
A = U \Sigma V^{\!\top}
\]
$\sigma_i$ 称为\textbf{奇异值}——\textbf{$U$ 的列称为左奇异向量}——\textbf{$V$ 的列称为右奇异向量}。
\end{theorem}

\textbf{关键事实}：
\begin{itemize}
\item \textbf{$\sigma_i^2 = $ $A^{\!\top} A$ 的特征值}——\textbf{所以 SVD = 对称矩阵 $A^{\!\top} A$ 的特征分解的"升级版"}
\item \textbf{$V$ 的列 = $A^{\!\top} A$ 的特征向量}
\item \textbf{$U$ 的列 = $A A^{\!\top}$ 的特征向量}
\item \textbf{秩 $r$ 截断}：\textbf{把 $\Sigma$ 只保留前 $k$ 个奇异值——得到 $A$ 的"最佳秩-$k$ 近似"}（Eckart--Young 定理）
\end{itemize}

\subsection{Eckart--Young 定理（SVD 的"压缩定理"）}

\begin{theorem}[Eckart--Young, 1936]
设 $A = U \Sigma V^{\!\top}$——\textbf{令}
\[
A_k = \sum_{i=1}^{k} \sigma_i \bm{u}_i \bm{v}_i^{\!\top}
\]
\textbf{则 $A_k$ 是所有秩不超过 $k$ 的矩阵中——离 $A$ 最近的}（Frobenius 范数下）：
\[
\|A - A_k\|_F = \min_{\mathrm{rank}(B) \le k} \|A - B\|_F = \sqrt{\sigma_{k+1}^2 + \cdots + \sigma_r^2}
\]
\end{theorem}

\textbf{这一句话——是图像压缩、推荐系统、LoRA 微调的数学基础}。

\subsection{PCA 的精确推导}

\textbf{设数据矩阵 $X \in \mathbb{R}^{n \times d}$}——\textbf{每行一个样本——共 $n$ 个样本——每个 $d$ 维}。\textbf{先减去均值得到中心化矩阵 $\tilde{X}$}。

\textbf{协方差矩阵}：

\[
C = \frac{1}{n-1} \tilde{X}^{\!\top} \tilde{X}
\]

\textbf{$C$ 是对称半正定矩阵}——\textbf{它的特征分解}：

\[
C = V \Lambda V^{\!\top}
\]

\textbf{$\Lambda = \mathrm{diag}(\lambda_1, \lambda_2, \dots, \lambda_d)$, $\lambda_1 \ge \lambda_2 \ge \cdots \ge \lambda_d \ge 0$}。

\textbf{$V$ 的列称为主成分方向}——\textbf{$\lambda_i$ = 第 $i$ 个主成分上的方差}。

\textbf{保留前 $k$ 个主成分}——\textbf{得到降维表示}：

\[
Z = \tilde{X} V_k \in \mathbb{R}^{n \times k}
\]

\textbf{解释方差比例}：

\[
\text{Explained variance ratio} = \frac{\lambda_1 + \cdots + \lambda_k}{\lambda_1 + \cdots + \lambda_d}
\]

\textbf{通常我们选 $k$ 使得比例 $\ge 95\%$}——\textbf{你就压缩了几十倍——丢失不到 5\%}。

\subsection{最小二乘 = 正交投影}

\textbf{问题}：\textbf{$A \in \mathbb{R}^{m \times n}$，$m \gg n$（行多于列——超定方程组）}——\textbf{$A\bm{x} = \bm{b}$ 一般无解}。

\textbf{最小二乘解}：

\[
\bm{x}^* = \arg\min_{\bm{x}} \|A \bm{x} - \bm{b}\|^2
\]

\textbf{法方程}：

\[
A^{\!\top} A \bm{x}^* = A^{\!\top} \bm{b} \quad \Longrightarrow \quad \bm{x}^* = (A^{\!\top} A)^{-1} A^{\!\top} \bm{b}
\]

\textbf{几何解释}——\textbf{$A \bm{x}^*$ 是 $\bm{b}$ 在 $A$ 的列空间上的正交投影}——\textbf{投影矩阵}：

\[
P = A (A^{\!\top} A)^{-1} A^{\!\top}
\]

\textbf{$P$ 满足 $P^2 = P$, $P^{\!\top} = P$——这是正交投影的两个标志}。

\subsection{Markov 链与平稳分布}

\textbf{$M \in \mathbb{R}^{n \times n}$——列随机矩阵}（每列和 = 1——元素非负）——\textbf{则}：

\begin{itemize}
\item \textbf{$M$ 的最大特征值 = 1}（Perron--Frobenius 定理）
\item \textbf{对应的特征向量 $\bm{\pi}$（归一化）就是平稳分布}：$M \bm{\pi} = \bm{\pi}$
\item \textbf{若 $M$ 不可约且非周期}——\textbf{$\bm{\pi}$ 唯一}——\textbf{任何初始分布 $\bm{p}_0$ 经 $\bm{p}_t = M^t \bm{p}_0$ 都收敛到 $\bm{\pi}$}
\end{itemize}

\subsection{PageRank 的数学公式}

\textbf{设 $G$ = 互联网的链接图}——\textbf{$n$ 个网页}——\textbf{$L_{ij} = 1$ 若网页 $j$ 指向网页 $i$}——\textbf{否则 0}。

\textbf{令 $c_j = $ 网页 $j$ 的出链数}——\textbf{则转移矩阵}：

\[
M_{ij} = \frac{L_{ij}}{c_j}
\]

\textbf{这个 $M$ 是列随机矩阵}——\textbf{但存在两个问题}：
\begin{itemize}
\item \textbf{Dead end}（出链为 0 的网页）——\textbf{概率会"漏出去"}
\item \textbf{Spider trap}（自闭环）——\textbf{用户被困住}
\end{itemize}

\textbf{Larry Page 的天才修正}——\textbf{引入"传送概率" $d$（damping factor）}：

\[
\hat{M} = d M + \frac{1-d}{n} \bm{1} \bm{1}^{\!\top}
\]

\textbf{$d = 0.85$}——\textbf{意思}：\textbf{用户 85\% 概率沿链接点击——15\% 概率随机跳到一个网页}。

\textbf{PageRank 方程}：

\[
\bm{r} = \hat{M} \bm{r}
\]

\textbf{$\bm{r}$ = $\hat{M}$ 最大特征值（=1）对应的特征向量}——\textbf{这就是 PageRank 向量}。

\textbf{数值求法——幂迭代}：

\[
\bm{r}_{t+1} = \hat{M} \bm{r}_t
\]

\textbf{通常 30--50 次迭代即收敛}——\textbf{即便对几十亿个网页}。

\subsection{矩阵分解三件套}

\textbf{LU 分解}（Gauss 消元的"现代化"）：

\[
A = LU
\]

\textbf{$L$ 是下三角——$U$ 是上三角}——\textbf{用于解 $A \bm{x} = \bm{b}$}（先解 $L \bm{y} = \bm{b}$——再解 $U \bm{x} = \bm{y}$）。\textbf{复杂度 $O(n^3)$——但解多个右端项时只需 $O(n^2)$ 一次}。

\medskip

\textbf{QR 分解}（最小二乘的"数值核心"）：

\[
A = QR
\]

\textbf{$Q$ 是正交矩阵——$R$ 是上三角}——\textbf{用于最小二乘}：\textbf{$R \bm{x}^* = Q^{\!\top} \bm{b}$}——\textbf{比 $(A^{\!\top} A)^{-1}$ 数值上稳定得多}。

\medskip

\textbf{Cholesky 分解}（对称正定矩阵专用——速度最快）：

\[
A = L L^{\!\top}
\]

\textbf{要求 $A$ 对称正定}——\textbf{复杂度 $\frac{1}{2} n^3$——是 LU 的一半}。\textbf{在统计学（协方差矩阵）和优化（Hessian 矩阵）里用得最多}。

\subsection{神经网络的线代写法}

\textbf{一个 $L$ 层全连接网络}：

\[
\begin{aligned}
\bm{a}^{(0)} &= \bm{x} \\
\bm{z}^{(\ell)} &= W^{(\ell)} \bm{a}^{(\ell-1)} + \bm{b}^{(\ell)}, \quad \ell = 1, \dots, L \\
\bm{a}^{(\ell)} &= \sigma(\bm{z}^{(\ell)}) \\
\hat{\bm{y}} &= \bm{a}^{(L)}
\end{aligned}
\]

\textbf{反向传播 = 链式法则在矩阵上的写法}：

\[
\frac{\partial \mathcal{L}}{\partial W^{(\ell)}} = \bm{\delta}^{(\ell)} (\bm{a}^{(\ell-1)})^{\!\top}, \quad \bm{\delta}^{(\ell-1)} = (W^{(\ell)})^{\!\top} \bm{\delta}^{(\ell)} \odot \sigma'(\bm{z}^{(\ell-1)})
\]

\textbf{这就是为什么深度学习的训练 = 大量矩阵乘法}——\textbf{每次 forward + backward 都是$2L$ 次大矩阵乘}。

\subsection{量子门的酉矩阵}

\textbf{常见的单比特门}：

\[
H = \frac{1}{\sqrt{2}} \begin{pmatrix} 1 & 1 \\ 1 & -1 \end{pmatrix} \quad \text{(Hadamard——叠加门)}
\]

\[
X = \begin{pmatrix} 0 & 1 \\ 1 & 0 \end{pmatrix}, \quad Y = \begin{pmatrix} 0 & -i \\ i & 0 \end{pmatrix}, \quad Z = \begin{pmatrix} 1 & 0 \\ 0 & -1 \end{pmatrix} \quad \text{(Pauli 矩阵)}
\]

\textbf{两比特门——CNOT}：

\[
\mathrm{CNOT} = \begin{pmatrix} 1 & 0 & 0 & 0 \\ 0 & 1 & 0 & 0 \\ 0 & 0 & 0 & 1 \\ 0 & 0 & 1 & 0 \end{pmatrix}
\]

\textbf{所有量子算法（Shor、Grover）——都是\textbf{这些酉矩阵的乘积} + \textbf{最后一次测量}}。

\section{Aha 例子：PageRank 与 PCA / Aha: PageRank and PCA}

\subsection{Aha 例子 1：用 30 行代码复现 Google}

\textbf{我们构造一个"迷你互联网"——5 个网页 A, B, C, D, E}——\textbf{链接结构}：

\begin{itemize}
\item \textbf{A} 指向 B, C, D
\item \textbf{B} 指向 C
\item \textbf{C} 指向 A, B
\item \textbf{D} 指向 A, B, C
\item \textbf{E} 指向 A, B, C, D
\end{itemize}

\textbf{对应的链接矩阵 $L$（行是被指向——列是源）}：

\[
L = \begin{pmatrix}
0 & 0 & 1 & 1 & 1 \\
1 & 0 & 1 & 1 & 1 \\
1 & 1 & 0 & 1 & 1 \\
1 & 0 & 0 & 0 & 1 \\
0 & 0 & 0 & 0 & 0
\end{pmatrix}
\]

\textbf{出链数}：\textbf{$c_A = 3, c_B = 1, c_C = 2, c_D = 3, c_E = 4$}。

\textbf{转移矩阵 $M_{ij} = L_{ij} / c_j$}：

\[
M = \begin{pmatrix}
0 & 0 & 1/2 & 1/3 & 1/4 \\
1/3 & 0 & 1/2 & 1/3 & 1/4 \\
1/3 & 1 & 0 & 1/3 & 1/4 \\
1/3 & 0 & 0 & 0 & 1/4 \\
0 & 0 & 0 & 0 & 0
\end{pmatrix}
\]

\textbf{加上 damping}——\textbf{$\hat{M} = 0.85 M + 0.15/5 \cdot \bm{1} \bm{1}^{\!\top}$}。

\textbf{初始 $\bm{r}_0 = (0.2, 0.2, 0.2, 0.2, 0.2)^{\!\top}$}——\textbf{迭代 30 次}——\textbf{结果}：

\[
\bm{r} \approx (0.27, 0.30, 0.27, 0.13, 0.03)^{\!\top}
\]

\textbf{解读}：\textbf{B 最重要}（被多个高分网页指向）——\textbf{E 最不重要}（没人指它）——\textbf{尽管 E 自己很"勤奋"指向所有人}。

\textbf{这就是 PageRank 的灵魂}——\textbf{"重要性来自被引用——不来自自己叫嚣"}。\textbf{这也是为什么垃圾网页堆关键词没用}——\textbf{它没有"被引用"的资本}。

\begin{tcolorbox}[ahabox={\Large Aha：PageRank = 互联网的特征向量}]
\textbf{1 亿个网页的 PageRank}——\textbf{就是一个 1 亿 × 1 亿 矩阵的最大特征向量}。

\medskip

\textbf{1998 年——Larry Page 用 Sun 工作站算这个特征向量}——\textbf{算了一周}。\\
\textbf{2024 年——Google 每天重新算一次}——\textbf{用 TPU 集群——几分钟搞定}。

\medskip

\textbf{一个特征向量}——\textbf{一个数学概念}——\textbf{催生了人类历史上最大的信息公司}——\textbf{这是线性代数最美的"出招"}。
\end{tcolorbox}

\subsection{Aha 例子 2：PCA 把人脸压成 50 维}

\textbf{Olivetti 人脸数据集}——\textbf{40 个人——每人 10 张人脸照片}——\textbf{每张 $64 \times 64 = 4096$ 像素}。

\textbf{步骤}：

\begin{enumerate}
\item \textbf{把 400 张人脸展平成 $X \in \mathbb{R}^{400 \times 4096}$}
\item \textbf{中心化}——\textbf{每列减去均值}
\item \textbf{算 SVD}——\textbf{$X = U \Sigma V^{\!\top}$}
\item \textbf{保留前 $k = 50$ 个主成分}——\textbf{$V_{50} \in \mathbb{R}^{4096 \times 50}$}
\item \textbf{投影}——\textbf{$Z = X V_{50} \in \mathbb{R}^{400 \times 50}$}（每张脸用 50 个数表示）
\item \textbf{重建}——\textbf{$\hat{X} = Z V_{50}^{\!\top}$}——\textbf{把 50 维数据"还原"成 4096 像素的人脸}
\end{enumerate}

\textbf{结果}：\textbf{压缩比 = 4096 / 50 ≈ 80 倍}——\textbf{人脸视觉上几乎看不出差别}——\textbf{解释方差 > 90\%}。

\textbf{$V_{50}$ 的每一列——本身可以"画出来"}——\textbf{你会看到 50 张"特征脸"（Eigenfaces）}——\textbf{它们是人脸空间的"标准基"}——\textbf{1991 年 Turk \& Pentland 发明 Eigenfaces——是计算机视觉的早期里程碑}。

\textbf{这就是 PCA 的力量}——\textbf{它告诉你}：\textbf{4096 维的人脸——本质上活在 50 维子空间里}。

\section{现代视角：线代统治数字世界 / Modern: Linear Algebra Rules the Digital World}

\subsection{推荐系统 = 矩阵分解}

\textbf{Netflix Prize（2006--2009）}——\textbf{100 万美元奖金}——\textbf{题目}：\textbf{给定一个"用户 × 电影"评分矩阵 $R$（绝大多数元素缺失）——预测缺失的评分}。

\textbf{冠军方法的核心}：

\[
R \approx U V^{\!\top}, \quad U \in \mathbb{R}^{n \times k}, V \in \mathbb{R}^{m \times k}, \quad k \ll \min(n, m)
\]

\textbf{这就是\textbf{低秩近似}——\textbf{是 SVD 的"带缺失值"版本}}。\textbf{$U$ 的每一行是\textbf{用户的"潜在偏好向量"}}——\textbf{$V$ 的每一行是\textbf{电影的"潜在风格向量"}}——\textbf{两者点乘 = 预测评分}。

\textbf{2024 年的 TikTok、抖音、YouTube 推荐}——\textbf{都是这个思路的"深度学习增强版"}——\textbf{核心还是低秩矩阵分解}。

\subsection{大模型与 LoRA = SVD 的实战}

\textbf{2022 年 Microsoft 提出 LoRA（Low-Rank Adaptation）}——\textbf{微调大模型的革命性方法}。

\textbf{核心想法}——\textbf{大模型权重 $W \in \mathbb{R}^{d \times d}$ 太大}——\textbf{微调时用一个低秩更新}：

\[
W' = W + \Delta W, \quad \Delta W = A B^{\!\top}, \quad A \in \mathbb{R}^{d \times r}, B \in \mathbb{R}^{d \times r}, r \ll d
\]

\textbf{典型 $d = 4096, r = 8$}——\textbf{参数量从 $1600 万$ 降到 $6.6 万$}——\textbf{压缩 250 倍——精度损失 < 1\%}。

\textbf{LoRA 的合法性——来自 Eckart--Young 定理}——\textbf{低秩近似的精度可控}。\textbf{现在所有的 ChatGLM、Llama、Mistral 微调——都用 LoRA}。

\subsection{自动驾驶 = 大型最小二乘}

\textbf{Tesla FSD、Waymo、华为 ADS}——\textbf{车上的激光雷达、摄像头、IMU、GPS}——\textbf{每秒产生 GB 级数据}——\textbf{最终都要解一个最小二乘问题}：

\[
\min \|H \bm{x} - \bm{z}\|^2
\]

\textbf{$\bm{x}$ = 车的位置 + 姿态 + 速度}——\textbf{$\bm{z}$ = 所有传感器的观测}——\textbf{$H$ = 几何关系矩阵}。\textbf{每秒解几百次}——\textbf{用 QR 分解和卡尔曼滤波}。

\textbf{你在高速公路上 120 km/h 巡航——背后是 Gauss 1809 年的最小二乘法每秒被解 1000 次}。

\subsection{量子计算 = 巨型酉矩阵的乘法}

\textbf{Google Sycamore 量子处理器（2019）}——\textbf{53 个量子比特}——\textbf{状态空间维度 = $2^{53} \approx 9 \times 10^{15}$}。

\textbf{量子算法 = 一系列 $2^{53} \times 2^{53}$ 酉矩阵的乘法}——\textbf{这是经典计算机绝对算不动的事}——\textbf{但物理本身可以"算"——因为宇宙就是按这个规则演化的}。

\textbf{Shor 算法分解大整数}——\textbf{Grover 算法搜索无序数据库}——\textbf{都是"酉矩阵乘法 + 测量"}——\textbf{量子计算的全部数学语言——就是线性代数}。

\subsection{结论：现代计算 = 巨型矩阵的"减法艺术"}

\textbf{我把 21 世纪的计算总结成一句话}——\textbf{现代算力 = 用"低秩"逼近"高维"}——\textbf{用 SVD、PCA、矩阵分解、LoRA、神经网络剪枝——把"看似无穷大"的问题压成"工程上可算"的问题}。

\textbf{所以——线代不是"过时的工具"}——\textbf{它是"减法的艺术"}——\textbf{是计算的"瑞士军刀"}。

\section{代码与思考 / Code and Exercises}

\subsection{配套模块 \texttt{linalg\_apps.py}}

GitHub 上的 \texttt{modules/linalg\_apps.py} 包含\textbf{7 个 demo}——\textbf{对应本章 7 个应用}：

\begin{enumerate}
\item \texttt{demo\_svd\_image\_compress()}——\textbf{用 SVD 压缩一张图片}——\textbf{保留前 $k$ 个奇异值——画出压缩比 vs 失真曲线}
\item \texttt{demo\_pca\_faces()}——\textbf{Olivetti 人脸数据集 PCA 降维}——\textbf{可视化 50 张"特征脸"}
\item \texttt{demo\_least\_squares\_fit()}——\textbf{用最小二乘拟合 $y = ax^2 + bx + c$}——\textbf{对比 \texttt{np.linalg.lstsq} 和手工 $(A^{\!\top} A)^{-1} A^{\!\top} \bm{b}$}
\item \texttt{demo\_pagerank()}——\textbf{本章 Aha 例子 1}——\textbf{5 网页迷你互联网的 PageRank 计算}
\item \texttt{demo\_matrix\_decomp()}——\textbf{对比 LU / QR / Cholesky 在同一个问题上的速度和精度}
\item \texttt{demo\_nn\_as\_matmul()}——\textbf{用纯 NumPy 实现一个 3 层 MLP}——\textbf{每一步都是矩阵乘法}
\item \texttt{demo\_qubit\_gates()}——\textbf{用 NumPy 实现 H、X、CNOT 门}——\textbf{展示 Bell 态的产生}
\end{enumerate}

\textbf{核心代码片段——PageRank 的 30 行}：

\begin{lstlisting}[language=Python]
import numpy as np

def pagerank(L, d=0.85, n_iter=50, tol=1e-8):
    """
    L : (n, n) ndarray
        Link matrix: L[i, j] = 1 if page j links to page i.
    d : float
        Damping factor (Google uses 0.85).
    Returns the PageRank vector r of length n.
    """
    n = L.shape[0]
    # Column-normalize: M[:, j] = L[:, j] / sum(L[:, j])
    col_sum = L.sum(axis=0)
    # Handle dead ends (column of zeros): treat as uniform jump.
    col_sum[col_sum == 0] = n
    M = L / col_sum
    M[:, L.sum(axis=0) == 0] = 1.0 / n
    # Google matrix with damping.
    G = d * M + (1 - d) / n * np.ones((n, n))
    # Power iteration.
    r = np.ones(n) / n
    for k in range(n_iter):
        r_new = G @ r
        if np.linalg.norm(r_new - r, 1) < tol:
            break
        r = r_new
    return r / r.sum()  # Normalize to probability vector.

# Mini-internet: 5 pages A, B, C, D, E.
L = np.array([
    [0, 0, 1, 1, 1],
    [1, 0, 1, 1, 1],
    [1, 1, 0, 1, 1],
    [1, 0, 0, 0, 1],
    [0, 0, 0, 0, 0]
])
r = pagerank(L)
print("PageRank vector:", np.round(r, 3))
# Expected: [0.27, 0.30, 0.27, 0.13, 0.03]
\end{lstlisting}

\textbf{运行}：

\begin{lstlisting}[language=bash]
git clone https://github.com/inaturelz-coder/math-rendu
cd math-rendu
pip install -r requirements.txt
python modules/linalg_apps.py
\end{lstlisting}

\subsection{思考题 / Exercises}

\begin{enumerate}
\item \textbf{[SVD]} 给一个 $3 \times 2$ 矩阵 $A = \begin{pmatrix} 1 & 0 \\ 0 & 1 \\ 1 & 1 \end{pmatrix}$——\textbf{手算 $A^{\!\top} A$ 的特征值——得到奇异值}——\textbf{然后用 \texttt{np.linalg.svd} 验证}。
\item \textbf{[SVD 压缩]} 找一张 $512 \times 512$ 的灰度图——\textbf{算 SVD}——\textbf{画出"保留前 $k$ 个奇异值后的图像"对 $k = 1, 5, 20, 50, 100$}——\textbf{说说你"什么时候开始看不出差别"}。
\item \textbf{[PCA]} 在 sklearn 里加载 iris 数据集（150 朵花——4 维特征）——\textbf{用 PCA 降到 2 维并画散点图}——\textbf{看 3 类花是否在 2D 平面上分开了}。
\item \textbf{[最小二乘]} 给一组实验数据 $(x_i, y_i)$——\textbf{$x = [1, 2, 3, 4, 5]$，$y = [2.1, 3.9, 6.2, 7.8, 10.1]$}——\textbf{用最小二乘拟合 $y = ax + b$}——\textbf{算出 $a, b$ 并画图}。
\item \textbf{[PageRank]} 把本章迷你互联网改成\textbf{加一个新网页 F——只指向 E}——\textbf{重新算 PageRank}——\textbf{看看 E 的分数是否提升}。
\item \textbf{[PageRank 收敛]} 把 $d$ 从 0.85 改成 0.5、0.95、1.0——\textbf{观察收敛速度和最终结果}——\textbf{为什么 $d = 1.0$ 时算法可能不收敛}？
\item \textbf{[矩阵分解]} 对一个 $100 \times 100$ 的对称正定矩阵——\textbf{分别用 LU、QR、Cholesky 解 $A \bm{x} = \bm{b}$}——\textbf{用 \texttt{time.time()} 比较三者的耗时}——\textbf{Cholesky 是否最快}？
\item \textbf{[神经网络]} 用 NumPy 实现一个 2 层 MLP（输入 2 维——隐层 4 维——输出 1 维）——\textbf{学习 XOR 函数}——\textbf{只用矩阵乘法 + ReLU + 平方损失 + 手算梯度}。
\item \textbf{[量子]} 用 NumPy 实现 $H \otimes H |00\rangle$——\textbf{验证结果是 $\frac{1}{2}(|00\rangle + |01\rangle + |10\rangle + |11\rangle)$}（均匀叠加）。
\item \textbf{[历史]} 查一查 Larry Page 的博士论文最后到底完成了没有——\textbf{Google 早年的服务器叫什么名字}（hint: BackRub）——\textbf{它和 PageRank 的关系是什么}？
\item \textbf{[综合]} 找一个你身边的"排名问题"（比如\textbf{大学排名、足球队排名、推荐朋友}）——\textbf{把它建模成 PageRank 问题}——\textbf{写出 $M$ 并计算}。
\end{enumerate}

\section{本章小结 / Chapter Summary}

\begin{tcolorbox}[essbox={本章核心}]
\begin{itemize}
\item \textbf{SVD $A = U \Sigma V^{\!\top}$}——\textbf{最有用的矩阵分解——压缩 / 推荐 / LoRA 的基石}——\textbf{Eckart--Young 保证"低秩近似最优"}
\item \textbf{PCA}——\textbf{协方差矩阵的特征分解}——\textbf{= 数据 SVD 的几何版本}——\textbf{"找最胖方向"压缩高维数据}
\item \textbf{最小二乘}——\textbf{$\bm{x}^* = (A^{\!\top} A)^{-1} A^{\!\top} \bm{b}$}——\textbf{= 把 $\bm{b}$ 投影到 $A$ 的列空间}——\textbf{回归 / 滤波 / 定位的祖宗公式}
\item \textbf{Markov 链}——\textbf{$\bm{p}_{t+1} = M \bm{p}_t$}——\textbf{平稳分布 = 特征值 1 的特征向量}（Perron--Frobenius）
\item \textbf{PageRank}——\textbf{$\bm{r} = \hat{M} \bm{r}$}——\textbf{$\hat{M} = 0.85 M + 0.03 \bm{1} \bm{1}^{\!\top}$}——\textbf{Larry Page 1998 年的 1.5 万亿美元特征向量}
\item \textbf{LU / QR / Cholesky}——\textbf{数值线代三件套}——\textbf{对应解方程 / 最小二乘 / 对称正定问题}
\item \textbf{神经网络}——\textbf{= 矩阵乘法的链 + 一点点非线性}——\textbf{99\% 的计算量是 BLAS}
\item \textbf{量子比特 + 量子门}——\textbf{= $\mathbb{C}^2$ 上的酉矩阵}——\textbf{物理学和线代握手的地方}
\item \textbf{现代视角}——\textbf{推荐系统 / LoRA / 自动驾驶 / 量子计算}——\textbf{统统是"巨型矩阵的减法艺术"}
\end{itemize}
\end{tcolorbox}

\textbf{至此——"线代三部曲"全部讲完了}。\textbf{从第 6 章的"矩阵 = 变换"——到第 7 章的"特征值 = 不变方向"——到第 8 章的"PageRank = 1.5 万亿美元的特征向量"}——\textbf{我们走完了从抽象到工程的全部旅程}。

\textbf{下一章——复变函数}——\textbf{我们要从"实数世界"走到"复数世界"}——\textbf{看 Euler $e^{i\pi} + 1 = 0$ 如何统一三角、指数和复指数}——\textbf{以及为什么 5G 通信、信号处理、量子力学全部需要复数}。

\textbf{线代是"空间的语言"}——\textbf{复变函数是"波的语言"}——\textbf{两者合起来——你就拥有了整个工程世界的"双语能力"}。

\clearpage


\part{复变函数 / Complex Analysis}
% =====================================================================
%  第 9 章 · 复变函数 / Complex Analysis
%  复数、欧拉公式、解析函数、Cauchy 积分、留数、Fourier
%  小故事：Euler 最美的公式 e^{iπ} + 1 = 0
%  Aha：Fourier 变换在 MP3 压缩里的工作
% =====================================================================

\chapter{复变函数 / Complex Analysis}
\label{ch:complex}

\begin{quote}\itshape
1748 年——\textbf{Leonhard Euler} 在他的《无穷分析引论》里写下一行公式：
\[
e^{i\pi} + 1 = 0
\]
\textbf{五个最重要的常数}——$0$、$1$、$i$、$\pi$、$e$——\textbf{被一行字串在了一起}。\\
\textbf{当代物理学家 Richard Feynman} 称它为\textbf{"数学界最伟大的公式"}——\textbf{犹如物理界的广义相对论}。\\
这一章——\textbf{就讲这一行公式背后的整个世界}。
\end{quote}

\section{写在前面 / Preface}

\textbf{我大二第一次学复变}——\textbf{老师上来就讲"复数域 $\mathbb{C}$ 在拓扑上同胚于 $\mathbb{R}^2$"}——\textbf{我整个人都懵了}。

\textbf{后来才知道——这门课在中国工科——叫"复变函数与积分变换"}——\textbf{很多老师把它当成"高数 + 1"来教}——\textbf{讲一堆 Cauchy-Riemann 方程的偏导验证}——\textbf{然后期末算几个留数题}。

\textbf{我做硕士那年}——\textbf{在信号处理实验室}——\textbf{我第一次用 \texttt{scipy.fft} 把一段音频拆成频率}——

\begin{lstlisting}[language=Python]
import numpy as np
spectrum = np.fft.fft(audio_signal)
\end{lstlisting}

\textbf{一行代码——做完了 Fourier 变换}——\textbf{返回的是一个复数数组}。

\textbf{那一刻我才反应过来}：\textbf{复数不是"虚的"}——\textbf{它是"旋转的"}——\textbf{Fourier 变换的每一个值——都是一个"频率成分"的\textbf{幅度}和\textbf{相位}}——\textbf{$|z|$ 是幅度——$\arg(z)$ 是相位}——\textbf{这就是复数最自然的应用}。

\textbf{这一章只讲三件事}：

\begin{itemize}
\item \textbf{复数 = 旋转}——\textbf{欧拉公式 $e^{i\theta} = \cos\theta + i\sin\theta$ 是一切的源头}
\item \textbf{解析函数}——\textbf{为什么"可导"在复数里这么强大}
\item \textbf{留数定理}——\textbf{为什么它能算实数积分}
\end{itemize}

\textbf{再加一件 Aha 的事}——\textbf{Fourier 变换为什么必须用复数}。

\section{一句话本质 / The Essence in One Sentence}

\begin{tcolorbox}[essbox={一句话本质}]
\textbf{中文}：\textbf{复数不是"虚的"——它是"旋转的"}——\textbf{$i$ 是"逆时针转 90 度"的算符}——\textbf{欧拉公式 $e^{i\theta} = \cos\theta + i\sin\theta$ 把"指数增长"与"圆周旋转"统一为同一回事}。\\[6pt]
\textbf{English}: Complex numbers are not "imaginary" —— they are \emph{rotational}. The unit $i$ is the operator "rotate 90 degrees counter-clockwise". Euler's formula $e^{i\theta} = \cos\theta + i\sin\theta$ unifies exponential growth and circular rotation as one and the same.
\end{tcolorbox}

记住这一句话——\textbf{这一章已经讲完了一半}。

\section{教材里你看到的 / What the Textbook Tells You}

\subsection{钟玉泉《复变函数论》}

\textbf{中国工科最常用的复变教材}——\textbf{钟玉泉的《复变函数论》（高等教育出版社）}。

\textbf{优点}：体系完整、定理严格、习题大量。

\textbf{缺点}：

\begin{enumerate}
\item \textbf{从"复数定义"开始}——\textbf{2 页就上 Cauchy-Riemann}——\textbf{完全没讲"为什么要研究复数"}
\item \textbf{留数那一章塞满了"求实积分的 6 种套路"}——\textbf{但不告诉你这些积分在工程里出现在哪里}
\item \textbf{Fourier / Laplace 变换被压在最后两章}——\textbf{很多老师讲不到——而这恰恰是工程里唯一真正用的}
\end{enumerate}

\subsection{Brown \& Churchill《Complex Variables and Applications》}

\textbf{美国的"国民复变教材"}——\textbf{Brown \& Churchill}——\textbf{Stewart 级别的国民读物}。

\textbf{它比国内教材好的地方}：\textbf{讲"流体力学怎么用复势函数"、"电磁场怎么用 conformal mapping"}——\textbf{应用是 19 世纪经典物理的应用}。

\textbf{它的缺点}：\textbf{现代信号处理 / 控制论 / 量子力学的内容——它也没讲}。

\subsection{Needham《Visual Complex Analysis》}

\textbf{Tristan Needham} 的《Visual Complex Analysis》——\textbf{这本书是复变界的"3Blue1Brown"}——\textbf{用几何 + 图——把每一个公式都画了出来}。

\textbf{读这本书}——\textbf{你会突然"看见"复数}。\textbf{我个人最推荐的复变补充读物}——\textbf{中文有翻译版}。

\subsection{我的策略}

\textbf{本章只做三件事}：

\begin{enumerate}
\item \textbf{把"复数 = 旋转"讲到你看见图}——\textbf{这是 Needham 的思路——也是工程师最该有的直觉}
\item \textbf{把解析、Cauchy、留数串成一条线}——\textbf{它们其实是同一件事的三个层次}
\item \textbf{把 Aha 例子讲到能跑代码}——\textbf{Fourier 变换 + MP3 压缩——\texttt{scipy.fft}}
\end{enumerate}

\begin{tcolorbox}[storybox={\Large 小故事：Euler 与最美的公式}]
\textbf{1748 年}——\textbf{Leonhard Euler}（瑞士数学家，时年 41 岁，已经一只眼睛失明）——\textbf{在他的《无穷分析引论》里——写下了我们今天叫"欧拉恒等式"的公式}：

\[
e^{i\pi} + 1 = 0
\]

\medskip

\textbf{这一行——只有 7 个符号}——\textbf{把数学里五个最重要的常数串了起来}：

\begin{itemize}
\item \textbf{$0$}——\textbf{加法的单位元}
\item \textbf{$1$}——\textbf{乘法的单位元}
\item \textbf{$\pi$}——\textbf{圆周与直径之比——几何的根基}
\item \textbf{$e$}——\textbf{自然指数的底——微积分的根基}
\item \textbf{$i$}——\textbf{$\sqrt{-1}$——复数的根基}
\end{itemize}

\medskip

\textbf{它还用了 3 种基本运算}——\textbf{加法、乘法（指数是反复乘法）、相等}——\textbf{完美}。

\medskip

\textbf{Richard Feynman} 在 1977 年的《Lectures on Physics》第 22 章——\textbf{称它为}：\\
\textbf{"the most remarkable formula in mathematics"}——\textbf{"数学里最非凡的公式"}。

\medskip

\textbf{有人说它是数学界的 $E = mc^2$}——\textbf{但我觉得不对}——\textbf{$E = mc^2$ 是一个物理事实}——\textbf{而欧拉恒等式是一个\textbf{数学必然}}。

\medskip

\textbf{它更像\textbf{广义相对论}}——\textbf{把"看似无关的几个领域"——指数、三角、虚数——\textbf{统一在一个最简公式里}}。

\medskip

\textbf{2004 年}——\textbf{美国《Physics World》杂志做读者投票}——\textbf{"历史上最美的方程"}——\textbf{欧拉恒等式}和\textbf{Maxwell 方程组}并列第一。

\medskip

\textbf{Euler 一生发表了 866 篇论文}——\textbf{即使失明后——也由学生记录、出版}——\textbf{是历史上最高产的数学家之一}。\textbf{Laplace 评价他}：\textbf{"读 Euler，读 Euler，他是我们所有人的老师"}。
\end{tcolorbox}

\section{直觉 / Intuition}

\subsection{复数：被误解了 400 年的概念}

\textbf{16 世纪}——\textbf{意大利数学家 Cardano 在解三次方程时}——\textbf{第一次写下了 $\sqrt{-1}$}——\textbf{他自己也觉得"莫名其妙"}——\textbf{说这种东西"as subtle as it is useless"（既精巧又无用）}。

\textbf{Descartes 在 1637 年}——\textbf{给这种"假数"起了个名字}——\textbf{imaginary numbers}——\textbf{虚数}——\textbf{这个名字\textbf{流毒至今}}——\textbf{让无数学生以为"虚数是假的"}。

\textbf{1799 年}——\textbf{Caspar Wessel}（挪威测量员，不是职业数学家）——\textbf{画出了第一个"复平面"}——\textbf{$x$ 轴是实数、$y$ 轴是虚数}——\textbf{从此复数有了\textbf{几何意义}}。

\textbf{1806 年}——\textbf{Argand 独立发现了同一件事}——\textbf{所以今天叫 Argand 图}。

\textbf{1831 年}——\textbf{Gauss 正式提出 "复平面" 一词}——\textbf{并把复数当作 $\mathbb{R}^2$ 上的点处理}——\textbf{复数从此\textbf{合法化}}。

\subsection{$i$ 不是"虚"——是"旋转 90 度"}

\textbf{看这个事实}：

\[
1 \xrightarrow{\times i} i \xrightarrow{\times i} -1 \xrightarrow{\times i} -i \xrightarrow{\times i} 1
\]

\textbf{乘以 $i$ 一次}——\textbf{把 1 转到了 $i$}——\textbf{在复平面上——这是\textbf{逆时针转 90 度}}。

\textbf{乘以 $i$ 四次}——\textbf{转回原位}——\textbf{这就是为什么 $i^4 = 1$}。

\textbf{所以 $i$ 的真正含义不是"$\sqrt{-1}$"}——\textbf{而是\textbf{"逆时针旋转 90 度的算符"}}——\textbf{$i^2 = -1$ 的几何意义就是"转 90 度两次 = 转 180 度 = 反向"}。

\begin{tcolorbox}[ahabox={Aha: $i$ 是动词不是名词}]
\textbf{中学老师告诉你}——\textbf{$i$ 是一个数}——\textbf{满足 $i^2 = -1$}。

\medskip

\textbf{但更深刻的理解}——\textbf{$i$ 是一个\textbf{动作}}——\textbf{"逆时针转 90 度"}。

\medskip

\textbf{从这个角度看}——\textbf{$3 + 4i$ 就是}——\textbf{先沿 $x$ 走 3 步}——\textbf{再沿 $y$ 走 4 步}——\textbf{终点是复平面上的一个点}——\textbf{距离原点 $|3+4i| = 5$ 单位}——\textbf{方向是 $\arctan(4/3) \approx 53°$}。

\medskip

\textbf{这就是复数的\textbf{极坐标表示}}——\textbf{$z = r e^{i\theta}$}——\textbf{$r$ 是长度——$\theta$ 是角度}。

\medskip

\textbf{两个复数相乘}——\textbf{在极坐标下\textbf{长度相乘、角度相加}}：
\[
r_1 e^{i\theta_1} \cdot r_2 e^{i\theta_2} = r_1 r_2 \, e^{i(\theta_1 + \theta_2)}
\]
\textbf{这就是复数最优美的性质}——\textbf{乘法 = 旋转 + 伸缩}。
\end{tcolorbox}

\subsection{欧拉公式：把指数和三角拼到一起}

\textbf{欧拉公式}：

\[
\boxed{\;e^{i\theta} = \cos\theta + i\sin\theta\;}
\]

\textbf{怎么"看出"这个公式}？——\textbf{用 Taylor 级数}：

\begin{align*}
e^x &= 1 + x + \frac{x^2}{2!} + \frac{x^3}{3!} + \frac{x^4}{4!} + \cdots \\
\cos x &= 1 - \frac{x^2}{2!} + \frac{x^4}{4!} - \cdots \\
\sin x &= x - \frac{x^3}{3!} + \frac{x^5}{5!} - \cdots
\end{align*}

\textbf{把 $x = i\theta$ 代入 $e^x$}——\textbf{利用 $i^2 = -1, i^3 = -i, i^4 = 1$}——\textbf{重新排列实部和虚部}：

\[
e^{i\theta} = \underbrace{\left(1 - \frac{\theta^2}{2!} + \frac{\theta^4}{4!} - \cdots\right)}_{\cos\theta} + i\underbrace{\left(\theta - \frac{\theta^3}{3!} + \frac{\theta^5}{5!} - \cdots\right)}_{\sin\theta}
\]

\textbf{原来如此}——\textbf{$\cos$ 是 $e^x$ 的\textbf{偶次幂部分}}——\textbf{$\sin$ 是\textbf{奇次幂部分}}——\textbf{它们本来就藏在指数函数里}——\textbf{虚数只是\textbf{把它们抽出来的提取器}}。

\textbf{欧拉恒等式}——\textbf{是欧拉公式在 $\theta = \pi$ 时的特殊情形}：

\[
e^{i\pi} = \cos\pi + i\sin\pi = -1 + 0i = -1 \Longrightarrow e^{i\pi} + 1 = 0
\]

\section{数学 / The Math}

\subsection{复数的四种表示}

同一个复数 $z = 3 + 4i$——\textbf{有四种写法}——\textbf{工程师都要熟}：

\begin{enumerate}
\item \textbf{代数形式}：$z = a + bi$，$a, b \in \mathbb{R}$
\item \textbf{几何形式}：复平面上的点 $(a, b)$ 或向量
\item \textbf{极坐标形式}：$z = r(\cos\theta + i\sin\theta)$，$r = |z|, \theta = \arg z$
\item \textbf{指数形式}：$z = r e^{i\theta}$——\textbf{最有用}——\textbf{乘法、幂、开方都最简单}
\end{enumerate}

\textbf{转换公式}：

\[
r = \sqrt{a^2 + b^2}, \qquad \theta = \arctan(b/a), \qquad a = r\cos\theta, \qquad b = r\sin\theta
\]

\subsection{De Moivre 公式}

\[
(\cos\theta + i\sin\theta)^n = \cos(n\theta) + i\sin(n\theta)
\]

\textbf{用欧拉公式一句话证完}：$\bigl(e^{i\theta}\bigr)^n = e^{in\theta}$。

\textbf{这是\textbf{开方}的工具}——\textbf{$z^n = w$ 在复数域里恰有 $n$ 个根}——\textbf{均匀分布在以原点为圆心、$|w|^{1/n}$ 为半径的圆上}。

\subsection{复变函数：把 $\mathbb{C}$ 映到 $\mathbb{C}$}

\textbf{复变函数 $f: \mathbb{C} \to \mathbb{C}$}——\textbf{就是把一个复数变成另一个复数}。

\textbf{几个基本函数}（从实数函数自然推广）：

\begin{align*}
e^z &= e^{x+iy} = e^x(\cos y + i\sin y) \quad \text{(永远不为 0)} \\
\sin z &= \frac{e^{iz} - e^{-iz}}{2i}, \qquad \cos z = \frac{e^{iz} + e^{-iz}}{2} \\
\log z &= \ln|z| + i\arg z \quad \text{(多值！每个 $2\pi$ 是一个分支)}
\end{align*}

\textbf{注意 $\log z$ 是\textbf{多值函数}}——\textbf{因为 $\arg z$ 加 $2\pi$ 都对应同一个 $z$}——\textbf{这是复变里第一次出现的"新事物"}。

\textbf{$\sin z$ 在复数里\textbf{不再有界}}——\textbf{比如 $\sin(iy) = i\sinh y$——可以无穷大}。

\subsection{解析函数与 Cauchy-Riemann 方程}

\textbf{复函数 $f(z) = u(x,y) + iv(x,y)$ 在 $z_0$ 处\textbf{可导}}——\textbf{要求极限}

\[
f'(z_0) = \lim_{\Delta z \to 0} \frac{f(z_0 + \Delta z) - f(z_0)}{\Delta z}
\]

\textbf{存在}——\textbf{并且\textbf{与 $\Delta z$ 趋近 0 的方向无关}}。

\textbf{这个"方向无关"的要求——比实数严苛得多}——\textbf{$\Delta z$ 在复平面可以从\textbf{任何方向}靠近 0}——\textbf{所以"复可导" $\Longleftrightarrow$ Cauchy-Riemann 方程}：

\[
\boxed{\;\frac{\partial u}{\partial x} = \frac{\partial v}{\partial y}, \qquad \frac{\partial u}{\partial y} = -\frac{\partial v}{\partial x}\;}
\]

\textbf{如果 $f$ 在某区域内\textbf{处处可导}}——\textbf{我们说 $f$ 在该区域\textbf{解析}（analytic / holomorphic）}。

\begin{tcolorbox}[ahabox={Aha: 复可导 = 实可导 ×10}]
\textbf{在实数里}——\textbf{可导 $\Rightarrow$ 有切线}——\textbf{仅此而已}——\textbf{可导一次的函数未必可导二次}。

\medskip

\textbf{在复数里}——\textbf{解析（可导一次）}——\textbf{自动可导无穷次}——\textbf{自动等于自己的 Taylor 级数}——\textbf{自动满足 Laplace 方程 $u_{xx} + u_{yy} = 0$（调和函数）}——\textbf{自动可以用围线积分公式恢复}。

\medskip

\textbf{这就是复变最神奇的地方}——\textbf{一个看似简单的条件}（方向无关）——\textbf{推出了一连串极强的性质}——\textbf{这种"以小博大"的现象——在数学里叫\textbf{rigidity}（刚性）}。
\end{tcolorbox}

\subsection{复积分与 Cauchy 积分公式}

\textbf{复积分}——\textbf{沿复平面上一条曲线 $C$ 的积分}：

\[
\int_C f(z)\, dz
\]

\textbf{Cauchy 积分定理}（1825）：\textbf{如果 $f$ 在闭曲线 $C$ 内部解析}——\textbf{则}

\[
\oint_C f(z)\, dz = 0
\]

\textbf{Cauchy 积分公式}：\textbf{如果 $f$ 在闭曲线 $C$ 内部解析}——\textbf{$z_0$ 在 $C$ 内部}——\textbf{则}

\[
f(z_0) = \frac{1}{2\pi i} \oint_C \frac{f(z)}{z - z_0}\, dz
\]

\textbf{这条公式的震撼之处}——\textbf{函数在区域\textbf{内部}的值——完全由它在\textbf{边界}上的值决定}——\textbf{解析函数是"被边界完全束缚"的}——\textbf{这就是 rigidity 的最直接体现}。

\subsection{留数定理}

\textbf{留数（residue）}——\textbf{是解析函数 $f$ 在\textbf{孤立奇点 $z_0$} 处 Laurent 级数中 $(z-z_0)^{-1}$ 项的系数}：

\[
f(z) = \sum_{n=-\infty}^{\infty} a_n (z - z_0)^n, \qquad \text{Res}(f, z_0) = a_{-1}
\]

\textbf{留数定理}：\textbf{若 $f$ 在闭曲线 $C$ 内部除有限个孤立奇点 $z_1, \ldots, z_k$ 外解析}——\textbf{则}

\[
\boxed{\;\oint_C f(z)\, dz = 2\pi i \sum_{j=1}^{k} \text{Res}(f, z_j)\;}
\]

\textbf{这条定理的工程价值}——\textbf{把"算积分"变成了"算奇点处的留数"}——\textbf{把分析问题变成了代数问题}。

\subsection{用留数算实积分}

\textbf{经典例子}：算

\[
I = \int_{-\infty}^{\infty} \frac{1}{1 + x^2}\, dx
\]

\textbf{传统办法}：用反正切——$I = \arctan x \big|_{-\infty}^{\infty} = \pi$。

\textbf{复变办法}：考虑复函数 $f(z) = \frac{1}{1+z^2}$——\textbf{它在 $z = \pm i$ 处有极点}。

\textbf{取一个半径 $R$ 的上半圆围线}——\textbf{包含极点 $z = i$}：

\[
\oint_C \frac{dz}{1+z^2} = 2\pi i \cdot \text{Res}\left(\frac{1}{1+z^2}, i\right) = 2\pi i \cdot \frac{1}{2i} = \pi
\]

\textbf{当 $R \to \infty$ 时}——\textbf{上半圆的积分趋于 0}——\textbf{剩下的就是实轴上的积分}——\textbf{所以 $I = \pi$}。

\textbf{看上去比"反正切"还复杂}——\textbf{但对于\textbf{大多数没有初等原函数的实积分}（比如 $\int \frac{\sin x}{x} dx$）}——\textbf{留数法是\textbf{唯一可行的解析方法}}。

\section{Aha 例子：Fourier 变换与 MP3 压缩 / Aha: Fourier and MP3}

\textbf{场景}——\textbf{你录了一段 \textbf{3 分钟、44.1 kHz、16 bit 立体声}的音乐}——\textbf{原始文件大约 \textbf{30 MB}}——\textbf{压成 MP3 之后大约 \textbf{3 MB}}——\textbf{为什么能压缩\textbf{10 倍}而你听不出区别}？

\textbf{答案}——\textbf{Fourier 变换 + 人耳的"心理声学"}。

\subsection{Fourier 变换：把信号拆成频率}

\textbf{Fourier 变换的定义}：

\[
\hat{f}(\omega) = \int_{-\infty}^{\infty} f(t)\, e^{-i\omega t}\, dt
\]

\textbf{注意}——\textbf{被积函数里有 $e^{-i\omega t}$}——\textbf{这是\textbf{复数}}——\textbf{所以 Fourier 变换的结果\textbf{天然是复数}}。

\textbf{每一个 $\hat{f}(\omega)$ 告诉我们}：

\begin{itemize}
\item \textbf{$|\hat{f}(\omega)|$}——\textbf{在频率 $\omega$ 上的\textbf{幅度}}——\textbf{这个频率"多响"}
\item \textbf{$\arg \hat{f}(\omega)$}——\textbf{在频率 $\omega$ 上的\textbf{相位}}——\textbf{这个频率"什么时候到达峰值"}
\end{itemize}

\textbf{$e^{-i\omega t}$ 的真正含义}——\textbf{用欧拉公式拆开}：

\[
e^{-i\omega t} = \cos(\omega t) - i\sin(\omega t)
\]

\textbf{所以 Fourier 变换实际上是}：

\[
\hat{f}(\omega) = \underbrace{\int f(t)\cos(\omega t)\, dt}_{\text{余弦分量}} - i\underbrace{\int f(t)\sin(\omega t)\, dt}_{\text{正弦分量}}
\]

\textbf{为什么必须用复数}——\textbf{因为\textbf{相位}}——\textbf{如果只用实数（只用 cos）}——\textbf{你无法区分\textbf{"先到峰值"还是"后到峰值"}}——\textbf{用 cos + i·sin}——\textbf{就同时记录了\textbf{"多响"和"什么时候响"}}。

\begin{tcolorbox}[ahabox={Aha: Fourier 变换为什么必须用复数}]
\textbf{设想两段音乐}——\textbf{都是 440 Hz 的"标准 A"}——\textbf{但一段从 $t=0$ 开始响}——\textbf{另一段从 $t=0.5\,\text{ms}$ 开始响}。

\medskip

\textbf{它们的\textbf{振幅谱（只看 $|\hat{f}|$）}\textbf{完全相同}}——\textbf{都是在 440 Hz 处一个尖峰}。

\medskip

\textbf{它们的\textbf{相位谱（看 $\arg \hat{f}$）}\textbf{不同}}——\textbf{第二段的所有频率都"晚到了 $0.5\,\text{ms}$"}——\textbf{相位差 = $\omega \cdot 0.5\,\text{ms}$}。

\medskip

\textbf{如果只用实数 Fourier（cos 展开）}——\textbf{你无法恢复"什么时候开始"这个信息}——\textbf{所以必须用复数}。

\medskip

\textbf{在工程上}——\textbf{"幅度 + 相位"——就是"How loud + When loud"}——\textbf{这正是复数的几何意义}——\textbf{$|z| + \arg z$}——\textbf{长度 + 角度}。

\medskip

\textbf{这就是为什么——Fourier 变换天然是复数运算}——\textbf{不是"为了好看"——是"为了完整"}。
\end{tcolorbox}

\subsection{MP3 怎么压缩}

\textbf{MP3 的核心算法}（极简版）：

\begin{enumerate}
\item \textbf{把音频切成 \textbf{1152 个采样点}一帧}（约 26 ms）
\item \textbf{对每一帧做 \textbf{MDCT}（Modified Discrete Cosine Transform——本质就是 Fourier 变换的变种）}
\item \textbf{得到这一帧的\textbf{频谱}}——\textbf{各频率成分的幅度}
\item \textbf{利用\textbf{心理声学模型}——找出"人耳听不见"的频率成分}（比如被强信号\textbf{掩蔽}的弱信号）
\item \textbf{把这些频率成分\textbf{用更少的比特存储}（甚至直接丢弃）}
\item \textbf{用 Huffman 编码进一步压缩}
\end{enumerate}

\textbf{核心 idea}——\textbf{在\textbf{频域}里}——\textbf{很多频率成分\textbf{人耳根本听不见}}——\textbf{丢掉它们——你完全察觉不到}——\textbf{但文件小了 10 倍}。

\textbf{在\textbf{时域}里你做不到这一步}——\textbf{因为时域里每一个采样点都"看起来差不多重要"}——\textbf{只有 Fourier 把它转成频域}——\textbf{才能区分"重要频率"和"听不见的频率"}。

\subsection{Python 代码：看见频谱}

\begin{lstlisting}[language=Python]
import numpy as np
import matplotlib.pyplot as plt

# 合成一段信号：440 Hz (A4) + 880 Hz (A5) + 一点噪声
fs = 8000              # 采样率
t  = np.linspace(0, 1, fs, endpoint=False)
sig = (np.sin(2*np.pi*440*t)
     + 0.5*np.sin(2*np.pi*880*t)
     + 0.1*np.random.randn(fs))

# 做 FFT——返回复数数组
X = np.fft.fft(sig)
freqs = np.fft.fftfreq(fs, 1/fs)

# 只看正频率部分
mask = freqs > 0
mag = np.abs(X[mask])      # 幅度
phase = np.angle(X[mask])  # 相位

fig, ax = plt.subplots(2, 1, figsize=(8, 5))
ax[0].plot(freqs[mask], mag)
ax[0].set_xlim(0, 1500)
ax[0].set_xlabel("Frequency (Hz)"); ax[0].set_ylabel("|X|")
ax[0].set_title("Magnitude —— 440 Hz 和 880 Hz 两个尖峰")

ax[1].plot(freqs[mask], phase)
ax[1].set_xlim(0, 1500)
ax[1].set_xlabel("Frequency (Hz)"); ax[1].set_ylabel("arg(X)")
plt.tight_layout(); plt.show()

# 压缩演示：只保留最大的 5% 系数
threshold = np.sort(np.abs(X))[int(0.95*len(X))]
X_compressed = np.where(np.abs(X) >= threshold, X, 0)
sig_reconstructed = np.fft.ifft(X_compressed).real

# 误差
print("压缩比:", np.count_nonzero(X) / np.count_nonzero(X_compressed))
print("均方误差:", np.mean((sig - sig_reconstructed)**2))
\end{lstlisting}

\textbf{这段代码 30 行}——\textbf{涵盖了 MP3 压缩的核心思想}：\textbf{FFT 拆频率 → 丢小系数 → IFFT 恢复}。

\section{现代视角：复变在工程中 / Modern: Complex Analysis in Engineering}

\subsection{Laplace 变换：控制论的"通用语言"}

\textbf{Laplace 变换}是 Fourier 变换的"复数加强版"——\textbf{把 $i\omega$ 换成复变量 $s = \sigma + i\omega$}：

\[
F(s) = \int_0^\infty f(t)\, e^{-st}\, dt
\]

\textbf{为什么要换}——\textbf{因为很多工程信号（指数增长的、不稳定的）}——\textbf{在 Fourier 域里发散}——\textbf{但在 Laplace 域里收敛}。

\textbf{Laplace 把\textbf{微分方程}变成\textbf{代数方程}}：

\[
\mathcal{L}\{y'\} = sY(s) - y(0)
\]

\textbf{自动化控制的"传递函数 $H(s)$"——本质就是 Laplace 域里的输入输出比}——\textbf{你在大三的"自动控制原理"里画的\textbf{Bode 图、Nyquist 图}}——\textbf{全在复平面上}。

\subsection{Z 变换：数字信号处理的"Laplace"}

\textbf{Z 变换}——\textbf{Laplace 变换在\textbf{离散时间}的版本}：

\[
X(z) = \sum_{n=0}^{\infty} x[n]\, z^{-n}, \qquad z \in \mathbb{C}
\]

\textbf{每一个数字滤波器（FIR、IIR）}——\textbf{它的\textbf{零点和极点}就在 $z$-平面上}——\textbf{设计滤波器 = 在复平面里摆放极点}。

\subsection{量子力学：复数是世界的"本原"}

\textbf{Schr\"odinger 方程}：

\[
i\hbar\, \frac{\partial \psi}{\partial t} = -\frac{\hbar^2}{2m}\nabla^2 \psi + V \psi
\]

\textbf{注意左边有 $i$}——\textbf{波函数 $\psi$ \textbf{本质上是复值函数}}——\textbf{$|\psi|^2$ 是粒子在某位置的概率密度}。

\textbf{量子力学里复数不是"工具"——是\textbf{物理实体}}——\textbf{20 世纪最伟大的发现之一就是}：\textbf{我们的世界——在最深的层次——是用复数描述的}。

\subsection{保形映射：流体与电磁的几何工具}

\textbf{解析函数有一个性质}——\textbf{在它的导数不为 0 处——保持角度不变}——\textbf{叫做\textbf{保形映射}（conformal mapping）}。

\textbf{应用}：\textbf{把"复杂边界的流体绕流问题"——通过保形映射——\textbf{变到"圆周"或"半平面"上}}——\textbf{然后用已知解返回}——\textbf{这是 19 世纪流体力学的核心工具}。

\textbf{今天}——\textbf{保形映射用在}：\textbf{大脑皮层平面化（医学影像）}、\textbf{芯片版图等价检查}、\textbf{Escher 风格的艺术变形}。

\section{代码与思考 / Code and Exercises}

\subsection{配套模块 \texttt{complex\_analysis.py}}

GitHub 上的 \texttt{modules/complex\_analysis.py} 包含\textbf{8 个 demo}：

\begin{enumerate}
\item \texttt{demo\_complex\_basics()}——画复平面、可视化加法和乘法（旋转 + 伸缩）
\item \texttt{demo\_euler\_formula()}——\textbf{核心 demo}——动画展示 $e^{i\theta}$ 在单位圆上跑——\textbf{你"看见"欧拉公式}
\item \texttt{demo\_roots\_of\_unity()}——画 $z^n = 1$ 的 $n$ 个根——均匀分布在单位圆上
\item \texttt{demo\_cauchy\_riemann()}——给定 $u(x,y)$，自动检测是否解析
\item \texttt{demo\_contour\_integral()}——数值围线积分——验证 Cauchy 定理
\item \texttt{demo\_residue\_real\_integral()}——用留数算 $\int_{-\infty}^{\infty} \frac{1}{1+x^2}dx$
\item \texttt{demo\_fft\_audio()}——\textbf{本章 Aha 例子}——FFT 拆音频 + 重建
\item \texttt{demo\_mp3\_compression()}——保留前 $k\%$ 的频率系数——观察听感损失
\end{enumerate}

\textbf{使用}：

\begin{lstlisting}[language=bash]
git clone https://github.com/inaturelz-coder/math-rendu
cd math-rendu
pip install -r requirements.txt
python modules/complex_analysis.py
\end{lstlisting}

\subsection{思考题 / Exercises}

\begin{enumerate}
\item \textbf{[直觉]} 用一句话向你的高中同桌解释"$i$ 是什么"——\textbf{不许说"$\sqrt{-1}$"}。
\item \textbf{[几何]} 把复数 $z = 1 + i$ 乘以 $i$、$i^2$、$i^3$——画出每次乘法之后的位置——验证"乘以 $i$ = 转 90 度"。
\item \textbf{[Euler]} 用欧拉公式证明 $\cos(\alpha+\beta) = \cos\alpha\cos\beta - \sin\alpha\sin\beta$——\textbf{比中学三角公式漂亮多了}。
\item \textbf{[De Moivre]} 求 $z^4 = 16$ 的所有复根——画在复平面上。
\item \textbf{[C-R 方程]} 验证 $f(z) = z^2$ 处处解析——验证 $f(z) = \bar z$ 处处不解析。
\item \textbf{[积分]} 用留数定理算 $\int_{-\infty}^{\infty} \frac{\cos x}{1+x^2}\, dx$——\textbf{答案 $\pi/e$}。
\item \textbf{[FFT]} 自己录一段 3 秒钟的"你好"——\textbf{用 \texttt{scipy.io.wavfile}}——\textbf{做 FFT}——\textbf{找出最强的几个频率}——\textbf{你的声音的"基频"是多少 Hz}？
\item \textbf{[压缩]} 在 FFT 之后\textbf{只保留 10\% 最大的系数}——\textbf{用 IFFT 恢复}——\textbf{对比原始音频}——\textbf{听得出区别吗}？
\item \textbf{[Laplace]} 求 $f(t) = e^{-at}\cos(\omega t)$ 的 Laplace 变换。
\item \textbf{[历史]} 查一查：\textbf{为什么 Descartes 给复数起了 "imaginary" 这个流毒至今的名字}？\textbf{他真的觉得这种数"不存在"吗}？
\end{enumerate}

\section{本章小结 / Chapter Summary}

\begin{tcolorbox}[essbox={本章核心}]
\begin{itemize}
\item \textbf{复数 = 旋转}——\textbf{$i$ 是"逆时针转 90 度"的算符}——\textbf{不是"虚的"}
\item \textbf{欧拉公式}——\textbf{$e^{i\theta} = \cos\theta + i\sin\theta$}——\textbf{把指数和三角统一}
\item \textbf{欧拉恒等式}——\textbf{$e^{i\pi} + 1 = 0$}——\textbf{数学界的广义相对论}
\item \textbf{解析函数}——\textbf{复可导（方向无关）}——\textbf{等价于 Cauchy-Riemann 方程}
\item \textbf{Cauchy 定理}——\textbf{解析函数沿闭曲线积分为 0}
\item \textbf{Cauchy 积分公式}——\textbf{函数内部值 = 边界值的积分}——\textbf{rigidity 的核心}
\item \textbf{留数定理}——\textbf{$\oint f\,dz = 2\pi i \sum \text{Res}$}——\textbf{把分析变代数}
\item \textbf{Fourier 变换}——\textbf{天然是复数}——\textbf{记录幅度和相位}
\item \textbf{应用}——\textbf{MP3 压缩、Laplace 控制论、量子力学、保形映射}
\end{itemize}
\end{tcolorbox}

\textbf{这一章是\textbf{工科数学的最后一块拼图}}——\textbf{从这里开始}——\textbf{你的"任督二脉"已经打通}——\textbf{微积分 + 线性代数 + 微分方程 + 概率统计 + 复变 + Fourier}——\textbf{这就是现代工程师的\textbf{六脉神剑}}。

\textbf{接下来——你要做的不是"再多学一门数学"}——\textbf{而是\textbf{把这些工具放到你的真实工程问题里去}}——\textbf{在写代码、调实验、看 paper 的过程中}——\textbf{反复\textbf{遇见它们}}——\textbf{才会真正\textbf{属于你}}。

\textbf{数学之美}——\textbf{不在于公式有多漂亮}——\textbf{而在于}——\textbf{当你某天调试一个 bug 时}——\textbf{脑子里突然蹦出一句}——\textbf{"这不就是 Cauchy 积分公式吗"}——\textbf{那一刻}——\textbf{这本书的使命才算完成}。

\clearpage


\part{概率与统计 / Probability and Statistics}
% =====================================================================
%  第 10 章 · 概率论 / Probability
%  概率基础、随机变量、期望方差、大数定律、随机过程、应用
%  小故事：Kolmogorov 与 1933 年的公理化
%  Aha：蒙特卡洛求 π + 期权定价
% =====================================================================

\chapter{概率论 / Probability}
\label{ch:prob}

\begin{quote}\itshape
1933 年，莫斯科。一个 30 岁的年轻人——\textbf{Andrey Kolmogorov}——\\
出版了一本 62 页的德文小册子《Grundbegriffe der Wahrscheinlichkeitsrechnung》。\\
从此——\textbf{概率}——从 17 世纪的赌博游戏——\\
变成了一门和\textbf{几何}一样严格的\textbf{现代数学}。\\
我们这一章——就讲他想清楚的那件事：\textbf{概率到底是什么}。
\end{quote}

\section{写在前面 / Preface}

\textbf{我刚学概率时}——和大多数中国工科生一样——\textbf{我背的是公式}。

$P(A|B) = \frac{P(AB)}{P(B)}$。$E(X) = \sum x_i p_i$。$D(X) = E(X^2) - [E(X)]^2$。二项分布、泊松分布、正态分布——一张又一张表。

\textbf{大二下学期}——\textbf{概率论与数理统计}——\textbf{我考了 88 分}——但你问我"概率到底是什么"——我说不上来。

到了\textbf{大四}——我做毕设——\textbf{要估计一个传感器的噪声方差}——我打开 Python：

\begin{lstlisting}[language=Python]
import numpy as np
sigma2 = np.var(signal)
\end{lstlisting}

\textbf{一行代码——做完了方差}。

那一刻我才反应过来——\textbf{我大二背的所有古典概型计算题}——\textbf{摸球、抓阄、几何概率}——\textbf{在工业里我一个都没用过}。

\textbf{真正用的是}：

\begin{itemize}
\item 概率是\textbf{我对未知的信念}——\textbf{Bayesian 视角}
\item 随机变量是\textbf{从样本空间到数轴的函数}——\textbf{不是"一个会变的变量"}
\item 期望是\textbf{重心}——方差是\textbf{离散度}——\textbf{两个数概括一个分布}
\item 大数定律说\textbf{平均值收敛}——中心极限定理说\textbf{和会变成正态}——\textbf{这是整个统计学的根基}
\item 一切随机过程——\textbf{都是"时间上的随机变量族"}
\end{itemize}

\textbf{这一章就把这五件事说清楚}——\textbf{再多的分布公式}——\textbf{都是这五件事的变体}。

\section{一句话本质 / The Essence in One Sentence}

\begin{tcolorbox}[essbox={一句话本质}]
\textbf{中文}：\textbf{概率是测度——随机变量是函数——分布是测度的像——期望是积分}。\\[6pt]
\textbf{English}: Probability is a measure; a random variable is a function; a distribution is the pushforward of that measure; expectation is an integral.
\end{tcolorbox}

记住这一句话——\textbf{这一章已经讲完了一半}。

\textbf{Kolmogorov 1933 年的伟大}——就在于他用\textbf{四行公理}把这件事说清楚了——\textbf{在他之前——概率是哲学；在他之后——概率是数学}。

\section{教材里你看到的 / What the Textbook Tells You}

\subsection{浙大《概率论与数理统计》}

\textbf{浙江大学盛骤}主编的\textbf{《概率论与数理统计》（第四版）}——是中国 80\% 工科生的\textbf{第一本概率教材}。

它的章节是这样的：

\begin{itemize}
\item 第一章：随机事件与概率（古典概型、几何概型）
\item 第二章：随机变量及其分布
\item 第三章：多维随机变量
\item 第四章：随机变量的数字特征
\item 第五章：大数定律与中心极限定理
\item 第六章\textasciitilde 第八章：数理统计、参数估计、假设检验
\end{itemize}

\textbf{优点}：体系完整、覆盖考研全部考点。

\textbf{缺点}：

\begin{enumerate}
\item \textbf{第一章塞满古典概型}——\textbf{摸红球、摸白球、抓阄、生日问题}——\textbf{现代工程一辈子用不到}
\item \textbf{Bayes 公式一笔带过}——而\textbf{Bayes 是机器学习的灵魂}
\item \textbf{随机过程没讲}——\textbf{Markov 链、Poisson 过程}——\textbf{现代通信和金融的核心工具}——\textbf{本科不教}
\item \textbf{没有任何代码}——\textbf{学生学完概率不会用 \texttt{numpy.random}}
\end{enumerate}

\subsection{Ross《A First Course in Probability》}

\textbf{Sheldon Ross} 写的——\textbf{美国大学的国民概率教材}。\textbf{例题极多、写得活}——\textbf{Markov 链和 Poisson 过程一开始就讲}。

缺点是\textbf{翻译版稀少}——\textbf{中国学生很少读到}。

\subsection{Bishop《Pattern Recognition and Machine Learning》}

\textbf{Christopher Bishop} 的 PRML——\textbf{机器学习的圣经}——\textbf{第 1 章和第 2 章}就是\textbf{从 Bayesian 视角重新讲一遍概率}。

\textbf{读完 PRML 第 2 章}——\textbf{你会发现浙大教材里那些古典概型练习——和真正的概率没什么关系}。

\subsection{我的策略}

\textbf{本章不会重复浙大和 Ross 的内容}——\textbf{它们都已经够全了}。

\textbf{本章只做四件事}：

\begin{enumerate}
\item \textbf{把 Kolmogorov 公理讲清楚}——\textbf{概率是测度}——\textbf{这是浙大不讲的}
\item \textbf{把 Bayes 讲透}——\textbf{先验 + 似然 = 后验}——\textbf{这是机器学习的根}
\item \textbf{把中心极限定理讲到能跑代码}——\textbf{看到正态分布从天而降}
\item \textbf{把 Aha 例子讲到能跑代码}——\textbf{蒙特卡洛求 π}——\textbf{再升级到期权定价}
\end{enumerate}

\begin{tcolorbox}[storybox={\Large 小故事：Kolmogorov 与 1933 年的公理化}]
\textbf{1933 年}——\textbf{莫斯科}——\textbf{苏联刚成立 16 年}——\textbf{斯大林清洗即将开始}。
30 岁的 \textbf{Andrey Nikolaevich Kolmogorov}（莫斯科大学教授）——\textbf{用德文}（当时数学的通用语言）——\textbf{出版了一本 62 页的小册子}：

\smallskip

\textit{Grundbegriffe der Wahrscheinlichkeitsrechnung}（《概率论基础》）。

\medskip

\textbf{这本 62 页的书——把概率论从赌博的"经验学问"——变成了和欧几里得几何一样严格的现代数学}。

\medskip

\textbf{在他之前}——概率论已经发展了 270 年（从 \textbf{Pascal 和 Fermat 1654 年通信讨论赌博}开始）——但是\textbf{"概率到底是什么"}——一直没有公认的答案：

\begin{itemize}
\item \textbf{Laplace 派}：概率是\textbf{等可能结果的比例}（古典概型）
\item \textbf{频率学派}：概率是\textbf{长期重复的频率}（von Mises）
\item \textbf{Bayes 派}：概率是\textbf{个人信念的强度}
\end{itemize}

\textbf{三派吵了 200 年——没人说服得了谁}。

\medskip

\textbf{Kolmogorov 的天才——就是绕开这个哲学问题}：

\textbf{"我不管你怎么解释概率——只要你的'概率'满足这四条公理——剩下的定理都能从公理推出。"}

\begin{enumerate}
\item $P(A) \geq 0$（非负性）
\item $P(\Omega) = 1$（规范性）
\item $P\left(\bigcup_{i=1}^\infty A_i\right) = \sum_{i=1}^\infty P(A_i)$（可数可加性）
\item 条件概率定义为 $P(A|B) = P(AB)/P(B)$
\end{enumerate}

\medskip

\textbf{就是这四行}——\textbf{概率论从此有了"宪法"}。

\textbf{Kolmogorov 还做了什么}：\textbf{建立了随机过程理论、建立了算法信息论（Kolmogorov 复杂度）、解决了 Hilbert 第 13 问题}——\textbf{他是 20 世纪最伟大的数学家之一}。

\medskip

\textbf{1933 年没有计算机}——\textbf{没有 Python}——\textbf{他用纯思想——把"什么是随机"这个困扰人类 270 年的问题——一锤定音}。

\medskip

你今天学的每一个 \texttt{numpy.random}——\textbf{背后都站着这位莫斯科的年轻人}。
\end{tcolorbox}

\section{直觉 / Intuition}

\subsection{§1 概率基础：从样本空间到条件概率}

\textbf{概率论的三件套——一句话}：

\begin{quote}
\textbf{"样本空间 $\Omega$——事件 $A$——概率 $P(A)$。"}\\[4pt]
\itshape "Sample space, events, probability measure."
\end{quote}

\textbf{样本空间} $\Omega$：所有可能结果的集合。掷骰子——$\Omega = \{1,2,3,4,5,6\}$。

\textbf{事件} $A \subseteq \Omega$：你关心的某些结果。"掷出偶数"——$A = \{2,4,6\}$。

\textbf{概率} $P: \mathcal{F} \to [0,1]$：给每个事件分配一个 $[0,1]$ 之间的数——\textbf{满足 Kolmogorov 四公理}。

\textbf{这就是 Kolmogorov 框架——浙大教材的第一章——本质上就这几行}。

\subsection{条件概率：信息一来——概率就变}

\textbf{条件概率} $P(A|B)$：\textbf{已知 $B$ 发生的前提下——$A$ 发生的概率}。

\begin{equation}
P(A|B) = \frac{P(AB)}{P(B)}, \quad P(B) > 0
\end{equation}

\textbf{直觉}：\textbf{把样本空间从 $\Omega$ 缩小到 $B$}——\textbf{在新的小世界里重新算 $A$ 的概率}。

\begin{example}[著名的"三门问题" / Monty Hall]
\textbf{三扇门——其中一扇后面是车——你选一扇——主持人打开剩下两扇里没车的那一扇——问你换不换}？

\textbf{直觉}：剩下两扇——\textbf{概率 50/50——换不换无所谓}。

\textbf{Bayes 算一下}：你最初选中车的概率 = 1/3。\textbf{所以换——赢的概率是 2/3——不换只有 1/3}。

\textbf{这道题让 Marilyn vos Savant 在 1990 年代被一万个 PhD 围攻}——\textbf{她是对的}。
\end{example}

\subsection{Bayes 定理：先验 + 似然 = 后验}

\textbf{Bayes 公式}——\textbf{机器学习的灵魂}——\textbf{一行话}：

\begin{equation}
\boxed{\,P(H|D) = \frac{P(D|H) \cdot P(H)}{P(D)}\,}
\end{equation}

\begin{itemize}
\item $P(H)$：\textbf{先验}（prior）——\textbf{看到数据之前你对 $H$ 的信念}
\item $P(D|H)$：\textbf{似然}（likelihood）——\textbf{如果 $H$ 是真的——观测到 $D$ 的概率}
\item $P(H|D)$：\textbf{后验}（posterior）——\textbf{看到数据之后你对 $H$ 的更新信念}
\item $P(D)$：\textbf{归一化常数}——\textbf{保证后验加起来等于 1}
\end{itemize}

\textbf{用一句话讲完 Bayes}：\textbf{"信念遇到证据——按比例更新"}。

\begin{example}[医学检测的"经典误判"]
某种病人群\textbf{发病率 0.1\%}。检测\textbf{准确率 99\%}（真阳率 99\%、假阳率 1\%）。

你检测呈阳性——\textbf{你真正得病的概率是多少}？

直觉答："99\% 准确——所以我 99\% 得病了"。\textbf{错}！

Bayes 算：
\[
P(\text{病}|+) = \frac{0.99 \times 0.001}{0.99 \times 0.001 + 0.01 \times 0.999} \approx 9\%
\]

\textbf{你只有 9\% 概率真的得病}——\textbf{91\% 概率是假阳性}！

\textbf{原因}：\textbf{基率（base rate）很低}——\textbf{先验淹没了似然}。

\textbf{这是医生、律师、AI 工程师天天犯的错}——\textbf{Daniel Kahneman 因为研究这个拿了诺贝尔经济学奖}。
\end{example}

\section{数学 / Mathematics}

\subsection{§2 随机变量与分布}

\textbf{浙大教材的常见错误}——把随机变量讲成\textbf{"一个会随机取值的变量"}。

\textbf{真正的定义}（Kolmogorov 框架）：

\begin{equation}
X: \Omega \to \mathbb{R} \quad \text{是一个可测函数}
\end{equation}

\textbf{随机变量 $X$ 是从样本空间 $\Omega$ 到实数轴 $\mathbb{R}$ 的函数}——\textbf{它本身一点也不"随机"}——\textbf{随机的是输入 $\omega$}。

\textbf{分布}是 $X$ 把 $\Omega$ 上的概率\textbf{推到} $\mathbb{R}$ 上的像：
\[
P_X(A) = P(X^{-1}(A)) = P(\{\omega: X(\omega) \in A\})
\]

\textbf{常见分布——一张表}：

\begin{center}
\begin{tabular}{l|l|l}
\textbf{分布} & \textbf{何时用} & \textbf{典型工程场景} \\
\hline
Bernoulli$(p)$ & 抛硬币 & 二分类标签 \\
Binomial$(n,p)$ & $n$ 次硬币 & A/B 测试转化人数 \\
Poisson$(\lambda)$ & 单位时间事件数 & 网络包到达 \\
Uniform$(a,b)$ & 等可能 & 随机种子 / 蒙特卡洛 \\
Normal$(\mu,\sigma^2)$ & 噪声、误差 & 几乎一切信号 \\
Exponential$(\lambda)$ & 等待时间 & 排队、寿命 \\
\end{tabular}
\end{center}

\textbf{正态分布的密度}（重要到要背）：
\begin{equation}
f(x) = \frac{1}{\sqrt{2\pi}\,\sigma}\exp\!\left(-\frac{(x-\mu)^2}{2\sigma^2}\right)
\end{equation}

\textbf{多元情形}——用\textbf{均值向量} $\bm{\mu}$ 和\textbf{协方差矩阵} $\bm{\Sigma}$：
\begin{equation}
f(\bm{x}) = \frac{1}{(2\pi)^{n/2}|\bm{\Sigma}|^{1/2}}\exp\!\left(-\tfrac{1}{2}(\bm{x}-\bm{\mu})^\top\bm{\Sigma}^{-1}(\bm{x}-\bm{\mu})\right)
\end{equation}

\textbf{这是机器学习里出现频率最高的密度函数}——\textbf{GMM、Kalman 滤波、Gaussian Process——都靠它}。

\subsection{§3 期望、方差、协方差}

\textbf{期望}——\textbf{随机变量的"重心"}：
\begin{equation}
E(X) = \int_\Omega X \, dP = \int_{\mathbb{R}} x \, dP_X(x)
\end{equation}

\textbf{方差}——\textbf{偏离重心的平方平均}：
\begin{equation}
D(X) = E\!\left[(X-E(X))^2\right] = E(X^2) - [E(X)]^2
\end{equation}

\textbf{协方差}（两个变量"共变"的程度）：
\begin{equation}
\mathrm{Cov}(X,Y) = E[(X - E(X))(Y - E(Y))]
\end{equation}

\textbf{相关系数}——\textbf{把协方差归一化到 $[-1, 1]$}：
\begin{equation}
\rho_{XY} = \frac{\mathrm{Cov}(X,Y)}{\sqrt{D(X)D(Y)}}
\end{equation}

\textbf{多变量情形}——\textbf{协方差矩阵} $\bm{\Sigma}$：
\begin{equation}
\bm{\Sigma}_{ij} = \mathrm{Cov}(X_i, X_j), \quad \bm{\Sigma} = E[(\bm{X}-\bm{\mu})(\bm{X}-\bm{\mu})^\top]
\end{equation}

\textbf{协方差矩阵是对称半正定的}——\textbf{它的特征值告诉你数据沿哪个方向方差最大}——\textbf{这就是 PCA 的核心}（见第 8 章）。

\subsection{§4 大数定律与中心极限定理}

\textbf{概率论中最重要的两个定理}——\textbf{统计学的两根柱子}。

\textbf{弱大数定律}（Khinchin）：\textbf{独立同分布样本的平均值——依概率收敛到真期望}。
\begin{equation}
\bar{X}_n = \frac{1}{n}\sum_{i=1}^n X_i \xrightarrow{P} \mu \quad \text{当}~ n \to \infty
\end{equation}

\textbf{一句话}：\textbf{"样本多了——平均值就稳了"}。这就是为什么民意调查抽 1000 人就够、为什么蒙特卡洛能逼近真值。

\textbf{中心极限定理}（CLT）：\textbf{独立同分布样本的和——标准化后——收敛到正态分布}。
\begin{equation}
\frac{\bar{X}_n - \mu}{\sigma/\sqrt{n}} \xrightarrow{d} \mathcal{N}(0,1)
\end{equation}

\textbf{这是整个统计学最神奇的事实}——\textbf{不管原始分布长什么样}（哪怕是骰子、Bernoulli、指数）——\textbf{足够多个独立样本加起来——一定变成正态}。

\textbf{这就是为什么正态分布无处不在}：测量误差、身高、考试分数、股票回报——\textbf{它们都是大量独立小因素叠加的结果}。

\subsection{§5 随机过程入门：Markov + Poisson}

\textbf{随机过程}——\textbf{"时间上的随机变量族"} $\{X_t\}_{t \in T}$。

\textbf{Markov 链}——\textbf{"未来只依赖于现在——不依赖于过去"}：
\begin{equation}
P(X_{n+1} = j \mid X_n = i, X_{n-1}, \dots, X_0) = P(X_{n+1} = j \mid X_n = i) = p_{ij}
\end{equation}

\textbf{转移矩阵} $\bm{P} = (p_{ij})$ 是\textbf{行随机矩阵}（每行和为 1）。\textbf{长期分布}是 $\bm{P}^\top$ 对应特征值 1 的特征向量——\textbf{这正是第 7 章的内容}。

\textbf{Markov 链的应用}：\textbf{Google PageRank}（网页是状态、链接是转移）、\textbf{语音识别 HMM}、\textbf{强化学习 MDP}、\textbf{ChatGPT 的下一个 token}。

\textbf{Poisson 过程}——\textbf{"单位时间内事件数服从 Poisson 分布——不重叠区间独立"}：
\begin{equation}
P(N(t) = k) = \frac{(\lambda t)^k}{k!} e^{-\lambda t}
\end{equation}

\textbf{典型应用}：\textbf{电话呼叫到达}、\textbf{网络包到达}、\textbf{地震发生}、\textbf{放射性衰变}——\textbf{凡是"小概率事件大量重复"——都用 Poisson}。

\section{Aha 时刻 / The Aha Moment}

\begin{tcolorbox}[ahabox={\Large Aha：蒙特卡洛求 $\pi$ 与期权定价}]
\textbf{你在浙大教材里学了 200 页概率}——\textbf{却从来没用过它干一件"工程"的事}。

\medskip

\textbf{这一节告诉你}——\textbf{概率最优雅的现代应用之一}：\textbf{用随机数计算确定的数}。

\textbf{蒙特卡洛方法}（Monte Carlo method）——\textbf{1940 年代}由 \textbf{Stanislaw Ulam 和 John von Neumann} 在洛斯阿拉莫斯研究原子弹时发明——\textbf{以摩纳哥赌场命名}（Ulam 的叔叔是赌徒）。

\medskip

\textbf{第一站：求 $\pi$}。

往边长为 2 的正方形里随机扔 $N$ 个点。落在内切圆里的点数为 $M$。

\[
\frac{M}{N} \approx \frac{\text{圆面积}}{\text{正方形面积}} = \frac{\pi}{4} \implies \pi \approx \frac{4M}{N}
\]

\textbf{大数定律保证了 $N \to \infty$ 时——这个估计收敛到真值}。

\medskip

\textbf{第二站：期权定价（金融工程）}。

\textbf{欧式看涨期权}：\textbf{今天买一张合同——$T$ 年后有权利按价格 $K$ 买入一只股票}。

\textbf{Black-Scholes 模型假设}：股票价格 $S_t$ 服从\textbf{几何 Brown 运动}：
\[
S_T = S_0 \exp\!\left((r - \tfrac{\sigma^2}{2})T + \sigma\sqrt{T} \cdot Z\right), \quad Z \sim \mathcal{N}(0,1)
\]

\textbf{期权今天的价格}：
\[
C = e^{-rT} \cdot E[\max(S_T - K, 0)]
\]

\textbf{这个期望——可以用 Black-Scholes 公式解析解}（1997 年诺贝尔经济学奖）——\textbf{也可以用蒙特卡洛模拟}：

\begin{enumerate}
\item 模拟 $N$ 条股票路径 $S_T^{(1)}, \dots, S_T^{(N)}$
\item 计算每条路径的支付 $\max(S_T^{(i)} - K, 0)$
\item 取平均、再贴现 $e^{-rT}$——\textbf{就是期权价格}
\end{enumerate}

\textbf{大数定律 + 中心极限定理——告诉你估计的精度是 $O(1/\sqrt{N})$}——\textbf{不依赖于维度}。

\textbf{这是蒙特卡洛在高维问题（多资产期权、信用衍生品）打败一切确定性算法的原因}——\textbf{华尔街每天烧 GPU 跑的就是这件事}。

\medskip

\textbf{看代码——只要 30 行}。
\end{tcolorbox}

\begin{lstlisting}[language=Python, caption={probability.py · 蒙特卡洛求 $\pi$ 与期权定价}]
import numpy as np
import matplotlib.pyplot as plt

# =============================================================
# 1. Monte Carlo for pi
# =============================================================
def monte_carlo_pi(N=1_000_000, seed=0):
    rng = np.random.default_rng(seed)
    x = rng.uniform(-1, 1, N)
    y = rng.uniform(-1, 1, N)
    inside = (x**2 + y**2) <= 1
    return 4.0 * inside.mean()

print(f"pi approx = {monte_carlo_pi(10**6):.6f}")
print(f"true pi   = {np.pi:.6f}")

# =============================================================
# 2. Central Limit Theorem in action
# =============================================================
def clt_demo(dist_sampler, n_samples=30, n_trials=10000):
    """The sum of any iid distribution converges to Normal."""
    sums = np.array([dist_sampler(n_samples).mean()
                     for _ in range(n_trials)])
    return sums

# Try with a very non-normal source: Bernoulli
rng = np.random.default_rng(42)
sums = clt_demo(lambda n: rng.binomial(1, 0.3, n))
print(f"sample mean = {sums.mean():.4f}, expected = 0.3")
print(f"sample std  = {sums.std():.4f}, expected = "
      f"{np.sqrt(0.3*0.7/30):.4f}")

# =============================================================
# 3. European call option pricing via Monte Carlo
# =============================================================
def mc_european_call(S0, K, T, r, sigma, N=200_000, seed=0):
    """Black-Scholes European call by Monte Carlo."""
    rng = np.random.default_rng(seed)
    Z = rng.standard_normal(N)
    ST = S0 * np.exp((r - 0.5*sigma**2)*T + sigma*np.sqrt(T)*Z)
    payoff = np.maximum(ST - K, 0.0)
    price = np.exp(-r*T) * payoff.mean()
    stderr = np.exp(-r*T) * payoff.std() / np.sqrt(N)
    return price, stderr

# Test against closed-form Black-Scholes
from scipy.stats import norm
def bs_call(S0, K, T, r, sigma):
    d1 = (np.log(S0/K) + (r+0.5*sigma**2)*T) / (sigma*np.sqrt(T))
    d2 = d1 - sigma*np.sqrt(T)
    return S0*norm.cdf(d1) - K*np.exp(-r*T)*norm.cdf(d2)

S0, K, T, r, sigma = 100, 100, 1.0, 0.05, 0.2
mc_price, se = mc_european_call(S0, K, T, r, sigma)
bs_price = bs_call(S0, K, T, r, sigma)
print(f"MC price  = {mc_price:.4f} (+/- {1.96*se:.4f})")
print(f"BS price  = {bs_price:.4f}")

# =============================================================
# 4. Bayes update for the medical test
# =============================================================
def bayes_update(prior, likelihood_pos_given_H,
                 likelihood_pos_given_notH):
    evidence = (likelihood_pos_given_H * prior +
                likelihood_pos_given_notH * (1 - prior))
    return likelihood_pos_given_H * prior / evidence

posterior = bayes_update(0.001, 0.99, 0.01)
print(f"P(disease | positive) = {posterior:.4f}")
\end{lstlisting}

\textbf{跑一遍——你会看到}：

\begin{itemize}
\item \texttt{pi approx = 3.141980}——\textbf{蒙特卡洛 100 万样本——精度到 3 位}
\item \texttt{MC price = 10.4503 (+/- 0.0653)}——\textbf{和 Black-Scholes 解析解 10.4506 几乎一致}
\item \texttt{P(disease | positive) = 0.0902}——\textbf{Bayes 算出的医学误判率 9\%}
\end{itemize}

\textbf{30 行代码——把蒙特卡洛、CLT、Black-Scholes、Bayes——全部跑通}。

\section{现代视角 / The Modern Lens}

\subsection{§6 应用一：信号与噪声}

\textbf{现代通信工程的核心模型}：\textbf{接收信号 = 发送信号 + 高斯白噪声}。
\begin{equation}
y(t) = s(t) + n(t), \quad n(t) \sim \mathcal{N}(0, \sigma^2)
\end{equation}

\textbf{为什么是高斯}？——\textbf{中心极限定理}：\textbf{热噪声是无数电子的独立扰动叠加——CLT 保证它服从正态}。

\textbf{为什么是"白"}？——\textbf{所有频率成分功率谱平等}——\textbf{白光的类比}。

\textbf{Kalman 滤波}（\textbf{阿波罗登月、特斯拉自动驾驶}的核心算法）——\textbf{本质上就是 Bayesian 更新}：\textbf{用先验状态 + 新的观测 → 更新后验状态}。

\subsection{§6 应用二：机器学习}

\textbf{机器学习的核心}——一句话——\textbf{"从数据中学一个分布"}。

\begin{itemize}
\item \textbf{监督学习}：学条件分布 $p(y|x)$
\item \textbf{生成模型}：学联合分布 $p(x,y)$ 或 $p(x)$
\item \textbf{Bayesian 推断}：把模型参数 $\bm{\theta}$ 也看成随机变量——\textbf{用 Bayes 更新}：$p(\bm{\theta}|D) \propto p(D|\bm{\theta}) p(\bm{\theta})$
\item \textbf{交叉熵损失} $= -\sum y \log p$ \textbf{本质上是似然的对数}——\textbf{深度学习训练就是 MLE}
\end{itemize}

\textbf{Diffusion 模型}（DALL-E、Stable Diffusion 的核心）：\textbf{学一个去噪过程——把高斯噪声逆向地恢复成图片}——\textbf{完全基于随机过程理论}。

\textbf{Transformer 的 attention 权重}——\textbf{softmax 输出是概率分布}——\textbf{下一个 token 的预测就是 $p(x_{t+1} | x_1, \dots, x_t)$}——\textbf{ChatGPT 本质上是一个超大的 Markov 链}。

\subsection{§6 应用三：量子}

\textbf{量子力学是 20 世纪最深刻的物理理论}——\textbf{它的根基是概率}。

\textbf{波函数} $\psi(\bm{x}, t)$ 的模平方 $|\psi|^2$ 是\textbf{粒子在位置 $\bm{x}$ 的概率密度}（Born 法则，1926）。

\textbf{Schrödinger 方程描述 $\psi$ 的演化}——\textbf{是确定性的}——\textbf{但测量瞬间——波函数塌缩——结果是随机的}。

\textbf{Einstein 不接受这件事}——\textbf{"God does not play dice"}——\textbf{但他错了}。\textbf{Bell 不等式实验}（1980 年代）\textbf{彻底排除了"隐变量"}——\textbf{量子的随机性是宇宙的基本属性}——\textbf{不是因为我们"不知道"}。

\textbf{量子计算}——\textbf{利用叠加态的概率振幅做并行计算}——\textbf{Shor 算法因式分解、Grover 算法搜索}——\textbf{都是"用概率振幅而非概率"做计算}的实例。

\section{代码与思考 / Code \& Reflection}

\textbf{把 \texttt{probability.py} 跑通——你已经会用概率"做事"了}。

\textbf{再做三件事——把直觉锁住}：

\begin{enumerate}
\item \textbf{改变蒙特卡洛样本数 $N$}——观察 $\pi$ 的估计误差——\textbf{画 log-log 图}——你会看到\textbf{斜率 $-1/2$}——这就是 \textbf{CLT 给的 $O(1/\sqrt{N})$ 收敛速率}。

\item \textbf{用 \texttt{rng.uniform / rng.exponential / rng.poisson} 替换 CLT 演示里的 Bernoulli}——观察样本均值的直方图——\textbf{每一种原始分布——都会变成正态钟形}——\textbf{亲眼看见 CLT 的魔法}。

\item \textbf{把 Bayes 更新里的先验从 0.001 改成 0.5}——再看后验——\textbf{你会发现后验是 99\%}——\textbf{这说明 Bayes 的结论严重依赖于先验}——\textbf{这就是为什么 base rate 是统计学第一坑}。
\end{enumerate}

\textbf{再问自己三个问题}：

\begin{itemize}
\item \textbf{为什么概率必须满足"可数可加性"}（公理 3）——\textbf{而不是有限可加性}？\\
\textit{答}：\textbf{为了能处理"无限多个事件"——比如"在 $[0,1]$ 上随机取点"——必须有可数可加才能定义 Lebesgue 测度——这是 Kolmogorov 1933 年的关键洞察}。

\item \textbf{为什么 ML 损失函数最常用的是"交叉熵"而不是"准确率"}？\\
\textit{答}：\textbf{交叉熵 = $-\log$ 似然——它对概率本身敏感、可微分、有 CLT 性质}——\textbf{准确率不可微——梯度下降走不动}。

\item \textbf{为什么期权定价不直接用真实概率——而用"风险中性概率"}？\\
\textit{答}：\textbf{这是金融工程最深的一课}——\textbf{无套利论证}（Harrison-Pliska 1981）\textbf{证明}：\textbf{在没有套利的市场——存在一个等价测度 $Q$——使得贴现价格是鞅}——\textbf{用 $Q$ 算期望、再贴现——就是无套利价格}——\textbf{这是 Black-Scholes 真正的数学基础}。
\end{itemize}

\textbf{这一章——是\textbf{打通工科数学任督二脉}的一个关键节点}：\textbf{从此你看到的世界——不是确定的——而是分布的}。

\section{本章小结 / Chapter Summary}

\begin{tcolorbox}[essbox={本章小结}]
\textbf{中文}：
\begin{enumerate}
\item \textbf{Kolmogorov 1933 年用四条公理把概率变成现代数学}——概率是测度
\item \textbf{Bayes 定理是"信念遇到证据按比例更新"}——医学误判 / ML 的根
\item \textbf{随机变量是函数 $X: \Omega \to \mathbb{R}$}——分布是测度的像
\item \textbf{期望是重心、方差是离散度、协方差矩阵是 PCA 的入口}
\item \textbf{大数定律保证平均收敛、CLT 解释了正态无处不在}——\textbf{统计学的两根柱子}
\item \textbf{Markov 链 = PageRank/HMM/RL 的核心；Poisson 过程 = 排队/通信的核心}
\item \textbf{蒙特卡洛 + Black-Scholes 让概率走进华尔街}——\textbf{现代金融工程的发动机}
\item \textbf{深度学习本质是从数据中学一个分布}——\textbf{交叉熵就是负对数似然}
\item \textbf{量子力学告诉我们}——\textbf{随机性是宇宙的基本属性}
\end{enumerate}
\textbf{English}: Kolmogorov's 1933 axioms made probability a measure; Bayes' rule updates belief in proportion to evidence; random variables are measurable functions and distributions are their pushforwards; expectation is integration, variance is spread, the covariance matrix powers PCA; the law of large numbers guarantees convergence of sample means and the central limit theorem explains the ubiquity of the normal distribution; Markov chains and Poisson processes are the backbones of modern stochastic modeling; Monte Carlo and Black–Scholes brought probability to Wall Street; deep learning is fundamentally learning a distribution, and quantum mechanics reminds us that randomness is a basic feature of reality itself.
\end{tcolorbox}

\textbf{下一章——我们将进入第 11 章：数理统计与推断}——\textbf{把这一章的"概率"用来"从数据里学规律"}——\textbf{这是现代数据科学真正的起点}。

\clearpage

% =====================================================================
%  第 11 章 · 数理统计 / Statistics
%  描述性统计、参数估计、假设检验、线性回归、方差分析、应用
%  小故事：R.A. Fisher 与 Rothamsted 农场
%  Aha：A/B 测试 + 推荐系统（线性回归 = ML 的最简形态）
% =====================================================================

\chapter{数理统计 / Statistics}
\label{ch:stats}

\begin{quote}\itshape
1919 年——英国 Rothamsted 农业试验站——\\
一个 29 岁、近视到几乎失明的年轻人 R.A. Fisher——\\
被分配去整理\textbf{75 年来积压的土壤、肥料、产量数据}。\\
他在那个农场——\textbf{独自一人}——\textbf{发明了今天我们叫"现代统计学"的几乎一切}：\\
\textbf{方差分析、最大似然、实验设计、p 值}……\\
统计学——\textbf{真的是从田里长出来的}。\\
这一章——我们就讲他在那个农场里想清楚的事。
\end{quote}

\section{写在前面 / Preface}

\textbf{第 10 章我们讲了概率论}——\textbf{概率是从模型推数据}——\textbf{已知一枚硬币是公平的、问扔 10 次有几次正面}。

\textbf{这一章讲的是反过来的事}——\textbf{统计是从数据推模型}——\textbf{我捡到一枚硬币、扔了 100 次、73 次正面}——\textbf{问这枚硬币是不是公平的}。

\textbf{这是两个方向相反的学科}——\textbf{但用的是同一套概率工具}。

\textbf{我刚学统计时}——\textbf{和大多数人一样}——\textbf{我学的是公式}：

样本均值 $\bar{X} = \frac{1}{n}\sum X_i$。样本方差 $S^2 = \frac{1}{n-1}\sum(X_i - \bar{X})^2$。t 统计量、$\chi^2$ 统计量、F 统计量……

\textbf{我考试考得很好}——\textbf{但你问我"什么时候用 t 检验、什么时候用 $\chi^2$"}——\textbf{我说不清楚}。

\textbf{真正让我懂统计的}——\textbf{是大三那年我做了一个推荐系统的项目}——\textbf{我要回答一个问题}：\textbf{"新版本的推荐算法——比旧版本好吗？"}

我打开 Python：

\begin{lstlisting}[language=Python]
from scipy import stats
t_stat, p_value = stats.ttest_ind(old_ctr, new_ctr)
print(f"p = {p_value:.4f}")
# p = 0.0023  -> 新版本显著更好
\end{lstlisting}

\textbf{三行代码——做完了一次 A/B 测试}。

那一刻我才反应过来——\textbf{我学的所有公式——背后是同一个思想}：\textbf{"观察到的差异——是真的吗？还是运气？"}

\textbf{这一章就把这件事讲清楚}——再多的统计公式——都是这一个问题的变体。

\section{一句话本质 / The Essence in One Sentence}

\begin{tcolorbox}[essbox={一句话本质}]
\textbf{中文}：\textbf{统计是从有限的数据——推断背后的真相——并量化这个推断有多可信}。\\[6pt]
\textbf{English}: Statistics is the art of inferring truth from limited data, while quantifying how much we can trust that inference.
\end{tcolorbox}

记住这一句话——\textbf{这一章已经讲完了一半}。

\textbf{"量化可信度"——这五个字是统计的灵魂}。\textbf{物理学说"光速是 $3 \times 10^8$ m/s"}——\textbf{统计学说"我 95\% 相信真值在 [2.99, 3.01]$\times 10^8$ 之间"}——\textbf{前者是断言、后者是带不确定性的断言}——\textbf{真实世界里——后者才有用}。

\section{教材里你看到的 / What the Textbook Tells You}

\subsection{浙大《概率论与数理统计》}

\textbf{浙江大学盛骤主编}的\textbf{《概率论与数理统计》}——\textbf{是中国工科考研指定的统计教材}——\textbf{一半概率、一半统计}。

\textbf{优点}：\textbf{严谨、覆盖考研所有考点}。

\textbf{缺点}：

\begin{enumerate}
\item \textbf{先讲一堆抽样分布}（$\chi^2$、$t$、$F$ 分布的密度函数）——\textbf{学生不知道这些分布是干什么的}
\item \textbf{假设检验讲了 8 种情况}（单总体、双总体、方差已知、方差未知……）——\textbf{学生记不住}
\item \textbf{应用部分薄弱}——\textbf{回归只讲一元、方差分析点到为止}——\textbf{真实工业用的是多元回归 + 大规模 A/B 测试}
\end{enumerate}

\subsection{Casella \& Berger《Statistical Inference》}

\textbf{美国研究生统计课的国民教材}——\textbf{从测度论一直讲到贝叶斯}——\textbf{非常严谨}——\textbf{600 页}——\textbf{不适合工科生}。

\subsection{Wasserman《All of Statistics》}

\textbf{卡内基梅隆 ML 系的统计教材}——\textbf{薄、快、现代}——\textbf{适合工科生重读}——\textbf{我强烈推荐}。

\subsection{我的策略}

\textbf{本章不会重复浙大教材的细节}——\textbf{它已经够全了}。

\textbf{本章只做三件事}：

\begin{enumerate}
\item \textbf{把统计的核心思想讲清楚}——\textbf{"观察 vs 真相"——"信号 vs 噪声"}
\item \textbf{把假设检验的"为什么"讲透}——\textbf{p 值不是你以为的那个意思}
\item \textbf{把线性回归讲到能跑代码——并指出它就是 ML 的最简形态}
\end{enumerate}

\begin{tcolorbox}[storybox={\Large 小故事：R.A. Fisher 与 Rothamsted 农场}]
\textbf{1919 年}——\textbf{Ronald Aylmer Fisher} 29 岁——\textbf{近视到几乎失明}（从小靠口算学数学——因为读书太费眼）——\textbf{找了好几年工作都被拒}。

\medskip

\textbf{Rothamsted 农业试验站}——\textbf{当时世界上最老的农业研究机构}（成立于 1843 年）——\textbf{积压了 75 年的肥料、土壤、产量数据}——\textbf{需要有人整理}。

\medskip

\textbf{他们以为请来了一个数据录入员}——\textbf{他们请来的是统计学的爱因斯坦}。

\medskip

\textbf{Fisher 在 Rothamsted 待了 14 年}——\textbf{独自一人}——\textbf{从那堆农业数据里}——\textbf{发明了今天我们说的"现代统计学"几乎所有概念}：

\begin{enumerate}
\item \textbf{方差分析（ANOVA）}——\textbf{为了判断不同肥料是否真的造成产量差异}
\item \textbf{最大似然估计（MLE）}——\textbf{为了从有限样本估计真实参数}
\item \textbf{实验设计（DoE）}——\textbf{随机化、区组、析因设计——所有现代实验的根基}
\item \textbf{p 值与显著性检验}——\textbf{虽然这个概念后来被滥用——但它的本意是"信号能否压过噪声"}
\item \textbf{Fisher 信息量}——\textbf{量化"一个数据点能告诉你多少关于参数的信息"}
\end{enumerate}

\medskip

\textbf{1925 年}——他写了\textbf{《Statistical Methods for Research Workers》}——\textbf{这本书}——\textbf{被翻译成几十种语言}——\textbf{全世界的农学家、生物学家、心理学家——人手一本}。

\medskip

\textbf{统计学不是从象牙塔里发明的}——\textbf{是 Fisher 蹲在 Rothamsted 的田边}——\textbf{为了回答"这块地比那块地高 5\% 的产量——是肥料的功劳——还是运气？"}——\textbf{一点一点想出来的}。

\medskip

\textbf{今天}——\textbf{Google、Facebook、字节跳动}——\textbf{每天跑成千上万次 A/B 测试}——\textbf{用的还是 Fisher 在那片麦田里想出来的方法}。

\medskip

\textbf{你今天学的每一个 p 值}——\textbf{背后都有那片麦田}。
\end{tcolorbox}

\section{直觉 / Intuition}

\subsection{描述性统计：把数据压缩成几个数}

\textbf{你有 1000 个学生的数学成绩}——\textbf{你不可能挨个看}——\textbf{你需要把这 1000 个数压缩成几个关键数字}。

\textbf{三个最重要的数}：

\begin{itemize}
\item \textbf{均值 mean} $\bar{X} = \frac{1}{n}\sum X_i$——\textbf{中心在哪}
\item \textbf{方差 variance} $S^2 = \frac{1}{n-1}\sum(X_i - \bar{X})^2$——\textbf{散得多开}
\item \textbf{相关系数 correlation} $r = \frac{\sum(X_i-\bar{X})(Y_i-\bar{Y})}{\sqrt{\sum(X_i-\bar{X})^2 \sum(Y_i-\bar{Y})^2}}$——\textbf{两个变量同涨同跌的程度}
\end{itemize}

\textbf{为什么方差要除以 $n-1$ 而不是 $n$？}——\textbf{因为你用了样本均值 $\bar{X}$ 代替真均值 $\mu$}——\textbf{这"偷"掉了一个自由度}——\textbf{要把它还回来}——\textbf{这叫 Bessel 校正}。

\begin{example}[相关 $\neq$ 因果]
\textbf{冰淇淋销量}和\textbf{溺水死亡人数}——\textbf{相关系数 $r = 0.9$}——\textbf{是因为吃冰淇淋会让人溺水吗}？

\textbf{不是}——\textbf{是因为夏天到了}——\textbf{冰淇淋卖得多、游泳的人也多、溺水自然多}。\textbf{"夏天"是它们共同的原因——叫混淆变量（confounder）}。

\textbf{统计能告诉你"两件事一起发生"——但不能告诉你"谁导致谁"}——\textbf{这是统计的本质局限}——\textbf{记住这一句话——能让你在数据时代少犯 90\% 的错}。
\end{example}

\subsection{参数估计：从样本猜真相}

\textbf{背景}：\textbf{真实分布 $f(x; \theta)$ 的参数 $\theta$ 是未知的}（比如均值 $\mu$、方差 $\sigma^2$）——\textbf{我们只有 $n$ 个样本 $X_1, \ldots, X_n$}——\textbf{怎么猜 $\theta$？}

\textbf{两种猜法}：

\textbf{(1) 点估计 Point Estimation}——\textbf{给一个数}：\textbf{"我猜 $\mu = 5.3$"}。

\textbf{最常用方法——最大似然估计（MLE）}：\textbf{选一个 $\theta$——让我观察到当前数据的概率最大}。

\begin{equation}
\hat{\theta}_{\text{MLE}} = \arg\max_\theta \prod_{i=1}^n f(X_i; \theta)
\end{equation}

\textbf{"哪个 $\theta$ 最能解释我看到的数据——就选哪个"}——\textbf{这是 Fisher 1922 年提出的}——\textbf{今天所有 ML 模型的训练——本质上都是 MLE}。

\textbf{(2) 区间估计 Interval Estimation}——\textbf{给一个范围 + 信心}：\textbf{"我 95\% 相信 $\mu \in [4.8, 5.8]$"}。

\textbf{这就是置信区间（confidence interval）}——\textbf{比点估计更有用}——\textbf{因为它带了不确定性}。

\begin{tcolorbox}[colback=yellow!10, colframe=orange!70!black, title={\textbf{警告：95\% 置信区间不是你以为的那个意思}}]
\textbf{常见的错误理解}：\textbf{"真实的 $\mu$ 有 95\% 的概率落在 [4.8, 5.8] 里"}。

\textbf{正确的理解}：\textbf{"如果我重复这个实验 100 次——每次都算一个区间——大约 95 个区间会包含真实的 $\mu$"}。

\textbf{区别在哪}？——\textbf{真实的 $\mu$ 是一个固定的数——它要么在 [4.8, 5.8] 里、要么不在——没有"概率"可言}。\textbf{有概率的是"我们这个区间"——它是随机的}。

\textbf{这是频率学派和贝叶斯学派吵了 100 年的话题}——\textbf{工业里——你不需要纠结——记住"95\% 的区间会捞到真值"就够了}。
\end{tcolorbox}

\subsection{假设检验：信号能压过噪声吗？}

\textbf{这是 Fisher 思想的核心}——\textbf{也是统计在工业里最常用的一招}。

\textbf{场景}：\textbf{新版推荐算法 vs 旧版}——\textbf{新版点击率 5.3\%、旧版 5.0\%}——\textbf{差 0.3\%}——\textbf{这是真的提升还是随机波动？}

\textbf{假设检验的逻辑}（\textbf{反证法}）：

\begin{enumerate}
\item \textbf{假设两者没区别（原假设 $H_0$）}：\textbf{$\mu_{\text{new}} = \mu_{\text{old}}$}
\item \textbf{算一个"如果 $H_0$ 真——观察到当前数据的概率"——这就是 p 值}
\item \textbf{p 值很小（一般 < 0.05）}——\textbf{说明"如果 $H_0$ 真——这数据太离谱了"}——\textbf{拒绝 $H_0$}——\textbf{相信新版真的更好}
\item \textbf{p 值不小}——\textbf{数据和 $H_0$ 不矛盾}——\textbf{不能拒绝 $H_0$}——\textbf{不代表 $H_0$ 真——只代表"目前证据不够"}
\end{enumerate}

\textbf{这是逻辑上的"宁可放过、不要冤枉"}——\textbf{法庭原则}——\textbf{Fisher 设计这一套的时候——脑子里想的就是法庭}。

\begin{tcolorbox}[colback=red!5, colframe=red!70!black, title={\textbf{p 值不是"原假设为真的概率"}}]
\textbf{p 值的真实含义}：\textbf{"假设原假设 $H_0$ 为真——观察到比当前数据更极端结果的概率"}。

\textbf{形式上}：\textbf{$p = P(\text{data} \mid H_0)$}——\textbf{不是} $P(H_0 \mid \text{data})$。

\textbf{这两个混淆——是科学界最大的统计误用}——\textbf{2016 年美国统计学会专门发了一个声明}——\textbf{警告大家"p < 0.05 不等于结论可信"}。

\textbf{工业里的实用规则}：\textbf{p < 0.01——挺可信}；\textbf{p < 0.001——很可信}；\textbf{p < 0.05——值得继续看、但别下定论}。
\end{tcolorbox}

\textbf{两类错误}：

\begin{itemize}
\item \textbf{第一类错误（弃真）}：\textbf{$H_0$ 真——你拒绝了}——\textbf{冤枉好人}——\textbf{概率 = $\alpha$（显著性水平）}
\item \textbf{第二类错误（取伪）}：\textbf{$H_0$ 假——你没拒绝}——\textbf{放过坏人}——\textbf{概率 = $\beta$}
\item \textbf{统计功效（power）= $1 - \beta$}——\textbf{"真有效果时——能检测出来的概率"}——\textbf{样本量越大、功效越高}
\end{itemize}

\subsection{常用检验：t 检验和 $\chi^2$ 检验}

\textbf{t 检验（t-test）}：\textbf{比较两组数的均值是否相同}。\textbf{两个版本}：

\begin{itemize}
\item \textbf{单样本 t 检验}：\textbf{这组数的均值是不是 $\mu_0$}？
\item \textbf{双样本 t 检验}：\textbf{A 组和 B 组的均值是否相同}？
\end{itemize}

\textbf{统计量}（双样本、方差未知、近似相等）：

\begin{equation}
t = \frac{\bar{X}_A - \bar{X}_B}{S_p \sqrt{\frac{1}{n_A} + \frac{1}{n_B}}}, \quad S_p^2 = \frac{(n_A-1)S_A^2 + (n_B-1)S_B^2}{n_A + n_B - 2}
\end{equation}

\textbf{它服从自由度 $n_A + n_B - 2$ 的 t 分布}——\textbf{查表或用 \texttt{scipy.stats.t} 算 p 值}。

\textbf{$\chi^2$ 检验（卡方检验）}：\textbf{比较"分类数据"的分布是否符合预期}。

\textbf{例子}：\textbf{骰子扔 600 次}——\textbf{每个点数应该出现约 100 次}——\textbf{实际是 95、110、98、102、99、96}——\textbf{这个骰子公平吗}？

\begin{equation}
\chi^2 = \sum_{i=1}^k \frac{(O_i - E_i)^2}{E_i}
\end{equation}

\textbf{$O_i$ 是观察值、$E_i$ 是期望值}——\textbf{这个统计量服从 $\chi^2$ 分布}——\textbf{越大——越说明观察值偏离期望}。

\subsection{线性回归：ML 的最简形态}

\textbf{这是本章最重要的一节}——\textbf{读三遍}。

\textbf{场景}：\textbf{房价 $y$ 和面积 $x$}——\textbf{你有 1000 套房子的数据}——\textbf{画在散点图上——大致是一条直线}——\textbf{你想用这条直线预测新房子的价格}。

\textbf{模型}：\textbf{$y = \beta_0 + \beta_1 x + \varepsilon$}——\textbf{$\beta_0$ 是截距、$\beta_1$ 是斜率、$\varepsilon$ 是噪声}。

\textbf{怎么选 $\beta_0, \beta_1$？}——\textbf{最小二乘法（OLS）}——\textbf{让所有点到直线的"垂直距离平方和"最小}：

\begin{equation}
\min_{\beta_0, \beta_1} \sum_{i=1}^n (y_i - \beta_0 - \beta_1 x_i)^2
\end{equation}

\textbf{求导 = 0——解出来}：

\begin{equation}
\hat{\beta}_1 = \frac{\sum (x_i - \bar{x})(y_i - \bar{y})}{\sum (x_i - \bar{x})^2}, \quad \hat{\beta}_0 = \bar{y} - \hat{\beta}_1 \bar{x}
\end{equation}

\textbf{多元回归}：\textbf{$y = \beta_0 + \beta_1 x_1 + \beta_2 x_2 + \cdots + \beta_p x_p + \varepsilon$}——\textbf{多个特征预测一个目标}。

\textbf{矩阵形式}：\textbf{$\mathbf{y} = X\boldsymbol{\beta} + \boldsymbol{\varepsilon}$}——\textbf{解是}：

\begin{equation}
\boxed{\hat{\boldsymbol{\beta}} = (X^T X)^{-1} X^T \mathbf{y}}
\end{equation}

\textbf{这个公式——叫"正规方程"（normal equation）——你必须背下来}。

\begin{tcolorbox}[colback=blue!5, colframe=blue!70!black, title={\textbf{核心洞察：线性回归 = 所有 ML 的种子}}]
\textbf{把线性回归学透——你就懂了 80\% 的监督学习}。\textbf{为什么}？

\begin{enumerate}
\item \textbf{逻辑回归（分类）}——\textbf{线性回归 + sigmoid 激活}——\textbf{$P(y=1) = \sigma(\boldsymbol{\beta}^T \mathbf{x})$}
\item \textbf{神经网络}——\textbf{多层线性回归 + 非线性激活叠加}——\textbf{每一层都是 $\mathbf{Wx} + \mathbf{b}$}
\item \textbf{支持向量机（SVM）}——\textbf{线性回归 + 最大化间隔 + 核技巧}
\item \textbf{岭回归、Lasso}——\textbf{线性回归 + 正则化}——\textbf{防过拟合}
\item \textbf{深度学习}——\textbf{说到底就是} \textbf{"$\mathbf{Wx} + \mathbf{b}$ 堆很多层 + 非线性 + 反向传播算梯度"}
\end{enumerate}

\textbf{你看}——\textbf{$\mathbf{Wx} + \mathbf{b}$ 这个形式}——\textbf{从 1805 年 Legendre 的最小二乘}——\textbf{活到了 2025 年 GPT 的 transformer}——\textbf{220 年没换过}。

\textbf{把线性回归理解到能闭眼推导 $(X^T X)^{-1} X^T y$ 的程度}——\textbf{你已经掌握了 ML 的根}。\textbf{后面所有的复杂模型——都是这个根上长出来的枝叶}。
\end{tcolorbox}

\subsection{方差分析（ANOVA）：Fisher 的招牌}

\textbf{场景}：\textbf{三种肥料 A/B/C}——\textbf{每种用在 10 块地上}——\textbf{产量不一样}——\textbf{这是肥料的差异、还是地之间的随机波动}？

\textbf{ANOVA 的核心思想}——\textbf{把数据的总方差分解}：

\begin{equation}
\underbrace{SS_{\text{total}}}_{\text{总变异}} = \underbrace{SS_{\text{between}}}_{\text{组间变异（信号）}} + \underbrace{SS_{\text{within}}}_{\text{组内变异（噪声）}}
\end{equation}

\textbf{F 统计量}：

\begin{equation}
F = \frac{SS_{\text{between}}/(k-1)}{SS_{\text{within}}/(n-k)} = \frac{\text{信号}}{\text{噪声}}
\end{equation}

\textbf{F 越大——信号越强}——\textbf{说明组与组之间真的有差异}。\textbf{F 接近 1——组间和组内差不多——说明所谓的"差异"就是随机波动}。

\textbf{实验设计（DoE）}——\textbf{Fisher 的另一项发明}——\textbf{三个原则}：

\begin{itemize}
\item \textbf{随机化（randomization）}——\textbf{随机分配处理}——\textbf{抵消混淆变量}
\item \textbf{重复（replication）}——\textbf{每个处理多做几次}——\textbf{估计噪声大小}
\item \textbf{区组（blocking）}——\textbf{把相似的实验单元放一起}——\textbf{减少噪声}
\end{itemize}

\textbf{这三条原则}——\textbf{今天的 A/B 测试、临床试验、新药审批}——\textbf{全部遵循}。

\section{Aha 例子：A/B 测试 + 推荐系统}

\subsection{背景}

\textbf{某互联网公司的推荐系统}——\textbf{旧版算法 A、新版算法 B}——\textbf{产品经理问}：\textbf{"B 比 A 好吗？能上线吗？"}

\textbf{这是字节跳动、Meta、Google 每天发生几百次的事}。

\subsection{方法：A/B 测试}

\textbf{步骤}：

\begin{enumerate}
\item \textbf{随机把用户分成两组}——A 组用旧算法、B 组用新算法
\item \textbf{跑一段时间（一般 1$\sim$2 周）}——\textbf{收集每个用户的点击率（CTR）}
\item \textbf{用双样本 t 检验}——\textbf{看两组的 CTR 是否显著不同}
\item \textbf{p < 0.05 + B 的均值高}——\textbf{决定上线 B}
\end{enumerate}

\subsection{Python 实现}

\begin{lstlisting}[language=Python, caption={A/B 测试 + 线性回归}]
import numpy as np
from scipy import stats
from sklearn.linear_model import LinearRegression

# === 1. A/B 测试：两组 CTR 数据 ===
np.random.seed(42)
ctr_A = np.random.normal(0.050, 0.02, 5000)  # 旧版 CTR
ctr_B = np.random.normal(0.053, 0.02, 5000)  # 新版 CTR

# 描述性统计
print(f"A: mean={ctr_A.mean():.4f}, std={ctr_A.std():.4f}")
print(f"B: mean={ctr_B.mean():.4f}, std={ctr_B.std():.4f}")
print(f"提升: {(ctr_B.mean()/ctr_A.mean() - 1)*100:.2f}%")

# 双样本 t 检验
t_stat, p_value = stats.ttest_ind(ctr_A, ctr_B)
print(f"\nt = {t_stat:.4f}, p = {p_value:.6f}")
if p_value < 0.05:
    print(">>> 显著差异 -> 上线 B")
else:
    print(">>> 不显著 -> 继续观察")

# 95% 置信区间（差值）
diff = ctr_B - ctr_A
ci = stats.t.interval(0.95, len(diff)-1,
                      loc=diff.mean(),
                      scale=stats.sem(diff))
print(f"95% CI for B-A: [{ci[0]:.5f}, {ci[1]:.5f}]")

# === 2. 线性回归：CTR vs 特征 ===
# 假设特征：用户活跃度、内容相关性、推送时段
n = 1000
X = np.random.randn(n, 3)
true_beta = np.array([0.8, 1.2, -0.5])
y = X @ true_beta + 0.3 + np.random.randn(n) * 0.5

model = LinearRegression().fit(X, y)
print(f"\n真实系数:  {true_beta}")
print(f"估计系数:  {model.coef_}")
print(f"R^2: {model.score(X, y):.4f}")

# 手算正规方程验证
X_aug = np.column_stack([np.ones(n), X])
beta_hat = np.linalg.inv(X_aug.T @ X_aug) @ X_aug.T @ y
print(f"正规方程: {beta_hat}")
\end{lstlisting}

\textbf{运行结果}：\textbf{t 检验告诉你 p < 0.001——上线 B}；\textbf{线性回归告诉你哪个特征对 CTR 影响最大}——\textbf{这就是真实推荐系统每天在做的事}。

\subsection{推荐系统的统计学全貌}

\begin{itemize}
\item \textbf{召回}——\textbf{协同过滤、矩阵分解}——\textbf{本质是\textbf{相关系数}和 \textbf{SVD}}
\item \textbf{排序}——\textbf{逻辑回归、GBDT、DNN}——\textbf{本质是\textbf{广义线性模型}}
\item \textbf{评估}——\textbf{AUC、precision、recall}——\textbf{本质是\textbf{二分类的假设检验}}
\item \textbf{上线}——\textbf{A/B 测试}——\textbf{本质是\textbf{双样本 t 检验或 $\chi^2$}}
\item \textbf{长期效应}——\textbf{反事实分析、Uplift 模型}——\textbf{本质是\textbf{因果推断}}
\end{itemize}

\textbf{看到了吗}——\textbf{整个推荐系统——从头到尾——都是统计}。

\section{Module: \texttt{statistics.py}}

\textbf{把本章的所有内容打包成一个 Python 模块}——\textbf{你可以直接用}。

\begin{lstlisting}[language=Python, caption={statistics.py}]
"""
statistics.py - 工科统计工具集
用法: from statistics import describe, ab_test, linreg
"""
import numpy as np
from scipy import stats
from sklearn.linear_model import LinearRegression


def describe(x):
    """描述性统计：均值、方差、四分位"""
    x = np.asarray(x)
    return {
        'n': len(x),
        'mean': x.mean(),
        'std': x.std(ddof=1),
        'min': x.min(),
        'q25': np.percentile(x, 25),
        'median': np.median(x),
        'q75': np.percentile(x, 75),
        'max': x.max(),
    }


def confidence_interval(x, conf=0.95):
    """单样本均值的置信区间"""
    x = np.asarray(x)
    n = len(x)
    return stats.t.interval(conf, n-1,
                            loc=x.mean(),
                            scale=stats.sem(x))


def ab_test(a, b, alpha=0.05):
    """双样本 t 检验，返回判定结果"""
    t_stat, p = stats.ttest_ind(a, b)
    lift = (np.mean(b) - np.mean(a)) / np.mean(a)
    return {
        'mean_A': np.mean(a),
        'mean_B': np.mean(b),
        'lift_pct': lift * 100,
        't_stat': t_stat,
        'p_value': p,
        'significant': p < alpha,
        'recommend': 'B' if (p < alpha and lift > 0) else 'A',
    }


def chi2_test(observed, expected):
    """卡方检验"""
    chi2, p = stats.chisquare(observed, expected)
    return {'chi2': chi2, 'p_value': p,
            'reject_H0': p < 0.05}


def linreg(X, y, intercept=True):
    """线性回归（多元），手算 + sklearn 双验证"""
    X = np.asarray(X)
    y = np.asarray(y)
    if X.ndim == 1:
        X = X.reshape(-1, 1)

    # sklearn 解
    model = LinearRegression(fit_intercept=intercept).fit(X, y)

    # 正规方程解（验证）
    if intercept:
        X_aug = np.column_stack([np.ones(len(X)), X])
    else:
        X_aug = X
    beta = np.linalg.pinv(X_aug.T @ X_aug) @ X_aug.T @ y

    # R^2
    y_pred = model.predict(X)
    ss_res = ((y - y_pred) ** 2).sum()
    ss_tot = ((y - y.mean()) ** 2).sum()
    r2 = 1 - ss_res / ss_tot

    return {
        'beta_sklearn': np.r_[model.intercept_, model.coef_] \
                        if intercept else model.coef_,
        'beta_normal_eq': beta,
        'r_squared': r2,
        'predict': model.predict,
    }


def anova_oneway(*groups):
    """单因素方差分析"""
    f_stat, p = stats.f_oneway(*groups)
    return {'F': f_stat, 'p_value': p,
            'significant': p < 0.05}


if __name__ == '__main__':
    # 演示
    np.random.seed(0)
    a = np.random.normal(50, 10, 200)
    b = np.random.normal(53, 10, 200)
    print('描述:', describe(a))
    print('CI:', confidence_interval(a))
    print('A/B:', ab_test(a, b))

    X = np.random.randn(100, 2)
    y = X @ [1.5, -2.0] + 0.5 + np.random.randn(100) * 0.3
    print('回归:', linreg(X, y)['beta_sklearn'])
\end{lstlisting}

\section{常见陷阱 / Common Pitfalls}

\begin{enumerate}
\item \textbf{p-hacking}——\textbf{反复试不同切片、不同指标——直到 p < 0.05}——\textbf{这是科学造假}——\textbf{学术界已经严打 10 年了}
\item \textbf{多重比较（multiple testing）}——\textbf{同时检验 20 个假设、5\% 的水平}——\textbf{平均有 1 个会"偶然显著"}——\textbf{要做 Bonferroni 校正}
\item \textbf{样本量太小}——\textbf{$n < 30$ 时 t 分布和正态偏离明显}——\textbf{结论不可靠}——\textbf{先用功效分析（power analysis）算最少样本量}
\item \textbf{相关 $\neq$ 因果}——\textbf{再强调一次}——\textbf{$r = 0.99$ 都可能是混淆变量}——\textbf{要因果——必须做随机化实验或反事实分析}
\item \textbf{异常值的处理}——\textbf{均值对异常值极度敏感}——\textbf{有时用中位数、修剪均值（trimmed mean）更稳健}
\item \textbf{过拟合}——\textbf{回归里特征越多、训练 $R^2$ 越高}——\textbf{但泛化变差}——\textbf{用调整后 $R^2$ 或交叉验证判断}
\end{enumerate}

\section{核心公式速查 / Cheat Sheet}

\begin{center}
\begin{tabular}{|l|l|}
\hline
\textbf{概念} & \textbf{公式} \\ \hline
样本均值 & $\bar{X} = \frac{1}{n}\sum X_i$ \\ \hline
样本方差 & $S^2 = \frac{1}{n-1}\sum(X_i - \bar{X})^2$ \\ \hline
相关系数 & $r = \frac{\sum(x_i-\bar{x})(y_i-\bar{y})}{\sqrt{\sum(x_i-\bar{x})^2 \sum(y_i-\bar{y})^2}}$ \\ \hline
95\% 置信区间 & $\bar{X} \pm t_{\alpha/2, n-1} \cdot \frac{S}{\sqrt{n}}$ \\ \hline
单样本 t & $t = \frac{\bar{X} - \mu_0}{S/\sqrt{n}}$ \\ \hline
双样本 t & $t = \frac{\bar{X}_A - \bar{X}_B}{S_p \sqrt{1/n_A + 1/n_B}}$ \\ \hline
$\chi^2$ 检验 & $\chi^2 = \sum \frac{(O_i - E_i)^2}{E_i}$ \\ \hline
线性回归 (一元) & $\hat{\beta}_1 = \frac{\sum(x_i-\bar{x})(y_i-\bar{y})}{\sum(x_i-\bar{x})^2}$ \\ \hline
线性回归 (多元) & $\hat{\boldsymbol{\beta}} = (X^T X)^{-1} X^T \mathbf{y}$ \\ \hline
$R^2$ & $R^2 = 1 - \frac{SS_{\text{res}}}{SS_{\text{tot}}}$ \\ \hline
F 统计量 (ANOVA) & $F = \frac{MS_{\text{between}}}{MS_{\text{within}}}$ \\ \hline
MLE & $\hat{\theta} = \arg\max_\theta \prod f(X_i; \theta)$ \\ \hline
\end{tabular}
\end{center}

\section{本章小结 / Summary}

\textbf{这一章我们讲了 6 件事}：

\begin{enumerate}
\item \textbf{描述性统计}——\textbf{均值、方差、相关}——\textbf{把数据压缩成几个关键数}
\item \textbf{参数估计}——\textbf{MLE 给点估计、置信区间给范围 + 不确定性}
\item \textbf{假设检验}——\textbf{t 检验比均值、$\chi^2$ 检验比分布}——\textbf{p < 0.05 的真正含义}
\item \textbf{线性回归}——\textbf{$\hat{\boldsymbol{\beta}} = (X^T X)^{-1} X^T \mathbf{y}$}——\textbf{ML 的根}
\item \textbf{方差分析}——\textbf{Fisher 的发明}——\textbf{F = 信号/噪声}
\item \textbf{应用}——\textbf{A/B 测试是 t 检验的工业版}——\textbf{推荐系统从头到尾都是统计}
\end{enumerate}

\textbf{最重要的一句话}——\textbf{读三遍}：

\begin{tcolorbox}[colback=green!5, colframe=green!50!black]
\textbf{线性回归是 ML 的种子}——\textbf{深度理解线性回归——就是理解了 80\% 的监督学习}。

\textbf{从 1805 年 Legendre 的最小二乘——到 2025 年 GPT 的 transformer}——\textbf{$\mathbf{Wx} + \mathbf{b}$ 这个形式——220 年没换过}。

\textbf{当别人在追新模型时——你回头把线性回归学透——你就赢了}。
\end{tcolorbox}

\textbf{下一章}——\textbf{我们讲数值方法}——\textbf{当方程没有解析解时——计算机怎么算}——\textbf{这是把数学变成代码的最后一公里}。

\clearpage


\part{计算与优化 / Computation and Optimization}
% =====================================================================
%  第 12 章 · 数值方法 / Numerical Methods
%  数值微分、数值积分、非线性方程、线性方程组、ODE、误差分析
%  小故事：ENIAC (1945) 与现代数值计算的诞生
%  Aha：Newton 法求开方 + GPS 定位算法
% =====================================================================

\chapter{数值方法 / Numerical Methods}
\label{ch:numerical}

\begin{quote}\itshape
1945 年——宾夕法尼亚大学地下室——\textbf{ENIAC}——\textbf{30 吨重}——\textbf{18000 个真空管}——\textbf{每秒 5000 次加法}。\\
冯·诺依曼坐在它面前——\textbf{第一次让数学公式}——\textbf{在机器里"跑"了起来}。\\
这一章——我们要讲的——是\textbf{把数学变成"能算的东西"}的艺术——\textbf{数值方法}。
\end{quote}

\section{写在前面 / Preface}

\textbf{我研究生第一年}——\textbf{导师扔给我一个题}：\textbf{求解一个 ODE 系统}——\textbf{描述化学反应釜的温度演化}。

我打开同济《高等数学》第七版——\textbf{从头翻到尾}——\textbf{没有一个公式能解这个方程}。

我打开 Zill《微分方程》——\textbf{找到了 5 种"解析解法"}——\textbf{我的方程一种都不适用}。

\textbf{它是非线性的、刚性的、还带随机扰动}。

\textbf{我去问师兄}——他笑了笑——\textbf{打开 Jupyter Notebook}：

\begin{lstlisting}[language=Python]
from scipy.integrate import solve_ivp
sol = solve_ivp(rhs, [0, T], y0, method='RK45')
\end{lstlisting}

\textbf{两行代码——5 秒钟——出图}。

那一刻我才反应过来——\textbf{工业界 99\% 的方程——没有解析解}——\textbf{真正在用的}——\textbf{是数值方法}。

\textbf{这一章}——\textbf{我把六类核心数值方法讲清楚}——\textbf{每一类都给一个 SciPy 的"现成轮子"}——\textbf{你以后再也不用怕"算不出来"}。

\section{一句话本质 / The Essence in One Sentence}

\begin{tcolorbox}[essbox={一句话本质}]
\textbf{中文}：\textbf{数值方法——把无法解析求解的数学问题——化成"有限次加减乘除"——让机器去算}。\\[6pt]
\textbf{English}: Numerical methods convert intractable mathematical problems into finite sequences of arithmetic operations that a machine can execute.
\end{tcolorbox}

记住这一句话——\textbf{数值方法的灵魂}——是\textbf{"近似 + 迭代 + 收敛"}。

\section{教材里你看到的 / What the Textbook Tells You}

\subsection{Burden \& Faires《Numerical Analysis》}

\textbf{Burden \& Faires} 是美国数值分析的国民教材——\textbf{900 页}——\textbf{覆盖了你能想到的所有方法}。

\textbf{优点}：严格、全面、每个方法都有收敛性证明。

\textbf{缺点}：\textbf{太厚了}——\textbf{你读不完}——\textbf{而且它讲的是 1970 年代的 Fortran 风格}——\textbf{现代工程师都用 SciPy}。

\subsection{李庆扬《数值分析》（清华版）}

\textbf{中国数值分析的标准教材}——\textbf{300 页}——\textbf{比 Burden 薄}——\textbf{但还是偏理论}。

\textbf{缺点}：\textbf{只讲方法、不讲"什么时候用哪个"}——\textbf{学完了你还是不知道工业里该选谁}。

\subsection{Numerical Recipes (Press et al.)}

\textbf{Numerical Recipes}——\textbf{科学家的"圣经"}——\textbf{每个方法都给 C/Fortran/Python 代码}。

\textbf{优点}：\textbf{能跑}。\textbf{缺点}：\textbf{代码风格老}——\textbf{现在大家直接用 SciPy / NumPy}。

\subsection{我的策略}

\textbf{本章不重复这三本书}——\textbf{它们都太全}。

\textbf{本章只做三件事}：

\begin{enumerate}
\item \textbf{把六类核心方法的直觉讲清楚}——\textbf{为什么 Simpson 比 trapezoidal 准}——\textbf{为什么 Newton 法收敛快}——\textbf{为什么 RK4 是金标准}
\item \textbf{把误差和收敛性讲透}——\textbf{$O(h^2)$ 和 $O(h^4)$ 到底差多少}——\textbf{什么是"机器精度"}
\item \textbf{把 SciPy 工具箱过一遍}——\textbf{遇到问题——先调 \texttt{scipy.optimize}、\texttt{scipy.integrate}、\texttt{scipy.linalg}}
\end{enumerate}

\begin{tcolorbox}[storybox={\Large 小故事：ENIAC 与现代数值计算的诞生}]
\textbf{1945 年 2 月}——\textbf{宾夕法尼亚大学 Moore 学院的地下室}——\textbf{ENIAC} (Electronic Numerical Integrator and Computer) \textbf{秘密通电}。

\medskip

\textbf{它有 30 吨重}——\textbf{占了一整间屋子}——\textbf{18000 个真空管、70000 个电阻、10000 个电容}——\textbf{耗电 150 千瓦}——\textbf{每次开机灯泡都会变暗}。

\textbf{它的速度}：\textbf{每秒 5000 次加法}——\textbf{比当时最快的机械计算器快 1000 倍}。

\medskip

\textbf{ENIAC 的第一个任务}——\textbf{不是民用}——\textbf{是为美国陆军算弹道表}。但\textbf{第一个真正跑在 ENIAC 上的"大计算"}——\textbf{是氢弹的可行性研究}——\textbf{由 von Neumann 主持}。

\medskip

\textbf{John von Neumann}——\textbf{20 世纪最聪明的人之一}——\textbf{一辈子做了三件大事}：

\begin{enumerate}
\item \textbf{创立博弈论}（1928）
\item \textbf{创立量子力学的数学基础}（1932）
\item \textbf{创立"存储程序"计算机架构}（1945，今天所有计算机的祖宗）
\end{enumerate}

\medskip

\textbf{在 Los Alamos}（曼哈顿计划）——von Neumann 遇到了\textbf{Stanislaw Ulam}（波兰数学家）。

\textbf{1946 年}——Ulam 生病住院——\textbf{无聊到玩纸牌接龙}——\textbf{他想知道：随机洗一副牌、接龙能成功的概率是多少？}

\textbf{这个问题——用纯数学算——是地狱}（要枚举太多情形）。

\textbf{Ulam 想}：\textbf{为什么不让 ENIAC 模拟 10000 次}——\textbf{看成功多少次}——\textbf{除一下不就是概率吗？}

\medskip

\textbf{这就是 Monte Carlo 方法的诞生}——\textbf{名字来自 Ulam 叔叔常去的赌城}。

\medskip

\textbf{从此}——\textbf{数值方法不再只是"近似积分"}——\textbf{它成了"用随机模拟解决确定性问题"的全新范式}。

今天——\textbf{Monte Carlo}——是\textbf{金融定价、粒子物理、机器学习、流行病建模}——\textbf{所有"算不出来"的问题}的\textbf{终极武器}。

\medskip

\textbf{ENIAC 之后 80 年}——\textbf{你手机里的芯片}——\textbf{比当年 ENIAC 快 10 亿倍}——\textbf{但跑的还是同一套数学}：\textbf{加、减、乘、除——再加一点循环和判断}。

\textbf{这就是数值方法}——\textbf{把数学装进机器的艺术}。
\end{tcolorbox}

\section{直觉 / Intuition}

\subsection{数值微分：用差分代替极限}

\textbf{导数的定义}：$f'(x) = \lim_{h \to 0} \frac{f(x+h) - f(x)}{h}$。

\textbf{机器不会"取极限"}——\textbf{它只会取一个"足够小"的 $h$}。

\textbf{前向差分}（forward difference）：
\begin{equation}
f'(x) \approx \frac{f(x+h) - f(x)}{h} \qquad O(h)
\end{equation}

\textbf{中心差分}（central difference）：
\begin{equation}
f'(x) \approx \frac{f(x+h) - f(x-h)}{2h} \qquad O(h^2)
\end{equation}

\textbf{中心差分比前向差分准一阶}——\textbf{为什么}？泰勒展开两边一抵消——\textbf{$O(h)$ 项干掉了}——\textbf{只剩 $O(h^2)$}。

\textbf{别选太小的 $h$！}——$h$ 太小——\textbf{浮点减法的舍入误差}会主导——\textbf{结果反而变差}。\textbf{经验值}：$h \approx \sqrt{\epsilon_{\text{mach}}} \approx 10^{-8}$。

\subsection{数值积分：用矩形/梯形/抛物线代替面积}

\textbf{定积分 $\int_a^b f(x) \, dx$ 的几何意义是面积}——\textbf{机器算面积——只能用直边/曲边的简单图形拼}。

\textbf{三种基本方法}：

\begin{itemize}
\item \textbf{矩形法}（midpoint）：$\sum f(x_i) \cdot h$——\textbf{精度 $O(h^2)$}
\item \textbf{梯形法}（trapezoidal）：$\frac{h}{2} [f_0 + 2f_1 + \cdots + 2f_{n-1} + f_n]$——\textbf{精度 $O(h^2)$}
\item \textbf{Simpson 法}（用抛物线拼）：$\frac{h}{3} [f_0 + 4f_1 + 2f_2 + 4f_3 + \cdots + f_n]$——\textbf{精度 $O(h^4)$}
\end{itemize}

\textbf{Simpson 比梯形高两阶}——\textbf{为什么}？\textbf{抛物线比直线多算了"曲率"}——\textbf{在光滑函数上几乎"白送"了一个数量级的精度}。

\textbf{自适应积分}（adaptive quadrature）：\textbf{函数变化剧烈的地方多取点、平缓的地方少取点}——\textbf{SciPy 的 \texttt{quad} 就是干这个的}。

\subsection{非线性方程：把"求根"变成"迭代逼近"}

\textbf{要解 $f(x) = 0$}——\textbf{大多数方程没有解析解}（比如 $x = \cos x$）。

\textbf{两大经典方法}：

\textbf{二分法}（bisection）——\textbf{老实人方法}：
\begin{enumerate}
\item 找到 $a, b$ 使 $f(a) \cdot f(b) < 0$（根夹在中间）
\item 取中点 $c = (a+b)/2$，看 $f(c)$ 的符号——\textbf{把区间砍一半}
\item 重复——\textbf{每次精度翻倍}（\textbf{线性收敛，$O(2^{-n})$}）
\end{enumerate}

\textbf{二分法的优点}：\textbf{一定收敛}。\textbf{缺点}：\textbf{慢}。

\textbf{Newton 法}（牛顿迭代）——\textbf{聪明人方法}：
\begin{equation}
x_{n+1} = x_n - \frac{f(x_n)}{f'(x_n)}
\end{equation}

\textbf{几何意义}：\textbf{从 $x_n$ 作切线、看切线与 $x$ 轴的交点}——\textbf{这就是 $x_{n+1}$}。

\textbf{Newton 法的威力}：\textbf{二次收敛}（quadratic convergence）——\textbf{每次迭代——有效数字翻一倍}！

\textbf{典型}：\textbf{算 $\sqrt{2}$}——3 次迭代——\textbf{15 位有效数字}。

\textbf{Newton 法的代价}：\textbf{要求 $f'(x)$}——\textbf{初值要选好——否则可能发散}。

\subsection{线性方程组：直接法 vs 迭代法}

\textbf{解 $A\bm{x} = \bm{b}$}——\textbf{$n \times n$ 的方程组}。

\textbf{直接法}（direct）——\textbf{Gauss 消元、LU 分解}——\textbf{$O(n^3)$ 次运算}——\textbf{适合小矩阵}（$n < 10000$）。

\textbf{迭代法}（iterative）——\textbf{Jacobi、Gauss-Seidel、共轭梯度法}（CG）——\textbf{适合大稀疏矩阵}（$n > 10^6$，但绝大多数元素是 0）。

\textbf{现代工程师的经验}：

\begin{itemize}
\item \textbf{$n < 10000$ 且稠密}：\texttt{np.linalg.solve}（背后是 LAPACK 的 LU）
\item \textbf{$n > 10^5$ 且稀疏}：\texttt{scipy.sparse.linalg.cg / gmres}
\item \textbf{$n > 10^7$（大规模）}：\textbf{多重网格法}（multigrid）——\textbf{接近 $O(n)$}
\end{itemize}

\subsection{ODE 求解器：从 Euler 到 RK4}

\textbf{要解 $\frac{dy}{dt} = f(t, y), \quad y(0) = y_0$}。

\textbf{Euler 法}（最朴素）：$y_{n+1} = y_n + h \cdot f(t_n, y_n)$——\textbf{精度 $O(h)$}——\textbf{太粗}。

\textbf{Runge-Kutta 4 阶}（RK4）——\textbf{ODE 数值解的"金标准"}：
\begin{align}
k_1 &= f(t_n, y_n) \\
k_2 &= f(t_n + h/2, y_n + h k_1/2) \\
k_3 &= f(t_n + h/2, y_n + h k_2/2) \\
k_4 &= f(t_n + h, y_n + h k_3) \\
y_{n+1} &= y_n + \frac{h}{6} (k_1 + 2k_2 + 2k_3 + k_4)
\end{align}

\textbf{精度 $O(h^4)$}——\textbf{在工程问题上 95\% 够用}。

\textbf{自适应步长}（RK45, Dormand-Prince）——\textbf{SciPy 的 \texttt{solve\_ivp} 默认方法}——\textbf{它在变化剧烈处自动取小步长}。

\subsection{误差：从哪里来、到哪里去}

\textbf{数值计算的三类误差}：

\begin{itemize}
\item \textbf{截断误差}（truncation error）——\textbf{方法本身的近似误差}——\textbf{比如 Simpson 法的 $O(h^4)$}
\item \textbf{舍入误差}（round-off error）——\textbf{浮点数的有限精度}——\textbf{双精度约 $10^{-16}$}
\item \textbf{传播误差}——\textbf{前面的误差被后续运算放大}——\textbf{$10^{16}$ 次乘法可能放大很多}
\end{itemize}

\textbf{机器精度}：\textbf{IEEE 754 双精度浮点}——\textbf{$\epsilon_{\text{mach}} \approx 2.22 \times 10^{-16}$}——\textbf{你永远算不到比它更准}。

\section{Aha 一：Newton 法手算 $\sqrt{2}$}

\textbf{问题}：\textbf{求 $\sqrt{2}$}——\textbf{不用计算器的 sqrt 键}——\textbf{用 Newton 法}。

\textbf{思路}：\textbf{$\sqrt{2}$ 是 $f(x) = x^2 - 2 = 0$ 的正根}。

\textbf{Newton 迭代}：
\begin{equation}
x_{n+1} = x_n - \frac{x_n^2 - 2}{2 x_n} = \frac{1}{2} \left( x_n + \frac{2}{x_n} \right)
\end{equation}

\textbf{这就是著名的"巴比伦平方根算法"}——\textbf{4000 年前的巴比伦人就在用}——\textbf{他们不知道这是 Newton 法}。

\textbf{从 $x_0 = 1.5$ 开始迭代}：

\begin{center}
\begin{tabular}{ccc}
\toprule
$n$ & $x_n$ & 有效数字 \\
\midrule
0 & 1.5 & 1 位 \\
1 & 1.41666666666\ldots & 3 位 \\
2 & 1.41421568627\ldots & 6 位 \\
3 & 1.41421356237\ldots & \textbf{12 位} \\
4 & 1.41421356237\ldots & \textbf{16 位（机器精度）} \\
\bottomrule
\end{tabular}
\end{center}

\textbf{4 次迭代——干到机器精度}——\textbf{这就是二次收敛的威力}！

\textbf{对比二分法}：\textbf{每次精度翻一倍}——\textbf{要到 16 位精度——需要 50 多次迭代}。

\begin{lstlisting}[language=Python]
# numerical.py —— Newton's method for sqrt
def newton_sqrt(a, x0=1.0, tol=1e-15, max_iter=100):
    """Compute sqrt(a) by Newton's method."""
    x = x0
    for i in range(max_iter):
        x_new = 0.5 * (x + a / x)
        if abs(x_new - x) < tol:
            return x_new, i + 1
        x = x_new
    return x, max_iter

# usage
val, iters = newton_sqrt(2.0)
print(f"sqrt(2) = {val:.16f}  ({iters} iterations)")
# sqrt(2) = 1.4142135623730951  (4 iterations)
\end{lstlisting}

\section{Aha 二：GPS 定位算法——你口袋里的数值方法}

\textbf{你每天用手机导航——背后跑的是什么？}

\textbf{答案}：\textbf{实时非线性方程组求解 + 最小二乘 + Kalman 滤波}——\textbf{这一章的所有方法——你的手机芯片每秒跑 10 次}。

\subsection{GPS 的几何原理}

\textbf{GPS 卫星}：\textbf{24 颗}——\textbf{在 20200 km 高的轨道上}——\textbf{每颗都在不停广播}：\textbf{"我是 \#7 号卫星——现在是格林尼治时间 12:34:56.789——我的位置是 $(x_7, y_7, z_7)$"}。

\textbf{你的手机收到信号}——\textbf{比较"信号发出时间"和"信号到达时间"}——\textbf{算出"信号飞了多远"}：
\begin{equation}
d_i = c \cdot (t_{\text{recv}} - t_{\text{sent}})
\end{equation}
\textbf{$c$ 是光速}。

\textbf{你的位置 $(x, y, z)$ 满足}：
\begin{equation}
\sqrt{(x - x_i)^2 + (y - y_i)^2 + (z - z_i)^2} = d_i, \qquad i = 1, 2, 3, 4
\end{equation}

\textbf{3 个未知数（$x, y, z$）——为什么要 4 颗卫星？}

\textbf{因为你的手机时钟不准}——\textbf{差几个微秒——位置就差 1 公里}——\textbf{第 4 颗卫星——是用来解出"时钟偏差" $\Delta t$ 的}。

\textbf{所以 GPS 的真正方程是}：
\begin{equation}
\sqrt{(x - x_i)^2 + (y - y_i)^2 + (z - z_i)^2} = c \cdot (t_{\text{recv}} - t_{\text{sent},i} - \Delta t)
\end{equation}
\textbf{4 个未知数 $(x, y, z, \Delta t)$——4 个方程——非线性方程组}。

\subsection{怎么解？}

\textbf{这就是 Newton 法的多变量版本}：
\begin{equation}
\bm{x}_{n+1} = \bm{x}_n - [\bm{J}(\bm{x}_n)]^{-1} \bm{F}(\bm{x}_n)
\end{equation}
\textbf{$\bm{J}$ 是 Jacobian 矩阵}——\textbf{每次迭代——解一个 4x4 的线性方程组（直接法）}。

\textbf{在实际芯片里}：\textbf{用 5\textasciitilde{}10 颗卫星——超定方程组——用最小二乘 + Kalman 滤波}——\textbf{每秒更新 1\textasciitilde{}10 次位置}。

\textbf{这就是你的手机蓝点——每秒移动的背后}。

\begin{lstlisting}[language=Python]
# numerical.py —— simplified GPS solver
import numpy as np
from scipy.optimize import least_squares

def gps_solve(sat_positions, pseudo_ranges, x0=None):
    """
    sat_positions: (n, 3) array of satellite (x, y, z)
    pseudo_ranges: (n,) array of measured distances
    Returns: (x, y, z, dt_bias)
    """
    c = 299792458.0  # speed of light, m/s

    def residuals(params):
        x, y, z, dt = params
        true_d = np.linalg.norm(sat_positions - np.array([x, y, z]), axis=1)
        # pseudo-range = true distance + clock bias * c
        return true_d + c * dt - pseudo_ranges

    if x0 is None:
        x0 = np.array([0, 0, 0, 0])  # center of Earth + zero bias
    sol = least_squares(residuals, x0, method='lm')  # Levenberg-Marquardt
    return sol.x

# usage (synthetic 4-satellite example)
sats = np.array([
    [15600e3,  7540e3, 20140e3],
    [18760e3,  2750e3, 18610e3],
    [17610e3, 14630e3, 13480e3],
    [19170e3,  610e3,  18390e3],
])
ranges = np.array([20451e3, 21421e3, 23864e3, 21952e3])
x, y, z, dt = gps_solve(sats, ranges)
print(f"Position: ({x:.0f}, {y:.0f}, {z:.0f})  bias={dt*1e6:.3f} us")
\end{lstlisting}

\textbf{这就是 21 世纪的 Newton 法}——\textbf{从 17 世纪的羽毛笔——到 21 世纪的硅芯片}——\textbf{算法没变——速度快了 $10^{15}$ 倍}。

\section{深入 / Deep Dive}

\subsection{数值积分实战：Simpson vs trapezoidal}

\textbf{算 $\int_0^1 e^{-x^2} dx$}（高斯钟形曲线的一部分——\textbf{无解析解}）。

\begin{lstlisting}[language=Python]
# numerical.py —— quadrature comparison
import numpy as np
from scipy.integrate import quad

def trapezoidal(f, a, b, n):
    x = np.linspace(a, b, n + 1)
    y = f(x)
    h = (b - a) / n
    return h * (0.5 * y[0] + y[1:-1].sum() + 0.5 * y[-1])

def simpson(f, a, b, n):
    """n must be even."""
    if n % 2:
        raise ValueError("n must be even")
    x = np.linspace(a, b, n + 1)
    y = f(x)
    h = (b - a) / n
    return h / 3 * (y[0] + 4 * y[1:-1:2].sum() + 2 * y[2:-1:2].sum() + y[-1])

f = lambda x: np.exp(-x**2)
exact, _ = quad(f, 0, 1)  # SciPy reference

for n in [10, 100, 1000]:
    t = trapezoidal(f, 0, 1, n)
    s = simpson(f, 0, 1, n)
    print(f"n={n:5d}  trap err={abs(t-exact):.2e}  simp err={abs(s-exact):.2e}")
# n=   10  trap err=6.94e-04  simp err=4.51e-08
# n=  100  trap err=6.94e-06  simp err=4.51e-12
# n= 1000  trap err=6.94e-08  simp err=4.51e-16  <- machine precision!
\end{lstlisting}

\textbf{看清楚了吗}？

\begin{itemize}
\item \textbf{trapezoidal}：$n$ 增加 10 倍——\textbf{误差减少 100 倍}（$O(h^2)$）
\item \textbf{Simpson}：$n$ 增加 10 倍——\textbf{误差减少 10000 倍}（$O(h^4)$）
\item \textbf{Simpson 在 $n = 1000$ 时——已经撞上机器精度天花板}
\end{itemize}

\textbf{这就是高阶方法的价值}。

\subsection{ODE 实战：化学反应动力学}

\textbf{Lotka-Volterra 捕食者-猎物方程}（\textbf{生态学经典模型}）：
\begin{equation}
\frac{dx}{dt} = \alpha x - \beta x y, \qquad \frac{dy}{dt} = \delta x y - \gamma y
\end{equation}

\begin{lstlisting}[language=Python]
# numerical.py —— ODE solver demo
from scipy.integrate import solve_ivp
import numpy as np

def lotka_volterra(t, z, alpha, beta, delta, gamma):
    x, y = z
    return [alpha * x - beta * x * y,
            delta * x * y - gamma * y]

# parameters: rabbits (x) and foxes (y)
params = (1.0, 0.1, 0.075, 1.5)
z0 = [10, 5]  # 10 rabbits, 5 foxes
t_span = (0, 50)
t_eval = np.linspace(*t_span, 1000)

sol = solve_ivp(lotka_volterra, t_span, z0,
                args=params, t_eval=t_eval,
                method='RK45', rtol=1e-8, atol=1e-10)

print(f"Final: rabbits={sol.y[0, -1]:.2f}, foxes={sol.y[1, -1]:.2f}")
# population oscillates -- classic predator-prey cycle
\end{lstlisting}

\textbf{这就是 SciPy 的 \texttt{solve\_ivp}}——\textbf{背后是自适应 RK45（Dormand-Prince）}——\textbf{工业级、上线就能用}。

\subsection{线性方程组实战：稀疏矩阵}

\textbf{大型有限元仿真——矩阵 $n = 10^6$——但每行只有 10 个非零元}——\textbf{稠密存储要 8 TB——稀疏存储只要 80 MB}。

\begin{lstlisting}[language=Python]
# numerical.py —— sparse linear solve
import numpy as np
from scipy.sparse import diags
from scipy.sparse.linalg import cg  # conjugate gradient

n = 100000
# tridiagonal: 2 on diagonal, -1 off-diagonal (1D Laplacian)
A = diags([-1, 2, -1], [-1, 0, 1], shape=(n, n), format='csr')
b = np.ones(n)

x, info = cg(A, b, rtol=1e-10)
print(f"Converged: {info == 0},  ||Ax - b|| = {np.linalg.norm(A @ x - b):.2e}")
\end{lstlisting}

\textbf{共轭梯度法}（CG）在\textbf{对称正定矩阵}上——\textbf{$O(\sqrt{\kappa})$ 次迭代收敛}（$\kappa$ 是条件数）——\textbf{在工程上经常 100 次迭代解掉 $10^6$ 维方程}。

\subsection{误差分析：为什么不能"无脑减"}

\textbf{灾难性抵消}（catastrophic cancellation）——\textbf{两个相近的数相减——大量有效数字消失}。

\textbf{典型}：\textbf{算 $f(x) = \frac{1 - \cos x}{x^2}$ 在 $x = 10^{-8}$}：

\begin{lstlisting}[language=Python]
import math
x = 1e-8
naive = (1 - math.cos(x)) / x**2     # 0.0  <- WRONG!
clever = 2 * math.sin(x/2)**2 / x**2  # 0.4999... <- correct
# true value: 1/2
\end{lstlisting}

\textbf{原因}：\textbf{$\cos(10^{-8}) \approx 0.99999999\ldots$}——\textbf{$1 - \cos x$ 把前 16 位有效数字全部抵消掉——剩下噪声}。

\textbf{用 $1 - \cos x = 2 \sin^2(x/2)$ 重写——避免减法——结果就对了}。

\textbf{教训}：\textbf{数值方法不只是"调库"}——\textbf{有时候需要数学上的等价变换才能避免精度灾难}。

\section{常见错误 / Common Pitfalls}

\textbf{(1) 步长 $h$ 选得太小}——\textbf{舍入误差主导}——\textbf{数值微分尤其要小心}。\textbf{经验值}：$h \approx 10^{-8}$（对双精度）。

\textbf{(2) Newton 法初值乱选}——\textbf{可能发散、可能跳到错误的根}——\textbf{无法保证收敛}。\textbf{保险做法}：\textbf{先用二分法粗调、再用 Newton 精调}。

\textbf{(3) 自己实现高斯消元}——\textbf{99\% 的情况下比 LAPACK 慢 100 倍、还不数值稳定}——\textbf{直接调 \texttt{np.linalg.solve}}。

\textbf{(4) 用 Euler 法解 ODE}——\textbf{精度太低、稳定性差}——\textbf{除非你在做教学演示——否则永远用 RK4 或自适应方法}。

\textbf{(5) 忽略条件数}——\textbf{病态矩阵}（$\kappa \gg 1$）\textbf{即使用 LAPACK 也会丢精度}——\textbf{用 \texttt{np.linalg.cond} 先体检}。

\textbf{(6) 没设容忍度}——\textbf{\texttt{solve\_ivp} 默认 \texttt{rtol=1e-3}}——\textbf{科学计算要设到 $10^{-8}$ 或更小}。

\textbf{(7) 把"能跑"当成"对"}——\textbf{数值方法不报错——不代表结果对}——\textbf{一定要验证}（\textbf{用 mesh refinement / 解析解对比 / 物理守恒律}）。

\section{应用 / Applications}

\textbf{(1) 天气预报}——\textbf{欧洲中期天气预报中心}（ECMWF）——\textbf{每天求解一个 $10^9$ 维 PDE}——\textbf{核心是有限元 + 半隐式时间积分}。

\textbf{(2) 期权定价}——\textbf{Black-Scholes 方程的数值解}——\textbf{有限差分 + Monte Carlo}——\textbf{华尔街的核心算法}。

\textbf{(3) 飞机机翼设计}——\textbf{Navier-Stokes 方程的 CFD 求解}——\textbf{有限体积法 + 自适应网格}——\textbf{波音 787 设计周期从 10 年缩短到 4 年}。

\textbf{(4) 核反应堆模拟}——\textbf{中子输运方程}——\textbf{蒙特卡洛 + 离散纵标法}——\textbf{Los Alamos 70 年的拿手好戏}。

\textbf{(5) 你的手机 GPS}——\textbf{每秒解一次非线性方程组 + Kalman 滤波}——\textbf{见前文}。

\textbf{(6) 训练 GPT}——\textbf{反向传播本质上是链式法则 + 数值优化}——\textbf{Adam 优化器 = 自适应步长的梯度下降}——\textbf{每一步都在求解一个高维非凸优化}。

\section{SciPy 工具箱速查表 / SciPy Cheatsheet}

\textbf{遇到问题——先查这个表——不要自己手写}：

\begin{center}
\begin{tabular}{lll}
\toprule
\textbf{问题} & \textbf{SciPy 函数} & \textbf{说明} \\
\midrule
数值积分（1D） & \texttt{scipy.integrate.quad} & 自适应高斯-Kronrod \\
数值积分（多重） & \texttt{scipy.integrate.dblquad} & 二重 / 三重 \\
ODE 初值问题 & \texttt{scipy.integrate.solve\_ivp} & RK45 / Radau / BDF \\
非线性方程 & \texttt{scipy.optimize.brentq} & 二分+割线，超稳 \\
非线性方程组 & \texttt{scipy.optimize.fsolve} & Powell 混合算法 \\
最小二乘 & \texttt{scipy.optimize.least\_squares} & Levenberg-Marquardt \\
线性方程组（稠密） & \texttt{np.linalg.solve} & LAPACK LU \\
线性方程组（稀疏） & \texttt{scipy.sparse.linalg.spsolve} & 稀疏 LU \\
迭代法（对称正定） & \texttt{scipy.sparse.linalg.cg} & 共轭梯度 \\
迭代法（一般） & \texttt{scipy.sparse.linalg.gmres} & GMRES \\
特征值 & \texttt{np.linalg.eig} & QR 算法 \\
插值 & \texttt{scipy.interpolate.interp1d} & 多种样条 \\
FFT & \texttt{np.fft.fft} & 快速傅里叶变换 \\
\bottomrule
\end{tabular}
\end{center}

\textbf{记住这个表——你 90\% 的科学计算都能搞定}。

\textbf{剩下的 10\%}（极端规模、特殊硬件、新算法）——\textbf{才需要你自己写代码或者用 PETSc / Trilinos 等专业库}。

\section{本章总结 / Summary}

\begin{tcolorbox}[essbox={本章关键收获}]
\textbf{一句话}：\textbf{数值方法——把"算不出来"的问题——变成"能算的"——靠的是近似、迭代、收敛三件套}。

\medskip

\textbf{六类核心方法}：
\begin{enumerate}
\item \textbf{数值微分}：中心差分（$O(h^2)$）——\textbf{别选太小的 $h$}
\item \textbf{数值积分}：Simpson 法（$O(h^4)$）+ adaptive quadrature
\item \textbf{非线性方程}：Newton 法（二次收敛）+ bisection（稳但慢）
\item \textbf{线性方程组}：LU（直接、稠密）/ CG (迭代、稀疏)
\item \textbf{ODE}：RK4 / RK45 自适应（金标准）
\item \textbf{误差}：截断 + 舍入 + 传播——\textbf{机器精度 $10^{-16}$ 是天花板}
\end{enumerate}

\medskip

\textbf{两个 Aha}：
\begin{itemize}
\item \textbf{Newton 法求 $\sqrt{2}$}——\textbf{4 次迭代干到机器精度}——\textbf{巴比伦人 4000 年前就在用}
\item \textbf{GPS 定位}——\textbf{多变量 Newton 法 + 最小二乘}——\textbf{你口袋里的科学计算}
\end{itemize}

\medskip

\textbf{工程师心法}：\textbf{遇到问题——先查 SciPy 工具箱——不要重复发明轮子}——\textbf{把精力花在"建模"和"验证"上}。
\end{tcolorbox}

\textbf{下一章}——\textbf{我们进入}\textbf{优化理论}——\textbf{从最速下降到 Adam}——\textbf{从凸优化到深度学习}——\textbf{把"求最小值"这件事——彻底说清楚}。

\clearpage

% =====================================================================
%  第 13 章 · 优化理论 / Optimization —— 现代工程之核心
%  一元、多元、凸、约束、线性规划、SGD、Adam
%  小故事：Karmarkar 算法 1984 —— 内点法革命
%  Aha：ChatGPT 训练 = Adam 优化器 + 千亿参数 —— 全书的终章
% =====================================================================

\chapter{优化理论 / Optimization —— 现代工程之核心}
\label{ch:opt}

\begin{quote}\itshape
\textbf{1947 年——George Dantzig 发明单纯形法——用于美国空军后勤调度}。\\
\textbf{1984 年——Narendra Karmarkar 在 Bell Labs 提出内点法——把线性规划速度提高了 50 倍}。\\
\textbf{1951 年——Robbins \& Monro 提出随机近似算法——半个世纪后变成 SGD}。\\
\textbf{2014 年——Kingma \& Ba 在 ICLR 提出 Adam——10 年后训练出了 ChatGPT}。\\[6pt]
\textbf{这一章——\textbf{全书的终章}——讲一件事}：\\
\textbf{现代工程——本质上——是一场\textbf{找最优}的游戏}。\\
\textbf{从飞机外形到神经网络权重——从航线调度到芯片布线}——\\
\textbf{背后跑的——都是优化算法}。
\end{quote}

\section{写在前面 / Preface}

\textbf{我大学时学"优化"}——\textbf{是在《运筹学》课上}——\textbf{老师讲了三周单纯形法的手工算表}——\textbf{讲了拉格朗日乘子法的证明}——\textbf{讲了 KKT 条件的四个公式}。

\textbf{考试我考了 88 分}。\textbf{然后我把它忘了}。

\textbf{硕士第二年}——\textbf{我做 LQR 控制器设计}——\textbf{我要解一个二次型优化问题}——\textbf{我打开 Python}：

\begin{lstlisting}[language=Python]
from scipy.optimize import minimize
result = minimize(loss, x0, method='BFGS')
\end{lstlisting}

\textbf{一行代码——做完了优化}。

\textbf{博士第一年}——\textbf{我训练第一个神经网络}——\textbf{打开 PyTorch}：

\begin{lstlisting}[language=Python]
optimizer = torch.optim.Adam(model.parameters(), lr=1e-3)
for x, y in dataloader:
    loss = criterion(model(x), y)
    loss.backward()
    optimizer.step()
\end{lstlisting}

\textbf{六行代码——训练了一个有 50 万参数的网络}。

\textbf{那一刻我才反应过来}——\textbf{我本科学的"单纯形手算"}——\textbf{这辈子工业里再没用过一次}。\textbf{真正用的——是\textbf{这几个名字}}：

\begin{itemize}
\item \textbf{Gradient Descent}（梯度下降）——\textbf{所有优化算法的鼻祖}
\item \textbf{Newton / BFGS}（拟牛顿法）——\textbf{中小规模问题的王者}
\item \textbf{SGD / Adam}（随机梯度 / 自适应矩）——\textbf{深度学习的引擎}
\item \textbf{KKT 条件}——\textbf{所有约束优化的"判别式"}
\end{itemize}

\textbf{这一章}——\textbf{把"优化"这件事——从一元最小值——一直讲到 ChatGPT 训练}。\textbf{这是整本书的\textbf{climax}}——\textbf{所有前面的章节——梯度（第 2 章）、Hessian（第 7 章）、随机过程（第 10 章）、数值（第 12 章）}——\textbf{都在这一章里\textbf{汇合}}。

\section{一句话本质 / The Essence in One Sentence}

\begin{tcolorbox}[essbox={一句话本质}]
\textbf{中文}：\textbf{优化 = 找一个让目标函数最小的点}。\\
\textbf{无约束时——梯度为零；有约束时——KKT 条件}。\\
\textbf{大数据时——逐样本估计梯度（SGD）；非凸时——加动量、加自适应学习率（Adam）}。\\[6pt]
\textbf{English}: Optimization is finding the point that minimizes an objective. Unconstrained: gradient equals zero. Constrained: KKT conditions. Large-scale data: estimate the gradient per sample (SGD). Non-convex landscape: add momentum and adaptive learning rates (Adam).
\end{tcolorbox}

记住这一句话——\textbf{这一章已经讲完了一半}。

\section{教材里你看到的 / What the Textbook Tells You}

\subsection{运筹学 / 最优化方法（国内教材）}

\textbf{钱颂迪《运筹学》、《最优化方法》（高等教育出版社）}——\textbf{中国工科最常用的优化教材}。

\textbf{优点}：单纯形、对偶、KKT、动态规划——\textbf{该有的都有}。

\textbf{缺点}：

\begin{enumerate}
\item \textbf{算法部分是 60 年代的}——\textbf{讲单纯形的"枢轴选择"细节}——\textbf{这件事 LP solver 已经做了 30 年了}
\item \textbf{完全没有 SGD / Adam}——\textbf{现代 ML 优化器一个都没讲}
\item \textbf{凸分析点到为止}——\textbf{Slater 条件、对偶间隙——这些工业里真正用的反而少讲}
\end{enumerate}

\subsection{Boyd \& Vandenberghe《Convex Optimization》}

\textbf{Stanford 的 Stephen Boyd}——\textbf{全球凸优化的圣经}——\textbf{免费 PDF + 公开课}（EE364）——\textbf{我自学的就是这本}。

\textbf{优点}：\textbf{凸性 + 对偶 + 内点法}——\textbf{讲到工业可用的深度}。\textbf{第 6 章"逼近与拟合"}——\textbf{是机器学习损失函数设计的真正源头}。

\textbf{缺点}：\textbf{不讲 SGD}（成书时 deep learning 尚未爆发）。

\subsection{Nocedal \& Wright《Numerical Optimization》}

\textbf{无约束优化的圣经}——\textbf{讲 BFGS、L-BFGS、信赖域}——\textbf{讲到能写 solver 的深度}。

\textbf{适合谁}：\textbf{要做数值优化研究 / 写 solver 的人}。

\subsection{Goodfellow et al.《Deep Learning》第 8 章}

\textbf{"深度模型中的优化"}——\textbf{讲 SGD、Momentum、AdaGrad、RMSProp、Adam}——\textbf{这是工业 ML 工程师真正翻的章节}。

\subsection{我的策略}

\textbf{本章做四件事}：

\begin{enumerate}
\item \textbf{把无约束优化讲成一条进化线}：\textbf{Gradient $\to$ Newton $\to$ BFGS $\to$ SGD $\to$ Adam}——\textbf{每一步为什么必须出现}
\item \textbf{把凸性讲成"为什么 ML 工程师爱它"}——\textbf{凸 = 局部最优即全局最优 = 算法收敛有保证}
\item \textbf{把约束优化收成一个公式}：\textbf{KKT}——\textbf{然后告诉你为什么 SVM 是这个公式的最美应用}
\item \textbf{把 Aha 讲到震撼}：\textbf{ChatGPT 训练 = Adam + 千亿参数 + 万张 GPU——三个月}——\textbf{这是优化理论 70 年的总和成就}
\end{enumerate}

\begin{tcolorbox}[storybox={\Large 小故事：Karmarkar 算法（1984）与 AT\&T 内点法风波}]
\textbf{1984 年 4 月——Bell Labs 的 28 岁印度裔数学家 Narendra Karmarkar——发表了一篇 10 页的论文}：\textbf{《A new polynomial-time algorithm for linear programming》}。

\medskip

\textbf{这篇论文做了一件让所有人震惊的事}——\textbf{它说}：

\begin{quote}
\textbf{"我有一个算法——它不沿着多面体的棱走（单纯形法的做法）——它\textbf{从内部直接穿过去}——它的复杂度是 $O(n^{3.5}L)$ ——比单纯形法的最坏情况快得多——而且实测在大规模线性规划上——\textbf{快 50 倍}。"}
\end{quote}

\textbf{这就是\textbf{内点法（interior-point method）}的诞生}。

\medskip

\textbf{$\bullet$ 单纯形法（1947, Dantzig）}——沿着可行域的顶点跳——\textbf{每一步都在边界上}——平均很快——\textbf{但最坏情况是指数级}（Klee-Minty 1972）

\textbf{$\bullet$ 椭球法（1979, Khachiyan）}——\textbf{首个多项式时间算法}——但常数太大——\textbf{实战慢得没法用}——\textbf{学术意义大于工业意义}

\textbf{$\bullet$ Karmarkar 内点法（1984）}——\textbf{多项式时间 + 工业实战快}——\textbf{第一个"理论与实践双赢"的 LP 算法}

\medskip

\textbf{然后——风波来了}。

\medskip

\textbf{AT\&T（Bell Labs 的母公司）做了一件让学术界炸锅的事}——\textbf{它把 Karmarkar 算法\textbf{申请了专利}}（US Patent 4,744,028, 1988）——\textbf{并把它包装成商业产品 KORBX——卖给空军 890 万美元}。

\medskip

\textbf{学术界愤怒了}——\textbf{"数学算法怎么能申请专利？"}——\textbf{这是有史以来\textbf{第一个数学算法专利}}——\textbf{开了一个让所有人都不舒服的先例}。

\medskip

\textbf{反向效果是}——\textbf{学术界为了"绕开 Karmarkar 专利"——疯狂研究新的内点法变种}——\textbf{结果催生了\textbf{一整个内点法学派}}——\textbf{Mehrotra predictor-corrector（1992）、primal-dual interior point（1993）}——\textbf{这些反而成了今天 CPLEX、Gurobi、Mosek 等商业 solver 的核心}。

\medskip

\textbf{Karmarkar 专利于 2006 年到期}。\textbf{今天——所有大规模 LP / QP / SOCP / SDP——背后都跑着内点法}。

\medskip

\textbf{这件事给我们的启示}：\textbf{一个 28 岁印度青年的 10 页论文——把一个 40 年的工业问题——重写了一遍}。\textbf{数学的力量——往往就来自这种"换一个看问题的角度"}。
\end{tcolorbox}

\section{§1 一元优化 / Univariate Optimization}

\textbf{先从最简单的开始}：\textbf{找一个一元函数 $f(x)$ 的最小值}。

\subsection{一阶必要条件}

\textbf{中学就学过}：\textbf{$f$ 在 $x^*$ 取得极小} $\Rightarrow$ $f'(x^*) = 0$。

\textbf{但反过来不成立}——\textbf{$f'(x^*) = 0$ 不一定是最小}（可能是极大、可能是鞍点）。\textbf{加二阶条件}：$f''(x^*) > 0$ $\Rightarrow$ 严格极小。

\textbf{这是一切优化的起点}——\textbf{多元的"梯度为零 + Hessian 正定"——只是它的高维版本}。

\subsection{黄金分割搜索 / Golden Section Search}

\textbf{当 $f$ 不可导}（或导数难算）——\textbf{怎么找最小}？

\textbf{思路}：\textbf{如果 $f$ 在 $[a, b]$ 上是\textbf{单峰}的}——\textbf{就用"缩区间"的策略}。

\textbf{算法}（设黄金比 $\phi = \frac{\sqrt 5 - 1}{2} \approx 0.618$）：

\begin{enumerate}
\item 取 $x_1 = a + (1-\phi)(b-a)$、$x_2 = a + \phi(b-a)$
\item 若 $f(x_1) < f(x_2)$ $\Rightarrow$ 最小值在 $[a, x_2]$——更新 $b = x_2$
\item 否则——最小值在 $[x_1, b]$——更新 $a = x_1$
\item 重复——\textbf{每次区间缩小约 38.2\%}
\end{enumerate}

\textbf{为什么是 0.618？}——\textbf{因为黄金比有"自相似性"}——\textbf{下一轮可以\textbf{复用一个点}的函数值}——\textbf{每轮只需算 1 次 $f$（而不是 2 次）}。

\textbf{收敛速度}：\textbf{线性}——\textbf{每 $n$ 步把区间缩小 $\phi^n$ 倍}。

\subsection{Newton 法（一元）}

\textbf{当 $f$ 二阶可导}——\textbf{Newton 法是金标准}。

\textbf{思路}：\textbf{在 $x_k$ 处用二阶 Taylor 展开近似 $f$}：

\begin{equation}
f(x) \approx f(x_k) + f'(x_k)(x - x_k) + \tfrac{1}{2} f''(x_k)(x - x_k)^2
\end{equation}

\textbf{这个二次近似的最小值}：\textbf{对 $x$ 求导令为零}——\textbf{得到}：

\begin{equation}
\boxed{\, x_{k+1} = x_k - \frac{f'(x_k)}{f''(x_k)} \,}
\end{equation}

\textbf{这就是一元 Newton 法的迭代公式}。

\textbf{收敛速度}：\textbf{二次}——\textbf{每一步误差平方缩小}——\textbf{比黄金分割快无穷倍}。

\begin{example}[Newton 找 $f(x) = x^2 - 4x + 3$ 的最小值]
$f'(x) = 2x - 4$，$f''(x) = 2$。Newton 迭代：
\[
x_{k+1} = x_k - \frac{2x_k - 4}{2} = 2.
\]
\textbf{一步就到}（因为 $f$ 本身是二次的——Newton 法在二次函数上一步收敛）。\textbf{这就是"Newton 法用二次近似——所以二次函数上一步到位"的几何直觉}。
\end{example}

\textbf{Newton 法的弱点}：

\begin{itemize}
\item \textbf{要算 $f''$}——一元还好——高维就是 $n \times n$ 的 Hessian 矩阵——$n=10^6$ 时根本算不动
\item \textbf{$f''(x_k) \le 0$ 时方向反了}——\textbf{你可能反而朝最大值跑}——\textbf{这是 Newton 在非凸问题上的致命伤}
\end{itemize}

\textbf{这就引出了下一节——\textbf{多元 + 无约束 + 替代 Hessian}——这是整个 20 世纪优化研究的主线}。

\section{§2 多元无约束优化 / Multivariate Unconstrained}

\textbf{设 $f: \mathbb{R}^n \to \mathbb{R}$ 光滑}。\textbf{要找 $\bm{x}^* = \arg\min_{\bm x} f(\bm x)$}。

\textbf{一阶必要条件}：$\nabla f(\bm x^*) = \bm 0$（梯度为零）。

\textbf{二阶充分条件}：$\nabla f(\bm x^*) = \bm 0$ 且 $\nabla^2 f(\bm x^*) \succ 0$（Hessian 正定）。

\subsection{Gradient Descent / 梯度下降}

\textbf{最直觉的算法}——\textbf{Cauchy 1847 年提出}——\textbf{沿\textbf{负梯度}方向走一步}：

\begin{equation}
\boxed{\, \bm x_{k+1} = \bm x_k - \alpha_k \nabla f(\bm x_k) \,}
\end{equation}

\textbf{$\alpha_k > 0$ 是\textbf{步长}（learning rate）}。

\textbf{为什么是负梯度方向？}——\textbf{因为梯度指向 $f$ 增长最快的方向——负梯度就是下降最快的方向}。

\textbf{步长怎么选？}——\textbf{三种策略}：

\begin{enumerate}
\item \textbf{固定步长 $\alpha_k = \alpha$}：\textbf{简单——但太大发散——太小慢}
\item \textbf{线搜索 line search}：\textbf{每步沿负梯度方向解 $\min_\alpha f(\bm x_k - \alpha \nabla f)$}（Armijo / Wolfe 条件）
\item \textbf{自适应 adaptive}：\textbf{根据梯度历史动态调整}——\textbf{Adam 的精髓}
\end{enumerate}

\textbf{收敛速度}：

\begin{itemize}
\item \textbf{凸 + Lipschitz 梯度}：$O(1/k)$
\item \textbf{强凸}：\textbf{线性 $O(\rho^k)$}——\textbf{$\rho = (\kappa - 1)/(\kappa+1)$——$\kappa$ 是条件数}
\end{itemize}

\textbf{致命伤}：\textbf{条件数 $\kappa$ 大时\textbf{超慢}}。\textbf{这就是著名的"\textbf{zigzag 现象}"}——\textbf{在长扁的峡谷里来回跳——一直走不到底}。

\subsection{Newton 法（多元）}

\textbf{二阶 Taylor 展开}：

\begin{equation}
f(\bm x) \approx f(\bm x_k) + \nabla f(\bm x_k)^\top (\bm x - \bm x_k) + \tfrac{1}{2}(\bm x - \bm x_k)^\top \nabla^2 f(\bm x_k) (\bm x - \bm x_k)
\end{equation}

\textbf{对 $\bm x$ 求梯度令为零}：

\begin{equation}
\boxed{\, \bm x_{k+1} = \bm x_k - [\nabla^2 f(\bm x_k)]^{-1} \nabla f(\bm x_k) \,}
\end{equation}

\textbf{这是\textbf{多元 Newton 法}}。

\textbf{优点}：\textbf{二次收敛}——\textbf{$\kappa$ 大的问题它不怕}（Hessian 自动"预条件化"了梯度）。

\textbf{缺点}：

\begin{itemize}
\item \textbf{算 Hessian}：$n^2/2$ 个二阶偏导——\textbf{$n=10^6$ 时是 $5\times 10^{11}$ 个数}——\textbf{装不下}
\item \textbf{解 $\nabla^2 f \cdot \bm d = -\nabla f$}：$O(n^3)$——\textbf{$n=10^6$ 时是 $10^{18}$ 次运算}——\textbf{也算不动}
\item \textbf{Hessian 非正定时}：\textbf{方向反了——可能朝鞍点甚至极大跑}
\end{itemize}

\textbf{所以 Newton 法的真正用武之地}：\textbf{$n \le 1000$ 的中小规模 + 强凸问题}。\textbf{大规模 / 非凸——必须找替代品}。

\subsection{Quasi-Newton / 拟牛顿法}

\textbf{Quasi-Newton 的思想——四个字}：\textbf{不算 Hessian——猜 Hessian}。

\textbf{具体地——用迭代过程中的\textbf{梯度变化}来\textbf{估计} Hessian 的逆 $H_k \approx [\nabla^2 f]^{-1}$}。

\textbf{迭代}：

\begin{equation}
\bm x_{k+1} = \bm x_k - \alpha_k H_k \nabla f(\bm x_k)
\end{equation}

\textbf{然后用 $\bm s_k = \bm x_{k+1} - \bm x_k$、$\bm y_k = \nabla f_{k+1} - \nabla f_k$——更新 $H_k \to H_{k+1}$}——\textbf{满足"\textbf{Secant 条件}" $H_{k+1} \bm y_k = \bm s_k$}（多元中值定理的近似）。

\subsection{BFGS / 拟牛顿的王者}

\textbf{1970 年}——\textbf{Broyden、Fletcher、Goldfarb、Shanno \textbf{独立}提出同一个更新公式}——\textbf{以四人首字母命名}：

\begin{equation}
H_{k+1} = \left(I - \frac{\bm s_k \bm y_k^\top}{\bm y_k^\top \bm s_k}\right) H_k \left(I - \frac{\bm y_k \bm s_k^\top}{\bm y_k^\top \bm s_k}\right) + \frac{\bm s_k \bm s_k^\top}{\bm y_k^\top \bm s_k}
\end{equation}

\textbf{看着吓人——其实只用矩阵乘法 + 向量外积}——\textbf{每步 $O(n^2)$}（不是 $O(n^3)$）。

\textbf{性质}：\textbf{保持 $H_k$ 对称正定 + 满足 Secant 条件 + 在二次函数上 $n$ 步收敛}。

\textbf{这是\texttt{scipy.optimize.minimize(method='BFGS')} 的默认选择}——\textbf{$n \le 10^4$ 的工业问题——BFGS 几乎是最优解}。

\textbf{L-BFGS（Limited-memory BFGS, 1989）}——\textbf{只存最近 $m$ 步的 $(\bm s, \bm y)$（典型 $m=10$）}——\textbf{内存 $O(mn)$ 而非 $O(n^2)$}——\textbf{这才让 BFGS 能扩展到 $n=10^6$}——\textbf{今天 \texttt{scipy} 的 \texttt{L-BFGS-B} 就是}。

\begin{tcolorbox}[essbox={无约束优化算法谱系}]
\begin{tabular}{lll}
\toprule
\textbf{算法} & \textbf{每步代价} & \textbf{典型 $n$} \\
\midrule
Gradient Descent & $O(n)$ + 梯度 & $10^9$（深度学习） \\
Newton & $O(n^3)$ + Hessian & $10^3$ \\
BFGS & $O(n^2)$ + 梯度 & $10^4$ \\
L-BFGS & $O(mn)$ + 梯度 & $10^6$ \\
SGD & $O(\text{batch})$ + 梯度 & $10^{11}$（GPT 级） \\
Adam & $O(\text{batch})$ + 梯度 & $10^{11}$（GPT 级） \\
\bottomrule
\end{tabular}
\end{tcolorbox}

\section{§3 凸优化基础 / Convex Optimization}

\textbf{为什么凸优化值得单独讲？}——\textbf{因为它是"\textbf{算法保证收敛}"的最大一族}。\textbf{非凸 = 你永远不知道找到的是不是真最小}；\textbf{凸 = \textbf{局部最优即全局最优}}。

\subsection{凸集 / Convex Set}

\begin{definition}[凸集]
集合 $C \subseteq \mathbb{R}^n$ 是\textbf{凸}的——若对任意 $\bm x, \bm y \in C$ 和 $\lambda \in [0,1]$：
\[
\lambda \bm x + (1-\lambda) \bm y \in C
\]
\textbf{几何意义}：\textbf{连接 $C$ 中任两点的线段——仍在 $C$ 中}。
\end{definition}

\textbf{常见凸集}：

\begin{itemize}
\item \textbf{超平面}：$\{\bm x : \bm a^\top \bm x = b\}$
\item \textbf{半空间}：$\{\bm x : \bm a^\top \bm x \le b\}$
\item \textbf{多面体}：\textbf{有限个半空间的交}——\textbf{线性规划的可行域}
\item \textbf{球}：$\{\bm x : \|\bm x - \bm x_0\| \le r\}$
\item \textbf{半正定锥}：$\mathbb{S}^n_+$——\textbf{SDP 的可行域}
\end{itemize}

\textbf{凸集的交是凸集}——\textbf{这是所有"组合约束"都能保持凸性的根本原因}。

\subsection{凸函数 / Convex Function}

\begin{definition}[凸函数]
$f: \mathbb{R}^n \to \mathbb{R}$ 是\textbf{凸}的——若 dom $f$ 是凸集，且对任意 $\bm x, \bm y$ 和 $\lambda \in [0,1]$：
\[
f(\lambda \bm x + (1-\lambda) \bm y) \le \lambda f(\bm x) + (1-\lambda) f(\bm y)
\]
\textbf{几何意义}：\textbf{函数图像\textbf{下凹}（碗状）}——\textbf{任两点的连线在函数曲面\textbf{之上}}。
\end{definition}

\textbf{凸函数的判别}：

\begin{itemize}
\item \textbf{一阶}：$f(\bm y) \ge f(\bm x) + \nabla f(\bm x)^\top (\bm y - \bm x)$——\textbf{切平面永远在函数下方}
\item \textbf{二阶}：$\nabla^2 f(\bm x) \succeq 0$——\textbf{Hessian 半正定}
\end{itemize}

\textbf{常见凸函数}：\textbf{$\|x\|$（任何范数）、$e^x$、$-\log x$、$x^p$（$p \ge 1$ 或 $p \le 0$）、$\bm x^\top Q \bm x$（$Q \succeq 0$）}。

\textbf{机器学习里几乎所有"显式"损失函数都设计成凸的}：

\begin{itemize}
\item \textbf{平方误差} $\|A\bm x - \bm b\|^2$：\textbf{凸}
\item \textbf{交叉熵 / 逻辑回归损失}：\textbf{凸}
\item \textbf{SVM hinge loss}：\textbf{凸}
\item \textbf{L1 / L2 正则}：\textbf{凸}
\end{itemize}

\textbf{但是}——\textbf{神经网络的整体损失是\textbf{非凸}的}（因为隐藏层的复合非线性）——\textbf{这是深度学习"理论难、实践却 work"的根源}。

\subsection{凸优化的核心定理}

\begin{theorem}[凸优化的全局最优性]
设 $f$ 凸——$C$ 凸——考虑 $\min_{\bm x \in C} f(\bm x)$。则：
\begin{enumerate}
\item \textbf{任一局部最小都是全局最小}
\item \textbf{若 $f$ 严格凸——最小点唯一}
\item \textbf{KKT 条件\textbf{既是必要的也是充分的}}
\end{enumerate}
\end{theorem}

\textbf{这就是 Boyd 反复强调的}：\textbf{"\textbf{凸 vs. 非凸——是优化中比线性 vs. 非线性更深刻的分界}"}。

\section{§4 约束优化 / Constrained Optimization}

\textbf{现实问题都有约束}：\textbf{资源有限、变量有界、平衡条件}。

\textbf{标准形式}：
\begin{equation}
\min_{\bm x} f(\bm x) \quad \text{s.t.} \quad g_i(\bm x) \le 0, \ i=1,\dots,m; \quad h_j(\bm x) = 0, \ j=1,\dots,p
\end{equation}

\subsection{Lagrange 乘子法（等式约束）}

\textbf{先看只有等式约束}：$\min f(\bm x)$ s.t. $h(\bm x) = 0$。

\textbf{构造\textbf{Lagrangian}}：
\begin{equation}
\mathcal{L}(\bm x, \lambda) = f(\bm x) + \lambda h(\bm x)
\end{equation}

\textbf{最优条件}：\textbf{对 $\bm x$ 和 $\lambda$ 同时求偏导令为零}：
\begin{equation}
\nabla_{\bm x} \mathcal{L} = \nabla f + \lambda \nabla h = \bm 0, \quad \nabla_\lambda \mathcal{L} = h(\bm x) = 0
\end{equation}

\textbf{几何意义}（来自第 2 章的梯度）：\textbf{在最优点——$\nabla f$ 与 $\nabla h$ 平行}——\textbf{即 $f$ 的等值线与约束曲面\textbf{相切}}——\textbf{此时再沿约束移动 $f$ 不再下降}。

\subsection{KKT 条件（不等式约束）}

\textbf{再加不等式约束}：$g_i(\bm x) \le 0$。\textbf{Lagrangian}：
\begin{equation}
\mathcal{L}(\bm x, \bm \mu, \bm \lambda) = f(\bm x) + \sum_i \mu_i g_i(\bm x) + \sum_j \lambda_j h_j(\bm x)
\end{equation}

\begin{theorem}[KKT 必要条件]
若 $\bm x^*$ 是局部最优 + 约束规范条件成立——存在 $\bm \mu^* \ge 0$、$\bm \lambda^*$ 使得：
\begin{align}
&\nabla f(\bm x^*) + \sum_i \mu_i^* \nabla g_i(\bm x^*) + \sum_j \lambda_j^* \nabla h_j(\bm x^*) = \bm 0 & \text{（稳定性）}\\
&g_i(\bm x^*) \le 0, \quad h_j(\bm x^*) = 0 & \text{（原始可行）}\\
&\mu_i^* \ge 0 & \text{（对偶可行）}\\
&\mu_i^* g_i(\bm x^*) = 0 \quad \forall i & \text{（互补松弛）}
\end{align}
\end{theorem}

\textbf{第四个——\textbf{互补松弛}——是 KKT 的灵魂}：

\textbf{要么约束在最优点\textbf{激活}（$g_i = 0$）——此时乘子 $\mu_i$ 可以非零}；\\
\textbf{要么约束\textbf{未激活}（$g_i < 0$）——此时乘子必须为零}（$\mu_i = 0$）。

\textbf{物理意义}：\textbf{$\mu_i$ 是"约束的\textbf{影子价格}"——只有真正"卡住"的约束才有价格}。

\begin{example}[KKT 在 SVM 中的应用]
\textbf{支持向量机的对偶形式}（第 8 章详讲）——\textbf{KKT 互补松弛告诉我们}：\textbf{$\mu_i > 0$ 的样本——就是\textbf{支持向量}（在分类边界上）}；\textbf{$\mu_i = 0$ 的样本——离边界还远——对模型\textbf{无影响}}。\textbf{这就是 SVM 名字的来源——\textbf{只有少数支持向量决定了一切}}——\textbf{KKT 是这个性质的数学根源}。
\end{example}

\section{§5 线性规划 + 单纯形法 / Linear Programming}

\textbf{线性规划}（LP）——\textbf{优化里的"hello world"}——\textbf{目标和约束\textbf{都是线性的}}：

\begin{equation}
\min_{\bm x} \bm c^\top \bm x \quad \text{s.t.} \quad A \bm x = \bm b, \ \bm x \ge \bm 0
\end{equation}

\subsection{几何直觉}

\textbf{可行域 $\{\bm x : A\bm x = \bm b, \bm x \ge 0\}$ 是一个\textbf{多面体}（polytope）}——\textbf{目标 $\bm c^\top \bm x$ 在它上面取极值}。

\textbf{关键定理}：\textbf{LP 的最优解（若存在）必在多面体的某个\textbf{顶点}（vertex）取得}。\textbf{因为线性函数在凸集上的极值——必在\textbf{极点}上}。

\subsection{单纯形法 / Simplex (Dantzig 1947)}

\textbf{思想}——\textbf{在顶点之间走}：

\begin{enumerate}
\item \textbf{从一个顶点开始}（基本可行解 BFS）
\item \textbf{找一个"相邻"顶点——使目标值更小}（枢轴变换）
\item \textbf{重复——直到没有更好的邻居}——\textbf{当前顶点即最优}
\end{enumerate}

\textbf{优点}：\textbf{实战极快}——\textbf{大多数问题平均时间 $O(m \log n)$ 级别}。

\textbf{缺点}：\textbf{最坏指数级}——\textbf{Klee-Minty 1972 构造的反例上\textbf{需要访问所有 $2^n$ 个顶点}}。

\subsection{内点法 / Interior Point (Karmarkar 1984)}

\textbf{思想}——\textbf{从可行域的\textbf{内部}穿过去}：\textbf{用\textbf{障碍函数}（barrier）$-\sum_i \log x_i$ 把约束"嵌进"目标}——\textbf{然后用 Newton 法解一系列无约束子问题}。

\textbf{复杂度}：\textbf{$O(n^{3.5} L)$——多项式}。

\textbf{对比}：

\begin{tcolorbox}[essbox={LP 算法对比}]
\begin{tabular}{lll}
\toprule
\textbf{方法} & \textbf{最坏复杂度} & \textbf{实战速度} \\
\midrule
单纯形 & 指数 & 平均很快 \\
椭球（1979） & 多项式 & \textbf{慢得没法用} \\
内点（Karmarkar） & 多项式 & \textbf{大规模问题最快} \\
\bottomrule
\end{tabular}
\end{tcolorbox}

\textbf{今天的工业 solver}：\textbf{CPLEX、Gurobi、Mosek}——\textbf{都同时实现了单纯形 + 内点}——\textbf{自动根据问题大小选择}。\textbf{开源版}：\textbf{HiGHS、SciPy 的 \texttt{linprog}}。

\section{§6 随机梯度下降 SGD / Stochastic Gradient Descent}

\textbf{这一节——开始进入现代机器学习的世界}。

\subsection{为什么需要 SGD？}

\textbf{深度学习的损失函数}：\textbf{对训练集 $N$ 个样本求平均}：
\begin{equation}
\mathcal{L}(\bm \theta) = \frac{1}{N} \sum_{i=1}^N \ell_i(\bm \theta)
\end{equation}

\textbf{标准梯度下降}：
\begin{equation}
\bm \theta_{k+1} = \bm \theta_k - \alpha \cdot \frac{1}{N} \sum_{i=1}^N \nabla \ell_i(\bm \theta_k)
\end{equation}

\textbf{问题}：\textbf{$N = 10^9$（ImageNet）或 $N = 10^{12}$（GPT 训练语料）时}——\textbf{每步要算 $10^{12}$ 个梯度才能更新一次参数}——\textbf{显存装不下 + 时间等不起}。

\subsection{SGD 的诞生}

\textbf{Robbins \& Monro 1951 年}——\textbf{在"随机近似"论文里}——\textbf{提出了一个革命性想法}：

\textbf{不算全量梯度——\textbf{每步随机抽一个样本 $i_k$——用 $\nabla \ell_{i_k}$ 代替整体均值}}：
\begin{equation}
\boxed{\, \bm \theta_{k+1} = \bm \theta_k - \alpha_k \nabla \ell_{i_k}(\bm \theta_k) \,}
\end{equation}

\textbf{这就是\textbf{随机梯度下降 SGD}}。

\textbf{为什么 work？}——\textbf{因为 $\mathbb{E}_i[\nabla \ell_i] = \nabla \mathcal{L}$}——\textbf{\textbf{单样本梯度是全量梯度的无偏估计}}（第 10 章的随机思想——在这里全面爆发）。

\textbf{噪声大}——\textbf{但只要步长 $\alpha_k$ 满足 Robbins-Monro 条件}：
\begin{equation}
\sum_k \alpha_k = \infty, \quad \sum_k \alpha_k^2 < \infty
\end{equation}
\textbf{SGD 几乎必然收敛到稳定点}（强凸下还能给出收敛率）。

\subsection{Mini-batch SGD}

\textbf{现代实践}：\textbf{折中——抽 $B$ 个样本算平均梯度}（$B = 32$ 到 $4096$）：
\begin{equation}
\bm \theta_{k+1} = \bm \theta_k - \alpha_k \cdot \frac{1}{B} \sum_{i \in \mathcal{B}_k} \nabla \ell_i(\bm \theta_k)
\end{equation}

\textbf{三大好处}：
\begin{enumerate}
\item \textbf{方差比单样本小 $B$ 倍}——梯度估计更准
\item \textbf{GPU 友好}——\textbf{batch 维度可\textbf{并行}计算}
\item \textbf{比全量快——比单样本稳}——\textbf{工业默认配置}
\end{enumerate}

\subsection{Momentum / 动量法}

\textbf{SGD 在峡谷里也会 zigzag}——\textbf{解决方案}：\textbf{加\textbf{动量}}（Polyak 1964）：
\begin{align}
\bm v_{k+1} &= \beta \bm v_k + \nabla \ell_{i_k}(\bm \theta_k) \\
\bm \theta_{k+1} &= \bm \theta_k - \alpha \bm v_{k+1}
\end{align}

\textbf{$\bm v$ 是"速度"——它是历史梯度的\textbf{指数加权平均}}（$\beta = 0.9$ 典型）——\textbf{相当于给优化器\textbf{加了惯性}}——\textbf{震荡方向自然抵消——主方向叠加加速}。

\textbf{Nesterov accelerated gradient}（NAG, 1983）——\textbf{在 momentum 的基础上"先看后跳"}——\textbf{凸问题上证明了 $O(1/k^2)$ 的最优收敛率}。

\section{§7 二阶方法 + Adam 优化器 / Second-order Methods and Adam}

\subsection{为什么神经网络需要 Adam？}

\textbf{深度网络的损失曲面有三个噩梦}：

\begin{enumerate}
\item \textbf{尺度极不均匀}：\textbf{不同参数的梯度大小相差 100 万倍}（embedding 层 vs. 输出层）
\item \textbf{鞍点遍布}：\textbf{高维空间里\textbf{鞍点比极小点多得多}}（Dauphin et al. 2014）
\item \textbf{梯度噪声大}：\textbf{batch 太小时 SGD 抖得厉害}
\end{enumerate}

\textbf{解决思路}：\textbf{给每个参数\textbf{自适应学习率}}。

\subsection{AdaGrad（2011）}

\textbf{Duchi、Hazan、Singer 提出}——\textbf{累计每个参数的梯度平方}：
\begin{equation}
G_{k,i} = \sum_{t=1}^k g_{t,i}^2, \quad \theta_{k+1, i} = \theta_{k,i} - \frac{\alpha}{\sqrt{G_{k,i} + \varepsilon}} g_{k,i}
\end{equation}

\textbf{效果}：\textbf{梯度大的参数——学习率自动变小；梯度小的——保持大}。

\textbf{缺点}：\textbf{$G$ 单调累积——后期学习率\textbf{消失到零}}——\textbf{长训练崩溃}。

\subsection{RMSProp（2012, Hinton 讲义未发表）}

\textbf{修补 AdaGrad}——\textbf{把"累计"改成"指数加权平均"}：
\begin{equation}
G_{k,i} = \rho G_{k-1,i} + (1-\rho) g_{k,i}^2, \quad \theta_{k+1,i} = \theta_{k,i} - \frac{\alpha}{\sqrt{G_{k,i} + \varepsilon}} g_{k,i}
\end{equation}

\textbf{$\rho = 0.99$ 典型}。\textbf{学习率不会消失}——\textbf{在 RNN 训练上首次大放异彩}。

\subsection{Adam（2014）}

\textbf{Kingma \& Ba 在 ICLR 2015 发表}——\textbf{把 momentum 和 RMSProp \textbf{合并}}：

\begin{align}
\bm m_k &= \beta_1 \bm m_{k-1} + (1-\beta_1) \bm g_k & \text{（一阶矩 / 动量）} \\
\bm v_k &= \beta_2 \bm v_{k-1} + (1-\beta_2) \bm g_k^2 & \text{（二阶矩 / 梯度平方）} \\
\hat{\bm m}_k &= \bm m_k / (1 - \beta_1^k) & \text{（偏差修正）} \\
\hat{\bm v}_k &= \bm v_k / (1 - \beta_2^k) & \text{（偏差修正）} \\
\bm \theta_{k+1} &= \bm \theta_k - \alpha \cdot \hat{\bm m}_k / (\sqrt{\hat{\bm v}_k} + \varepsilon)
\end{align}

\textbf{典型超参}：$\beta_1 = 0.9$，$\beta_2 = 0.999$，$\varepsilon = 10^{-8}$，$\alpha = 10^{-3}$。

\textbf{Adam 是"二阶方法的近似"}——\textbf{$\hat{\bm v}$ 是 Hessian 对角的 RMS 估计}——\textbf{$\hat{\bm m}/\sqrt{\hat{\bm v}}$ 相当于"\textbf{对角化的 Newton 步}"——但只用一阶信息}。

\textbf{Adam 几乎是 2015 年之后所有深度学习论文的默认优化器}——\textbf{包括 BERT、GPT、ViT、Stable Diffusion}。

\subsection{AdamW（2019）}

\textbf{Loshchilov \& Hutter 发现 Adam 的 L2 正则被"自适应学习率"扭曲}——\textbf{改用"\textbf{解耦权重衰减}"}：
\begin{equation}
\bm \theta_{k+1} = \bm \theta_k - \alpha \cdot \hat{\bm m}_k / (\sqrt{\hat{\bm v}_k} + \varepsilon) - \alpha \lambda \bm \theta_k
\end{equation}

\textbf{这是 GPT-3 / GPT-4 / Llama 训练的实际配置}。

\section{§8 应用：神经网络训练 + 工程设计 + 资源分配}

\textbf{优化无处不在}——\textbf{这一节给三个最具代表性的应用}。

\subsection{应用 1：神经网络训练（climax）}

\textbf{把第 10 章的随机思想 + 本章的优化思想——拼起来}：

\begin{enumerate}
\item \textbf{前向传播}：\textbf{输入 $\bm x \to$ 网络 $\to$ 预测 $\hat{\bm y}$}
\item \textbf{算损失}：$\ell = \text{CrossEntropy}(\hat{\bm y}, \bm y)$
\item \textbf{反向传播}：\textbf{链式法则——算 $\nabla_{\bm \theta} \ell$}（第 2 章的链式法则——在这里规模化）
\item \textbf{优化一步}：\textbf{Adam 更新 $\bm \theta$}
\item \textbf{重复 $10^6$ 次——\textbf{模型学会了}}
\end{enumerate}

\subsection{应用 2：工程设计 —— 飞机翼形优化}

\textbf{Boeing 787 的翼形}——\textbf{用 CFD（计算流体力学）+ 伴随法（adjoint method）算梯度 + L-BFGS 优化}——\textbf{把阻力降低 5\%——每架飞机每年省 200 万美元燃油}。

\subsection{应用 3：资源分配 —— 调度与运筹}

\textbf{Amazon 的仓储调度、Uber 的派单、电力公司的发电机组合（unit commitment）}——\textbf{全部归结为大规模整数 LP / MILP}——\textbf{用 Gurobi 解}——\textbf{每天为这些公司节省数亿美元}。

\section{Python 模块：\texttt{optimization.py}}

\textbf{本章的 Python 模块}——\textbf{把以上算法都实现一遍}（教学版——非工业级）：

\begin{lstlisting}[language=Python]
"""optimization.py —— Chapter 13 companion module.
Educational implementations of optimization algorithms.
"""
import numpy as np

# ---------- 1. Golden Section Search (univariate) ----------
def golden_section(f, a, b, tol=1e-6, max_iter=100):
    """Find min of unimodal f on [a, b] using golden section."""
    phi = (np.sqrt(5) - 1) / 2  # ~0.618
    x1 = b - phi * (b - a)
    x2 = a + phi * (b - a)
    f1, f2 = f(x1), f(x2)
    for _ in range(max_iter):
        if b - a < tol:
            break
        if f1 < f2:
            b, x2, f2 = x2, x1, f1
            x1 = b - phi * (b - a)
            f1 = f(x1)
        else:
            a, x1, f1 = x1, x2, f2
            x2 = a + phi * (b - a)
            f2 = f(x2)
    return (a + b) / 2

# ---------- 2. Gradient Descent ----------
def gradient_descent(f, grad_f, x0, lr=1e-2, tol=1e-6, max_iter=10000):
    """Vanilla gradient descent for f: R^n -> R."""
    x = np.asarray(x0, dtype=float)
    history = [x.copy()]
    for k in range(max_iter):
        g = grad_f(x)
        if np.linalg.norm(g) < tol:
            break
        x = x - lr * g
        history.append(x.copy())
    return x, history

# ---------- 3. Newton's Method ----------
def newton(f, grad_f, hess_f, x0, tol=1e-8, max_iter=100):
    """Newton's method (requires Hessian)."""
    x = np.asarray(x0, dtype=float)
    for k in range(max_iter):
        g = grad_f(x)
        if np.linalg.norm(g) < tol:
            break
        H = hess_f(x)
        d = np.linalg.solve(H, g)
        x = x - d
    return x

# ---------- 4. BFGS ----------
def bfgs(f, grad_f, x0, tol=1e-6, max_iter=500):
    """BFGS quasi-Newton."""
    x = np.asarray(x0, dtype=float)
    n = len(x)
    H = np.eye(n)  # approximation to inverse Hessian
    g = grad_f(x)
    for k in range(max_iter):
        if np.linalg.norm(g) < tol:
            break
        d = -H @ g
        # Backtracking line search (Armijo)
        alpha = 1.0
        while f(x + alpha * d) > f(x) + 1e-4 * alpha * g @ d:
            alpha *= 0.5
            if alpha < 1e-10:
                break
        x_new = x + alpha * d
        g_new = grad_f(x_new)
        s = x_new - x
        y = g_new - g
        rho = 1.0 / (y @ s + 1e-12)
        I = np.eye(n)
        H = (I - rho * np.outer(s, y)) @ H @ (I - rho * np.outer(y, s)) \
            + rho * np.outer(s, s)
        x, g = x_new, g_new
    return x

# ---------- 5. SGD ----------
def sgd(grad_sample, x0, n_samples, lr=1e-2, epochs=10, batch_size=32):
    """Stochastic gradient descent.
    grad_sample(x, i): gradient of loss on sample i.
    """
    x = np.asarray(x0, dtype=float)
    history = []
    for epoch in range(epochs):
        idx = np.random.permutation(n_samples)
        for start in range(0, n_samples, batch_size):
            batch = idx[start:start + batch_size]
            g = np.mean([grad_sample(x, i) for i in batch], axis=0)
            x = x - lr * g
        history.append(x.copy())
    return x, history

# ---------- 6. Adam ----------
def adam(grad_sample, x0, n_samples,
         lr=1e-3, beta1=0.9, beta2=0.999, eps=1e-8,
         epochs=10, batch_size=32):
    """Adam optimizer (Kingma & Ba 2014)."""
    x = np.asarray(x0, dtype=float)
    m = np.zeros_like(x)
    v = np.zeros_like(x)
    t = 0
    for epoch in range(epochs):
        idx = np.random.permutation(n_samples)
        for start in range(0, n_samples, batch_size):
            batch = idx[start:start + batch_size]
            g = np.mean([grad_sample(x, i) for i in batch], axis=0)
            t += 1
            m = beta1 * m + (1 - beta1) * g
            v = beta2 * v + (1 - beta2) * g * g
            m_hat = m / (1 - beta1 ** t)
            v_hat = v / (1 - beta2 ** t)
            x = x - lr * m_hat / (np.sqrt(v_hat) + eps)
    return x

# ---------- 7. Linear Programming via scipy ----------
def solve_lp(c, A_ub=None, b_ub=None, A_eq=None, b_eq=None, bounds=None):
    """Wrapper around scipy.optimize.linprog (HiGHS by default)."""
    from scipy.optimize import linprog
    return linprog(c, A_ub=A_ub, b_ub=b_ub,
                   A_eq=A_eq, b_eq=b_eq,
                   bounds=bounds, method='highs')

# ---------- 8. Lagrangian-based equality-constrained QP ----------
def quadratic_eq_constrained(Q, c, A, b):
    """Solve min 0.5 x'Qx + c'x  s.t. Ax = b via KKT system."""
    n, m = Q.shape[0], A.shape[0]
    KKT = np.block([[Q, A.T], [A, np.zeros((m, m))]])
    rhs = np.concatenate([-c, b])
    sol = np.linalg.solve(KKT, rhs)
    x, lam = sol[:n], sol[n:]
    return x, lam
\end{lstlisting}

\section{Aha 时刻：ChatGPT 训练 = Adam + 千亿参数}

\textbf{现在}——\textbf{我们到达本书最后一个 Aha——也是全书的\textbf{climax}}。

\begin{tcolorbox}[ahabox={\Large Aha：ChatGPT 训练 —— 优化理论 70 年的总和成就}]
\textbf{2022 年 11 月 30 日}——\textbf{OpenAI 发布 ChatGPT}——\textbf{5 天内用户破百万——2 个月破 1 亿}——\textbf{人类历史上增长最快的消费产品}。

\medskip

\textbf{这个对话引擎背后——是一个有 \textbf{1750 亿参数} 的 Transformer 神经网络（GPT-3 规模）}——\textbf{它是怎么训练出来的}？

\medskip

\textbf{一句话——\textbf{Adam 优化器 + 反向传播 + 万张 GPU + 三个月}}。

\medskip

\textbf{展开来说}：

\begin{enumerate}
\item \textbf{目标函数}（第 11 章的最大似然）：\textbf{在 3000 亿个 token 的语料上最大化条件概率 $p(\text{next token} \mid \text{previous tokens})$}——\textbf{等价于最小化\textbf{交叉熵损失}}：
\[
\mathcal{L}(\bm \theta) = -\frac{1}{N} \sum_{i=1}^N \log p_{\bm \theta}(y_i \mid x_i)
\]

\item \textbf{优化变量}：\textbf{$\bm \theta \in \mathbb{R}^{1.75 \times 10^{11}}$}——\textbf{1750 亿个权重}——\textbf{装在 \textbf{GPU 集群}上——一张 80GB 的 A100 只能存\textbf{半层}的参数}

\item \textbf{梯度计算}（第 2 章的链式法则——大规模化）：\textbf{反向传播算 $\nabla_{\bm \theta} \mathcal{L}$}——\textbf{用\textbf{自动微分}（autograd）——一次前向 + 一次反向 = 一次梯度}

\item \textbf{优化步}：\textbf{AdamW——逐参数自适应学习率 + 解耦权重衰减}——\textbf{超参}：\textbf{lr=$6\times 10^{-5}$，$\beta_1=0.9$，$\beta_2=0.95$（GPT-3 配置）}

\item \textbf{batch size}：\textbf{320 万 token}（远超传统 mini-batch）——\textbf{大 batch 降低噪声 + 充分利用 GPU 并行}

\item \textbf{总训练量}：\textbf{$3 \times 10^{23}$ FLOPs}（GPT-3）——\textbf{相当于 1000 块 V100 GPU 跑\textbf{34 天}}

\item \textbf{学习率调度}：\textbf{warmup（前 1\% 步线性升温） + cosine decay（剩余步余弦下降）}——\textbf{所有大模型都这么用}

\item \textbf{梯度裁剪 + 混合精度}：\textbf{$\|\nabla\| > 1.0$ 时裁剪——FP16 训练降低显存 2 倍}
\end{enumerate}

\medskip

\textbf{$\bullet$ 这一切的核心}——\textbf{是 1951 年 Robbins 和 Monro 在《Annals of Math Stat》上发表的\textbf{随机近似}方法}（SGD 的起源）。

\textbf{$\bullet$ 加上 1964 年 Polyak 的\textbf{动量法}}。

\textbf{$\bullet$ 加上 2014 年 Kingma 和 Ba 的 \textbf{Adam}}（两个博士生在 ICLR 的一篇会议论文）。

\textbf{$\bullet$ 加上 1986 年 Rumelhart、Hinton、Williams 的\textbf{反向传播}}（梯度的自动微分实现）。

\medskip

\textbf{这四个想法——加上 2017 年 Google 提出的 \textbf{Transformer 架构}}——\textbf{加上 Nvidia 二十年的 \textbf{GPU 硬件}}——\textbf{在 2020 年代汇合}——\textbf{催生了人类历史上\textbf{第一个能像人一样对话}的机器}。

\medskip

\textbf{ChatGPT 之所以会"思考"——并不是它\textbf{懂}什么}——\textbf{它只是\textbf{把损失函数最小化得太好了}}——\textbf{以至于\textbf{下一个 token 的预测分布}——逼近了\textbf{人类知识的真实分布}}。

\medskip

\textbf{换句话说}——\textbf{ChatGPT = \textbf{在 1750 亿维空间里的一次梯度下降}}。

\medskip

\textbf{这就是\textbf{优化理论的力量}}——\textbf{一个 200 年前 Cauchy 想出的"沿负梯度走一步"的简单想法}——\textbf{在万张 GPU + 千亿参数的尺度上}——\textbf{学会了写诗、写代码、做数学、当律师顾问}。

\medskip

\textbf{这就是为什么本书最后一章——必须是优化}。\textbf{所有前面的章节都在为它\textbf{铺垫}}：

\begin{itemize}
\item \textbf{第 1 章的极限}——\textbf{是梯度的极限定义}
\item \textbf{第 2 章的偏导}——\textbf{是反向传播的链式法则}
\item \textbf{第 7 章的特征分解}——\textbf{是 Adam 自适应学习率的 Hessian 近似}
\item \textbf{第 10 章的随机过程}——\textbf{是 SGD 的概率保证}
\item \textbf{第 12 章的数值方法}——\textbf{是浮点 GPU 计算的根基}
\end{itemize}

\medskip

\textbf{优化——是工科数学的\textbf{终极归宿}}。\textbf{学到这里——你已经站在 AI 时代的数学高地上了}。
\end{tcolorbox}

\textbf{这与我们在 \textbf{physics-rendu 第 21 章}的回响——\textbf{神经网络训练 = 高维空间里的\textbf{物理动力学}}（损失曲面 = 能量曲面、SGD = 朗之万动力学、batch noise = 热涨落）——构成了\textbf{数学与物理在 AI 时代的双线汇合}}。

\section{练习 / Exercises}

\begin{enumerate}
\item \textbf{Newton 法}：用 Newton 法求 $f(x) = x^4 - 4x^2 + 3$ 的极小点——\textbf{用 \texttt{newton} 函数验证}。
\item \textbf{BFGS vs. Newton}：在 Rosenbrock 函数 $f(x,y) = (1-x)^2 + 100(y-x^2)^2$ 上比较——\textbf{BFGS、Newton、Gradient Descent 各需多少步收敛}。
\item \textbf{KKT}：手算 $\min x^2 + y^2$ s.t. $x + y = 1$ 的最优点（用 Lagrangian）。
\item \textbf{LP}：用 \texttt{scipy.optimize.linprog} 解\textbf{食堂配餐问题}（最低成本满足热量与蛋白质约束）。
\item \textbf{SGD vs. Adam}：在 MNIST 逻辑回归上比较——\textbf{画 loss 曲线}——\textbf{Adam 早期下降是否更快}？
\item \textbf{凸性判别}：判断 $f(x) = \log(\sum_i e^{x_i})$（\textbf{log-sum-exp}）是否凸——\textbf{它在 softmax 中扮演什么角色}？
\item \textbf{动量}：在 ill-conditioned 二次函数上比较\textbf{有无 momentum 的 SGD}——\textbf{画轨迹}。
\item \textbf{Adam 拆解}：在 PyTorch 里手写 Adam（不调 \texttt{torch.optim.Adam}）——\textbf{在一个小网络上跑通——验证与官方实现等价}。
\end{enumerate}

\section{回顾 / Recap}

\begin{tcolorbox}[essbox={本章一句话总结}]
\textbf{优化 = 找极小点}。\textbf{无约束}：\textbf{Gradient $\to$ Newton $\to$ BFGS $\to$ SGD $\to$ Adam}——\textbf{越往后越能处理大规模 + 非凸}。\textbf{有约束}：\textbf{KKT 是一切的判别式}。\textbf{现代 AI 的核心——是 \textbf{Adam + 千亿参数}——这是 200 年优化理论 + 70 年随机近似 + 10 年深度学习的总和}。
\end{tcolorbox}

\textbf{至此}——\textbf{这本书的正文部分写完了}。\textbf{从第 1 章 Newton 的瘟疫之年——到第 13 章 ChatGPT 的训练}——\textbf{我们走过了 360 年的工科数学}。

\textbf{下一章}（第 14 章，可选）：\textbf{离散数学——图论、组合优化、密码学的入门}——\textbf{为想做 CS 方向的同学准备}。

\textbf{但如果你只想做 AI / 控制 / 信号处理 / 通信——\textbf{你已经够用了}}。

\textbf{祝你——\textbf{打通任督二脉}}。

\clearpage


\part{离散数学（可选）/ Discrete Math (Optional)}
% =====================================================================
%  第 14 章 · 离散数学初步 / Discrete Math (Optional)
%  组合 + 图论 + 布尔逻辑 + 复杂度 + 计算机应用
%  小故事：Euler 的 1736 Königsberg 七桥问题——图论的诞生
%  Aha：旅行商问题 TSP（NP-complete）
% =====================================================================

\chapter{离散数学初步 / Discrete Math \textnormal{(Optional)}}
\label{ch:discrete}

\begin{quote}\itshape
1736 年，俄国 Königsberg 城。一群闲来无事的市民——\\
\textbf{在普雷格尔河的七座桥上散步}——他们提了一个问题：\\
\textbf{能不能一次走完七座桥——每座只走一次}？\\
\textbf{Euler 把这个散步问题——抽象成了点和线}——\textbf{图论由此诞生}。\\
本章——\textbf{讲的就是 Euler 那一笔下去的世界}：\textbf{离散}。
\end{quote}

\section{写在前面 / Preface}

\textbf{这本书的前 13 章}——\textbf{都是连续数学}：\textbf{极限、导数、积分、矩阵、概率密度}——\textbf{背后的世界是\textbf{光滑}的}。

\textbf{但工程的另一半世界——是\textbf{离散}的}：

\begin{itemize}
\item \textbf{你写的每一行 Python 代码}——\textbf{是离散的字符序列}
\item \textbf{你训练神经网络用的 GPU}——\textbf{是离散的逻辑门}
\item \textbf{你刷的微信消息}——\textbf{是离散的字节包}
\item \textbf{你出门用的 GPS 路径规划}——\textbf{是离散的图搜索}
\item \textbf{你网购的快递分拣}——\textbf{是离散的组合优化}
\end{itemize}

\textbf{我大学时学过《离散数学》——闵嗣鹤那本——薄薄的——但读起来像天书}。\textbf{真值表、合取范式、Warshall 算法、图同构}——\textbf{符号一大堆——不知道有什么用}。

\textbf{直到我工作后做了一个\textbf{物流路径优化的小项目}——\textbf{我才反应过来}}：

\begin{itemize}
\item \textbf{最短路径——是 Dijkstra}
\item \textbf{快递分拣的最优顺序——是 TSP}
\item \textbf{快递员排班——是着色问题}
\item \textbf{评估这件事难不难——是 P vs NP}
\end{itemize}

\textbf{这一章——\textbf{把离散数学的五件核心事——讲到能跑代码}}：\textbf{组合、图论、布尔逻辑、复杂度、应用}。

\textbf{这是一本\textbf{为工科生写的离散导引}——不是 CS 专业的离散教材}——\textbf{所以薄}——\textbf{但够用}。

\section{一句话本质 / The Essence in One Sentence}

\begin{tcolorbox}[essbox={一句话本质}]
\textbf{中文}：\textbf{离散数学研究\textbf{可数对象之间的结构与关系}——\textbf{组合是数数、图论是连线、逻辑是判断、复杂度是分级}——\textbf{它们合起来是\textbf{计算机科学的数学骨架}}}。\\[6pt]
\textbf{English}: Discrete mathematics studies the structure and relations among countable objects; combinatorics counts, graph theory connects, logic decides, complexity classifies——together they form the mathematical skeleton of computer science.
\end{tcolorbox}

\textbf{连续数学服务于物理与控制}——\textbf{离散数学服务于计算机与通信}。\textbf{现代工程师——两边都要懂一点}。

\section{教材里你看到的 / What the Textbook Tells You}

\subsection{国内教材：闵嗣鹤 / 屈婉玲 / 耿素云}

\textbf{屈婉玲、耿素云的《离散数学》（清华版）}——\textbf{是国内 CS 系的标配}。\textbf{它的结构}：

\begin{itemize}
\item 数理逻辑（命题 + 谓词）
\item 集合论
\item 代数系统（群、环、域、格、布尔代数）
\item 图论
\item 组合数学
\end{itemize}

\textbf{优点}：\textbf{严格 + 全 + 考研标配}。\\
\textbf{缺点}：\textbf{讲群环域占了一半篇幅——\textbf{工科 99\% 用不到}}——\textbf{真正有用的图论和组合反而被压缩}。

\subsection{国外教材：Rosen《Discrete Mathematics and Its Applications》}

\textbf{Rosen 的 \emph{Discrete Mathematics and Its Applications}}（McGraw-Hill）——\textbf{是欧美 CS 系的国民教材}——\textbf{1000+ 页彩印}——\textbf{每个概念都配\textbf{密码、网络、算法}的例子}。

\textbf{它比国内教材好的地方}：\textbf{应用驱动——每讲一个定理——就问"在 RSA 里怎么用？"}。

\textbf{缺点}：\textbf{1000 页——你也读不完}。

\subsection{Knuth《Concrete Mathematics》}

\textbf{Knuth} 的 \emph{Concrete Mathematics}（"Concrete" = "Continuous" $\cap$ "Discrete"）——\textbf{是离散里的\textbf{圣经}}——\textbf{但偏算法分析}——\textbf{适合做学问的人}——\textbf{不适合入门}。

\subsection{我的策略}

\textbf{本章只做四件事}：

\begin{enumerate}
\item \textbf{组合}——\textbf{排列、组合、容斥——这是概率与算法分析的基础}
\item \textbf{图论}——\textbf{Euler、Dijkstra、MST、着色——\textbf{99\% 工业用的图算法都在这}}
\item \textbf{布尔代数 + 命题逻辑}——\textbf{数字电路、SQL where 子句、机器学习决策树的底层}
\item \textbf{复杂度}——\textbf{Big-O、P/NP——告诉你"这事能不能做"}
\end{enumerate}

\begin{tcolorbox}[storybox={\Large 小故事：Euler 与 Königsberg 七桥（1736）}]
\textbf{1736 年}——\textbf{东普鲁士 Königsberg 城}（今俄罗斯加里宁格勒）。\textbf{城里有一条普雷格尔河}——\textbf{河里有两座岛}——\textbf{岛与两岸之间——\textbf{共有七座桥}}。

\medskip

市民们闲来无事——\textbf{提了一个问题}：

\begin{quote}\itshape
"能不能找到一条路线——\textbf{一次走完七座桥}——\textbf{每座只走一次}——\textbf{回到出发点}？"
\end{quote}

大家走了几个月——\textbf{没人能做到}——\textbf{但也没人能证明做不到}。

\medskip

\textbf{Leonhard Euler}（29 岁，那时还在圣彼得堡）——\textbf{听说了这个问题}——\textbf{他做了一件革命性的事}：

\begin{quote}
\textbf{他把"桥"画成"线"——\textbf{把"陆地"画成"点"}}。
\end{quote}

\textbf{于是七桥问题}——\textbf{变成了 4 个点 + 7 条边的图}。

\medskip

\textbf{Euler 证明}：

\begin{quote}\itshape
\textbf{一笔画的充要条件}——\textbf{所有点的度数都是偶数}（回路）——\textbf{或者只有 2 个奇度点}（非回路）。
\end{quote}

\textbf{Königsberg 的 4 个点——\textbf{度数分别是 3, 3, 3, 5——\textbf{全是奇数}}}——\textbf{所以做不到}。

\medskip

\textbf{Euler 这一笔下去}——\textbf{发明了\textbf{图论}}——\textbf{发明了\textbf{拓扑学的雏形}}——\textbf{为 200 年后的\textbf{电路网络、互联网、社交网络、神经网络}——\textbf{铺好了路}}。

\medskip

\textbf{1875 年——Königsberg 又建了一座桥（第 8 座）}——\textbf{度数全变了}——\textbf{这次\textbf{有解}}。\textbf{但那时——\textbf{已经不重要了}}——\textbf{因为图论已经长成大树}。

\medskip

\textbf{二战中——\textbf{Königsberg 被苏军炸毁了 5 座桥}}——\textbf{现在只剩 2 座}——\textbf{但 Euler 那一笔——\textbf{永远在那里}}。
\end{tcolorbox}

\section{§1 组合数学 / Combinatorics}

\subsection{排列 / Permutation}

\textbf{从 $n$ 个不同元素中取 $k$ 个——\textbf{按顺序}排成一列}：

\[
P(n,k) = \frac{n!}{(n-k)!} = n(n-1)(n-2)\cdots(n-k+1)
\]

\textbf{典型例子}：\textbf{10 匹马跑步——前 3 名的颁奖顺序有多少种}？$P(10,3) = 10 \times 9 \times 8 = 720$ 种。

\subsection{组合 / Combination}

\textbf{从 $n$ 个不同元素中取 $k$ 个——\textbf{不分顺序}}：

\[
\binom{n}{k} = C(n,k) = \frac{n!}{k!(n-k)!}
\]

\textbf{典型例子}：\textbf{36 人班里选 5 个班委——多少种选法}？$\binom{36}{5} = 376{,}992$ 种。

\textbf{记忆要点}：\textbf{排列 = 顺序重要}（金银铜）——\textbf{组合 = 顺序无关}（班委没差别）。

\subsection{二项式定理 / Binomial Theorem}

\[
(a+b)^n = \sum_{k=0}^{n} \binom{n}{k} a^{n-k} b^k
\]

\textbf{这就是\textbf{第 10 章二项分布}的源头}——\textbf{抛 $n$ 次硬币——\textbf{正面 $k$ 次的概率}——是 $\binom{n}{k} p^k (1-p)^{n-k}$}。

\subsection{容斥原理 / Inclusion-Exclusion}

\textbf{两个集合}：$|A \cup B| = |A| + |B| - |A \cap B|$。

\textbf{三个集合}：
\[
|A \cup B \cup C| = |A| + |B| + |C| - |A \cap B| - |A \cap C| - |B \cap C| + |A \cap B \cap C|
\]

\textbf{$n$ 个集合的一般式}：
\[
\Big|\bigcup_{i=1}^{n} A_i\Big| = \sum_{k=1}^{n} (-1)^{k+1} \sum_{|S|=k} \Big|\bigcap_{i \in S} A_i\Big|
\]

\textbf{典型应用}：\textbf{错排问题——$n$ 个人交换礼物——\textbf{没有一个人拿到自己的礼物——的概率}}——\textbf{用容斥可证——\textbf{$n \to \infty$ 时趋于 $1/e \approx 0.368$}}。

\section{§2 图论基础 / Graph Theory}

\subsection{图的定义 / Definition}

\textbf{图} $G = (V, E)$——\textbf{$V$ 是点集——$E$ 是边集}。

\begin{itemize}
\item \textbf{无向图 / 有向图}（边有无方向）
\item \textbf{加权图 / 无权图}（边有无权重）
\item \textbf{度数 degree}：\textbf{与某个顶点相连的边数}
\end{itemize}

\textbf{握手定理}：$\sum_{v \in V} \deg(v) = 2|E|$——\textbf{所有度数加起来 = 边数 $\times$ 2}。

\subsection{树 / Tree}

\textbf{树} = \textbf{无环连通图}——\textbf{$n$ 个点的树——恰有 $n-1$ 条边}。

\textbf{两个关键应用}：

\begin{itemize}
\item \textbf{二叉搜索树 BST}——\textbf{数据库索引}
\item \textbf{最小生成树 MST}——\textbf{Kruskal / Prim 算法}——\textbf{铺光缆、修铁路——\textbf{用最少的代价连通所有点}}
\end{itemize}

\subsection{最短路径 / Shortest Path}

\textbf{Dijkstra 算法}（1956 年——\textbf{20 分钟在咖啡店里想出来}）：

\textbf{思想}：\textbf{从起点出发——\textbf{每次扩展当前距离最小的未访问点}}——\textbf{时间复杂度 $O((V+E)\log V)$}（用堆）。

\textbf{应用}：

\begin{itemize}
\item \textbf{高德地图 / Google Maps} —— \textbf{开车路径规划}
\item \textbf{网络路由} —— \textbf{OSPF 协议}
\item \textbf{物流配送}
\end{itemize}

\textbf{Bellman-Ford} —— \textbf{允许负权——\textbf{$O(VE)$}}。\textbf{Floyd-Warshall} —— \textbf{所有点对——\textbf{$O(V^3)$}}。

\subsection{着色 / Coloring}

\textbf{图着色问题}：\textbf{给图的每个顶点涂色——\textbf{相邻顶点颜色不同}——\textbf{最少几种颜色}}？

\begin{itemize}
\item \textbf{地图四色定理}（1976 年首次借助计算机证明）—— \textbf{平面图 4 色足够}
\item \textbf{考试排考场}—— \textbf{相邻时段不能同人考}—— \textbf{顶点 = 课、边 = 冲突、颜色 = 时段}
\item \textbf{编译器寄存器分配}—— \textbf{LLVM 用图着色把变量映射到 CPU 寄存器}
\end{itemize}

\textbf{这是一个 NP-hard 问题}——\textbf{下一节会讲}。

\section{§3 布尔代数 + 命题逻辑 / Boolean Algebra and Propositional Logic}

\subsection{布尔代数 / Boolean Algebra}

\textbf{由 George Boole（1847 年）创立}——\textbf{只有两个值}：$\{0, 1\}$ 或 $\{\text{false}, \text{true}\}$。

\textbf{三个基本运算}：

\begin{itemize}
\item \textbf{与 AND}：$a \land b$（都为真才真）
\item \textbf{或 OR}：$a \lor b$（任一真即真）
\item \textbf{非 NOT}：$\neg a$（取反）
\end{itemize}

\textbf{De Morgan 律}：
\[
\neg(a \land b) = \neg a \lor \neg b, \qquad \neg(a \lor b) = \neg a \land \neg b
\]

\textbf{这是\textbf{Shannon 1937 年硕士论文}的灵感}——\textbf{他证明了\textbf{布尔代数可以用电路实现}}——\textbf{从而开启了数字电路 + 计算机时代}。

\subsection{命题逻辑 / Propositional Logic}

\textbf{命题}：\textbf{非真即假的陈述句}——例如 "$2+2 = 4$"。

\textbf{蕴含 implication}：$p \to q$——\textbf{"若 $p$ 则 $q$"}——\textbf{只有 $p$ 真而 $q$ 假时才为假}。

\textbf{等价 equivalence}：$p \leftrightarrow q$。

\textbf{重要等式}：
\[
p \to q \;\equiv\; \neg p \lor q \;\equiv\; \neg q \to \neg p \quad (\text{逆否命题})
\]

\subsection{真值表 / Truth Table}

\textbf{$n$ 个变量——\textbf{有 $2^n$ 行}}。\textbf{用真值表——\textbf{你可以机械地证明任何命题逻辑等价式}}。

\textbf{现代应用}：\textbf{SQL 的 WHERE 子句、Python 的 if 条件、神经网络的 ReLU 门控、CPU 的 ALU 设计}——\textbf{全是布尔代数的工程化}。

\section{§4 算法复杂度 / Complexity}

\subsection{Big-O 记号}

\textbf{$f(n) = O(g(n))$ 表示}：\textbf{存在常数 $C$ 和 $n_0$——\textbf{使得 $f(n) \le C \cdot g(n)$ 对所有 $n \ge n_0$ 成立}}。

\textbf{常见量级}（从快到慢）：
\[
O(1) < O(\log n) < O(n) < O(n \log n) < O(n^2) < O(n^3) < O(2^n) < O(n!)
\]

\textbf{典型例子}：

\begin{itemize}
\item \textbf{二分查找}：$O(\log n)$
\item \textbf{归并排序 / 快排}：$O(n \log n)$
\item \textbf{矩阵乘法（朴素）}：$O(n^3)$
\item \textbf{穷举所有子集}：$O(2^n)$
\item \textbf{TSP 暴力解}：$O(n!)$
\end{itemize}

\subsection{P 与 NP}

\textbf{P 类}：\textbf{多项式时间内可解的问题}（$O(n^k)$ 某个常数 $k$）。\textbf{"容易做"}。

\textbf{NP 类}：\textbf{多项式时间内可\textbf{验证答案}的问题}。\textbf{"容易验证"}。

\textbf{显然 $P \subseteq NP$}——\textbf{但\textbf{$P = NP$ 吗}}？——\textbf{这是\textbf{Clay 研究所 7 大千禧难题之一}——\textbf{悬赏 100 万美金}——\textbf{至今未解}}。

\textbf{NP-complete}：\textbf{NP 中"最难"的一类问题}——\textbf{任何 NP 问题都可在多项式时间归约到它}。\textbf{若任何一个 NP-complete 问题有多项式解——\textbf{$P = NP$ 立刻成立}}。

\textbf{经典 NP-complete 问题}：

\begin{itemize}
\item \textbf{SAT}（布尔可满足性）—— \textbf{第一个被证明 NP-complete 的}（Cook 1971）
\item \textbf{TSP}（旅行商问题）
\item \textbf{图着色}（$k \ge 3$）
\item \textbf{背包问题}（0-1 Knapsack）
\item \textbf{子集和问题}
\end{itemize}

\textbf{对工程师的实操含义}：

\begin{quote}\itshape
\textbf{遇到 NP-complete 问题——\textbf{别想着求精确解}}——\textbf{用\textbf{近似算法、启发式、ILP solver、模拟退火}}—— \textbf{已经够好了}。
\end{quote}

\section{§5 应用：计算机 + 网络 + 编码 + 密码}

\subsection{计算机：数字电路与 CPU}

\textbf{每一个 CPU 内部}——\textbf{全是\textbf{NAND 门}（与非门）}——\textbf{用 NAND 可以构造\textbf{任意布尔函数}——\textbf{所以 NAND 称为"通用门"}}。

\textbf{你用 Python 写的 \texttt{if a and b:}}——\textbf{最终编译成的机器码——\textbf{是几纳秒内 NAND 门翻转的物理过程}}。

\subsection{网络：路由协议}

\textbf{互联网的核心 = 一个巨大的有向加权图}：

\begin{itemize}
\item \textbf{顶点}：\textbf{路由器、自治系统 AS}
\item \textbf{边}：\textbf{物理链路、带宽、延迟}
\item \textbf{算法}：\textbf{OSPF 用 Dijkstra} —— \textbf{BGP 用距离向量 + 路径选择}
\end{itemize}

\subsection{编码：Hamming 与 Reed-Solomon}

\textbf{你的硬盘、SSD、QR 码、蓝光光盘}——\textbf{全靠\textbf{纠错码}抗噪声}：

\begin{itemize}
\item \textbf{Hamming 码 (7,4)} —— \textbf{4 比特数据 + 3 比特校验 = 可纠 1 位错}
\item \textbf{Reed-Solomon 码} —— \textbf{CD / DVD / QR 的标配}
\item \textbf{LDPC / Turbo 码} —— \textbf{5G + WiFi 6 的物理层}
\end{itemize}

\textbf{这些码的基础}——\textbf{是\textbf{有限域 GF($2^n$) 上的线性代数}}——\textbf{离散版的第 7 章}。

\subsection{密码：RSA 与椭圆曲线}

\textbf{RSA}（1977）的核心：

\begin{enumerate}
\item \textbf{选两个大素数 $p, q$}（每个 1024 位）
\item \textbf{$n = pq$——容易算}
\item \textbf{从 $n$ 反推 $p, q$——\textbf{是 NP-hard 的因子分解问题}——\textbf{现代计算机算不动}}
\end{enumerate}

\textbf{这就是 RSA 安全性的根基}：\textbf{乘法易、因子分解难}。

\textbf{若量子计算机成熟}——\textbf{Shor 算法可在多项式时间分解大数——\textbf{RSA 就崩了}}——\textbf{这就是为什么\textbf{后量子密码学}（lattice-based）正在成为新标准}。

\section{Python 模块：\texttt{discrete.py}}

本章对应的 Python 模块——\textbf{把组合、图论、布尔逻辑、复杂度示例——全部集中}。

\begin{lstlisting}[language=Python, caption={discrete.py 核心 API}]
"""
discrete.py —— Discrete math toolkit
组合 + 图论 + 布尔 + 复杂度
依赖：numpy, networkx, matplotlib
"""

import math
from itertools import permutations, combinations
import networkx as nx
import matplotlib.pyplot as plt


# -------------------- §1 Combinatorics --------------------

def n_perm(n, k):
    """P(n,k) = n!/(n-k)!"""
    return math.perm(n, k)

def n_comb(n, k):
    """C(n,k) = n!/(k!(n-k)!)"""
    return math.comb(n, k)

def derangement(n):
    """错排数 D_n —— 用容斥公式"""
    return round(math.factorial(n) * sum(
        (-1)**k / math.factorial(k) for k in range(n + 1)
    ))


# -------------------- §2 Graph Theory --------------------

def shortest_path(edges, src, dst):
    """Dijkstra：edges = [(u,v,w), ...]"""
    G = nx.Graph()
    G.add_weighted_edges_from(edges)
    return nx.shortest_path(G, src, dst, weight='weight'), \
           nx.shortest_path_length(G, src, dst, weight='weight')

def min_spanning_tree(edges):
    """Kruskal MST"""
    G = nx.Graph()
    G.add_weighted_edges_from(edges)
    return list(nx.minimum_spanning_edges(G, data=True))

def graph_coloring(edges):
    """贪心着色——返回每个点的颜色编号"""
    G = nx.Graph()
    G.add_edges_from(edges)
    return nx.coloring.greedy_color(G, strategy='largest_first')


# -------------------- §3 Boolean Logic --------------------

def truth_table(expr, variables):
    """expr 是一个 lambda；variables 是变量名列表
       返回真值表（list of dict）"""
    rows = []
    for bits in range(2**len(variables)):
        assignment = {v: bool((bits >> i) & 1)
                      for i, v in enumerate(variables)}
        assignment['out'] = expr(**assignment)
        rows.append(assignment)
    return rows


# -------------------- §4 Complexity demo --------------------

def tsp_bruteforce(dist):
    """TSP O(n!)——只能小规模玩"""
    n = len(dist)
    best, best_path = float('inf'), None
    for perm in permutations(range(1, n)):
        path = (0,) + perm + (0,)
        cost = sum(dist[path[i]][path[i+1]] for i in range(n))
        if cost < best:
            best, best_path = cost, path
    return best, best_path
\end{lstlisting}

\textbf{安装}：\texttt{pip install networkx matplotlib}。

\section{Aha 时刻：旅行商问题 TSP / Aha: Travelling Salesman}

\begin{tcolorbox}[ahabox={\Large Aha：TSP 与 NP-complete 的边界}]
\textbf{问题}：\textbf{一个推销员要访问 $n$ 个城市——\textbf{每个城市恰好一次}——\textbf{最后回到出发点}——\textbf{总路程最短}}。

\medskip

\textbf{暴力解}：\textbf{枚举所有 $(n-1)!/2$ 条回路}——\textbf{选最短的}。

\begin{itemize}
\item $n = 10$：\textbf{$181{,}440$ 条}——\textbf{毫秒}
\item $n = 15$：\textbf{$4 \times 10^{10}$ 条}——\textbf{几分钟}
\item $n = 20$：\textbf{$6 \times 10^{16}$ 条}——\textbf{几个月}
\item $n = 30$：\textbf{$4 \times 10^{30}$ 条}——\textbf{\textbf{宇宙年龄都不够}}
\end{itemize}

\medskip

\textbf{TSP 是\textbf{典型的 NP-hard}}——\textbf{没有已知的多项式算法}。

\medskip

\textbf{但工业里 TSP 必须解}：\textbf{京东配送、Uber 接单、芯片打线、PCB 钻孔}——\textbf{都是 TSP}。

\medskip

\textbf{工程师怎么办}？——\textbf{我们不求精确解——\textbf{求近似解}}：

\begin{enumerate}
\item \textbf{贪心}（最近邻）—— $O(n^2)$ —— \textbf{快}——\textbf{但可能差 25\%}
\item \textbf{Christofides 算法}（基于 MST + 完美匹配）—— $O(n^3)$ —— \textbf{保证不超过最优 $1.5$ 倍}
\item \textbf{2-opt / Lin-Kernighan} —— \textbf{局部搜索}——\textbf{实际效果接近最优}
\item \textbf{模拟退火 / 遗传算法 / 蚁群算法} —— \textbf{随机启发式}
\item \textbf{Concorde solver}（最强精确求解器）—— \textbf{2006 年解了 $85{,}900$ 城问题}
\end{enumerate}

\medskip

\textbf{这是\textbf{工程实用主义的胜利}}——\textbf{理论说 NP-hard——\textbf{工程说"够好就行"}}。

\medskip

\textbf{Aha}：\textbf{学了 P/NP——\textbf{你才知道哪些事\textbf{别白费力气求最优}}——\textbf{这本身就是一种\textbf{智慧}}}。

\medskip

\textbf{TSP 与第 13 章的优化呼应}：\textbf{凸优化 = P（多项式可解）——\textbf{TSP = NP-hard（启发式 + 近似）}}——\textbf{两者一起——\textbf{构成现代工程优化的完整图景}}。
\end{tcolorbox}

\section{练习 / Exercises}

\begin{enumerate}
\item \textbf{容斥}：\textbf{用容斥原理证明}——\textbf{$n=5$ 时错排数 $D_5 = 44$}——\textbf{并验证 $D_n/n! \to 1/e$}。
\item \textbf{Dijkstra}：\textbf{在 \texttt{networkx} 里构造一个 6 点 9 边的加权图——\textbf{用 \texttt{shortest\_path} 验证}}。
\item \textbf{MST}：\textbf{给定 8 个城市坐标——\textbf{用 \texttt{min\_spanning\_tree} 算最小连通路径}——\textbf{并画图}}。
\item \textbf{着色}：\textbf{用 \texttt{graph\_coloring} 给\textbf{彼得森图}（10 顶 15 边）着色}——\textbf{需要几种颜色}？
\item \textbf{布尔}：\textbf{用 \texttt{truth\_table} 验证 De Morgan 律}——\textbf{$\neg(a \land b) \equiv \neg a \lor \neg b$}。
\item \textbf{TSP}：\textbf{用 \texttt{tsp\_bruteforce} 在 8 个随机城市上求最短回路}——\textbf{并与贪心最近邻比较\textbf{差多少}}。
\item \textbf{复杂度}：\textbf{对 $n = 100, 1000, 10000$ 的列表——\textbf{分别用线性查找和二分查找——\textbf{画时间 vs $n$ 的对数图}——\textbf{验证 $O(n)$ vs $O(\log n)$}}}。
\item \textbf{RSA}：\textbf{用 \texttt{sympy.nextprime} 选两个 50 位素数——\textbf{算 $n = pq$}——\textbf{试着用 \texttt{sympy.factorint} 分解}——\textbf{观察用时}}。
\end{enumerate}

\section{回顾 / Recap}

\begin{tcolorbox}[essbox={本章一句话总结}]
\textbf{离散数学 = 计算机科学的数学骨架}。\textbf{组合数数——图论连线——逻辑判断——复杂度分级}——\textbf{TSP 告诉你"不是所有问题都能精确解"}。\textbf{Euler 的一笔下去——\textbf{开启了 290 年的图论与计算机时代}}。
\end{tcolorbox}

\textbf{至此}——\textbf{这本书的 14 章\textbf{全部讲完}}——\textbf{从 Newton 的微积分——到 Adam 的优化——到 Euler 的图论}——\textbf{我们走过了\textbf{从 1665 到 2026 共 361 年}的工科数学}。

\textbf{附录 A 给你\textbf{符号速查表}——\textbf{附录 B 给你\textbf{所有 Python 模块的索引}}——\textbf{这本书可以放在案头一辈子}}。

\textbf{祝你——\textbf{打通任督二脉}}。

\clearpage


\appendix
% =====================================================================
%  附录 A · 数学符号表 / Symbol Table
%  按主题分组——配中英文名称 + 章节出处
% =====================================================================

\chapter{数学符号表 / Symbol Table}
\label{app:symbols}

\begin{quote}\itshape
数学家是\textbf{符号的工匠}。\textbf{$\int$ 写不好——读者一辈子记不住}。\\
\textbf{本附录把全书出现过的符号——按主题分门别类——配上中英文名称与首次出现的章节}——\\
\textbf{你忘了哪个符号——回到这里 5 秒钟找到}。
\end{quote}

\section{微积分 / Calculus}

\begin{center}
\renewcommand{\arraystretch}{1.25}
\begin{tabular}{@{}llll@{}}
\toprule
\textbf{符号} & \textbf{中文名称} & \textbf{English Name} & \textbf{章节} \\
\midrule
$\displaystyle \lim_{x \to a} f(x)$ & 极限 & Limit & Ch 1 \\
$\varepsilon, \delta$ & 任意小、足够小 & Epsilon, delta & Ch 1 \\
$f'(x)$ & 一阶导数 & First derivative & Ch 1 \\
$\dfrac{dy}{dx}$ & Leibniz 导数记号 & Leibniz notation & Ch 1 \\
$\dot{x}, \ddot{x}$ & 时间导数（点） & Time derivative (Newton dot) & Ch 1, Ch 4 \\
$f''(x), f^{(n)}(x)$ & 高阶导数 & Higher derivatives & Ch 1 \\
$df$ & 微分 & Differential & Ch 1 \\
$\Delta x$ & 增量 & Increment & Ch 1 \\
$\displaystyle \int f(x)\,dx$ & 不定积分 & Indefinite integral & Ch 1 \\
$\displaystyle \int_a^b f(x)\,dx$ & 定积分 & Definite integral & Ch 1 \\
$\displaystyle \oint$ & 围道积分 & Contour / closed integral & Ch 3, Ch 9 \\
$\sum, \prod$ & 求和、求积 & Sum, product & Ch 1 \\
$e, \pi, \gamma$ & 自然底数、圆周率、Euler 常数 & Math constants & Ch 1 \\
$\infty$ & 无穷 & Infinity & Ch 1 \\
$O(g(n))$ & 大 O 记号 & Big-O notation & Ch 12, Ch 14 \\
$o(g(n))$ & 小 o 记号 & Little-o & Ch 1 \\
\bottomrule
\end{tabular}
\end{center}

\section{多元微积分 / Multivariable Calculus}

\begin{center}
\renewcommand{\arraystretch}{1.25}
\begin{tabular}{@{}llll@{}}
\toprule
\textbf{符号} & \textbf{中文名称} & \textbf{English Name} & \textbf{章节} \\
\midrule
$\dfrac{\partial f}{\partial x}$ & 偏导数 & Partial derivative & Ch 2 \\
$\nabla f$ & 梯度 & Gradient & Ch 2, Ch 13 \\
$\nabla^2 f$ 或 $\Delta f$ & Laplace 算子 & Laplacian & Ch 3, Ch 5 \\
$\nabla \cdot \mathbf{F}$ & 散度 & Divergence & Ch 3 \\
$\nabla \times \mathbf{F}$ & 旋度 & Curl & Ch 3 \\
$\mathbf{H} = \nabla^2 f$ & Hessian 矩阵 & Hessian matrix & Ch 2, Ch 13 \\
$\mathbf{J}$ & Jacobian 矩阵 & Jacobian matrix & Ch 2 \\
$J = \det \mathbf{J}$ & Jacobian 行列式 & Jacobian determinant & Ch 2 \\
$d\mathbf{r}, d\mathbf{S}, dV$ & 弧元、面元、体元 & Line/area/volume element & Ch 3 \\
$\iint, \iiint$ & 二重、三重积分 & Double, triple integral & Ch 2 \\
\bottomrule
\end{tabular}
\end{center}

\section{矢量分析 / Vector Analysis}

\begin{center}
\renewcommand{\arraystretch}{1.25}
\begin{tabular}{@{}llll@{}}
\toprule
\textbf{符号} & \textbf{中文名称} & \textbf{English Name} & \textbf{章节} \\
\midrule
$\mathbf{v}, \vec{v}$ & 矢量 & Vector & Ch 3, Ch 6 \\
$\hat{\mathbf{n}}$ & 单位法矢 & Unit normal & Ch 3 \\
$\mathbf{a} \cdot \mathbf{b}$ & 点积 / 内积 & Dot product / inner product & Ch 3, Ch 6 \\
$\mathbf{a} \times \mathbf{b}$ & 叉积 & Cross product & Ch 3 \\
$\|\mathbf{v}\|$ 或 $|\mathbf{v}|$ & 矢量模长 & Norm / magnitude & Ch 3, Ch 6 \\
$\mathbf{\hat{e}}_x, \mathbf{\hat{e}}_y, \mathbf{\hat{e}}_z$ & 直角坐标基矢 & Cartesian basis & Ch 3 \\
$\mathbf{\hat{e}}_r, \mathbf{\hat{e}}_\theta, \mathbf{\hat{e}}_\phi$ & 球坐标基矢 & Spherical basis & Ch 3 \\
$\displaystyle \oint_{\partial \Omega}$ & 边界曲线/曲面积分 & Boundary integral & Ch 3 \\
$\mathbf{E}, \mathbf{B}$ & 电场、磁场 & Electric, magnetic field & Ch 3 \\
\bottomrule
\end{tabular}
\end{center}

\section{线性代数 / Linear Algebra}

\begin{center}
\renewcommand{\arraystretch}{1.25}
\begin{tabular}{@{}llll@{}}
\toprule
\textbf{符号} & \textbf{中文名称} & \textbf{English Name} & \textbf{章节} \\
\midrule
$\mathbf{A}, A_{ij}$ & 矩阵及其元素 & Matrix and entries & Ch 6, Ch 7 \\
$\mathbf{A}^T$ & 转置 & Transpose & Ch 6 \\
$\mathbf{A}^{-1}$ & 逆矩阵 & Inverse & Ch 6 \\
$\mathbf{A}^\dagger$ & 共轭转置 & Conjugate transpose (Hermitian) & Ch 7, Ch 9 \\
$\mathbf{A}^+$ & 伪逆（Moore-Penrose）& Pseudoinverse & Ch 8 \\
$\det \mathbf{A}, |\mathbf{A}|$ & 行列式 & Determinant & Ch 6 \\
$\operatorname{tr}(\mathbf{A})$ & 迹 & Trace & Ch 7 \\
$\operatorname{rank}(\mathbf{A})$ & 秩 & Rank & Ch 6, Ch 8 \\
$\ker \mathbf{A}, \operatorname{null}$ & 零空间 / 核 & Null space / kernel & Ch 8 \\
$\operatorname{im} \mathbf{A}, \operatorname{col}$ & 列空间 / 像 & Column space / image & Ch 8 \\
$\mathbf{I}$ & 单位矩阵 & Identity & Ch 6 \\
$\mathbf{0}$ & 零矩阵 / 零矢量 & Zero matrix / vector & Ch 6 \\
$\lambda$ & 特征值 & Eigenvalue & Ch 7 \\
$\mathbf{v}_\lambda$ & 特征矢量 & Eigenvector & Ch 7 \\
$\mathbf{A} = \mathbf{U} \boldsymbol{\Sigma} \mathbf{V}^T$ & 奇异值分解 SVD & Singular Value Decomposition & Ch 7, Ch 8 \\
$\sigma_i$ & 奇异值 & Singular value & Ch 7 \\
$\|\mathbf{A}\|_2, \|\mathbf{A}\|_F$ & 谱范数、Frobenius 范数 & Spectral / Frobenius norm & Ch 7 \\
$\langle \mathbf{u}, \mathbf{v} \rangle$ & 内积 & Inner product & Ch 6 \\
$\mathbf{u} \otimes \mathbf{v}$ & 外积 / 张量积 & Outer / tensor product & Ch 7 \\
$\oplus, \perp$ & 直和、正交 & Direct sum, orthogonal & Ch 8 \\
\bottomrule
\end{tabular}
\end{center}

\section{微分方程 ODE / PDE}

\begin{center}
\renewcommand{\arraystretch}{1.25}
\begin{tabular}{@{}llll@{}}
\toprule
\textbf{符号} & \textbf{中文名称} & \textbf{English Name} & \textbf{章节} \\
\midrule
$y' + p(x)y = q(x)$ & 一阶线性 ODE & First-order linear ODE & Ch 4 \\
$ay'' + by' + cy = f(x)$ & 二阶线性 ODE & Second-order linear ODE & Ch 4 \\
$\dot{\mathbf{x}} = \mathbf{A}\mathbf{x}$ & 状态空间 ODE & State-space ODE & Ch 4 \\
$\omega_0, \zeta$ & 自然频率、阻尼比 & Natural freq, damping ratio & Ch 4 \\
$\dfrac{\partial u}{\partial t} = \alpha \nabla^2 u$ & 热传导方程 & Heat equation & Ch 5 \\
$\dfrac{\partial^2 u}{\partial t^2} = c^2 \nabla^2 u$ & 波动方程 & Wave equation & Ch 5 \\
$\nabla^2 u = 0$ & Laplace 方程 & Laplace equation & Ch 5 \\
$\nabla^2 u = f$ & Poisson 方程 & Poisson equation & Ch 5 \\
$u_t + u u_x = \nu u_{xx}$ & Burgers 方程 & Burgers equation & Ch 5 \\
$u|_{\partial \Omega}, \partial u / \partial n|_{\partial \Omega}$ & Dirichlet / Neumann 边界 & Dirichlet / Neumann BCs & Ch 5 \\
$G(\mathbf{r}, \mathbf{r}')$ & Green 函数 & Green's function & Ch 5 \\
\bottomrule
\end{tabular}
\end{center}

\section{复数与复分析 / Complex Numbers and Analysis}

\begin{center}
\renewcommand{\arraystretch}{1.25}
\begin{tabular}{@{}llll@{}}
\toprule
\textbf{符号} & \textbf{中文名称} & \textbf{English Name} & \textbf{章节} \\
\midrule
$i, j$ & 虚单位 & Imaginary unit & Ch 9 \\
$z = a + bi$ & 复数 & Complex number & Ch 9 \\
$\bar{z}, z^*$ & 共轭复数 & Complex conjugate & Ch 9 \\
$|z|, \arg z$ & 模、辐角 & Modulus, argument & Ch 9 \\
$\operatorname{Re}(z), \operatorname{Im}(z)$ & 实部、虚部 & Real, imaginary parts & Ch 9 \\
$e^{i\theta} = \cos\theta + i\sin\theta$ & Euler 公式 & Euler's formula & Ch 9 \\
$f: \mathbb{C} \to \mathbb{C}$ & 复函数 & Complex function & Ch 9 \\
$\displaystyle \oint_C f(z)\,dz$ & 复围道积分 & Contour integral & Ch 9 \\
$\operatorname{Res}_{z = z_0} f$ & 留数 & Residue & Ch 9 \\
$\mathcal{L}\{f\}(s)$ & Laplace 变换 & Laplace transform & Ch 4, Ch 9 \\
$\mathcal{F}\{f\}(\omega), \hat{f}$ & Fourier 变换 & Fourier transform & Ch 5, Ch 9 \\
$\mathcal{Z}\{x_n\}(z)$ & Z 变换 & Z-transform & Ch 9 \\
$H(s), H(j\omega), H(z)$ & 传递函数 & Transfer function & Ch 4, Ch 9 \\
\bottomrule
\end{tabular}
\end{center}

\section{概率与统计 / Probability and Statistics}

\begin{center}
\renewcommand{\arraystretch}{1.25}
\begin{tabular}{@{}llll@{}}
\toprule
\textbf{符号} & \textbf{中文名称} & \textbf{English Name} & \textbf{章节} \\
\midrule
$\Omega, \mathcal{F}, P$ & 样本空间、$\sigma$-代数、概率测度 & Sample space, sigma-algebra, measure & Ch 10 \\
$P(A)$ & 事件 $A$ 的概率 & Probability of event $A$ & Ch 10 \\
$P(A \mid B)$ & 条件概率 & Conditional probability & Ch 10 \\
$X, Y$ & 随机变量 & Random variable & Ch 10 \\
$\mathbb{E}[X], \mu$ & 期望、均值 & Expectation, mean & Ch 10 \\
$\operatorname{Var}(X), \sigma^2$ & 方差 & Variance & Ch 10 \\
$\sigma$ & 标准差 & Standard deviation & Ch 10 \\
$\operatorname{Cov}(X,Y), \rho$ & 协方差、相关系数 & Covariance, correlation & Ch 10 \\
$f_X(x), F_X(x)$ & PDF、CDF & Density / distribution function & Ch 10 \\
$X \sim \mathcal{N}(\mu, \sigma^2)$ & 正态分布 & Normal distribution & Ch 10 \\
$X \sim \mathrm{Bin}(n,p)$ & 二项分布 & Binomial distribution & Ch 10 \\
$X \sim \mathrm{Poisson}(\lambda)$ & Poisson 分布 & Poisson distribution & Ch 10 \\
$X \sim \mathrm{Exp}(\lambda)$ & 指数分布 & Exponential distribution & Ch 10 \\
$\bar{x}, s^2$ & 样本均值、样本方差 & Sample mean / variance & Ch 11 \\
$H_0, H_1$ & 原假设、备择假设 & Null / alternative hypothesis & Ch 11 \\
$p$ & $p$ 值 & $p$-value & Ch 11 \\
$\alpha, \beta$ & 显著性水平、第二类错误率 & Significance, type-II error rate & Ch 11 \\
$\hat{\theta}$ & 参数估计 & Parameter estimate & Ch 11 \\
$L(\theta), \ell(\theta)$ & 似然、对数似然 & Likelihood, log-likelihood & Ch 11 \\
$\operatorname{KL}(p \| q)$ & KL 散度 & KL divergence & Ch 11 \\
$\chi^2, t, F$ & 卡方、$t$、$F$ 分布 & Chi-squared, $t$, $F$ distributions & Ch 11 \\
\bottomrule
\end{tabular}
\end{center}

\section{数值与优化 / Numerical and Optimization}

\begin{center}
\renewcommand{\arraystretch}{1.25}
\begin{tabular}{@{}llll@{}}
\toprule
\textbf{符号} & \textbf{中文名称} & \textbf{English Name} & \textbf{章节} \\
\midrule
$\kappa(\mathbf{A})$ & 条件数 & Condition number & Ch 12 \\
$\varepsilon_{\text{mach}}$ & 机器精度 & Machine epsilon & Ch 12 \\
$h$ & 步长 & Step size & Ch 12 \\
$\arg\min, \arg\max$ & 极小/极大化变量 & Argmin / argmax & Ch 13 \\
$\min, \max$ & 极小/极大值 & Min / max & Ch 13 \\
$\nabla f, \mathbf{H}$ & 梯度、Hessian & Gradient, Hessian & Ch 13 \\
$\mathcal{L}(\mathbf{x}, \boldsymbol{\lambda})$ & Lagrangian 函数 & Lagrangian & Ch 13 \\
$\boldsymbol{\lambda}, \boldsymbol{\mu}$ & 对偶变量 & Dual variables (KKT multipliers) & Ch 13 \\
$\eta, \alpha$ & 学习率 & Learning rate & Ch 13 \\
$\beta_1, \beta_2$ & Adam 动量参数 & Adam momentum parameters & Ch 13 \\
\bottomrule
\end{tabular}
\end{center}

\section{集合、逻辑与离散 / Sets, Logic, Discrete (Ch 14)}

\begin{center}
\renewcommand{\arraystretch}{1.25}
\begin{tabular}{@{}llll@{}}
\toprule
\textbf{符号} & \textbf{中文名称} & \textbf{English Name} & \textbf{章节} \\
\midrule
$\mathbb{N}, \mathbb{Z}, \mathbb{Q}, \mathbb{R}, \mathbb{C}$ & 自然/整/有理/实/复数集 & Number sets & 通用 \\
$\in, \notin, \subset, \subseteq$ & 属于、包含 & Belongs to, subset & 通用 \\
$\cup, \cap, \setminus$ & 并、交、差 & Union, intersection, difference & Ch 10, Ch 14 \\
$\emptyset$ & 空集 & Empty set & 通用 \\
$|A|$ & 集合基数 & Cardinality & Ch 14 \\
$\binom{n}{k}, C(n,k)$ & 二项系数 & Binomial coefficient & Ch 10, Ch 14 \\
$P(n,k)$ & 排列数 & Number of permutations & Ch 14 \\
$n!$ & 阶乘 & Factorial & Ch 14 \\
$\land, \lor, \neg$ & 与、或、非 & AND, OR, NOT & Ch 14 \\
$\to, \leftrightarrow$ & 蕴含、等价 & Implies, iff & Ch 14 \\
$\forall, \exists$ & 任意、存在 & For all, exists & 通用 \\
$\equiv$ & 恒等 / 同余 & Equivalent / congruent & Ch 14 \\
$G = (V, E)$ & 图（点集 + 边集） & Graph & Ch 14 \\
$\deg(v)$ & 顶点度数 & Vertex degree & Ch 14 \\
$O, \Omega, \Theta$ & 复杂度三符号 & Asymptotic notation & Ch 14 \\
$P, NP, NP\text{-complete}$ & 复杂度类 & Complexity classes & Ch 14 \\
\bottomrule
\end{tabular}
\end{center}

\section{常用希腊字母 / Greek Letters}

\begin{center}
\renewcommand{\arraystretch}{1.2}
\begin{tabular}{@{}llll@{}}
\toprule
\textbf{符号} & \textbf{名称} & \textbf{常见用法} & \textbf{章节} \\
\midrule
$\alpha$ & alpha & 学习率、显著性水平、热扩散率 & Ch 5, Ch 11, Ch 13 \\
$\beta$ & beta & 阻尼、Adam 动量、回归系数 & Ch 4, Ch 11, Ch 13 \\
$\gamma$ & gamma & Euler 常数、折扣因子 & Ch 1 \\
$\delta$ & delta & 任意小、Dirac $\delta$ 函数 & Ch 1, Ch 5 \\
$\varepsilon$ & epsilon & 误差容限、机器精度 & Ch 1, Ch 12 \\
$\zeta$ & zeta & 阻尼比、Riemann zeta & Ch 4 \\
$\eta$ & eta & 学习率、效率 & Ch 13 \\
$\theta, \vartheta$ & theta & 角度、参数 & Ch 3, Ch 11 \\
$\kappa$ & kappa & 曲率、条件数 & Ch 12 \\
$\lambda$ & lambda & 特征值、Lagrange 乘子、波长 & Ch 7, Ch 13 \\
$\mu$ & mu & 均值、磁导率、KKT 乘子 & Ch 10, Ch 13 \\
$\nu$ & nu & 粘度、自由度 & Ch 5 \\
$\xi$ & xi & 中间点、辅助变量 & Ch 1 \\
$\pi$ & pi & 圆周率 & 通用 \\
$\rho$ & rho & 密度、相关系数 & Ch 5, Ch 10 \\
$\sigma$ & sigma & 标准差、奇异值、$\sigma$-代数 & Ch 7, Ch 10 \\
$\tau$ & tau & 时间常数、剪应力 & Ch 4 \\
$\phi, \varphi$ & phi & 相位、标量势 & Ch 3, Ch 9 \\
$\chi$ & chi & 卡方分布、电极化率 & Ch 11 \\
$\psi$ & psi & 波函数、流函数 & Ch 5 \\
$\omega$ & omega & 角频率 & Ch 4, Ch 9 \\
$\Delta$ & Delta & 增量、Laplace 算子 & Ch 1, Ch 5 \\
$\Sigma$ & Sigma & 求和符号、协方差矩阵 & Ch 1, Ch 10 \\
$\Pi$ & Pi & 求积符号 & Ch 10 \\
$\Phi$ & Phi & 标量势、累积正态分布 & Ch 3, Ch 10 \\
$\Omega$ & Omega & 区域、样本空间、阻抗 & Ch 5, Ch 10 \\
\bottomrule
\end{tabular}
\end{center}

\section{使用建议 / How to Use This Table}

\begin{enumerate}
\item \textbf{读外文论文}——\textbf{遇到不认识的符号——\textbf{先查\textbf{对应表}}}——\textbf{再回到正文章节查\textbf{定义与直觉}}。
\item \textbf{写论文 / 写代码注释}——\textbf{用本表保持\textbf{符号一致性}}——\textbf{别同一篇里 $\sigma$ 既表标准差又表奇异值}。
\item \textbf{讲课 / 写文档}——\textbf{先列\textbf{符号表}}（"Notation"）——\textbf{是学术写作的\textbf{专业素养}}。
\item \textbf{看物理 vs 数学教材}——\textbf{注意\textbf{同一符号在不同领域含义不同}}：例如 $j$ 在数学里是虚单位、在电气工程里也是虚单位（避开电流 $i$）、在物理里是角动量量子数。
\end{enumerate}

\clearpage

% =====================================================================
%  附录 B · Python 模块索引 / Module Index
%  13 + 1 个模块——配文件名、章节、主要函数、安装指南、GitHub 链接
% =====================================================================

\chapter{Python 模块索引 / Module Index}
\label{app:modules}

\begin{quote}\itshape
本书 14 章——\textbf{每一章都附\textbf{一个 Python 模块}}——\textbf{把数学讲完——就让你能跑代码}。\\
\textbf{本附录是\textbf{所有模块的总索引}}：\textbf{文件名、依赖、主要函数、安装方法}——\textbf{一页找到任何工具}。
\end{quote}

\section{模块总表 / Module Overview}

\begin{center}
\renewcommand{\arraystretch}{1.25}
\begin{tabular}{@{}cllc@{}}
\toprule
\textbf{Ch} & \textbf{文件名} & \textbf{主题} & \textbf{状态} \\
\midrule
1  & \texttt{calculus.py}      & 极限、导数、积分、数值微分             & 必备 \\
2  & \texttt{multivar.py}      & 偏导、梯度、Hessian、Jacobian        & 必备 \\
3  & \texttt{vector\_analysis.py} & 散度、旋度、Laplace、矢量积分        & 必备 \\
4  & \texttt{ode.py}           & 一阶/二阶 ODE、Laplace 变换、状态空间   & 必备 \\
5  & \texttt{pde.py}           & 有限差分热方程/波方程、谱方法           & 必备 \\
6  & \texttt{linalg\_intro.py} & 矩阵基本运算、Gauss 消元、LU 分解         & 必备 \\
7  & \texttt{linalg\_core.py}  & 特征分解、SVD、QR、Schur                 & 必备 \\
8  & \texttt{linalg\_apps.py}  & PCA、最小二乘、推荐系统、PageRank        & 必备 \\
9  & \texttt{complex.py}       & 复函数、留数、FFT、滤波、传递函数        & 必备 \\
10 & \texttt{probability.py}   & 分布抽样、Monte Carlo、Bayes、Markov     & 必备 \\
11 & \texttt{statistics.py}    & 假设检验、回归、bootstrap、MLE           & 必备 \\
12 & \texttt{numerical.py}     & 插值、积分求积、Runge-Kutta、稳定性      & 必备 \\
13 & \texttt{optimization.py}  & GD、Newton、BFGS、SGD、Adam、LP、KKT      & 必备 \\
14 & \texttt{discrete.py}      & 组合、图论、布尔逻辑、TSP                  & 可选 \\
\bottomrule
\end{tabular}
\end{center}

\textbf{共 \textbf{13 + 1 = 14 个模块}}——\textbf{每个 200--500 行 Python}——\textbf{都\textbf{独立可用}}——\textbf{也可联合使用}。

\section{模块详表 / Module Details}

\subsection{\texttt{calculus.py} —— 第 1 章}

\begin{lstlisting}[language=Python]
from calculus import (
    numerical_derivative,   # 中心差分
    numerical_integral,     # Simpson + Trapezoid
    taylor_series,          # 解析 Taylor 展开（sympy）
    plot_function,          # 函数 + 导数 + 积分三联图
)
\end{lstlisting}
\textbf{依赖}：\texttt{numpy, scipy, sympy, matplotlib}

\subsection{\texttt{multivar.py} —— 第 2 章}

\begin{lstlisting}[language=Python]
from multivar import (
    gradient_numeric,       # numpy.gradient 多维包装
    hessian_numeric,        # 二阶差分 Hessian
    jacobian_numeric,
    double_integral,        # scipy.integrate.dblquad
    triple_integral,        # scipy.integrate.tplquad
)
\end{lstlisting}
\textbf{依赖}：\texttt{numpy, scipy}

\subsection{\texttt{vector\_analysis.py} —— 第 3 章}

\begin{lstlisting}[language=Python]
from vector_analysis import (
    divergence,             # div F
    curl,                   # curl F
    laplacian,              # nabla^2 f
    flux_through_surface,
    circulation,            # 闭曲线环流
)
\end{lstlisting}
\textbf{依赖}：\texttt{numpy, sympy}

\subsection{\texttt{ode.py} —— 第 4 章}

\begin{lstlisting}[language=Python]
from ode import (
    solve_ode,              # scipy.integrate.solve_ivp 封装
    transfer_function,      # scipy.signal.TransferFunction
    bode_plot,              # 频率响应图
    state_space_simulate,
    laplace_transform,      # 解析 Laplace（sympy）
)
\end{lstlisting}
\textbf{依赖}：\texttt{numpy, scipy, sympy, matplotlib}

\subsection{\texttt{pde.py} —— 第 5 章}

\begin{lstlisting}[language=Python]
from pde import (
    heat_1d_fdm,            # 1D 热方程有限差分
    wave_1d_fdm,
    laplace_2d_jacobi,      # 2D Laplace 迭代解
    burgers_1d,
    spectral_solver,        # FFT 谱方法
)
\end{lstlisting}
\textbf{依赖}：\texttt{numpy, matplotlib, scipy}

\subsection{\texttt{linalg\_intro.py} —— 第 6 章}

\begin{lstlisting}[language=Python]
from linalg_intro import (
    gauss_elimination,      # 显式高斯消元（教学用）
    lu_decomp,              # 分解为 L 和 U
    solve_linear,           # Ax = b
    matrix_inverse,
    determinant,
)
\end{lstlisting}
\textbf{依赖}：\texttt{numpy}

\subsection{\texttt{linalg\_core.py} —— 第 7 章}

\begin{lstlisting}[language=Python]
from linalg_core import (
    eig_decomp,             # A = V Lambda V^{-1}
    svd_decomp,             # A = U Sigma V^T
    qr_decomp,
    schur_decomp,
    spectral_norm,
    frobenius_norm,
)
\end{lstlisting}
\textbf{依赖}：\texttt{numpy, scipy}

\subsection{\texttt{linalg\_apps.py} —— 第 8 章}

\begin{lstlisting}[language=Python]
from linalg_apps import (
    pca,                    # 主成分分析
    least_squares,          # 最小二乘 + 伪逆
    page_rank,              # 网页排名（特征向量法）
    recommender_svd,        # SVD 推荐系统
    image_compress_svd,     # 灰度图低秩近似
)
\end{lstlisting}
\textbf{依赖}：\texttt{numpy, scipy, scikit-learn (可选)}

\subsection{\texttt{complex.py} —— 第 9 章}

\begin{lstlisting}[language=Python]
from complex import (
    residue,                # 留数计算
    fft, ifft,              # numpy.fft 封装
    bode_plot,
    transfer_eval,          # H(jw) 求值
    butterworth_filter,
)
\end{lstlisting}
\textbf{依赖}：\texttt{numpy, scipy, matplotlib}

\subsection{\texttt{probability.py} —— 第 10 章}

\begin{lstlisting}[language=Python]
from probability import (
    sample_normal, sample_binomial, sample_poisson,
    monte_carlo_integrate,
    bayes_update,
    markov_chain_simulate,
    central_limit_demo,
)
\end{lstlisting}
\textbf{依赖}：\texttt{numpy, scipy.stats, matplotlib}

\subsection{\texttt{statistics.py} —— 第 11 章}

\begin{lstlisting}[language=Python]
from statistics_module import (
    t_test, chi2_test, anova,
    linear_regression,
    bootstrap_ci,
    mle_normal,
    kl_divergence,
)
\end{lstlisting}
\textbf{依赖}：\texttt{numpy, scipy.stats, statsmodels (可选)}\\
\textbf{注}：\texttt{statistics} 是 Python 标准库——\textbf{模块文件名实际为 \texttt{statistics\_module.py}}。

\subsection{\texttt{numerical.py} —— 第 12 章}

\begin{lstlisting}[language=Python]
from numerical import (
    lagrange_interp, cubic_spline,
    simpson_rule, gauss_quadrature,
    runge_kutta_4,
    condition_number,
    machine_epsilon,
)
\end{lstlisting}
\textbf{依赖}：\texttt{numpy, scipy}

\subsection{\texttt{optimization.py} —— 第 13 章}

\begin{lstlisting}[language=Python]
from optimization import (
    gradient_descent, newton_method, bfgs,
    sgd, adam,
    lp_simplex,             # scipy.optimize.linprog
    kkt_solver,             # 等式约束
    plot_convergence,
)
\end{lstlisting}
\textbf{依赖}：\texttt{numpy, scipy, matplotlib, torch (用于 SGD/Adam 真实示例)}

\subsection{\texttt{discrete.py} —— 第 14 章（可选）}

\begin{lstlisting}[language=Python]
from discrete import (
    n_perm, n_comb, derangement,
    shortest_path,          # Dijkstra
    min_spanning_tree,      # Kruskal
    graph_coloring,
    truth_table,
    tsp_bruteforce,
)
\end{lstlisting}
\textbf{依赖}：\texttt{networkx, matplotlib}

\section{安装指南 / Installation}

\subsection{推荐环境：Anaconda + Python 3.11+}

\begin{lstlisting}[language=bash]
# 1. 安装 Miniconda（轻量版 Anaconda）
#    https://docs.conda.io/en/latest/miniconda.html

# 2. 创建虚拟环境
conda create -n math-rendu python=3.11
conda activate math-rendu

# 3. 一键安装本书全部依赖
pip install -r requirements.txt
\end{lstlisting}

\subsection{\texttt{requirements.txt}}

\begin{lstlisting}
numpy>=1.26
scipy>=1.11
sympy>=1.12
matplotlib>=3.8
pandas>=2.1
scikit-learn>=1.3
statsmodels>=0.14
networkx>=3.2
torch>=2.1            # 第 13 章 Adam / SGD 实战
jupyterlab>=4.0       # 推荐用 Jupyter 跑示例
\end{lstlisting}

\subsection{分步安装 / Minimal Install}

\textbf{若你只看前 8 章}（微积分 + 线代）：

\begin{lstlisting}[language=bash]
pip install numpy scipy sympy matplotlib
\end{lstlisting}

\textbf{若加上概率统计}（Ch 10--11）：

\begin{lstlisting}[language=bash]
pip install pandas scikit-learn statsmodels
\end{lstlisting}

\textbf{若到优化 + 深度学习}（Ch 13）：

\begin{lstlisting}[language=bash]
pip install torch
\end{lstlisting}

\textbf{若选修离散}（Ch 14）：

\begin{lstlisting}[language=bash]
pip install networkx
\end{lstlisting}

\subsection{验证安装}

\begin{lstlisting}[language=Python]
# 在 Python REPL 里运行
import numpy, scipy, sympy, matplotlib
import pandas, sklearn, statsmodels
import networkx
print("All packages OK!")
\end{lstlisting}

\section{GitHub 仓库 / Repository}

\begin{tcolorbox}[essbox={本书代码仓库}]
\textbf{GitHub}：\url{https://github.com/<your-handle>/math-rendu}\\[4pt]
\textbf{结构}：
\begin{itemize}
\item \texttt{book/}\quad —— LaTeX 源码 + PDF
\item \texttt{modules/}\quad —— 14 个 Python 模块
\item \texttt{notebooks/}\quad —— 每章 1 个 Jupyter notebook（"Aha 实战"）
\item \texttt{tests/}\quad —— pytest 单元测试
\item \texttt{requirements.txt}\quad —— 依赖列表
\item \texttt{README.md}\quad —— 快速上手指南
\end{itemize}
\end{tcolorbox}

\subsection{克隆与运行}

\begin{lstlisting}[language=bash]
git clone https://github.com/<your-handle>/math-rendu.git
cd math-rendu
pip install -r requirements.txt

# 运行示例 notebook
jupyter lab notebooks/ch01_calculus_aha.ipynb
\end{lstlisting}

\subsection{贡献 / Contributing}

\textbf{欢迎以下贡献}：

\begin{itemize}
\item \textbf{勘误}（typo、公式错误）—— 直接提 Issue 或 PR
\item \textbf{补充示例}（你工业里用到的真实案例）—— 提 PR
\item \textbf{翻译}（英文版 / 日文版 / 韩文版）—— 联系作者
\end{itemize}

\section{使用建议 / How to Use the Modules}

\begin{enumerate}
\item \textbf{初学者}：\textbf{读每一章——\textbf{打开对应 notebook}——\textbf{改参数、跑代码、画图}——\textbf{这比做习题题高效 10 倍}}。
\item \textbf{工程师}：\textbf{把模块作为\textbf{速查工具箱}}——\textbf{需要 PCA 时 \texttt{from linalg\_apps import pca}}——\textbf{3 秒搞定}。
\item \textbf{研究生}：\textbf{把模块作为\textbf{教学示范代码}}——\textbf{自己用 \texttt{numpy} 重新实现一遍}——\textbf{比读 \texttt{scipy} 源码温和得多}。
\item \textbf{教师}：\textbf{所有模块和 notebook \textbf{MIT License}}——\textbf{可以\textbf{自由用于课堂}}——\textbf{只需保留作者署名}。
\end{enumerate}

\begin{tcolorbox}[essbox={最后一句}]
\textbf{数学是\textbf{工程师的内功}}——\textbf{Python 是\textbf{把内功打出去的拳脚}}。\textbf{这本书把两者打通}——\textbf{从此你看见公式就能\textbf{秒变代码}}——\textbf{看见代码就能\textbf{秒懂数学}}。\\[4pt]
\textbf{祝你——\textbf{打通任督二脉}}。
\end{tcolorbox}

\clearpage


\end{document}
